\documentclass[a4paper,11pt]{report}

% --- 1. CONFIGURAZIONE LINGUA E CODIFICA ---
\usepackage[utf8]{inputenc}
\usepackage[T1]{fontenc}
\usepackage[italian]{babel}
\usepackage{lmodern} 

% --- 2. LAYOUT E FORMATTAZIONE ---
\usepackage[a4paper, margin=1in]{geometry}
\usepackage{fancyhdr}
\usepackage{titlesec}
\usepackage{tocloft}
\usepackage[bottom]{footmisc} 
\usepackage{float}
\usepackage{calc}
\usepackage{multicol}
\usepackage[table,xcdraw]{xcolor} % Caricato con opzioni tabella
\usepackage{enumitem}
\usepackage{booktabs} 
\usepackage{caption}
\usepackage{subcaption}
\usepackage{chngcntr} 

% --- 3. MATEMATICA E FISICA ---
\usepackage{amsmath}
\usepackage{physics}
\usepackage{mathtools}
\usepackage{amsfonts}
\usepackage{amssymb}
\usepackage{bm}      
\usepackage{siunitx} 
\usepackage{xparse}  

% --- 4. GRAFICA E TIKZ ---
\usepackage{graphicx}
\usepackage{tikz}
\usepackage{pgfplots}
\usepackage{tikz-3dplot}
\usetikzlibrary{circuits.ee.IEC, arrows.meta, calc, positioning, 3d, decorations.markings, angles, quotes, decorations.pathmorphing}
\pgfplotsset{compat=1.18}

% --- 5. NAVIGAZIONE ---
\usepackage{hyperref}
\hypersetup{
    colorlinks=true,
    linkcolor=black, % Nero per la stampa, cambia in blue se solo digitale
    urlcolor=blue,
    pdftitle={Dispense di Fisica 2},
    pdfauthor={Politecnico di Torino}
}

% --- 6. COMANDI PERSONALIZZATI & STILE ---
\renewcommand{\contentsname}{Indice}
\renewcommand{\tablename}{Tabella}
\pagecolor{white}

% Stile Sezioni con linea orizzontale
\NewDocumentCommand{\fancysectiontitle}{m}{%
  \noindent
  \textbf{\thesection\quad #1}%
  \hspace{0.5em}%
  \raisebox{0.6ex}{\rule{\linewidth - \widthof{\textbf{\thesection\quad #1}} - 0.6em}{0.4pt}}%
}

\titleformat{\section}
  {\normalfont\Large\bfseries}
  {}
  {0pt}
  {\fancysectiontitle}

\setcounter{tocdepth}{1}

% --- CONFIGURAZIONE HEADER/FOOTER (Stile Fancy) ---
\pagestyle{fancy}
\fancyhf{} 
\setlength{\headheight}{15pt} 

\lhead{\itshape \nouppercase{\rightmark}} 
\rhead{\itshape \nouppercase{\leftmark}}                                   
\lfoot{}                         
\cfoot{}                                   
\rfoot{\thepage}                           

\renewcommand{\headrulewidth}{0.4pt}     
\renewcommand{\footrulewidth}{0.4pt}     

% Marcatori
\renewcommand{\sectionmark}[1]{\markright{\thesection\quad #1}}
\renewcommand{\chaptermark}[1]{\markboth{#1}{}}

% Gestione Note a piè pagina (reset per capitolo)
\counterwithin*{footnote}{chapter} 
\renewcommand{\thefootnote}{\thechapter.\arabic{footnote}} 


% --- 7. BOX APPROFONDIMENTO (PROFESSIONALE - SFONDO GRIGIO VISIBILE) ---
\usepackage[most]{tcolorbox} 

% Definizione colori
\definecolor{formalblue}{RGB}{0, 51, 102} 
% Scurito al 10% di grigio (0.90) invece del 5% (0.95) per renderlo visibile
\definecolor{formalbackground}{gray}{0.95} 

\newtcolorbox{approfondimento}[1][]{
    enhanced,
    breakable,
    frame hidden,                       % Nasconde il bordo (box aperto)
    interior style={color=formalbackground}, % Forza il riempimento dello sfondo
    borderline west={3pt}{0pt}{formalblue},  % Barra laterale
    coltitle=formalblue,                
    fonttitle=\bfseries\sffamily\large, 
    title={Approfondimento},            
    sharp corners,
    detach title,                       
    before upper={\tcbtitle\par\vspace{0.3em}}, 
    left=1em, right=1em, top=1em, bottom=1em, 
    parbox=false,                       
    #1
}


% ==========================================
% INIZIO DOCUMENTO
% ==========================================
\begin{document}

% --- FRONT MATTER (Numeri Romani) ---
\pagenumbering{Roman}   % Attiva numerazione i, ii, iii...
\pagestyle{plain}       % Stile pulito (niente header, solo numero in basso)

% Indice
\tableofcontents
\clearpage

% Qui puoi inserire la Prefazione o Introduzione alle dispense
% \chapter*{Prefazione}
% Testo della prefazione...
% \clearpage

% --- MAIN MATTER (Numeri Arabi) ---
\cleardoublepage        % Assicura che il cap. 1 inizi su pagina destra (opzionale)
\pagenumbering{arabic}  % Riparte da 1, 2, 3...
\pagestyle{fancy}

\chapter{Elettrostatica}

Il dibattito sulla cosiddetta "azione a distanza'' ha segnato profondamente l'evoluzione del pensiero scientifico. Accanto all'interazione gravitazionale, storicamente modellata come forza istantanea, l'osservazione dei fenomeni elettrici ha imposto un radicale cambio di paradigma, rendendo necessaria l'introduzione formale del concetto di \emph{campo}.
Così come la \emph{massa} è la sorgente dell'interazione gravitazionale, in questo capitolo si introdurrà la \emph{carica elettrica}, proprietà fondamentale della materia e sorgente dell'interazione elettrostatica.

\section{La forza elettrostatica}

L'evidenza sperimentale mostra che, dati due corpi puntiformi carichi $q_1$ e $q_2$ posti nel vuoto a distanza reciproca $r$, si manifesta un'interazione\footnote{Si assume l'approssimazione $r \gg d$, dove $d$ rappresenta la dimensione caratteristica dei corpi. La carica $q$ è scalare e dotata di segno: $q \in \mathbb{R}$.} dinamica lungo la retta congiungente i due corpi.
Tale forza risulta direttamente proporzionale al prodotto delle cariche e inversamente proporzionale al quadrato della distanza relativa. In termini vettoriali:
\begin{equation}
    \bm{F}_{1 \to 2} \propto \frac{q_1 q_2}{|\bm{r}_1-\bm{r}_2|^2} \hat{\bm{u}}_r = \frac{q_1 q_2}{r^2} \hat{\bm{u}}_r,
    \label{eq:proporzionalita_coulomb}
\end{equation}
dove $\hat{\bm{u}}_r$ è il versore diretto dalla carica sorgente alla carica esploratrice e, evidentemente, $r = |\bm{r}_1-\bm{r}_2|$.
Si osserva un'immediata \textit{analogia formale} con la Legge di gravitazione universale di Newton ($F \propto m_1 m_2 / r^2$). Tuttavia, sussiste una differenza sostanziale di natura fisica: mentre la massa è una grandezza definita positiva (rendendo l'interazione gravitazionale esclusivamente \textit{attrattiva}), la carica elettrica può assumere valori positivi o negativi. Di conseguenza, la forza elettrostatica (anche conosciuta come \textit{forza di Coulomb}) descrive un'interazione duale: \textit{repulsiva} per cariche omonime (segno concorde, prodotto $q_1q_2 > 0$) e \textit{attrattiva} per cariche eteronome (segno discorde, prodotto $q_1q_2 < 0$).


\begin{figure}[h]
    \centering
    % --- FIGURA 1: REPULSIONE (Segni Concordi) ---
    \begin{minipage}{0.45\textwidth}
        \centering
        \begin{tikzpicture}[>=Stealth, scale=0.7]
            % Definizioni stili
            \tikzstyle{charge}=[circle, draw, minimum size=0.8cm, thick]
            
            % Cariche (Positiva - Positiva)
            \node[charge, fill=red!10] (q1) at (0,0) {$+q_1$};
            \node[charge, fill=red!10] (q2) at (4,0) {$+q_2$};
            
            % Vettori Forza (Repulsivi - verso l'esterno)
            \draw[->, thick, red] (-0.6,0) -- (-1.5, 0) node[above] {$\vec{F}_{21}$};
            \draw[->, thick, red] (4.6,0) -- (5.5, 0) node[above] {$\vec{F}_{12}$};
            
            % Linea distanza
            \draw[dashed] (0,-0.8) -- (0,-1.2);
            \draw[dashed] (4,-0.8) -- (4,-1.2);
            \draw[<->] (0,-1) -- (4,-1) node[midway, below] {$r$};
            
            % Didascalia interna (opzionale)
            \node[below=1.5cm] at (2,0) {\small \textbf{Repulsione} ($q_1 q_2 > 0$)};
        \end{tikzpicture}
    \end{minipage}
    \hfill
    % --- FIGURA 2: ATTRAZIONE (Segni Discordi) ---
    \begin{minipage}{0.45\textwidth}
        \centering
        \begin{tikzpicture}[>=Stealth, scale=0.7]
            % Definizioni stili
            \tikzstyle{charge}=[circle, draw, minimum size=0.8cm, thick]
            
            % Cariche (Positiva - Negativa)
            \node[charge, fill=red!10] (q1) at (0,0) {$+q_1$};
            \node[charge, fill=blue!10] (q2) at (4,0) {$-q_2$};
            
            % Vettori Forza (Attrattivi - verso l'interno)
            \draw[->, thick, blue!80!black] (0.6,0) -- (1.5, 0) node[above] {$\vec{F}_{21}$};
            \draw[->, thick, blue!80!black] (3.4,0) -- (2.5, 0) node[above] {$\vec{F}_{12}$};
            
            % Linea distanza
            \draw[dashed] (0,-0.8) -- (0,-1.2);
            \draw[dashed] (4,-0.8) -- (4,-1.2);
            \draw[<->] (0,-1) -- (4,-1) node[midway, below] {$r$};
            
             % Didascalia interna (opzionale)
            \node[below=1.5cm] at (2,0) {\small \textbf{Attrazione} ($q_1 q_2 < 0$)};
        \end{tikzpicture}
    \end{minipage}
    \label{fig:coulomb_forces}
    \caption{Rappresentazione vettoriale della forza di Coulomb. Si noti l'analogia con la forza gravitazionale, evidente nel caso 2, con cariche discordi.}
\end{figure}

\begin{approfondimento}[ title={La costante di proporzionalità: \texorpdfstring{$\bm{\varepsilon_0}$}{epsilon zero}}]
\noindent
Come ogni grandezza in fisica, la forza elettrostatica ha bisogno di essere definita ``totalmente'', valutando il valore delle costanti di proporzionalità. Si definisce inizialmente l'unità di misura della carica, il \textit{Coulomb}: $[q] = \text{C}$.
Si consideri che $\SI{1}{\coulomb}$ corrisponde alla carica che, posta a distanza di $\SI{1}{\meter}$ da una carica uguale nel vuoto, risente di una forza elettrostatica di modulo $|\bm{F}| = F = c^2 \cdot 10^{-7} \, \si{\newton}$, con $c$ la velocità della luce.
\begin{equation}
    F = k \frac{q_1q_2}{r^2} = k \frac{\SI{1}{\coulomb}\cdot \SI{1}{\coulomb}}{(\SI{1}{\meter})^2} \implies k = c^2 \cdot 10^{-7} \, [\si{\newton\meter\squared\per\coulomb\squared}],
\end{equation}
dove si è esplicitata la proporzionalità tramite la costante $k$. Si ricorre ora a una relazione introdotta da Maxwell che verrà successivamente dimostrata:
\begin{equation}
    c = \frac{1}{\sqrt{\varepsilon_0 \mu_0}}.
\end{equation}
Ricordando che $\mu_0 = 4\pi \times 10^{-7} \, \si{\tesla\meter\per\ampere}$, si ottiene:
\begin{equation}
    c^2 = \frac{1}{\varepsilon_0\cdot 4\pi \times 10^{-7}} \rightarrow c^2 \cdot 10^{-7} = \frac{1}{4\pi\varepsilon_0 } = k.
\end{equation}
La forza elettrostatica si scrive quindi come:
\begin{equation}
    \boxed{\bm{F} = \frac{1}{4 \pi \varepsilon_0} \frac{q_1q_2}{r^2}\hat{\bm{u}}_r}. \label{eq:coulomb_completa}
\end{equation}
    \end{approfondimento}

\subsection{Conservatività}
Si vuole ora valutare il lavoro della forza elettrostatica esercitata da una carica puntiforme fissa $q_s$ su una carica di prova $q_p$ in moto lungo una traiettoria $\gamma$ qualsiasi, dalla posizione di partenza $A$ a quella di arrivo $B$. 
\begin{figure}[H]
    \centering
    % Impostiamo una vista 3D "simulata"
    \begin{tikzpicture}[x={(-0.5cm,-0.4cm)}, y={(1cm,0cm)}, z={(0cm,1cm)}, scale=1, >=Stealth]

        % --- 1. GLI ASSI CARTESIANI ---
        \coordinate (O) at (0,0,0);
        \draw[->, black] (O) -- (4,0,0) node[anchor=north east] {$x$};
        \draw[->, black] (O) -- (0,5,0) node[anchor=north west] {$y$};
        \draw[->, black] (O) -- (0,0,4) node[anchor=south] {$z$};

        % --- 2. LA CARICA SORGENTE ---
        \shade[ball color=red!20] (O) circle (0.15cm);
        \node at (O) [above left=0.1cm] {$q_S$};

        % --- 3. LA TRAIETTORIA GAMMA ---
        \coordinate (A) at (4,3,1);
        \coordinate (B) at (1,4,4);
        
        \draw[thick, densely dashed, black] (A) .. controls (3,3,1) and (2,4,2) .. (B) node[midway, below right] {$\gamma$};
        
        \node[circle, fill, inner sep=1pt, label=below:$A$] at (A) {};
        \node[circle, fill, inner sep=1pt, label=above:$B$] at (B) {};

        % --- 4. IL PUNTO P ---
        \coordinate (P) at (2.15, 3.40, 1.7);
        \node[circle, fill=black, inner sep=1.5pt] at (P) {};
        \node[above left] at (P) {$q_p$};

        % --- 5. I VETTORI ---
        
        % Vettore posizione r
        \draw[->, thick, black!80!black] (O) -- (P) node[midway, above] {$\bm{r}$};
        
        % Vettore Forza F
        \coordinate (F_tip) at ($(O)!1.45!(P)$);
        \draw[->, very thick, violet] (P) -- (F_tip) node[right] {$\bm{F}$};
        
        % Vettore spostamento dl
        \coordinate (dl_tip) at ($(P) + (-0.3, 0.5, 0.8)$);
        \draw[->, thick, black] (P) -- (dl_tip) node[anchor=south west] {$d\bm{l}$};

        % --- 6. L'ANGOLO THETA ---
        \pic [draw, "$\theta$", angle radius=0.6cm, angle eccentricity=1.3] {angle = F_tip--P--dl_tip};

        % --- 7. ELEMENTI DI SUPPORTO ---
        \draw[dotted, thin, red] (O) -- (P);

    \end{tikzpicture}

    \label{fig:lavoro_campo_elettrico}
    \caption{Schema per il calcolo del lavoro del campo elettrostatico. La forza è radiale, il lavoro infinitesimo dipende solo dalla variazione radiale $dr$.}
\end{figure}


\noindent
Senza perdita di generalità, si assume che la carica sorgente $q_S$ coincida con l'origine del sistema di riferimento.
Considerando uno spostamento infinitesimo $d\bm{l}$ lungo la traiettoria $\gamma$, il lavoro elementare compiuto dalla forza elettrostatica è:
\begin{equation}
    \delta W = \bm{F} \cdot d\bm{l} = \frac{1}{4 \pi \varepsilon_0} \frac{q_S q_p}{r^2} (\hat{\bm{u}}_r \cdot d\bm{l}).
\end{equation}
L'osservazione geometrica fondamentale risiede nel prodotto scalare: $\hat{\bm{u}}_r \cdot d\bm{l} = dl \cos\theta = dr$, dove $dr$ rappresenta la variazione infinitesima della coordinata radiale.
Integrando lungo l'intero percorso $A \to B$:
\begin{equation}
    W_{AB} = \int_{r_A}^{r_B} \frac{1}{4 \pi \varepsilon_0} \frac{q_S q_p}{r^2} dr = \frac{q_S q_p}{4 \pi \varepsilon_0} \left[ -\frac{1}{r} \right]_{r_A}^{r_B} = \frac{q_S q_p}{4 \pi \varepsilon_0} \left(\frac{1}{r_A}-\frac{1}{r_B}\right).
    \label{eq:lavoro_integrato}
\end{equation}
Il risultato mostra che il lavoro non dipende dalla traiettoria $\gamma$ ma esclusivamente dalle posizioni radiali iniziale ($r_A$) e finale ($r_B$).
Tale proprietà identifica il campo vettoriale della forza elettrostatica come un \textit{campo conservativo}: è pertanto possibile definire una funzione scalare di stato $U$, detta \textit{energia potenziale elettrostatica}, per cui valga:
\begin{equation}
\boxed{\bm{F} = -\nabla U}.
\label{eq:conservatività}
\end{equation}
Esplicitando l'operatore differenziale gradiente in coordinate cartesiane\footnote{La trattazione assume come dominio lo spazio euclideo tridimensionale $\mathbb{R}^3$.}, la relazione si scrive componente per componente come:
\begin{equation}
\bm{F} = -\left( \frac{\partial U}{\partial x}\hat{\bm{u}}_x + \frac{\partial U}{\partial y}\hat{\bm{u}}_y + \frac{\partial U}{\partial z}\hat{\bm{u}}_z \right).
\end{equation}

\subsubsection{L'energia potenziale elettrostatica}
In piena analogia con il caso gravitazionale, la funzione $U$ rappresenta l'energia immagazzinata nel sistema in virtù della posizione relativa delle cariche. 
Dal punto di vista analitico, poiché $U$ è ottenuta tramite integrazione, essa è definita a meno di una costante arbitraria. Per interazioni che decadono come $r^{-2}$ (come nel caso coulombiano), è convenzione standard fissare lo \emph{zero} del potenziale all'infinito ($\lim_{r \to \infty} U = 0$) per semplificare i calcoli. Esemplificativamente, nel caso elementare di un sistema costituito da due cariche puntiformi $q_1$ e $q_2$ poste nel vuoto a distanza $r$, l'energia potenziale di interazione è data da:
\begin{equation}
U(\bm{r}) = \frac{1}{4 \pi \varepsilon_0}\frac{q_1q_2}{|\bm{r}_1-\bm{r}_2|}.
\end{equation}

\noindent
Fondamentale conseguenza della conservatività è la modalità di calcolo del lavoro. Sfruttando la relazione locale $\bm{F} = -\nabla U$ e la definizione di \textit{differenziale totale} di una funzione scalare ($dU = \nabla U \cdot d\bm{l}$), si ottiene:
\begin{equation}
    W_{A \to B} = \int_A^B \bm{F} \cdot d\bm{l} = - \int_A^B \nabla U \cdot d\bm{l} = - \int_A^B dU = -\Bigl[ U(B) - U(A) \Bigr] = -\Delta U.
    \label{eq:teorema_lavoro_energia}
\end{equation}
Il lavoro compiuto dalle forze del campo corrisponde dunque all'opposto della variazione dell'energia potenziale.


\section{I principi dell'elettrostatica}
In questa sezione si procede a enucleare i due essenziali fondamenti dell'elettrostatica: il \textit{Principio di conservazione della carica} e il \textit{Principio di sovrapposizione degli effetti}. Si tratta in ambedue i casi di risultati provenienti da evidenze sperimentali.

\subsection{Principio di conservazione della carica}
\begin{quote}
\textit{La carica netta totale in un sistema isolato è costante, a prescindere dai processi che avvengono al suo interno.}
\end{quote}
Si può formalizzare matematicamente questa asserzione in questo modo: la somma algebrica delle singole cariche (con il loro proprio segno) in un sistema isolato deve rimanere costante. Da cui,
\begin{equation}
    \sum_i q_i = \text{cost.}
\end{equation}

\subsubsection{Estensione al continuo}
Il Principio di conservazione mantiene validità anche nel caso in cui le cariche non siano discrete ma distribuite con continuità nello spazio. In tal caso, la sommatoria nel limite del continuo viene sostituita da un integrale.
Data una distribuzione volumetrica di carica descritta dalla densità $\rho(\bm{r})$ confinata in una regione di spazio $\Omega$:
\begin{equation}
     \int_\Omega \rho (\bm{r}) \, dV = \text{cost.}
\end{equation}

\subsection{Principio di sovrapposizione}
\begin{quote}
    \textit{La forza risultante agente su una carica di prova è data dalla somma vettoriale delle singole forze esercitate da ciascuna delle cariche sorgente presenti nel sistema, considerate come se agissero singolarmente.}
\end{quote}
In un sistema discreto costituito da $N$ cariche puntiformi $q_1, \dots, q_N$, la forza totale $\bm{F}_i$ agente sulla carica $i$-esima è:
\begin{equation}
    \bm{F}_i = \sum_{\substack{j=1 \\ j \neq i}}^{N} \bm{F}_{j \to i} = \frac{q_i}{4\pi\varepsilon_0} \sum_{\substack{j=1 \\ j \neq i}}^{N} \frac{q_j}{r_j^2} \bm{u_{r_j}}.
    \label{eq:sovrapposizione_discreta}
\end{equation}


\begin{figure}[htbp]
    \centering
    \begin{tikzpicture}[scale=1, >=Stealth]
        
        % --- 1. DEFINIZIONE DELLE COORDINATE ---
        \coordinate (Q1) at (0, 0);      % Carica 1
        \coordinate (Q2) at (5, 0.5);    % Carica 2
        \coordinate (P) at (2.5, 3.5);   % Punto P di osservazione

        % --- 2. LINEE DI CONGIUNZIONE (r1, r2) ---
        % Tratteggiate grigie dalle cariche al punto P
        \draw[dashed, gray, thin] (Q1) -- (P) node[midway, left] {$r_1$};
        \draw[dashed, gray, thin] (Q2) -- (P) node[midway, right] {$r_2$};

        % --- 3. DEFINIZIONE DEI VETTORI CAMPO (E1, E2) ---
        % Calcolo manuale o "a occhio" per una buona resa grafica
        % E1: Prolungamento della linea Q1 -> P
        % P - Q1 = (2.5, 3.5). Scaliamo per fare il vettore: (1, 1.4)
        \coordinate (E1_tip) at ($(P) + (1, 1.4)$); 
        
        % E2: Prolungamento della linea Q2 -> P
        % P - Q2 = (-2.5, 3). Scaliamo: (-1.25, 1.5)
        \coordinate (E2_tip) at ($(P) + (-1.25, 1.5)$);

        % --- 4. CALCOLO VETTORE SOMMA (Etot) ---
        % Usiamo la sintassi calc di TikZ: (P) + (Vettore1) + (Vettore2)
        % Vettore1 = E1_tip - P
        % Vettore2 = E2_tip - P
        \coordinate (Etot_tip) at ($(E1_tip) + (E2_tip) - (P)$);

        % --- 5. COSTRUZIONE DEL PARALLELOGRAMMA ---
        \draw[dashed, blue!40] (E1_tip) -- (Etot_tip);
        \draw[dashed, blue!40] (E2_tip) -- (Etot_tip);

        % --- 6. DISEGNO DEI VETTORI ---
        
        % Vettori Componenti (Blu)
        \draw[->, thick, blue] (P) -- (E1_tip) node[right] {$\vec{F}_1$};
        \draw[->, thick, blue] (P) -- (E2_tip) node[left] {$\vec{F}_2$};
        
        % Vettore Totale (Viola, più spesso)
        \draw[->, very thick, violet] (P) -- (Etot_tip) node[above] {$\vec{F}_{tot}$};

        % --- 7. CARICHE E PUNTO P ---
        
        % Punto P
        \fill (P) circle (2pt) node[below=0.3cm] {$P$};
        
        % Carica q1 (Positiva)
        \node[circle, draw=red, fill=red!10, thick, inner sep=3pt] (cq1) at (Q1) {$+q_1$};
        
        % Carica q2 (Positiva)
        \node[circle, draw=red, fill=red!10, thick, inner sep=3pt] (cq2) at (Q2) {$+q_2$};

    \end{tikzpicture}
    \caption{Principio di sovrapposizione per la forza elettrica. La forza totale $\vec{F}_{tot}$ nel punto $P$ è la somma vettoriale delle forze $\vec{E}_1$ e $\vec{E}_2$ generate dalle singole cariche puntiformi.}
    \label{fig:sovrapposizione_vettoriale}
\end{figure}

\subsubsection{Estensione al continuo}
Il Principio di sovrapposizione mantiene validità anche nel caso in cui le cariche non siano discrete ma distribuite con continuità nello spazio. In tal caso, la sommatoria viene sostituita da un integrale.
Data una distribuzione volumetrica di carica descritta dalla densità $\rho(\bm{r})$ confinata in una regione di spazio $\Omega$, la forza agente su una carica puntiforme $q_0$ posta in $\bm{r}$ è:
\begin{equation}
    \bm{F}(\bm{r}) = \frac{q_0}{4\pi\varepsilon_0} \int_\Omega \frac{\rho(\bm{r}')}{|\bm{r} - \bm{r}'|^2} \bm{u_r}  \, dV,
\end{equation}
dove $\bm{r'}$ rappresenta la posizione dell'$i$-esimo infinitesimo di carica che induce una forza in $\bm{r}'$. 

\section{Il campo elettrostatico}
Una carica elettrica non interagisce direttamente con un'altra carica remota, bensì altera le proprietà fisiche dello spazio circostante generando un \textit{campo elettrostatico}. Una seconda carica immersa in tale spazio risentirà di una forza locale dovuta all'interazione con il campo in quel punto.
Si definisce vettore campo elettrostatico $\bm{E}$ in un punto $P$ il rapporto limite tra la forza elettrica esercitata su una \textit{carica di prova} $q_p$ (posta in $P$) e la carica stessa. L'operazione di limite è necessaria per garantire che $q_0$ sia sufficientemente piccola da non perturbare la distribuzione delle cariche sorgente:
\begin{equation}
    \boxed{\bm{E}(\bm{r}) = \lim_{q_p \to 0} \frac{\bm{F}}{q_p}} \quad [\si{\newton\per\coulomb}] = [\si{\volt\per\meter}].
\end{equation}
Nel caso particolare di una singola carica sorgente puntiforme $q$, applicando la legge di Coulomb, l'espressione del campo diviene:
\begin{equation}
    \bm{E}(\bm{r}) = \frac{1}{4 \pi \varepsilon_0}\frac{q}{r^2}\hat{\bm{u}}_r.
\end{equation}
Si osserva che il campo $\bm{E}$ è una funzione vettoriale di punto: esso eredita la natura vettoriale della forza ma, a differenza di quest'ultima, è una proprietà intrinseca dello spazio e delle sorgenti, indipendente dall'esistenza o dal valore della carica di prova.

\subsection{Principio di sovrapposizione degli effetti}
Il Principio di sovrapposizione è estendibile anche al caso del campo elettrostatico. Si consideri una distribuzione di $N$ cariche sorgenti $q_1, \dots, q_N$. Ponendo una {carica di prova} $q_p$ in un punto $P$ dello spazio, la forza totale agente su di essa è data dalla somma vettoriale delle singole forze esercitate da ciascuna carica $q_j$ (come visto nella \eqref{eq:sovrapposizione_discreta}).
Dalla definizione di campo elettrico segue:
\begin{equation}
    \bm{E}_{tot}(P) = \lim_{q_p \to 0} \frac{\bm{F}_{tot}}{q_p} = \lim_{q_p \to 0} \frac{1}{q_p} \left( \sum_{j=1}^{N} \bm{F}_{j \to p} \right).
\end{equation}
Data la linearità delle operazioni di somma e limite, gli operatori commutano:
\begin{equation}
    \bm{E}_{tot}(P) = \sum_{j=1}^{N} \left( \lim_{q_p \to 0} \frac{\bm{F}_{j \to p}}{q_p} \right).
\end{equation}
Riconoscendo nel termine tra parentesi la definizione del campo elettrico generato dalla singola carica $j$-esima, si ottiene il principio di sovrapposizione per il campo elettrostatico:
\begin{equation}
    \bm{E}_{tot} = \sum_{j=1}^{N} \bm{E}_j.
    \label{eq:sovrapposizione_campo}
\end{equation}
In altre parole: \textit{il campo elettrico totale in un punto è la somma vettoriale dei campi prodotti singolarmente dalle cariche sorgenti in quel punto.}


\section{Il potenziale elettrostatico}
Come discusso in precedenza, essendo il campo elettrico e la forza elettrica legate da una costante moltiplicativa, $\bm{E}$ deve ereditare tutte le caratteristiche vettoriali che contraddistinguono la forza, fra cui la \emph{conservatività}. Applicando allora la definizione di campo elettrico alla condizione di conservatività trovata per la forza elettrostatica \eqref{eq:conservatività}, sfruttando una carica di prova $q_p$ e la linearità dell'operazione di limite, si ottiene:
\begin{equation}
\bm{E} = \lim_{q_p \rightarrow 0} \frac{\bm{F}}{q_p} = -\lim_{q_p \rightarrow 0} \frac{1}{q_p} \nabla U.
\end{equation}
L'operatore $\nabla$ opera rispetto alla variabile posizione della funzione $U$, il limite rispetto alla carica di prova: i due operatori commutano. Si ottiene così la relazione costituiva del campo elettrico in funzione del suo potenziale:
\begin{equation} 
\boxed{\bm{E} = -\nabla (\lim_{q_p \rightarrow 0}  \frac{U}{q_p}) = -\nabla V}
\label{eq:potenziale}
\end{equation}
Il campo scalare $V$ è detto \textit{potenziale elettrostatico}. 
La ricerca di un'espressione per la funzione potenziale è possibile, ancora una volta, a meno di un'operazione di integrazione: {il potenziale elettrostatico}, così come l'energia potenziale, {è definito a meno di una costante} che può essere scelta arbitrariamente. 

\subsubsection{Il legame fra potenziale elettrostatico e energia potenziale}
Così come il campo elettrostatico rappresenta la forza normalizzata rispetto alla carica di prova ($\bm{E} = \bm{F}/q_p$), è possibile definire una grandezza analoga per l'energia. Si introduce allora il precedentemente discusso \textit{potenziale elettrostatico} $V$ come l'energia potenziale per unità di carica.
Dalla \eqref{eq:potenziale}, detta $q_p$ una carica di prova positiva, si definisce:
\begin{equation}
    V = \lim_{q_p \rightarrow 0} \frac{U}{q_p} \quad [\si{\volt}].
\end{equation}
In termini operativi, l'energia potenziale acquisita da una particella di carica $q$ posta in un punto a potenziale $V$ è data semplicemente da:
\begin{equation}
    U = qV.
    \label{eq:relazione_energia_potenziale}
\end{equation}
Sfruttando il Teorema dell'energia cinetica e la conservatività del campo \eqref{eq:teorema_lavoro_energia}, il lavoro compiuto dalla forza elettrostatica su una carica $q$ che si sposta tra due punti $A$ e $B$ si scrive:
\begin{equation}
    W_{A \to B} = -\Delta U = -[U(B)-U(A)] = -q[V(B)-V(A)] = -q\Delta V.
\end{equation}
Questa relazione evidenzia come il lavoro dipenda linearmente dalla differenza di potenziale (o \textit{tensione}) ai capi della traiettoria.


\begin{approfondimento}[title={Unità di energia: L'Elettronvolt}]
\noindent
Dalla relazione \eqref{eq:relazione_energia_potenziale} emerge che il Joule, unità SI per l'energia, risulta spesso sproporzionato per la fisica microscopica.
    Si definisce l'\textit{elettronvolt} (\si{\electronvolt}) come l'energia cinetica acquisita da un elettrone libero ($e \approx \SI{1.602e-19}{\coulomb}$) accelerato attraverso una differenza di potenziale di \SI{1}{\volt}.
    \begin{equation}
        \SI{1}{\electronvolt} = (\SI{1.602e-19}{\coulomb})(\SI{1}{\volt}) = \SI{1.602e-19}{\joule}.
    \end{equation}
\end{approfondimento}


\subsection{Principio di sovrapposizione}
Il Principio di sovrapposizione per il campo elettrico può essere esteso al potenziale elettrostatico $V$, definito a partire dall'applicazione dell'operatore $\nabla$, lineare in quanto derivatore. Dal momento che il campo elettrico risultante in un punto dello spazio è dato dalla combinazione lineare dei campi prodotti dalle varie sorgenti disposte nella regione (sovrapposizione), l'applicazione lineare $\nabla$ non modifica la proprietà di combinazione. Se ne deriva dunque che:
\begin{equation}
\begin{aligned}
   -\nabla V_{tot} = \bm{E}_{tot} =& \, \sum_{j=1}^{N} \bm{E}_j = -\sum_{j=1}^{N}\nabla V_{j} = -\nabla \sum_{j=1}^{N} V_j \\
    V_{tot} =& \sum_{j=1}^{N} V_j.
    \end{aligned}
    \label{eq:sovrapposizione_potenziale}
\end{equation}

\begin{approfondimento}[title={Superfici (o Linee) Equipotenziali}]
\noindent
Dalla definizione di potenziale come campo scalare, risulta naturale introdurre un ente geometrico utile alla visualizzazione dell'andamento energetico del campo: le \textit{superfici equipotenziali}.
Si definisce superficie equipotenziale il luogo geometrico dei punti dello spazio in cui il potenziale elettrostatico assume lo stesso valore costante:
\begin{equation}
    V(x, y, z) = \text{cost}.
\end{equation}
Nel caso bidimensionale (o nelle rappresentazioni su piano), si parla di \textit{linee equipotenziali}. Un esempio comune è quello della cartografia topografica: le linee equipotenziali sono analoghe alle curve di livello (isoipse) che uniscono i punti alla stessa quota altimetrica. 

% --- INIZIO FIGURA NON FLOTTANTE ---
\begin{center}
    \begin{tikzpicture}[scale=1.2, >=Stealth]
        % Definizione stile freccia a metà linea
        \tikzset{midarrow/.style={decoration={markings, mark=at position 0.6 with {\arrow{Stealth}}}, postaction={decorate}}}
        
        % 1. Linee di campo (Rosse)
        \foreach \angle in {0, 30, ..., 330} { 
            \draw[midarrow, red!80!black, thin] (0,0) -- (\angle:2.2cm); 
        }
        
        % 2. Equipotenziali (Blu)
        \foreach \r in {0.8, 1.3, 1.8} { 
            \draw[densely dashed, blue!80!black, thick] (0,0) circle (\r cm); 
        }
        
        % 3. Carica centrale
        \node[circle, draw=red, fill=red!10, inner sep=2pt, minimum size=0.6cm] at (0,0) {\large $+$};
        \node[below=0.4cm, red!80!black] at (0,0) {\scriptsize Sorgente};
        
        % 4. Dimostrazione ortogonalità (Focus geometrico)
        % Coordinate dinamiche
        \def\ang{45} % Angolo di analisi
        \def\rad{1.8} % Raggio equipotenziale
        \coordinate (P) at (\ang:\rad);
        
        % Vettori
        \draw[->, very thick, red] (P) -- ++(\ang:0.8) node[right] {$\bm{E}$};
        \draw[->, thick, blue] (P) -- ++(\ang+90:0.8) node[above left] {$d\bm{l}$};
        
        % Angolo Retto (Metodo pulito con libreria 'calc')
        \draw[thin, black] ($(P)!0.15cm!(\ang:3cm)$) -- 
                           ($(P)!0.15cm!(\ang:3cm)!0.15cm!90:(\ang:3cm)$) -- 
                           ($(P)!0.15cm!(\ang+90:3cm)$);

        % Punto P
        \node[circle, fill, inner sep=1pt] at (P) {};
        \node[below right] at (P) {$P$};
        
    \end{tikzpicture}
    % Didascalia specifica per contenuti NON flottanti
    \captionof{figure}{Rappresentazione di campo generato da carica puntiforme: linee di campo (rosse) e superfici equipotenziali (blu). Si nota l'ortogonalità locale in $P$.}
    \label{fig:equipotenziali}
\end{center}
% --- FINE FIGURA ---

\bigskip % Spaziatura verticale per separare dal testo successivo

\subsubsection{Proprietà fondamentali}
Le superfici equipotenziali godono di due proprietà fisiche cruciali:

\begin{enumerate}
    \item \textit{lavoro nullo:} per definizione, spostare una carica $q$ lungo una superficie equipotenziale non richiede lavoro. Infatti, se $V_A = V_B$, allora:
    \begin{equation}
        W_{A \to B} = -q \Delta V = -q(V_B - V_A) = 0;
    \end{equation}
    
    \item \textit{ortogonalità alle linee di campo:} in ogni punto, il vettore campo elettrico $\bm{E}$ è perpendicolare alla superficie equipotenziale passante per quel punto. Considerato infatti uno spostamento infinitesimo $d\bm{l}$ tangente alla superficie ($dV=0$), il lavoro deve essere nullo:
    \begin{equation}
        \delta W = \bm{F} \cdot d\bm{l} = q \bm{E} \cdot d\bm{l} = 0.
    \end{equation}
    Affinché il prodotto scalare sia nullo con $\bm{E}, d\bm{l} \neq 0$, i vettori devono essere ortogonali:
    \begin{equation}
        \bm{E} \perp d\bm{l} \quad \forall \, d\bm{l} \in \text{superficie}.
    \end{equation}
\end{enumerate}
\end{approfondimento}

\section{La legge di Gauss}
In questa sezione viene presentata la \textit{Legge di Gauss} nelle sue due formulazioni, integrale e differenziale. L'utilità della stessa, anche nota come \textit{prima equazione di Maxwell}, si ritrova nel calcolo del flusso del campo elettrico attraverso una generica superficie chiusa.

\subsection{Forma integrale}
Sia $\Sigma$ una superficie chiusa che racchiude un volume $\Omega$\footnote{La superficie $\Sigma$ si configura pertanto come il bordo della regione $\Omega$: $\Sigma = \partial \Omega$.}, all'interno del quale sono distribuite $N$ cariche puntiformi $q_1, \dots, q_N$.
In virtù del {Principio di sovrapposizione}, il campo elettrico totale in ogni punto dello spazio è la somma vettoriale dei campi generati dalle singole cariche. Di conseguenza, il flusso totale attraverso $\Sigma$ è pari alla somma dei flussi dei singoli campi.
È dunque sufficiente analizzare, in prima istanza, il contributo di una singola carica puntiforme $q$ posta internamente alla superficie.
Il flusso elementare del campo elettrico attraverso un elemento infinitesimo di superficie $d\bm{S} = \hat{\bm{u}}_n \, dS$ vale:
\begin{equation}
    d\Phi = \bm{E} \cdot d\bm{S} = \left( \frac{1}{4\pi \varepsilon_0} \frac{q}{r^2} \hat{\bm{u}}_r \right) \cdot (\hat{\bm{u}}_n \, dS).
\end{equation}
Sviluppando il prodotto scalare, si evidenzia la dipendenza geometrica dall'orientamento relativo tra il campo e la superficie:
\begin{equation}
    d\Phi = \frac{q}{4\pi \varepsilon_0 r^2} (\hat{\bm{u}}_r \cdot \hat{\bm{u}}_n) \, dS = \frac{q}{4\pi \varepsilon_0 r^2} \cos\theta \, dS,
    \label{eq:flusso_diff}
\end{equation}
dove $\theta$ è l'angolo compreso tra la direzione radiale uscente dalla carica ($\hat{\bm{u}}_r$) e la normale alla superficie ($\hat{\bm{u}}_n$).
Il termine $dS_{\perp} = dS \cos\theta$ rappresenta la proiezione dell'area $dS$ sulla superficie sferica di raggio $r$ centrata nella carica. Il rapporto tra tale proiezione e il quadrato della distanza definisce l'\textit{angolo solido} elementare $d\Omega$ sotto cui l'area $dS$ è vista dalla carica:
\begin{equation}
   d\Omega{G} = \frac{dS \cos\theta}{r^2} = \frac{\hat{\bm{u}}_r \cdot d\bm{S}}{r^2}.
\end{equation}
Sostituendo questa definizione nella \eqref{eq:flusso_diff} ed estendendo l'integrale all'intera superficie chiusa $\Sigma$, si ottiene il flusso totale generato dalla singola carica interna. Poiché la superficie avvolge completamente la carica, l'integrale dell'angolo solido copre l'intera sfera unitaria ($4\pi \, \si{\steradian}$):
\begin{equation}
    \Phi_{\Sigma}(\bm{E}_q) = \frac{q}{4\pi \varepsilon_0} \oint_{\Sigma} d\Omega= \frac{q}{4\pi \varepsilon_0} (4\pi) = \frac{q}{\varepsilon_0}.
\end{equation}

\subsubsection{Estensione a sistemi di cariche}
Si consideri ora il caso generale con $N$ cariche, di cui solo una parte è contenuta nel volume racchiuso da $\Sigma$.
È fondamentale osservare come le cariche \textit{esterne} alla superficie non contribuiscano al flusso netto\footnote{L'intuizione può a questo punto suggerire che il flusso del campo elettrico sia in qualche modo legato alle \emph{sorgenti} che lo generano. Il flusso attraverso una superficie generica chiusa $\Sigma$ che racchiude punti da cui il campo "sgorga" (generiche cariche puntiformi positive) o "muore" (generiche cariche puntiformi negative) deve essere diverso da zero (a meno che la carica netta nella regione entro $\Sigma$ sia nulla), altrimenti identicamente nullo. Questa idea viene formalizzata nel passaggio alla forma differenziale (o \emph{locale}) della Legge di Gauss.\label{pensiero}}. Per una carica esterna, infatti, il cono di angolo solido elementare interseca la superficie chiusa sempre in un numero pari di punti (entrata e uscita). I contributi al flusso ($\bm{E}\cdot\hat{\bm{u}}_n$) hanno segni opposti in ingresso ($\cos\theta < 0$, angolo ottuso) e in uscita ($\cos\theta > 0$, angolo acuto), annullandosi globalmente nell'integrazione.
Pertanto, sommando i contributi delle sole cariche interne ($q_i \in V$), si perviene alla formulazione generale della \textit{Legge di Gauss}:
\begin{equation}
    \boxed{\Phi_{\Sigma}(\bm{E}) = \oint_{\Sigma} \bm{E} \cdot d\bm{S} = \frac{1}{\varepsilon_0} \sum_{i} q_{i, \text{int}} = \frac{Q_{\text{int}}}{\varepsilon_0}}.
    \label{eq:teorema_gauss}
\end{equation}
Il Teorema di Gauss stabilisce che il flusso del campo elettrostatico attraverso una superficie chiusa dipende \textit{esclusivamente} dalla somma algebrica delle cariche interne, risultando indipendente dalla loro distribuzione spaziale e dalla forma geometrica della superficie stessa.


\begin{figure}[H]
\centering
\begin{tikzpicture}[scale=1.8, >=Stealth]
    % --- 1. DEFINIZIONI ---
    \coordinate (O) at (0,0); % Posizione della carica q
    \def\R{3.5} % Raggio della sfera di riferimento
    \def\angCenter{35} % Angolo centrale del cono
    \def\angWidth{12}  % semi-ampiezza del cono

    % --- 2. CONO DI ANGOLO SOLIDO ---
    % Riempimento del cono
    \fill[red!5, opacity=0.8] (O) -- (\angCenter-\angWidth:\R) arc (\angCenter-\angWidth:\angCenter+\angWidth:\R) -- cycle;
    % Linee delimitatrici del cono (tratteggiate)
    \draw[red!40, dashed, thin] (O) -- (\angCenter-\angWidth:4);
    \draw[red!40, dashed, thin] (O) -- (\angCenter+\angWidth:4);
    % Etichetta Angolo Solido
    \node[red!60!black] at (\angCenter:1.8) {$d\Omega$};

    % --- 3. SUPERFICIE SFERICA DI RIFERIMENTO (Proiezione) ---
    % Arco di cerchio che rappresenta la calotta sferica perpendicolare al raggio
    \draw[red!80!black, thick] (\angCenter-\angWidth:\R) arc (\angCenter-\angWidth:\angCenter+\angWidth:\R);
    % Etichetta della proiezione
    \draw[<-] (\angCenter+5:\R) -- ++(\angCenter+5:0.6) node[above, red!80!black, scale=0.9] {$dS_\perp = dS \cos\theta$};
    % Linea di raggio di riferimento
    \draw[gray, dashed] (O) -- (\angCenter:\R) node[midway, below] {$r$};


    % --- 4. SUPERFICIE ARBITRARIA dS (Inclinata) ---
    \coordinate (P) at (\angCenter:\R); % Punto di applicazione dei vettori
    % Disegno una "patch" inclinata. Usiamo una curva di Bezier per farla sembrare 3D
    \filldraw[fill=blue!20, draw=blue!50!black, thick, opacity=0.7, rotate around={-25:(P)}] 
        ($(P)+(-0.3,-0.5)$) to[out=80, in=260] ($(P)+(-0.3,0.5)$) -- 
        ($(P)+(0.3,0.5)$) to[out=280, in=100] ($(P)+(0.3,-0.5)$) -- cycle;
    \node[blue!50!black, right] at ($(P)+(0.4,0)$) {$d\bm{S}$};


    % --- 5. VETTORI E ANGOLI (Nel punto P) ---
    % Vettore unitario radiale u_r (rosso, perpendicolare alla sfera)
    \draw[->, very thick, red] (P) -- ++(\angCenter:1) coordinate (Ur) node[below] {$\hat{\bm{u}}_r$};
    
    % Vettore normale n (blu, perpendicolare alla superficie arbitraria)
    % Deve essere inclinato di -25 gradi rispetto a u_r
    \draw[->, very thick, blue!80!black] (P) -- ++(\angCenter-25:1) coordinate (N) node[right] {$\hat{\bm{n}}$};

    % Angolo Theta
    \pic [draw, "$\theta$", angle radius=0.6cm, angle eccentricity=1.3, thick] {angle = N--P--Ur};

    % --- 6. CARICA SORGENTE ---
    \shade[ball color=red!30] (O) circle (0.2cm);
    \node at (O) {\textbf{+}};
    \node[left=0.3cm] at (O) {$q$};

\end{tikzpicture}
\caption{Interpretazione geometrica del flusso. Il cono (rosso chiaro) definisce l'angolo solido $d\Omega$. La superficie arbitraria $d\bm{S}$ (blu) è inclinata. Il flusso dipende dalla proiezione $dS_\perp$ (arco rosso spesso) di $d\bm{S}$ sulla sfera di riferimento. Tale proiezione è legata all'angolo $\theta$ tra la normale $\hat{\bm{n}}$ e il raggio $\hat{\bm{u}}_r$.}
\label{fig:flusso_angolo_solido_v2}
\end{figure}

\subsection{Forma differenziale (locale)}
Per descrivere il comportamento del campo elettrico punto per punto nello spazio, svincolandosi dalla scelta di una specifica superficie di integrazione, è necessario passare alla \textit{forma locale} (o differenziale) della Legge di Gauss.
Si consideri un volume generico $\Omega$ delimitato da una superficie chiusa $\Sigma = \partial \Omega$.
Applicando il \textit{Teorema della Divergenza} (o di Gauss-Green), il flusso di un campo vettoriale attraverso una superficie chiusa è pari all'integrale della divergenza del campo esteso al volume racchiuso\footnote{Si assume che il campo vettoriale sia di classe $C^1$ (continuo con derivate prime continue) e che il dominio $\Omega$ sia regolare a tratti.}:
\begin{equation}
    \oint_{\partial \Omega} \bm{E} \cdot d\bm{S} = \int_\Omega (\nabla \cdot \bm{E}) \, dV.
    \label{eq:teorema_divergenza}
\end{equation}
La carica interna $Q_{\text{int}}$ è esprimibile come l'integrale della densità volumetrica di carica $\rho(\bm{r})$ sul volume $\Omega$:
\begin{equation}
    Q_{\text{int}} = \int_\Omega \rho(\bm{r}) \, dV.
\end{equation}
Sostituendo tali espressioni nella forma integrale ($\Phi = Q_{\text{int}}/\varepsilon_0$) e sfruttando la linearità dell'integrale, si ottiene:
\begin{equation}
    \int_\Omega (\nabla \cdot \bm{E}) \, dV = \frac{1}{\varepsilon_0} \int_\Omega \rho \, dV \implies \int_\Omega \left( \nabla \cdot \bm{E} - \frac{\rho}{\varepsilon_0} \right) dV = 0.
\end{equation}
Affinché tale integrale sia nullo per \textit{qualsiasi} volume arbitrario $\Omega$, è necessario\footnote{Come risultato del \textit{Lemma di localizzazione}.} che l'integrando sia identicamente nullo in ogni punto dello spazio. Si giunge così alla forma differenziale della Legge di Gauss:
\begin{equation}
    \boxed{\nabla \cdot \bm{E} = \frac{\rho}{\varepsilon_0}}.
    \label{eq:gauss_differenziale}
\end{equation}

\subsubsection{Significato fisico della Divergenza}
In questo contesto l'intuizione di cui si discuteva in precedenza (in (\ref{pensiero})) risulta finalmente formalizzata. L'operatore divergenza, infatti, quantifica la "fecondità" di un punto dello spazio in termini di campo, ovvero la sua capacità di generare o assorbire linee di campo. L'equazione \eqref{eq:gauss_differenziale} stabilisce una corrispondenza puntuale biunivoca tra la presenza di carica e la divergenza del campo:

\begin{description}
    \item[$\nabla \cdot \bm{E} > 0$] (\emph{sorgente}): in presenza di densità di carica positiva ($\rho > 0$), il punto agisce come una sorgente da cui "sgorgano" le linee di forza;
    \item[$\nabla \cdot \bm{E} < 0$] (\emph{pozzo}): In presenza di densità di carica negativa ($\rho < 0$), il punto agisce come un pozzo in cui le linee di forza convergono e terminano;
    \item[$\nabla \cdot \bm{E} = 0$]: nel vuoto ($\rho = 0$), il campo non ha né sorgenti né pozzi: le linee di campo che entrano in un volume infinitesimo devono necessariamente uscirne.
\end{description}

\begin{approfondimento}[title={Le Linee di Campo: Realtà Fisica o Artificio Matematico?}]
\noindent
    Introdotte da Michael Faraday (che le chiamava "linee di forza"), le linee di campo costituiscono un potente strumento \textit{euristico} per visualizzare la struttura vettoriale dello spazio. Tuttavia, è fondamentale non confondere la mappa con il territorio: le linee sono una rappresentazione discreta di una realtà fisica continua.

    \subsubsection{Definizione geometrica}
    Una linea di campo è una curva orientata costruita in modo tale che, in ogni suo punto, la retta tangente sia parallela al vettore campo elettrico $\bm{E}$ in quel punto. Matematicamente, se $d\bm{l}$ è uno spostamento infinitesimo lungo la linea:
    \begin{equation}
        d\bm{l} \parallel \bm{E}(\bm{r}) \implies d\bm{l} \times \bm{E} = 0.
    \end{equation}
    
    \begin{center}
        \begin{tikzpicture}[scale=1.2, >=Stealth]
            % Disegno delle linee di campo curve (Semplificato, senza 'name path')
            \draw[thick, blue!80!black] (0,0) .. controls (1,1) and (3,1) .. (4,0);
            \draw[thick, blue!80!black] (0,-1) .. controls (1,0) and (3,0) .. (4,-1);
            
            \node[blue!80!black, right] at (4,-0.2) {$\gamma$};
            
            % Coordinate manuali per i punti tangenti
            \coordinate (P1) at (0.5, 0.35); 
            \coordinate (P2) at (2, 0.75);
            \coordinate (P3) at (3.5, 0.35);
    
            % Disegno vettori tangenti (rossi)
            \draw[->, thick, red] (P1) -- ++(30:0.8) node[above] {$\bm{E}_1$};
            \draw[->, thick, red] (P2) -- ++(0:0.8) node[above] {$\bm{E}_2$};
            \draw[->, thick, red] (P3) -- ++(-30:0.8) node[above] {$\bm{E}_3$};
    
            % Punti neri
            \fill (P1) circle (1.5pt);
            \fill (P2) circle (1.5pt);
            \fill (P3) circle (1.5pt);
            
            % Annotazione densità
            \draw[<->, thin, gray] (2, 0.75) -- (2, -0.25);
            \node[gray, right, scale=0.8] at (0.3, 0.05) {Densità $\propto |\bm{E}|$};
            
        \end{tikzpicture}
        % ATTENZIONE: Assicurati di avere \usepackage{caption} nel preambolo
        \captionof{figure}{Relazione tra linea di campo (blu) e vettori campo elettrico (rossi). Il vettore $\bm{E}$ è sempre tangente alla linea.}
    \end{center}

    \subsubsection{Regole di tracciamento e interpretazione}
    Perché questa rappresentazione sia fisicamente coerente, occorre rispettare tre principi:
    \begin{enumerate}
        \item \textit{densità e Intensità}: il campo elettrico è definito ovunque, anche "tra" le linee disegnate. La densità spaziale delle linee (numero di linee che attraversano un'area unitaria perpendicolare) è proporzionale al modulo del campo $|\bm{E}|$. Dove le linee si addensano, il campo è intenso; dove si diradano, è debole;
        \item \textit{sorgenti e pozzi}: le linee nascono dalle cariche positive (o all'infinito) e muoiono nelle cariche negative (o all'infinito). Il numero di linee generate/assorbite è proporzionale alla quantità di carica $Q$;
        \item \textit{non intersecabilità}: due linee di campo non possono mai intersecarsi. Se lo facessero, nel punto di intersezione il campo $\bm{E}$ dovrebbe avere due direzioni distinte contemporaneamente, violando l'unicità della funzione vettoriale (il campo in un punto è univocamente determinato dalla somma delle forze).
    \end{enumerate}
\end{approfondimento}

\section{La circuitazione del campo elettrico}
In questa sezione si analizza la proprietà di circuitazione del campo elettrostatico, introducendo la forma statica della \textit{terza equazione di Maxwell}. 
Si segnala che tale relazione costituisce un caso particolare della più generale \textit{Legge di Faraday-Neumann-Lenz}, valida in regime variabile, che verrà trattata successivamente. L'obiettivo è valutare la circuitazione\footnote{Fisicamente, la circuitazione corrisponde al lavoro compiuto dal campo sulla carica unitaria lungo un percorso chiuso. Matematicamente, è l'integrale curvilineo di seconda specie del campo vettoriale lungo una curva chiusa $\gamma$.} $\Gamma_\gamma(\bm{E})$ lungo una generica curva \emph{chiusa}.

\subsection{Forma integrale}
Si è precedentemente discusso come il campo elettrostatico sia intrinsecamente conservativo. Si intende ora formalizzare questa proprietà per via analitica.
Sia $\gamma$ una curva chiusa e regolare\footnote{Perché l'operazione di circuitazione sia fattibile, è necessario richiedere la \emph{regolarità} della curva lungo cui si valuta (ovvero il fatto che $\gamma' \neq \bm{0}$).} nello spazio. La circuitazione del campo $\bm{E}$ lungo $\gamma$ è definita come:
\begin{equation}
    \Gamma_{\gamma}(\bm{E}) = \oint_{\gamma} \bm{E} \cdot d\bm{l}.
\end{equation}
Ricordando la relazione locale tra campo e potenziale ($\bm{E} = -\nabla V$), l'integrando può essere riscritto sfruttando la definizione di differenziale esatto di una funzione scalare ($dV = \nabla V \cdot d\bm{l}$):
\begin{equation}
    \oint_{\gamma} \bm{E} \cdot d\bm{l} = -\oint_{\gamma} \nabla V \cdot d\bm{l} = -\oint_{\gamma} dV.
\end{equation}
Per il \textit{Teorema fondamentale del calcolo per integrali di linea}, l'integrale del differenziale esatto dipende esclusivamente dagli estremi del percorso. Poiché la curva $\gamma$ è chiusa per ipotesi, il punto iniziale $A$ coincide con il punto finale $B$ ($A \equiv B$). Ne consegue che la differenza di potenziale è nulla:
\begin{equation}
    -\oint_{\gamma} dV = -[V(A) - V(A)] = 0.
\end{equation}
Si ottiene dunque la seguente relazione\footnote{Anche conosciuta come \textit{Legge di Kirchhoff delle tensioni} in ambito circuitale.}:
\begin{equation}
    \boxed{\oint_{\gamma} \bm{E} \cdot d\bm{l} = 0}.
    \label{eq:circuitazione_nulla}
\end{equation}

% --- INSERIMENTO FIGURA ---
\begin{figure}[H]
    \centering
    \begin{tikzpicture}[scale=1.2, >=Stealth]
        % Definizione della curva chiusa (blob)
        \draw[thick, blue!80!black, decoration={markings, mark=at position 0.15 with {\arrow{>}}}, postaction={decorate}] 
            plot [smooth cycle, tension=0.7] coordinates {(0,0) (2,0.5) (2.5,2) (1,2.5) (-0.5,1.5)};
        
        \node[blue!80!black] at (-0.6, 1.5) {$\gamma$};
        
        % Punto P sulla curva
        \coordinate (P) at (2.5, 2);
        \fill (P) circle (1.5pt);
        \node[below right] at (P) {$P$};
        
        % Vettore tangente dl
        \draw[->, thick, blue] (P) -- ++(110:0.8) node[left] {$d\bm{l}$};
        
        % Vettore Campo E (generico, non parallelo)
        \draw[->, thick, red] (P) -- ++(45:1) node[right] {$\bm{E}$};
        
        % Angolo
        \draw[thin, gray] (P) ++(110:0.4) arc (110:45:0.4);
        
    \end{tikzpicture}
    \caption{Rappresentazione geometrica della circuitazione. L'integrale somma i contributi scalari $\bm{E} \cdot d\bm{l}$ lungo il percorso chiuso $\gamma$. In elettrostatica, la somma è sempre nulla.}
    \label{fig:circuitazione}
\end{figure} 

\subsection{Forma differenziale}
L'analisi locale delle proprietà rotazionali del campo richiede il passaggio alla forma differenziale.
Lo strumento matematico necessario è il \textit{Teorema di Stokes} (o del Rotore), il quale stabilisce che la circuitazione di un campo vettoriale lungo una curva chiusa $\gamma$ equivale al flusso del rotore del campo calcolato attraverso una qualsiasi superficie aperta $\Sigma$ che abbia $\gamma$ come bordo (ovvero $\gamma = \partial \Sigma$).
In formule:\footnote{Si presti attenzione alla notazione: l'integrale a membro destro non presenta il simbolo di chiusura ($\oint$) poiché il flusso è calcolato attraverso una superficie \textit{aperta}, sebbene il percorso di circuitazione sia chiuso.}
\begin{equation}
    \oint_{\gamma = \partial \Sigma} \bm{E} \cdot d\bm{l} = \int_{\Sigma} (\nabla \times \bm{E}) \cdot d\bm{S}.
\end{equation}
Ricordando che per il campo elettrostatico la circuitazione è sempre nulla (per la \eqref{eq:circuitazione_nulla}), si ha:
\begin{equation}
    \int_{\Sigma} (\nabla \times \bm{E}) \cdot d\bm{S} = 0.
\end{equation}
Affinché questa relazione sia verificata per \textit{ogni} possibile superficie $\Sigma$ delimitata da $\gamma$, l'integrando deve essere identicamente nullo in ogni punto dello spazio:
\begin{equation}
    \boxed{\nabla \times \bm{E} = 0}.
    \label{eq:circuitazione_diff}
\end{equation}


\subsubsection{Significato fisico}
Il rotore di un campo vettoriale misura la tendenza del campo a "ruotare" attorno ad un punto. La \eqref{eq:circuitazione_diff} permette di stabilire che il campo elettrostatico è in generale un campo \textit{irrotazionale}. Il risultato appena ottenuto consente inoltre di caratterizzare in maniera differente i campi conservativi come quello elettrostatico\footnote{Si noti come, quando si discute della \emph{conservatività}, non si faccia riferimento al campo elettrico generale ma al campo elettro\emph{statico}. Un campo elettrico tempo variabile, infatti, risulta intrinsecamente non conservativo, come verrà successivamente discusso.} adducendo agli stessi la proprietà dell'\textit{irrotazionalità}\footnote{Si segnala che \emph{conservatività} $\implies$ \emph{irrotazionalità} ma, in generale, non vale il viceversa. Perché un campo irrotazionale sia sicuramente conservativo, infatti, è necessario che sia definito su un dominio \emph{semplicemente connesso} ("senza buchi": un esempio di dominio non semplicemente connesso è $\mathbb{R}^2 \setminus \{0,0\}$), in accordo al \emph{Lemma di Poincaré}. }.   

\section{L'equazione di Poisson}
Conseguenza della conservatività del campo elettrostatico e della Legge di Gauss è \textit{l'equazione di Poisson}. Valgono la \eqref{eq:potenziale} e la \eqref{eq:gauss_differenziale}, da cui:
\begin{equation}
\begin{aligned}
\frac{\rho}{\varepsilon_0} = \nabla \cdot \bm{E} &= \nabla \cdot (-\nabla V) = -\nabla^2 V \\
\Aboxed{\nabla^2 V &= -\frac{\rho}{\varepsilon_0}}.
\end{aligned}
\label{eq:poisson}
\end{equation}
Nel caso particolare in cui ci si trovi nel vuoto o in una regione priva di cariche ($\rho = 0$), l'equazione si riduce all'\textit{Equazione di Laplace}:
\begin{equation}
    \nabla^2 V = 0.
    \label{eq:laplace}
\end{equation}

\subsection{Interpretazione geometrica: Curvatura e Sorgenti}
Per cogliere il significato profondo dell'equazione, occorre analizzare l'azione dell'operatore Laplaciano.
In una dimensione ($1D$), il Laplaciano coincide con la derivata seconda $\frac{d^2 V}{dx^2}$, la quale misura la \textit{concavità} locale della funzione:
\begin{itemize}
    \item $\frac{d^2 V}{dx^2} > 0$ (concavità $\cup$): la funzione presenta un minimo locale (struttura "a scodella");
    \item $\frac{d^2 V}{dx^2} < 0$ (concavità $\cap$): la funzione presenta un massimo locale (struttura "a collina").
\end{itemize}
Estendendo il concetto a tre dimensioni, l'equazione di Poisson \eqref{eq:poisson} fornisce una potente chiave di lettura geometrica:
\begin{quote}
    \textit{La densità di carica $\rho$ in un punto agisce come sorgente della "curvatura'' del potenziale $V$ in quel punto.}
\end{quote}

\subsubsection{Topologia del Potenziale}
Il segno meno nell'equazione di Poisson determina la tipologia di curvatura indotta dalle cariche:
\begin{itemize}
    \item {carica positiva ($\rho > 0$):} implica $\nabla^2 V < 0$ (concavità verso il basso).
    Fisicamente, le cariche positive generano dei \textit{massimi} locali di potenziale ("picchi" o "colline"). Le linee di campo "scendono" giù dalla collina allontanandosi dalla carica;
    
    \item {carica negativa ($\rho < 0$):} implica $\nabla^2 V > 0$ (concavità verso l'alto).
    Fisicamente, le cariche negative generano dei \textit{minimi} locali di potenziale ("buche" o "pozzi"). Le linee di campo "cadono" dentro il pozzo convergendo sulla carica.
\end{itemize}


\begin{figure}[H]
    \centering
    \begin{subfigure}[b]{0.45\textwidth}
        \centering
        \begin{tikzpicture}[scale=0.8]
            \begin{axis}[
                colormap/hot,
                hide axis,
                view={25}{30},
                z buffer=sort,
            ]
            % Collina di potenziale (Carica +)
            \addplot3 [surf, shader=interp, domain=-2:2, y domain=-2:2, samples=40] 
                {2 * exp(-(x^2+y^2)/0.8)};
            \end{axis}
            \node[below, red!80!black] at (3.4,0.5) {\textbf{Carica } $\bm{+}$};
            \node[below, scale=0.8] at (3.4,0.1) {$\nabla^2 V < 0$ (Collina)};
        \end{tikzpicture}
    \end{subfigure}
    \hfill
    \begin{subfigure}[b]{0.45\textwidth}
        \centering
        \begin{tikzpicture}[scale=0.8]
            \begin{axis}[
                colormap/bluered, % Se non hai colormap personalizzate usa 'cool' o 'viridis'
                hide axis,
                view={25}{30},
                z buffer=sort,
            ]
            % Pozzo di potenziale (Carica -)
            \addplot3 [surf, shader=interp, domain=-2:2, y domain=-2:2, samples=40] 
                {-2 * exp(-(x^2+y^2)/0.8)};
            \end{axis}
            \node[below, blue!80!black] at (3.4,0.5) {\textbf{Carica } $\bm{-}$};
            \node[below, scale=0.8] at (3.4,0.1) {$\nabla^2 V > 0$ (Pozzo)};
        \end{tikzpicture}
    \end{subfigure}
    \caption{Rappresentazione topologica del potenziale elettrostatico generato da una carica positiva (sinistra) e negativa (destra). La curvatura della superficie è determinata dal segno del Laplaciano.}
    \label{fig:potential_landscape}
\end{figure}

Si noti come l'operatore di Laplace $\nabla^2$ rappresenti la traccia della \textit{matrice hessiana}: 


\subsubsection{L'analogia della membrana elastica}
Un'analogia fisica estremamente efficace è quella di un telo elastico teso orizzontalmente:
\begin{enumerate}
    \item \textbf{Regime di Laplace ($\rho = 0$):} Se non si applicano pesi sul telo, esso si tende in modo liscio per raccordare i bordi (eventualmente fissati ad altezze diverse), senza creare avvallamenti o picchi locali. In questa configurazione, l'altezza del telo in un punto è esattamente la media delle altezze dei punti circostanti.
    \item \textbf{Regime di Poisson ($\rho \neq 0$):} Se si appoggiano delle sfere pesanti (cariche) sul telo, questo si deforma. Una massa (sorgente) forza il telo a curvarsi localmente, creando un avvallamento che altera la tensione dell'intera superficie.
\end{enumerate}

\newpage
\section{Esercizi}
\subsection{Pendolo elettrostatico}






\chapter{Sorgenti e campi elettrici notevoli}
Il presente capitolo è dedicato alla determinazione analitica del campo elettrico generato da configurazioni di sorgenti di particolare rilevanza fisica. Verrà spesso sfruttato il \textit{Teorema di Gauss} in forma integrale, il concetto di \textit{superficie gaussiana}\footnote{Si definisce superficie gaussiana una superficie chiusa arbitraria nello spazio tridimensionale, utilizzata per il calcolo del flusso del vettore campo elettrico. La scelta opportuna di tale superficie consente di semplificare l'applicazione della Legge di Gauss per ricavare il campo generato da distribuzioni di carica dotate di elevata simmetria. } nonché le particolari proprietà di simmetria dei sistemi analizzati.

\section{Distribuzione piana infinita di carica}
Sia $\sigma$ la densità superficiale di carica distribuita uniformemente su di un piano di estensione infinita.
Per ragioni di simmetria, il campo elettrico generato da tale distribuzione in un punto $P$ deve necessariamente:
\begin{itemize}
    \item giacere sulla direzione perpendicolare al piano. Infatti, ogni infinitesimo di carica $dq$ che non si trova sulla proiezione di $P$ sul piano ha il suo simmetrico rispetto a tale proiezione: le componenti del campo parallele al piano dunque si elidono a vicenda;
       \item essere uscente da entrambe le facce del piano (se $\sigma > 0$) o entrante (se $\sigma < 0$).
\end{itemize}
Per calcolare il modulo del campo, si sceglie come superficie gaussiana un cilindro retto (detto anche \textit{scatoletta gaussiana}) di area di base $A$ e altezza $2z$, disposto in modo tale che le basi siano parallele al piano carico e che questo lo tagli esattamente a metà (si veda figura (\ref{fig:piano_infinito})).
\begin{figure}[htbp]
    \centering
    \begin{tikzpicture}[scale=1.2, >=Stealth]
        % Piano Infinito (rappresentato come romboide)
        \fill[gray!10] (-2,-1) -- (2,-1) -- (3,1) -- (-1,1) -- cycle;
        \draw[gray!50] (-2,-1) -- (2,-1) -- (3,1) -- (-1,1) -- cycle;
        \node at (2.5,0.8) {$\sigma$};
        
        % Cariche sul piano
        \foreach \x in {-0.5, 0, 0.5} {
            \foreach \y in {-0.3, 0, 0.3} {
                \node[red!50, font=\tiny] at (\x + \y*0.5, \y) {$+$};
            }
        }

        % --- CILINDRO GAUSSIANO ---
        % Parte Sotto (Nascosta o tratteggiata)
        \draw[dashed] (0,0) ellipse (0.4 and 0.15); % Cerchio sul piano
        \draw[thick] (-0.4, -1.5) -- (-0.4, 0);
        \draw[thick] (0.4, -1.5) -- (0.4, 0);
        
        % Base inferiore
        \draw[thick, fill=white] (0, -1.5) ellipse (0.4 and 0.15);
        \draw[->, blue, thick] (0, -1.5) -- (0, -2.2) node[right] {$\bm{E}$};
        \draw[->, black, thin] (-0.2, -1.5) -- (-0.2, -2.0) node[left, font=\tiny] {$d\bm{S}$};

        % Parte Sopra
        \draw[thick] (-0.4, 0) -- (-0.4, 1.5);
        \draw[thick] (0.4, 0) -- (0.4, 1.5);
        
        % Base superiore
        \draw[thick, fill=white] (0, 1.5) ellipse (0.4 and 0.15);
        \draw[->, blue, thick] (0, 1.5) -- (0, 2.2) node[right] {$\bm{E}$};
        \draw[->, black, thin] (-0.2, 1.5) -- (-0.2, 2.0) node[left, font=\tiny] {$d\bm{S}$};
        
        % Asse e quote
        \draw[dotted] (0, -1.5) -- (0, 1.5);
        \draw[<->, thin] (-0.8, 0) -- (-0.8, 1.5) node[midway, left] {$z$};
        \draw[->, thin] (-0.8, -1) -- (-0.8, -1.5) node[midway, left] {$z$};
        \draw[<-, dashed] (-0.8, 0) -- (-0.8, -1);
        
        % Vettore normale superficie laterale
        \draw[->, black, thin] (0.4, 0.75) -- (0.8, 0.75) node[right, font=\tiny] {$d\bm{S}_{lat}$};

    \end{tikzpicture}
    \caption{Superficie gaussiana cilindrica che attraversa un piano infinito carico. Il flusso attraverso la superficie laterale è nullo perché il campo è radente alla superficie.}
    \label{fig:piano_infinito}
\end{figure}
Il flusso totale attraverso la superficie cilindrica chiusa $\Sigma_{cil}$ è la somma dei flussi attraverso le due basi ($S_{\text{alto}}, S_{\text{basso}}$) e la superficie laterale ($S_{\text{lat}}$):
\begin{equation}
    \Phi(\bm{E}) = \int_{S_{\text{alto}}} \bm{E} \cdot d\bm{S} + \int_{S_{\text{basso}}} \bm{E} \cdot d\bm{S} + \int_{S_{\text{lat}}} \bm{E} \cdot d\bm{S}.
\end{equation}
Analizzando i tre contributi:
\begin{itemize}
    \item sulla {superficie laterale}, il vettore campo $\bm{E}$ (verticale) è perpendicolare al vettore superficie $d\bm{S}$ (a sua volta perpendicolare alla superficie stessa). Il prodotto scalare è quindi nullo: $\Phi_{\text{lat}} = 0$;
    \item sulle {due basi}, il vettore campo $\bm{E}$ è parallelo ed equiverso al vettore superficie $d\bm{S}$. Essendo il campo uniforme a distanza $z$ fissa, il flusso su ciascuna base è semplicemente $E \cdot A$.
\end{itemize}
Il flusso totale è dunque:
\begin{equation}
    \Phi_{\text{tot}}(\bm{E}) = EA + EA = 2EA.
\end{equation}
La carica totale $Q_{\text{int}}$ racchiusa dal cilindro è la porzione di piano intercettata dal cilindro stesso, che ha area $A$. Dunque:
\begin{equation}
    Q_{\text{int}} = \sigma A.
\end{equation}
Applicando la legge di Gauss ($\Phi = Q_{\text{int}}/\varepsilon_0$):
\begin{equation}
    2EA = \frac{\sigma A}{\varepsilon_0} \implies  
    E = \frac{\sigma}{2\varepsilon_0}.
    \label{eq:campo_piano_scalare}
\end{equation}
Il risultato è quindi indipendente dall'area della superficie gaussiana scelta. In forma vettoriale, indicando con $\bm{u}_n$ il versore normale al piano uscente dalla superficie:
\begin{equation}
    {\bm{E} = \frac{\sigma}{2\varepsilon_0} \bm{u}_n}.
    \label{eq:campo_piano}
\end{equation}
Si noti che il campo generato da un piano infinito è \textbf{uniforme}: non dipende dalla distanza $z$ dal piano stesso.


\section{Distribuzione rettilinea infinita di carica}
Sia \(\lambda\) la densità lineare di carica uniforme su un filo rettilineo infinito. 
Per simmetria cilindrica, il campo elettrico deve essere radiale e dipendere solo dalla distanza \(r\) dal filo.
Si scelga come superficie gaussiana \(\Sigma\) un cilindro coassiale al filo, di raggio \(r\) e altezza \(h\). 

\begin{center}
\begin{tikzpicture}[
    scale=1.2,
    line cap=round,
    line join=round,
    >=Stealth,
    charge/.style={circle, inner sep=1pt, fill=red},
    field/.style={->, thick, blue},
    gauss/.style={dashed, thick, black!70},
    dim/.style={<->, black!50, thin},
    axis/.style={->, thin, black!50}
]


% Coordinate system
\draw[axis] (-2.5,0,0) -- (2.5,0,0) node[right] {$x$};
\draw[axis] (0,-2.5,0) -- (0,2.5,0) node[above] {$y$};
\draw[axis] (0,0,-2.5) -- (0,0,2.5) node[ left] {$z$};

% Infinite wire along z-axis
\draw[thick, red!80!black] (0,0,-2.5) -- (0,0,2.5);


% Add charge symbols along the wire
\foreach \z in {-2,-1.5,-1,-0.5,0,0.5,1,1.5,2} {
    \filldraw[charge] (0,0,\z) circle (1.5pt);
}

% Gaussian cylinder
\draw[gauss] (1.5,0,-2) -- (1.5,0,2);
\draw[gauss] (-1.5,0,-2) -- (-1.5,0,2);
\draw[gauss] (0,1.5,-2) -- (0,1.5,2);
\draw[gauss] (0,-1.5,-2) -- (0,-1.5,2);

% Top and bottom circles of cylinder
\draw[gauss] (0,0,2) ellipse (1.5 and 1.5);
\draw[gauss] (0,0,-2) ellipse (1.5 and 1.5);

% Electric field vectors (radial)
\foreach \angle in {0,45,90,135,180,225,270,315} {
    \foreach \z in {-1.5,0,1.5} {
        \draw[field] ({\angle}:1.5) -- ++({\angle}:0.5);
    }
}

% Labels
\node[red!80!black] at (-0.9,-0.7,3) {$\lambda$ (densità lineare)};
\node at (1.7,-1.7,0) {$\mathbf{E}$};
\node at (0,3,0) {$\mathbf{E}$};

% Dimensions
\draw[dim] (0,0,-2) -- (1.5,0,-2) node[midway, below] {$r$};
\draw[dim] (0,0,-2.5) -- (0,0,2.5) node[midway, right] {$h$};

% Coordinate point
\filldraw (1.5,0,0) circle (1.5pt) node[above right] {$P$};
\draw[dashed, thin] (0,0,0) -- (1.5,0,0) node[midway, above] {$r$};

% Gaussian surface label
\node[black!70] at (-2,2.2,0) {Superficie gaussiana $\Sigma$};

% Add symmetry explanation
\draw[->, green!70!black, thick] (2.5,1.5,0) -- (1.8,1,0);
\node[green!70!black, align=center, text width=3cm] at (3.5,1.8,0) 
    {Simmetria cilindrica:\\$\mathbf{E}$ è radiale e\\dipende solo da $r$};

\end{tikzpicture}
\end{center}

Il flusso attraverso le basi è nullo perché \(\bm{E}\) è radiale e quindi perpendicolare a \(d\bm{S}_{\text{base}}\). 
Attraverso la superficie laterale:
\begin{equation}
    \Phi_{\Sigma}(\bm{E}) = \oint_{\Sigma} \bm{E} \cdot d\bm{S} = E \cdot (2\pi r h),
\end{equation}
poiché \(\bm{E}\) è costante in modulo su tale superficie e parallelo a \(d\bm{S}\).
La carica interna è \(Q_{\text{int}} = \lambda h\). Per la legge di Gauss:
\begin{equation}
    E \cdot 2\pi r h = \frac{\lambda h}{\varepsilon_0} \quad\Rightarrow\quad E = \frac{\lambda}{2\pi\varepsilon_0 r}.
\end{equation}
In forma vettoriale, indicando con \(\bm{u}_r\) il versore radiale uscente dal filo:
\begin{equation}
    \bm{E} = \frac{\lambda}{2\pi\varepsilon_0 r} \bm{u}_r.
\end{equation}
\textit{Osservazione}: il campo decresce con l'inverso della distanza, a differenza del caso puntiforme (\(1/r^2\)) o del piano infinito (costante). 
Il potenziale associato è \(V(r) \propto \ln r\), definito a meno di una costante.

\section{Dipolo elettrico}

Un sistema fisico di fondamentale importanza nello studio dell'elettrostatica e della struttura della materia è il \textit{dipolo elettrico}. Esso è costituito da due cariche puntiformi di uguale modulo $q$ ma segno opposto ($+q$ e $-q$), poste a una distanza $d$ fissa tra loro.
Si definisce innanzitutto il \textit{momento di dipolo elettrico} il vettore $\bm{p}$:
\begin{equation}
    \bm{p} = q \bm{d},
\end{equation}
dove il vettore distanza $\bm{d}$ è diretto, per convenzione, dalla carica negativa verso quella positiva.
Il sistema non possiede evidentemente carica netta ($Q_{tot} = 0$) tuttavia, poiché le cariche non sono coincidenti, il campo elettrico non è nullo, ma possiede una particolare geometria detta \textit{anisotropa}\footnote{Con ciò a significare che varia a seconda della direzione in cui viene misurato.}.

Si consideri un sistema di riferimento con origine nel centro del dipolo e asse $z$ parallelo allo stesso, orientato come $\bm{p}$: le cariche hanno coordinate dunque rispettivamente $(0,0,d/2), q_+$ e $(0,0,-d/2), q_-$. Sia $P$ un punto individuato dal vettore $\bm{r}$ e siano $\bm{r}_1$ e $\bm{r}_2$ le distanze rispettivamente fra la carica positiva, la carica negativa e il punto $P$. Detto $\theta$ l'angolo fra $\bm{r}$ e l'asse $z$, applicando il Teorema del coseno ai triangoli $\widehat{q_+OP}$ e $\widehat{q_-OP}$  si ha:
\begin{equation}
\begin{cases}
    r_1 =\sqrt{ \left(\frac{d}{2}\right)^2 + r^2 - \frac{d}{2}r\cos\theta}  \\
    r_2 = \sqrt{\left(\frac{d}{2}\right)^2 + r^2 + \frac{d}{2}r\cos\theta} 
\end{cases}
,
\end{equation}
Assumendo che $r \gg d$, si può utilizzare l'espansione di McLaurin al primo ordine per $d/r$ infinitesimo: 
\begin{equation}
\begin{cases}
    r_1 \approx r\left(1-\frac{d}{2r}\cos \theta \right) \\
    r_2  \approx r\left(1+\frac{d}{2r}\cos \theta \right)
\end{cases}
.
\end{equation}
\begin{center}
    % Imposta l'angolo di visuale: (elevazione, azimut)
    \tdplotsetmaincoords{70}{120} 
    
    \begin{tikzpicture}[scale=1.3, tdplot_main_coords, >=Stealth]
        
        % --- 1. ASSI CARTESIANI ---
        \draw[->, gray] (0,0,0) -- (3,0,0) node[anchor=north east]{$x$};
        \draw[->, gray] (0,0,0) -- (0,3,0) node[anchor=north west]{$y$};
        \draw[->, gray] (0,0,0) -- (0,0,3.5) node[anchor=south]{$z$};

        % --- 2. IL DIPOLO (Sull'asse Z) ---
        \coordinate (O) at (0,0,0);
        \coordinate (Pos) at (0,0,1);
        \coordinate (Neg) at (0,0,-1);
        
        % Vettore momento di dipolo p
        \draw[ultra thick, ->, purple] (Neg) -- (Pos) node[midway, left] {$\bm{p}$}; ;
        
        % Le cariche (disegnate come sfere)
        \shade[ball color=red!20] (Pos) circle (0.15cm) node[anchor=south east] {$+q$};
        \shade[ball color=blue!20] (Neg) circle (0.15cm) node[anchor=north east] {$-q$};

        % --- 3. IL PUNTO P NELLO SPAZIO ---
        % Definisco P in coordinate sferiche (r=3.5, theta=50, phi=45)
        \tdplotsetcoord{P}{8}{50}{60}
        \tdplotsetcoord{Q}{9}{53}{60}
        
       
        \draw[->, thick, brown!95!black] (P) -- (Q) node[anchor=north] {$\bm{E}$};
        
        % Vettore r (Dall'origine a P)
        \draw[->, thick, black] (O) -- (P) node[anchor=south] {$P(r, \theta)$};
        
        % Linee di proiezione (per dare senso di profondità)
        \draw[dashed, gray] (P) -- (Pxy); % Proiezione su piano XY
        \draw[dashed, gray] (Pxy) -- (O);
        
        % --- 4. GLI ANGOLI ---
        
        % Angolo Theta (Tra asse Z e r)
        % Ruotiamo il sistema di coordinate per disegnare l'arco nel piano corretto
        \tdplotsetthetaplanecoords{60} % Ruota sul piano dell'azimut di P
        \tdplotdrawarc[tdplot_rotated_coords, ->, color=blue]{(0,0,0)}{0.8}{0}{50}{anchor=south west}{$\theta$}
        
        % Angolo Phi (Azimutale sul piano XY - opzionale ma bello)
        \tdplotdrawarc[->, color=gray]{(0,0,0)}{0.6}{0}{60}{anchor=north}{$\phi$}

        % --- 5. DETTAGLI GEOMETRICI (Distanze r1, r2) ---
        % Linee tratteggiate dalle cariche a P
        \draw[densely dashed, red!50!black] (Pos) -- (P) node[midway, above left, font=\tiny] {$r_1$};
        \draw[densely dashed, blue!50!black] (Neg) -- (P) node[midway, below right, font=\tiny] {$r_2$};

    \end{tikzpicture}
    \captionof{figure}{Rappresentazione 3D del dipolo elettrico. Il potenziale in $P$ dipende dall'angolo polare $\theta$ tra il momento di dipolo $\bm{p}$ e il vettore posizione $\bm{r}$.}
    \label{fig:dipolo_3d}
\end{center}
Per il principio di sovrapposizione, scelto lo $0$ a infinito, il potenziale nel punto $P$ risulta:
\begin{equation}
    V(P) = \frac{1}{4\pi\varepsilon_0}\frac{q}{r_1} - \frac{1}{4\pi\varepsilon_0}\frac{q}{r_2} = \frac{q}{4\pi\varepsilon_0r}\left(\frac{1}{1-\frac{d}{2r}\cos \theta} - \frac{1}{1+\frac{d}{2r}\cos \theta}\right)
\end{equation}
\begin{equation}
    = \frac{qd\cos\theta}{4\pi\varepsilon_0r^2} = \frac{q\bm{d}\cdot\bm{u}_r}{4\pi\varepsilon_0r^2},
\end{equation}
dove si sono trascurati infinitesimi di ordine superiore a $\frac{d}{r}$. 
Il campo elettrostatico, in coordinate sferiche, dunque:
\begin{equation}
    \bm{E} = -\nabla V = -\left(\frac{\partial V}{\partial r},\frac{1}{r} \frac{\partial V}{\partial \theta}, \frac{1}{r\sin\theta}\frac{\partial V}{\partial \phi}\right) = \left(\frac{qd\cos\theta}{2\pi \varepsilon_0r^3}\bm{u}_r, \frac{qd\sin\theta}{4\pi \varepsilon_0r^3}\bm{u}_{\theta}, 0\right).
\end{equation}








\chapter{Campi elettrici nella materia}

L'interazione fra campi elettrostatici e materia è argomento di fondamentale importanza, anche in termini prettamente applicativi. In questo capitolo vengono analizzate alcune proprietà del campo elettrostatico indotto in materiale e “generato" dagli stessi (in certe condizioni). 

\section{Conduttori e equilibrio elettrostatico}

Si definiscono \textit{conduttori} quei materiali che possiedono portatori di carica liberi di muoversi macroscopicamente all'interno del materiale stesso.
Nel caso più comune, i \textit{conduttori metallici}, i portatori di carica sono gli elettroni di valenza che, essendo debolmente legati ai nuclei atomici, si delocalizzano formando un "gas di elettroni" libero di scorrere attraverso il reticolo cristallino.

\subsection{Condizioni di equilibrio elettrostatico}
Si dice che un conduttore è in \textit{equilibrio elettrostatico} quando non vi è moto ordinato netto di cariche al suo interno (corrente nulla). In tali condizioni, valgono le seguenti proprietà fondamentali:

\begin{enumerate}
    \item \textbf{Il campo elettrico interno è nullo.}
    \begin{equation}
        \bm{E}_{\text{int}} = 0.
    \end{equation}
    \textit{Dimostrazione:} Se il campo interno non fosse nullo ($\bm{E} \neq 0$), sui portatori di carica agirebbe una forza $\bm{F} = q\bm{E}$ che li metterebbe in moto. Questo violerebbe l'ipotesi di equilibrio statico. Le cariche si ridistribuiscono istantaneamente fino ad annullare il campo interno.

    \item \textbf{Il potenziale elettrico è costante su tutto il volume.}
    Essendo $\bm{E} = -\nabla V$, se il campo è nullo in tutto il volume, il gradiente del potenziale è nullo. Di conseguenza, il potenziale $V$ è costante in ogni punto del conduttore.
    \begin{equation}
        V(P) = \text{costante} \quad \forall P \in \text{Conduttore}.
    \end{equation}
    Il conduttore è, a tutti gli effetti, una superficie equipotenziale.

    \item \textbf{La carica netta si dispone solo sulla superficie.}
    Applicando la legge di Gauss a una superficie chiusa $\Sigma$ interna ad un conduttore \textit{chiuso}:
    \begin{equation}
        \oint_{\Sigma} \bm{E} \cdot d\bm{S} = \frac{Q_{\text{int}}}{\varepsilon_0} = 0.
    \end{equation}
    Poiché $\bm{E}=0$ ovunque all'interno, il flusso è nullo e quindi $Q_{\text{int}} = 0$. Eventuali eccessi di carica devono necessariamente risiedere sulla superficie esterna.

    \item \textbf{Il campo elettrico esterno è perpendicolare alla superficie.}
    Sulla superficie del conduttore, il campo elettrico deve essere ortogonale alla superficie stessa ($\bm{E} \perp d\bm{S}$). Se avesse una componente tangenziale, le cariche superficiali scorrerebbero lungo la superficie, violando l'equilibrio.
\end{enumerate}

\subsection{Teorema di Coulomb}
Il modulo del campo elettrico nelle immediate vicinanze della superficie di un conduttore è direttamente proporzionale alla densità di carica superficiale locale $\sigma$:
\begin{equation}
    \bm{E}_{\text{sup}} = \frac{\sigma}{\varepsilon_0} \bm{n},
\end{equation}
dove $\bm{n}$ è il versore normale uscente dalla superficie. Questo implica che il campo è più intenso dove la densità di carica è maggiore (es. sulle punte o zone a forte curvatura, fenomeno noto come \textit{potere delle punte}).

Con buona approssimazione, una superficie sufficientemente regolare (\textit{differenziabile}) può essere ben approssimata localmente in un punto $P$ da un piano tangente ad essa in $P$. Da qui deriva dunque l'asserzione del Teorema di Coulomb. 


\section{Il condensatore a facce piane parallele}

Il \textit{condensatore} è un sistema fisico costituito da due conduttori separati da un isolante (dielettrico o vuoto), preposto all'accumulo di carica elettrica e, conseguentemente, di energia potenziale elettrostatica.

Il modello più semplice e fondamentale è il \textit{condensatore a facce piane parallele}.
Esso è costituito da due lastre metalliche piane (dette \textit{armature}) di area $A$, poste parallelamente una di fronte all'altra a una distanza $d$.

\subsection{Modello fisico e approssimazioni}
Per trattare analiticamente questo sistema, si introduce l'\textbf{approssimazione delle facce piane indefinite}.
Si assume che la distanza $d$ tra le armature sia molto piccola rispetto alle dimensioni lineari delle armature stesse ($d \ll \sqrt{A}$).

Sotto questa ipotesi, è possibile trascurare gli \textit{effetti di bordo} (ovvero la distorsione delle linee di campo ai bordi delle lastre) e trattare il campo elettrico nella regione interna come se fosse generato da due piani infiniti uniformemente carichi.

\subsection{Analisi del campo elettrico}
Si supponga di collegare le armature a una differenza di potenziale $V$, in modo che su un'armatura si accumuli una carica $+Q$ e sull'altra $-Q$. La densità superficiale di carica sarà uniforme e pari a $\sigma = Q/A$.

Il campo elettrico prodotto da un piano infinito segue la \eqref{eq:campo_piano_scalare}; analizzando le tre regioni dello spazio (una interna al condensatore e due esterne):

\begin{enumerate}
    \item \textbf{Regione esterna (sopra e sotto):} I campi elettrici generati dalla lastra positiva (uscente) e da quella negativa (entrante) hanno stesso modulo ma verso opposto. La somma vettoriale è nulla.
    \[ \bm{E}_{ext} = \bm{E}_+ + \bm{E}_- = 0 \]
    \item \textbf{Regione interna:} I campi elettrici hanno lo stesso verso (dalla lastra positiva alla negativa). I moduli si sommano:
    \[ E_{int} = E_+ + E_- = \frac{\sigma}{2\varepsilon_0} + \frac{\sigma}{2\varepsilon_0} = \frac{\sigma}{\varepsilon_0}. \]
\end{enumerate}

Ne consegue una proprietà fondamentale: \textit{all'interno di un condensatore piano ideale, il campo elettrico è \textbf{uniforme} (costante in modulo, direzione e verso) e diretto dall'armatura positiva a quella negativa. Esternamente è nullo.}

\subsection{Capacità del condensatore piano}
Essendo il campo uniforme diretto unicamente lungo la direzione perpendicolare alle facce del condensatore, la differenza di potenziale $V$ tra le armature è legata al campo dalla relazione:
\begin{equation}
    V = \int_-^+ \bm{E} \cdot d\bm{l} = E \cdot d = \frac{\sigma d}{\varepsilon_0} = \frac{Q d}{\varepsilon_0 A}.
\end{equation}
Scrivendo la carica in funzione della $ddp$, si ottiene:
\begin{equation}
    Q = \varepsilon_0\frac{A}{d}V.
\end{equation}
Si definisce $capacità$ di un condensatore la costante di proporzionalità fra $Q$ e $V$:
\begin{equation}
   Q = CV \implies  C = \varepsilon_0\frac{A}{d}.
\end{equation}
Questa formula evidenzia un concetto fisico cruciale: la capacità di un condensatore non dipende dalla carica o dal potenziale, ma \textbf{esclusivamente dalle caratteristiche geometriche} del sistema (area $A$ e distanza $d$).

\begin{figure}[htbp]
    \centering
    \begin{tikzpicture}[scale=1.2, >=Stealth]
        % Definizione parametri
        \def\plateW{4}   % Larghezza armatura
        \def\plateH{0.3} % Spessore armatura
        \def\dist{2}     % Distanza tra le armature
        
        % --- ARMATURA SUPERIORE (+Q) ---
        \draw[fill=gray!20, thick] (-2, \dist/2) rectangle (2, \dist/2 + \plateH);
        \node at (2.3, \dist/2 + \plateH/2) {$+Q$};
        % Cariche positive
        \foreach \x in {-1.8, -1.4, ..., 1.8} {
            \node[red, font=\small] at (\x, \dist/2 - 0.2) {$+$};
        }
        
        % --- ARMATURA INFERIORE (-Q) ---
        \draw[fill=gray!20, thick] (-2, -\dist/2 - \plateH) rectangle (2, -\dist/2);
        \node at (2.3, -\dist/2 - \plateH/2) {$-Q$};
        % Cariche negative
        \foreach \x in {-1.8, -1.4, ..., 1.8} {
            \node[blue, font=\small] at (\x, -\dist/2 + 0.2) {$-$};
        }
        
        % --- LINEE DI CAMPO ELETTRICO (Uniforme) ---
        \foreach \x in {-1.6, -1.2, ..., 1.6} {
            \draw[->, very thick, orange!90!black] (\x, \dist/2 - 0.4) -- (\x, -\dist/2 + 0.4);
        }
        % Etichetta Campo Elettrico
        \node[orange!90!black, fill=white, inner sep=1pt] at (0,0) {$\bm{E} = \frac{\sigma}{\varepsilon_0}\bm{u}_y$};
        
        % --- DISTANZA d ---
        \draw[<->, thin] (-2.5, \dist/2) -- (-2.5, -\dist/2) node[midway, left] {$d$};
        \draw[dashed, gray] (-2, \dist/2) -- (-2.6, \dist/2);
        \draw[dashed, gray] (-2, -\dist/2) -- (-2.6, -\dist/2);
        
        % --- EFFETTI DI BORDO (Opzionale - didattico) ---
        % Linee curve ai lati per mostrare la realtà fisica
        \draw[->, orange, thin] (1.9, \dist/2 -0.1) .. controls (2.5, 0) .. (1.9, -\dist/2 +0.1);
        \draw[->, orange, thin] (-1.9, \dist/2 -0.1) .. controls (-2.5, 0) .. (-1.9, -\dist/2 +0.1);
        
    \end{tikzpicture}
    \caption{Schema in sezione di un condensatore a facce piane parallele. Nella regione centrale il campo elettrico $\bm{E}$ è uniforme e perpendicolare alle armature. Ai margini si notano i cosiddetti ``effetti di bordo'', trascurabili se $d \ll \sqrt{A}$.}
    \label{fig:condensatore_piano}
\end{figure}

\subsection{Energia immagazzinata in un condensatore}
Ora si tenta di rispondere alla domanda: quanto vale il lavoro necessario per carica un condensatore?
La necessità è quella di trasportare ogni carica (positiva o negativa) sulla corrispondente faccia del condensatore. Comunque avvenga questo "trasporto" di cariche, il lavoro richiesto deve essere lo stesso essendo il campo elettrostatico conservativo. Il lavoro sulla carica infinitesima $dq$, quindi:
\begin{equation}
    \delta W = Vdq = Vd(CV) = CVdV,
\end{equation}
dove si è sfruttato il fatto che la capacità del condensatore è una sua costante caratteristica. Integrando si ottiene:
\begin{equation}
    W = \int_0^{Q_\text{tot}} Vdq  = \int_0^{V} CV'dV = \frac{1}{2}CV^2 = \frac{1}{2}Q_{\text{tot}}V.
\end{equation}
Si noti come non si sono fatte assunzioni riguardo la forma del condensatore. Sfruttando ora invece le relazioni trovate per un condensatore a facce piane parallele
\begin{equation}
    W = \frac{1}{2}\varepsilon_0\frac{A}{d}(Ed)^2 = \frac{1}{2}\varepsilon_0E^2 Ad,
\end{equation}
dove $Ad$ è il volume della regione compresa fra le armature del condensatore, immersa nel campo elettrico $E$. 

\begin{figure}[htbp]
    \centering
    \begin{tikzpicture}[scale=1.2, >=Stealth]
        % Assi
        \draw[->] (0,0) -- (4,0) node[right] {$Q$ (Carica)};
        \draw[->] (0,0) -- (0,3) node[above] {$V$ (Potenziale)};
        
        % Linea V = Q/C
        \draw[thick, blue] (0,0) -- (3,2.5) node[right] {$V(q) = q/C$};
        
        % Area Sottesa
        \fill[blue!10] (0,0) -- (3,2.5) -- (3,0) -- cycle;
        
        % Linee tratteggiate finale
        \draw[dashed] (3,2.5) -- (3,0) node[below] {$Q_{tot}$};
        \draw[dashed] (3,2.5) -- (0,2.5) node[left] {$V_{fin}$};
        
        % Etichetta Area
        \node at (2, 0.8) {Energia $U_e$};
        \node[font=\footnotesize] at (2, 0.4) {(Area = $\frac{1}{2}QV$)};
    \end{tikzpicture}
    \caption{Grafico della tensione in funzione della carica durante il processo di carica. L'energia immagazzinata corrisponde all'area sottesa dalla curva (triangolo), che rappresenta il lavoro compiuto.}
    \label{fig:lavoro_carica}
\end{figure}

\subsubsection{Densità di energia del campo elettrico}
La quantità 
\begin{equation}
    w_e = \frac{1}{2}\varepsilon_0E^2 
    \label{eq:densita_energia}
\end{equation}
corrisponde dunque all'energia spesa per unità di volume per spostare le cariche sulle armature e, dunque, generare il campo elettrico. Se ne deriva che il campo elettrico ha "immagazzinata" una densità volumetrica di energia $w$.  

Sebbene questa formula sia stata derivata per il caso specifico del condensatore piano, essa ha validità \textit{generale}: ovunque nello spazio (nel vuoto) sia presente un campo elettrico $\bm{E}$, a tale campo è associata una densità di energia data dalla \eqref{eq:densita_energia}.




\begin{approfondimento}[colback= white, colframe=green!35!black, no shadow, title={Validità generale della densità di energia elettrostatica}]
Nelle sezioni precedenti è stata ricavata l'espressione per la densità di energia elettrostatica $w_e = \frac{1}{2}\varepsilon_0 E^2$ analizzando il caso particolare del condensatore piano. Si dimostra ora che tale risultato ha validità \textit{generale} per qualsiasi configurazione di campo elettrico statico nel vuoto.
\subsubsection{Lavoro di formazione di una distribuzione di carica}
Si consideri una distribuzione generica di carica volumetrica $\rho(\bm{r})$ in un volume $V$. Il lavoro $W$ necessario per assemblare tale distribuzione portando le cariche dall'infinito è dato da:
\begin{equation}
    W = \frac{1}{2} \int_V \rho(\bm{r}) V(\bm{r}) \, d\tau,
    \label{eq:lavoro_cariche}
\end{equation}
dove $V(\bm{r})$ è il potenziale elettrostatico nel punto $\bm{r}$ e $d\tau$ è l'elemento di volume.
Per esprimere l'energia in funzione del solo campo, si utilizza la prima equazione di Maxwell (Legge di Gauss differenziale), $\rho = \varepsilon_0 \nabla \cdot \bm{E}$, sostituendola nella \eqref{eq:lavoro_cariche}:
\begin{equation}
    W = \frac{\varepsilon_0}{2} \int_V (\nabla \cdot \bm{E}) V(\bm{r}) \, d\tau.
\end{equation}

\subsubsection{Analisi vettoriale}
Si sfrutta ora l'identità vettoriale per la divergenza del prodotto di uno scalare per un vettore:
\begin{equation}
    \nabla \cdot (V \bm{E}) = V (\nabla \cdot \bm{E}) + \bm{E} \cdot (\nabla V).
\end{equation}
Isolando il termine di interesse $V (\nabla \cdot \bm{E})$ e sostituendo nell'integrale, si ottiene:
\begin{equation}
    W = \frac{\varepsilon_0}{2} \left[ \int_V \nabla \cdot (V \bm{E}) \, d\tau - \int_V \bm{E} \cdot (\nabla V) \, d\tau \right].
\end{equation}

Il primo termine, per il \textit{Teorema della Divergenza}, può essere trasformato in un integrale di superficie sulla frontiera $S$ che racchiude il volume $V$:
\begin{equation}
    \int_V \nabla \cdot (V \bm{E}) \, d\tau = \oint_S V \bm{E} \cdot d\bm{S}.
\end{equation}
Se si estende il dominio di integrazione a tutto lo spazio ($V \to \infty$), la superficie $S$ si sposta all'infinito. Per distribuzioni di carica localizzate, a grandi distanze il potenziale decade come $1/r$ e il campo come $1/r^2$. L'integranda decade dunque come $1/r^3$, mentre l'area della superficie sferica cresce come $r^2$. Il flusso totale tende a zero come $1/r$, annullando il primo termine.
Rimane solo il secondo termine. Ricordando la relazione tra campo e potenziale $\bm{E} = -\nabla V$, si ha:
\begin{equation}
    - \int_V \bm{E} \cdot (\nabla V) \, d\tau = - \int_V \bm{E} \cdot (-\bm{E}) \, d\tau = \int_V E^2 \, d\tau.
\end{equation}

\subsubsection{Conclusione}
L'energia totale risulta quindi:
\begin{equation}
    W = \int_{\text{spazio}} \left( \frac{1}{2} \varepsilon_0 E^2 \right) \, d\tau.
\end{equation}
Poiché l'integrale rappresenta l'energia totale sommando i contributi di ogni elemento di volume, l'argomento dell'integrale deve necessariamente rappresentare la \textit{densità di energia locale}:
\begin{equation}
    w_e = \frac{1}{2} \varepsilon_0 E^2.
\end{equation}
Questo conferma che l'espressione è universale e non dipende dalla geometria delle sorgenti, ma è una proprietà intrinseca del campo elettrico.
\end{approfondimento}







\section{I materiali dielettrici}

Un \textit{dielettrico} (o isolante) è un materiale che, a differenza dei conduttori, non possiede cariche libere di muoversi macroscopicamente attraverso il suo volume. Gli elettroni sono infatti strettamente legati ai rispettivi nuclei atomici.
Tuttavia, quando un dielettrico viene immerso in un campo elettrico esterno $\bm{E}_0$, sebbene non si generi una corrente, avviene una ridistribuzione microscopica delle cariche nota come \textbf{polarizzazione}.

\subsection{Meccanismo microscopico: la polarizzazione}
A livello microscopico, le molecole del dielettrico reagiscono alla presenza del campo esterno in due modi principali, a seconda della loro natura chimica:

\begin{enumerate}
    \item \textbf{Polarizzazione per deformazione (molecole apolari):} Se le molecole sono neutre e simmetriche, il campo esterno sposta leggermente il baricentro della nube elettronica (negativa) rispetto al nucleo (positivo), inducendo un \textit{momento di dipolo elettrico} microscopico parallelo al campo.
    \item \textbf{Polarizzazione per orientamento (molecole polari):} Se le molecole possiedono già un momento di dipolo intrinseco (come l'acqua, $H_2O$), ma sono orientate casualmente per l'agitazione termica, il campo esterno tende ad allinearle lungo la direzione delle linee di forza.
\end{enumerate}

In entrambi i casi, l'effetto macroscopico è la comparsa di un momento di dipolo netto per unità di volume (vettore \textbf{polarizzazione} $\bm{P}$).
Ciò comporta la formazione di densità di carica superficiali sulle facce del materiale. Queste cariche dette \textit{cariche di polarizzazione} $q_p$, generano un campo elettrico interno $\bm{E}_p$ che si oppone al campo esterno $\bm{E}_0$.



\section{I materiali dielettrici}

Un \textit{dielettrico} (o isolante ideale) è un materiale caratterizzato dall'assenza di portatori di carica liberi di muoversi su distanze macroscopiche. A differenza dei conduttori, gli elettroni rimangono vincolati ai rispettivi nuclei atomici.
Tuttavia, i dielettrici non sono inerti dal punto di vista elettrico. Quando vengono immersi in un campo elettrico esterno $\bm{E}_0$, la struttura atomica subisce una distorsione microscopica che porta alla formazione di momenti di dipolo orientati. Questo fenomeno è noto come \textbf{polarizzazione}.

\subsection{Meccanismo microscopico}
La risposta del materiale varia a seconda della struttura molecolare, ma conduce al medesimo effetto macroscopico. Si distinguono due meccanismi principali:

\begin{enumerate}
    \item \textbf{Polarizzazione per deformazione (molecole apolari):} In atomi o molecole simmetriche, il campo esterno sposta il baricentro della nube elettronica (negativa) rispetto a quello del nucleo (positivo). Si induce così un momento di dipolo temporaneo, parallelo e concorde al campo esterno.
    \item \textbf{Polarizzazione per orientamento (molecole polari):} Nelle molecole dotate di momento di dipolo intrinseco (es. $H_2O$), l'agitazione termica mantiene solitamente un orientamento casuale. Il campo esterno esercita una coppia torcente che tende ad allineare i dipoli preesistenti lungo le linee di forza.
\end{enumerate}

L'effetto collettivo di questi processi è la comparsa di un \textbf{vettore polarizzazione} $\bm{P}$ (momento di dipolo per unità di volume) e la formazione di densità di carica localizzate sulle superfici del materiale, denominate \textit{cariche di polarizzazione} $q_p$.

\subsection{La costante dielettrica relativa}
Le cariche di polarizzazione generano a loro volta un campo elettrico interno $\bm{E}_p$ che si oppone vettorialmente al campo esterno $\bm{E}_0$.
Il campo elettrico risultante $\bm{E}$ all'interno del mezzo è dato dalla sovrapposizione dei due contributi:
\begin{equation}
    \bm{E} = \bm{E}_0 + \bm{E}_p.
    \label{eq:campo_dielettrico}
\end{equation}
Poiché $\bm{E}_p$ e $\bm{E}_0$ hanno verso opposto, il modulo del campo risultante risulta \textbf{attenuato} rispetto a quello che si avrebbe nel vuoto. Tale attenuazione è descritta dalla relazione:
\begin{equation}
    E = \frac{E_0}{\varepsilon_r},
\end{equation}
dove $\varepsilon_r$ è la \textbf{costante dielettrica relativa} del mezzo. Essendo un fattore di attenuazione, vale sempre la condizione:
\begin{equation}
    \varepsilon_r \ge 1 \quad (\text{con } \varepsilon_r = 1 \text{ solo nel vuoto}).
\end{equation}
Si definisce inoltre la \textit{permittività elettrica assoluta} del mezzo come $\varepsilon = \varepsilon_r \varepsilon_0$.

\subsection{Applicazione ai condensatori}
L'effetto macroscopico più evidente si osserva inserendo un dielettrico tra le armature di un condensatore carico e isolato (carica $Q$ costante).

\begin{figure}[htbp]
    \centering
    \begin{tikzpicture}[scale=1, >=Stealth]
        % --- ARMATURE ---
        \draw[fill=gray!20] (-3, 2) rectangle (3, 2.3);
        \node at (0, 2.6) {$+Q_{\text{libera}}$};
        \foreach \x in {-2.5, -1.5, ..., 2.5} \node[red] at (\x, 2.15) {$+$};

        \draw[fill=gray!20] (-3, -2.3) rectangle (3, -2);
        \node at (0, -2.6) {$-Q_{\text{libera}}$};
        \foreach \x in {-2.5, -1.5, ..., 2.5} \node[blue] at (\x, -2.15) {$-$};

        % --- DIELETTRICO ---
        \draw[thick, fill=blue!5] (-2.8, -1.9) rectangle (2.8, 1.9);
        \node[right, gray, font=\small] at (2.9, -1) {$\varepsilon_r$};

        % --- DIPOLI MICROSCOPICI ---
        \foreach \x in {-2, -1, 0, 1, 2} {
            \foreach \y in {-1.2, 0, 1.2} {
                \draw[fill=white, thin] (\x, \y) ellipse (0.3 and 0.5);
                \node[red, font=\tiny] at (\x, \y-0.25) {$+$};
                \node[blue, font=\tiny] at (\x, \y+0.25) {$-$};
                \draw[->, purple, thin] (\x, \y+0.1) -- (\x, \y-0.1);
            }
        }

        % --- CAMPI ELETTRICI ---
        % E0 (Campo nel vuoto)
        \draw[->, very thick] (-4, 1.5) -- (-4, -1.5) node[midway, left] {$\bm{E}_0$};
        % Ep (Campo di polarizzazione)
        \draw[->, thick, red] (-3.5, -0.8) -- (-3.5, 0.8) node[midway, right] {$\bm{E}_p$};
        % E (Campo risultante)
        \draw[->, very thick, blue] (3.5, 1.5) -- (3.5, -0.5) node[midway, right] {$\bm{E} = \frac{\bm{E}_0}{\varepsilon_r}$};

        % --- CARICHE DI POLARIZZAZIONE ---
        \foreach \x in {-2, -1, ..., 2} \node[blue, font=\small] at (\x, 1.7) {$-$};
        \node[blue, right, font=\footnotesize] at (2.1, 1.7) {$-q_p$};
        
        \foreach \x in {-2, -1, ..., 2} \node[red, font=\small] at (\x, -1.7) {$+$};
        \node[red, right, font=\footnotesize] at (2.1, -1.7) {$+q_p$};

    \end{tikzpicture}
    \caption{Rappresentazione schematica della polarizzazione. Il campo esterno $\bm{E}_0$ induce l'allineamento dei dipoli. Sulle facce del dielettrico appaiono le cariche di polarizzazione $q_p$ che generano un campo opposto $\bm{E}_p$, riducendo il campo netto interno.}
    \label{fig:dielettrico_polarizzazione}
\end{figure}

Il campo netto diminuisce ($E = E_0 / \varepsilon_r$), e di conseguenza diminuisce la differenza di potenziale tra le armature ($V = V_0 / \varepsilon_r$). Poiché la capacità è definita come $C = Q/V$, si ottiene un aumento della stessa:
\begin{equation}
    C = \varepsilon_r C_0 = \varepsilon_r \varepsilon_0 \frac{A}{d}.
\end{equation}
L'utilizzo dei dielettrici è dunque fondamentale nella tecnologia dei condensatori per aumentare la carica accumulabile a parità di tensione e per fornire isolamento galvanico tra le armature.

\subsection{La suscettibilità elettrica}
Per quantificare la risposta del materiale al campo applicato, si introduce il concetto di suscettibilità.

\subsubsection{Definizione e relazione lineare}
Nella maggior parte dei dielettrici comuni (dielettrici lineari) e per campi non eccessivamente intensi, il grado di polarizzazione è direttamente proporzionale al campo elettrico \textit{totale} presente nel mezzo.
Si definisce \textbf{suscettibilità elettrica} la costante adimensionale $\chi_e$ tale che:
\begin{equation}
    \bm{P} = \varepsilon_0 \chi_e \bm{E}.
    \label{eq:suscettibilita}
\end{equation}
È fondamentale notare che la polarizzazione dipende dal campo risultante $\bm{E}$, non dal solo campo esterno $\bm{E}_0$.
La costante $\chi_e$ misura la "facilità" con cui il materiale si polarizza: $\chi_e=0$ corrisponde al vuoto, valori elevati indicano materiali fortemente polarizzabili.

\subsubsection{Relazione generale tra suscettibilità e permettività}
Per ottenere una relazione di validità generale, indipendente dalla geometria del sistema, è necessario introdurre il vettore ausiliario \textit{induzione elettrica} (o spostamento elettrico) $\bm{D}$.
La definizione generale di $\bm{D}$, che tiene conto sia del contributo del vuoto che della risposta del mezzo materiale, è:
\begin{equation}
    \bm{D} = \varepsilon_0 \bm{E} + \bm{P}.
    \label{eq:definizione_D}
\end{equation}

Si consideri un dielettrico lineare, omogeneo e isotropo. In tale approssimazione, il vettore polarizzazione $\bm{P}$ è parallelo e proporzionale al campo elettrico macroscopico $\bm{E}$ secondo la relazione costitutiva definita dalla suscettibilità elettrica $\chi_e$:
\begin{equation}
    \bm{P} = \varepsilon_0 \chi_e \bm{E}.
\end{equation}
Sostituendo tale espressione nella \eqref{eq:definizione_D} e raccogliendo i termini comuni, si ottiene:
\begin{equation}
    \bm{D} = \varepsilon_0 \bm{E} + \varepsilon_0 \chi_e \bm{E} = \varepsilon_0 (1 + \chi_e) \bm{E}.
    \label{eq:D_micro}
\end{equation}

D'altra parte, dal punto di vista fenomenologico macroscopico, la relazione costitutiva che lega l'induzione elettrica al campo elettrico è definita tramite la \textit{costante dielettrica relativa} $\varepsilon_r$:
\begin{equation}
    \bm{D} = \varepsilon \bm{E} = \varepsilon_0 \varepsilon_r \bm{E}.
    \label{eq:D_macro}
\end{equation}
Confrontando le equazioni \eqref{eq:D_micro} e \eqref{eq:D_macro}, ed essendo la relazione valida per un generico campo $\bm{E} \neq 0$, si ricava l'identità fondamentale che lega la grandezza macroscopica ($\varepsilon_r$) alla risposta microscopica del materiale ($\chi_e$):
\begin{equation}
    \varepsilon_r = 1 + \chi_e.
\end{equation}
Tale risultato conferma che la costante dielettrica relativa rappresenta il fattore di scala complessivo della permettività del mezzo rispetto al vuoto, inglobando in sé l'effetto della polarizzazione elettronica e molecolare.

\subsection{Classificazione dei materiali}
La relazione semplice $\bm{P} = \varepsilon_0 \chi_e \bm{E}$ è valida solo sotto specifiche ipotesi. 

\begin{figure}[htbp]
    \centering
    \begin{tikzpicture}[scale=1.2, >=Stealth]
        % Assi
        \draw[->] (0,0) -- (5,0) node[right] {$|\bm{E}|$ (Campo Elettrico)};
        \draw[->] (0,0) -- (0,3.5) node[above] {$|\bm{P}|$ (Polarizzazione)};
        
        % Comportamento Lineare
        \draw[thick, blue] (0,0) -- (3, 2.5) node[midway, above, rotate=40] {Lineare ($\chi_e = \text{cost}$)};
        
        % Saturazione (Non lineare)
        \draw[thick, blue, dashed] (3, 2.5) .. controls (4, 3.3) .. (5, 3.4) node[right] {Saturazione};
        
        % Linee tratteggiate
        \draw[dotted] (3, 0) -- (3, 2.5);
        \draw[dotted] (0, 2.5) -- (3, 2.5);
        
        % Annotazione sulla pendenza
        \draw[->, red] (1.5, 0.5) to[out=90, in=-90] (1.5, 1.2);
        \node[red, font=\footnotesize] at (1.5, 0.3) {Pendenza $\propto \chi_e$};

    \end{tikzpicture}
    \caption{Andamento del modulo della polarizzazione in funzione del campo elettrico applicato. Per campi deboli, la relazione è lineare. Per campi molto intensi, si verifica la saturazione (tutti i dipoli sono ormai allineati).}
    \label{fig:suscettibilita_grafico}
\end{figure}

Un materiale si dice:
\begin{enumerate}
    \item \textbf{Lineare:} Se $\chi_e$ non dipende dal modulo del campo elettrico $E$ (altrimenti si hanno fenomeni di isteresi o saturazione).
    \item \textbf{Omogeneo:} Se $\chi_e$ è costante in ogni punto del materiale (non dipende dalla posizione $\bm{r}$).
    \item \textbf{Isotropo:} Se $\chi_e$ è uno scalare, ovvero le proprietà elettriche sono uguali in tutte le direzioni. Se il materiale fosse \textit{anisotropo} (come certi cristalli), la polarizzazione $\bm{P}$ potrebbe non essere parallela a $\bm{E}$ e $\chi_e$ diventerebbe un \textit{tensore}.
\end{enumerate}



\section{Conduttori fuori equilibrio}



\subsection{Corrente e densità di corrente}
Se internamente ad un conduttore c'è un moto ordinato di carica libera, si registra una corrente elettrica $I$, definita come la quantità di carica che attraversa una data superficie nell'unità di tempo:
\begin{equation}
   I = \frac{dq}{dt}.
\end{equation}

Per descrivere il fenomeno puntualmente, si introduce il vettore \textit{densità di corrente} $\bm{J}$.
Se $n$ è il numero di portatori di carica per unità di volume e $q$ la carica di ciascun portatore (densità di carica volumetrica $\rho = nq$), e se questi si muovono con velocità media di deriva $\bm{v}_d$, allora:
\begin{equation}
   \bm{J} = n q \bm{v}_d.
   \label{eq:j}
\end{equation}
La corrente macroscopica $I$ attraverso una superficie $S$ (come la sezione di un filo) è il flusso del vettore densità di corrente attraverso tale superficie:
\begin{equation}
   I = \int_S \bm{J} \cdot d\bm{S} = nq\int_S \bm{v}_d \cdot d\bm{S}.
\end{equation}
Nel caso semplice di un conduttore di sezione $A$ costante e ortogonale alla velocità dei portatori (che è uniforme), la relazione si semplifica in:
\begin{equation}
   I = J A = n q v_d A.
   \label{eq:corrente}
\end{equation}

\subsection{La legge di Ohm microscopica (Modello di Drude)}
In regime stazionario, la corrente è costante. Ciò implica che la velocità media dei portatori ($\bm{v}_d$) deve essere costante, nonostante essi siano sottoposti alla forza elettrica accelerante $\bm{F} = q\bm{E}$.
Questo apparente paradosso si risolve considerando il modello microscopico (Drude): 
\begin{figure}[htbp]
    \centering
    \begin{tikzpicture}[scale=1.2, >=Stealth]
        % Atomi del reticolo
        \foreach \x in {0, 1.5, 3, 4.5, 6} {
            \foreach \y in {0, 1.5, 3} {
                \shade[ball color=gray!30] (\x,\y) circle (0.2);
            }
        }
        
        % Percorso Elettrone (Zig Zag)
        \draw[thick, blue, ->] (0.2, 1.7) -- (1.3, 2.8);
        \draw[thick, blue, ->] (1.3, 2.8) -- (2.8, 1.7);
        \draw[thick, blue, ->] (2.8, 1.7) -- (3.2, 0.2);
        \draw[thick, blue, ->] (3.2, 0.2) -- (4.3, 1.3);
        \draw[thick, blue, ->] (4.3, 1.3) -- (5.8, 1.5) node[below] {$\bm{v}_{reale}$};
        
        % Vettore Deriva (Drift)
        \draw[ultra thick, red, ->] (0, 1.5) -- (6, 1.5) node[above] {$\bm{v}_d$ (Drift)};
        
        % Campo Elettrico (va opposto agli elettroni se q è negativa, ma qui semplifichiamo il verso)
        \draw[ultra thick, orange, ->] (0, -0.5) -- (2, -0.5) node[right] {$\bm{E}$};
        \node[orange] at (1, -0.8) {Forza $\bm{F} = q\bm{E}$};

    \end{tikzpicture}
    
    \caption{Modello di Drude per la conduzione. Il moto reale (blu) è caotico a causa degli urti; il campo elettrico sovrappone una piccola velocità di deriva costante (rossa).}
    \label{fig:drude_model}
\end{figure}
gli elettroni urtano continuamente contro gli ioni del reticolo cristallino:
\begin{enumerate}
    \item tra un urto e l'altro, l'elettrone è accelerato dal campo;
    \item durante l'urto, l'elettrone cede l'energia cinetica acquisita al reticolo (dissipazione per \textit{effetto Joule}) e la sua velocità media si azzera (o si randomizza).
\end{enumerate}
Detto $\tau$ il \textit{tempo medio di libero cammino} (tempo medio tra due urti), la velocità di deriva raggiunta è quella per cui la forza elettrica bilancia una "forza d'attrito"\footnote{Nel modello di Drude si vogliono gli urti contro il reticolo come una sorta di forza di attrito viscoso: raggiunta una certa velocità limite, l'accelerazione si annulla così come la forza risultante agente sulla carica.} media:
\begin{equation}
   m \frac{v_d}{\tau} = qE \implies v_d = \left( \frac{q\tau}{m} \right) E.
\end{equation}
La quantità tra parentesi è definita \textit{mobilità} $\mu$:
\begin{equation}
    \mu = \frac{q\tau}{m}.
\end{equation}
Sostituendo nella definizione di densità di corrente, otteniamo la \textit{legge di Ohm microscopica}:
\begin{equation}
    \bm{J} = nq\bm{v}_d =  nq (\mu \bm{E}) = (nq\mu) \bm{E} = \sigma \bm{E},
    \label{eq:ohm_micro}
\end{equation}
dove $\sigma = nq\mu$ è la \textit{conducibilità elettrica} del materiale e corrisponde all'inverso della \textit{resistività} $\rho$. La \eqref{eq:ohm_micro} si può quindi anche scrivere come:
\begin{equation}
    \rho\bm{J} = \bm{E}.
    \label{eq:ohm_micro_rho}
\end{equation}

\subsection{La Legge di Ohm macroscopica}
Alle estremità di un campione di materiale $\xi$ lungo $dl$ e di sezione $S$, in cui fluisce una densità di corrente $\bm{J}$, si registra una differenza di potenziale $\Delta V$ che si può ricavare come dalla \eqref{eq:ohm_micro_rho}:
\begin{equation}
	\Delta V = -\int_\xi \bm{E} \cdot d\bm{l} = -\int_\xi \rho\bm{J} \cdot d\bm{l}.
\end{equation}
Nella usuale condizione di densità di corrente uniforme e parallela al vettore superficie, si riscrive l'equazione come:
\begin{equation}
	\Delta V = -\int_\xi \rho\frac{I}{S}  d{l} = - I \int_\xi \rho\frac{dl}{S},
\end{equation}
dove, stanti le approssimazioni discusse in precedenza, si può trattare $I$ come una costante. La costante di proporzionalità fra la differenza di potenziale e la corrente è detta \textit{resistenza} $R$ del materiale, in accordo alla \textit{prima Legge di Ohm}:
\begin{equation}
	\Delta V = -RI
\end{equation}
con 
\begin{equation}
dR = \rho\frac{dl}{S}
\end{equation}
la resistenza (\textit{seconda Legge di Ohm}). Si segnala che il segno $-$ nella formulazione della prima legge è dovuta ad una semplice convenzione (\textit{dei generatori}). Per ovviarlo si fa uso della \textit{convenzione degli utilizzatori}.


\chapter{Magnetostatica}

Sin dall'antichità era stata osservata una particolare interazione fra sbarrette di \textit{magnetite} (ossido di ferro, $\text{Fe}_3\text{O}_4$), simile a quella fra dipoli elettrici: attrazione (o repulsione) avvicinando gli estremi. 
Si è verificato che l'interazione non fosse di genere elettrostatico o gravitazionale\footnote{Una carica $q$ di massa $m$, $ferma$ nelle vicinanze della sbarretta non era minimamente interessata dalla stessa.}, concludendo che doveva esistere una nuova proprietà della materia responsabile del comportamento delle sbarrette. 
Per analogia con la legge di Coulomb, nel XVIII secolo venne ipotizzata l'esistenza della cosiddetta \textit{massa magnetica} (o carica magnetica), concentrata nei poli dei magneti, che avrebbe dovuto giustificare le interazioni, appunto, \textit{magnetiche}.

Tuttavia, questa ipotesi presentava un limite insormontabile: l'impossibilità di isolare un monopolo magnetico. Si notò che qual volta che si provava a dividere un \textit{dipolo magnetico} in $n$ parti, si venivano a costituire $n$ nuovi dipoli, più piccoli in dimensioni, ma apparentemente con le stesse caratteristiche $magnetiche$ del genitore. Si teorizzò dunque l'\textit{inseparabilità dei poli}.

\section{L'esperimento di Oersted}
La svolta nella comprensione della natura fisica del magnetismo avvenne nel 1820 grazie al fisico danese Hans Christian Ørsted. 
Egli osservò che un ago magnetico (una bussola), libero di ruotare, subiva una deviazione quando posto nelle vicinanze di un filo conduttore percorso da corrente elettrica.

\begin{figure}[htbp]
    \centering
    \begin{tikzpicture}[scale=0.9, >=Stealth]
        % Filo conduttore
        \draw[ultra thick, gray] (0,-2) -- (0,2) node[above] {Filo};
        
        % Corrente I
        \draw[->, ultra thick, orange] (0,-1.5) -- (0,1.5) node[midway, right] {$I$};
        
        % Linea di campo magnetico (Cerchio)
        % Disegnamo un'ellisse per la prospettiva
        \draw[thick, blue, dashed] (0,0) ellipse (2 and 0.8);
        
        % Vettori Campo B
        % Davanti
        \draw[->, thick, blue] (0, -0.8) -- (-0.1, -0.8) node[below] {$\bm{B}$};
        % Dietro (accennato)
        %\draw[->, thin, blue] (0, 0.8) -- (0.1, 0.8);
        % Destra
        \draw[->, thick, blue] (2, 0) -- (2, 0.1);
        
        % L'Ago Magnetico (Bussola) posto su una linea di campo
        % Lo posizioniamo "davanti" al filo
        \begin{scope}[shift={(1.5, -0.5)}, rotate=160] 
            % Ago
            \draw[fill=red!80] (0,0) -- (0.15, 0.5) -- (-0.15, 0.5) -- cycle; % Nord (Rosso)
            \draw[fill=blue!30] (0,0) -- (0.15, -0.5) -- (-0.15, -0.5) -- cycle; % Sud (Blu)
            \node[font=\tiny, above] at (0, 0.5) {N};
            \node[font=\tiny, below] at (0, -0.5) {S};
            \fill (0,0) circle (1.5pt); % Perno
        \end{scope}
        
        % Etichetta ago
        \node[right, font=\small] at (2, -1) {Ago magnetico};

    \end{tikzpicture}
    \caption{Esperimento di Oersted. Il passaggio di corrente $I$ nel filo genera un campo magnetico $\bm{B}$ le cui linee di forza sono circonferenze concentriche al filo. L'ago magnetico si allinea tangenzialmente a tali linee.}
    \label{fig:oersted}
\end{figure}

L'esperimento dimostrò inequivocabilmente che:
\begin{enumerate}
    \item le cariche elettriche in moto (corrente) generano un campo magnetico. Se nel conduttore non fluiva corrente, infatti, l'ago magnetico non subiva alcuna deviazione;
    \item la forza magnetica agisce su una direzione perpendicolare sia alla direzione della corrente che alla distanza dal filo;
    \item l'ago si disponeva lungo curve chiuse attorno al filo.
\end{enumerate}
Si concluse che doveva esistere un nuovo $campo$ le cui linee di campo erano "guida" per l'ago magnetico. Venne così definito il \textit{campo di induzione magnetica} (o più semplicemente \textit{campo magnetico}) $\bm{B}$. Similmente al caso del campo elettrostatico, si scelse per convenzione che le sue linee di campo nascessero dal polo positivo di un'eventuale dipolo e morissero in quello negativo\footnote{Nel caso di un campo magnetico generato da una sorgente diversa da un materiale, la convenzione è differente e verrà esplorata in seguito}. 

\section{La forza magnetica e il campo magnetico}

Una particella carica puntiforme $q$ in moto con velocità $\bm{v}$ in una regione immersa in un campo magnetico $\bm{B}$ risente di una forza sperimentalmente osservata le cui caratteristiche sono:
\begin{itemize}
    \item proporzionalità diretta al valore della carica $q$ e al modulo del campo $B$;
    \item proporzionalità alla componente di velocità perpendicolare alle linee di campo ($v_{\perp} = v \sin\theta$).
\end{itemize}
In formule scalari:
\begin{equation}
    F \propto q v B \sin\theta.
\end{equation}
Riscrivendo la relazione in termini vettoriali, si ottiene la definizione della \textit{forza magnetica}:
\begin{equation}
    \bm{F} = q \bm{v} \times \bm{B}.
    \label{eq:forza_lorentz}
\end{equation}

\begin{figure}[htbp]
    \centering
    \begin{tikzpicture}[scale=1.2, >=Stealth]
        % Coordinate
        \coordinate (O) at (0,0,0);
        
        % Vettori
        \draw[->, thick, blue] (O) -- (2,0,0) node[anchor=north east] {$\bm{v}$};
        \draw[->, thick, blue] (O) -- (0.5, 2, 0) node[anchor=north west] {$\bm{B}$};
        
        % Forza (Prodotto vettoriale v x B per q>0)
        % v=(2,0,0), B=(0.5, 2, 0). Il prodotto vettoriale è lungo Z.
        \draw[->, ultra thick, red] (O) -- (0,0,2.5) node[anchor=south east] {$\bm{F} = q \bm{v} \times \bm{B}$};
        
        % Piano v-B
        \fill[blue!5, opacity=0.5] (0,0,0) -- (2,0,0) -- (2.5, 2, 0) -- (0.5, 2, 0) -- cycle;
        
        % Angolo theta
        \draw[->, gray] (0.8,0,0) arc (0:76:0.8);
        \node at (1, 0.6, 0) {$\theta$};
        
        % Simbolo perpendicolare su F
        \draw[thin] (0,0,0.4) -- (0.2,0,0.4) -- (0.2,0,0);
        
        % Nota q>0
        \node[red, font=\small] at (0,0,2.8) {($q > 0$)};

    \end{tikzpicture}
    \caption{Rappresentazione vettoriale della forza di Lorentz. La forza $\bm{F}$ è sempre perpendicolare al piano individuato dai vettori velocità $\bm{v}$ e campo magnetico $\bm{B}$, secondo la regola della mano destra (per $q>0$).}
    \label{fig:regola_mano_destra}
\end{figure}


\subsection{Lavoro della forza magnetica}
Una conseguenza fondamentale della struttura vettoriale della forza magnetica è l'essere sempre ortogonale alla velocità istantanea della particella ($\bm{F} \perp \bm{v}$).
Pertanto, il lavoro fatto dalla forza magnetica su una carica in moto è identicamente nullo:
\begin{equation}
    W = \bm{F} \cdot d\bm{s} = \bm{F} \cdot \frac{d\bm{s}}{dt} dt = \bm{F} \cdot \bm{v} dt = 0.
\end{equation}
\textit{Il campo magnetico statico non compie lavoro}: esso può modificare la \textit{direzione} della velocità (il verso di $\bm{v}$), ma non può modificarne il \textit{modulo}: l'energia cinetica della particella non è influenzata dalle interazioni di natura magnetostatiche.

    \subsection{Definizione operativa di \texorpdfstring{$\bm{B}$}{B}}
La relazione \eqref{eq:forza_lorentz} permette di fornire una \textit{definizione operativa} del vettore induzione magnetica $\bm{B}$. In un punto dello spazio:
\begin{enumerate}
    \item la direzione di $\bm{B}$ è quella lungo la quale una carica di prova lanciata con velocità $\bm{v}$ \textit{non} subisce alcuna forza magnetica (caso $\bm{v} \parallel \bm{B}$);
    \item il verso e il modulo sono determinati misurando la forza massima $\bm{F}_{max}$ che si ottiene quando la velocità è perpendicolare a tale direzione.
\end{enumerate}

\begin{approfondimento}[colback= white, colframe=green!35!black, no shadow, title={Moto di cariche in regioni di campo magnetico uniforme}]
\textbf{Caso 1: velocità perpendicolare al campo} \\
Si analizza inizialmente il caso di una particella di carica $q$ e massa $m$ in moto in una regione di campo uniforme il cui vettore velocità è costante e perpendicolare a quello di induzione magnetica. La forza magnetica, stanti le considerazioni precedenti, si configura in tutto e per tutto come una \textit{forza centripeta}, costringendo $q$ a un moto circolare uniforme:
\begin{equation}
    \bm{F} = m\bm{a_c} = q \bm{v} \times \bm{B}
\end{equation}
che in forma scalare diventa
\begin{equation}
    m\frac{v^2}{R} = q v  B \implies R = \frac{mv}{qB}.
\end{equation}

\begin{center}
        \begin{tikzpicture}[scale=1, >=Stealth]
            % Campo magnetico (Entrante - croci)
            \foreach \x in {0, 1, ..., 4} {
                \foreach \y in {0, 1, ..., 3} {
                    \node[blue!50] at (\x, \y) {$\times$};
                }
            }
            \node[blue, right] at (4.2, 1.5) {$\bm{B}$ (entrante)};
    
            % Traiettoria circolare
            \draw[thick, dashed, red] (2, 1.5) circle (1.2);
            \node[red] at (2, 1.5) {$\bullet$}; % Centro
            
            % Particella in un punto
            \coordinate (P) at (3.2, 1.5); % A destra
            \shade[ball color=red] (P) circle (0.15) node[above right] {$q$};
            
            % Vettori v e F
            \draw[->, thick, black] (P) -- (3.2, 2.5) node[right] {$\bm{v}$};
            \draw[->, very thick, red] (P) -- (2.2, 1.5) node[below right] {$\bm{F}$};
            
            % Raggio
            \draw[->, thin] (2, 1.5) -- (2.85, 2.35) node[ above, left ] {$R$};
        \end{tikzpicture}
        \captionof{figure}{Moto circolare uniforme generato dalla forza di Lorentz.}
    \end{center}

    Il valore di $R$ è particolarmente interessante in applicazioni come lo \textit{spettrometro di massa}, uno strumento tramite cui si misura la massa di particelle. A parità di carica, infatti, il raggio di curvatura del moto di particelle in una regione di campo uniforme e perpendicolare alla loro velocità è direttamente proporzionale alla loro massa. Si segnala che, in generale, la spettrometria di massa fornisce informazioni interessanti circa il rapporto $\frac{q}{m}$, vastamente utilizzato in ambito teorico e applicativo. \\
    
    \textbf{Caso 2: velocità con componente parallela al campo non nulla} \\
    Nel caso in cui il vettore velocità della particella non sia perpendicolare al campo si ha: 
    \begin{equation}
        \bm{F} = q \bm{v} \times \bm{B} = q(\bm{v}_\perp + \bm{v}_\parallel)\times \bm{B} = q\bm{v}_\perp \times \bm{B} + q\bm{v}_\parallel \times \bm{B} = q\bm{v}_\perp \times \bm{B}.
        \end{equation}
       La forza agisce solo sul piano perpendicolare a $\bm{B}$, causando un moto circolare con velocità $v_\perp$. Lungo la direzione del campo, invece, non agisce alcuna forza ($a_\parallel = 0$), quindi la particella avanza di moto rettilineo uniforme con velocità $v_\parallel$.
    
    La composizione di un moto circolare uniforme e di un moto rettilineo uniforme perpendicolare al piano del cerchio genera una traiettoria \textbf{elicoidale}.
    
    \begin{center}
        \begin{tikzpicture}[x={(-0.6cm,-0.4cm)}, y={(1cm,0cm)}, z={(0cm,1cm)}, scale=1.2, >=Stealth]
            % Asse Z (Campo B)
            \draw[->, blue, thick] (0,0,0) -- (0,0,4.5) node[right] {$\bm{B}$};
            
            % Traiettoria Elicoidale
            \draw[red, thick, samples=100, domain=0:12.5] plot ({0.8*cos(\x r)}, {0.8*sin(\x r)}, {0.3*\x});
            
            % Particella in un punto
            \coordinate (P) at ({0.8*cos(10 r)}, {0.8*sin(10 r)}, {3});
            \shade[ball color=red] (P) circle (0.12);
            
            % Componenti velocità
            \draw[->, dashed, black] (P) -- ++(0,0,0.8) node[right, font=\tiny] {$\bm{v}_\parallel$};
            % Tangente al cerchio
            \draw[->, dashed, black] (P) -- ++({-0.8*sin(10 r)}, {0.8*cos(10 r)}, 0) node[above, font=\tiny] {$\bm{v}_\perp$};
            
            % Passo dell'elica
            \draw[|<->|, gray] (1.2, 0, 0) -- (1.2, 0, {0.3*2*pi}) node[midway, left, font=\footnotesize] {Passo $p$};
        \end{tikzpicture}
        \captionof{figure}{Moto elicoidale: la particella ruota attorno alle linee di campo mentre avanza lungo di esse.}
    \end{center}
    \end{approfondimento}
   

\subsection{Principio di sovrapposizione degli effetti}
Per il campo magnetico, come nel caso del campo elettrico, vale il principio di sovrapposizione degli effetti: infatti, il campo totale $\bm{B}(\bm{r})$ generato in un punto $P$ da un sistema di più sorgenti è pari alla \textit{somma vettoriale} dei campi che ciascuna sorgente genererebbe in quel punto se agisse singolarmente.
Se il campo è generato da $N$ sorgenti distinte, il campo magnetico totale nel punto $\bm{r}$ è:
\begin{equation}
    \bm{B}_{tot}(\bm{r}) = \sum_{i=1}^{N} \bm{B}_i(\bm{r}),
\end{equation}
dove $\bm{B}_i$ è il campo generato dalla sola $i$-esima sorgente.
    
  \section{Origine del campo}
Dall'esperimento di Oersted si deduce che la sorgente fisica del campo magnetico è la carica elettrica in movimento. Sebbene storicamente le leggi siano state formulate per correnti macroscopiche stazionarie (flussi continui di cariche nei conduttori, studiati da Biot, Savart e Ampère), è fondamentale definire il campo generato dalla sorgente elementare: la singola carica puntiforme in moto. Sfruttando l'equivalenza formale tra un elemento di corrente e una carica in moto ($I \propto qn$), è possibile scrivere la legge che descrive il campo B generato da una particella.

\subsection{Campo magnetico prodotto dalla singola carica}
Sulla base di evidenze sperimentali fornite da Ampere, il campo magnetico prodotto da una carica puntiforme in moto con velocità $\bm{v}$, inizialmente in posizione $\bm{r}'$, in un punto dello spazio identificato dal vettore posizione $\bm{r}$:
\begin{equation}
\bm{B}(P) \propto \frac{q\bm{v} \times \bm{u}_r}{|\bm{r}-\bm{r}'|^2},
\end{equation}
dove $\bm{u}_r$ rappresenta il versore giacente sulla congiungente fra la carica e il punto $P$ e $|\bm{r}-\bm{r}'|^2$ rappresenta la loro distanza. 

\begin{figure}[htbp]
    \centering
    \begin{tikzpicture}[x={(0.8cm,-0.3cm)}, y={(0.6cm,0.5cm)}, z={(0cm,1cm)}, scale=1.5, >=Stealth]
        
        % --- 1. DEFINIZIONE COORDINATE ---
        \coordinate (O) at (0,0,0);       % Posizione Carica q
        \coordinate (P) at (3,0,0);       % Punto P
        \coordinate (V_end) at (0,0,2.5); % Fine vettore v
        \coordinate (Proj) at (0,0,0);    % Proiezione di P sulla retta (in questo caso coincide con O per semplicità geometrica relativa, ma concettualmente è sull'asse z)
        
        % Spostiamo la carica un po' "sotto" per mostrare la linea di moto intera
        \coordinate (Q_start) at (0,0,-1.5);
        \coordinate (Q_pos) at (0,0,0);
        \coordinate (Line_end) at (0,0,3);

        % --- 2. IL PIANO (v, ur) ---
        % Il piano generato da v (asse z) e ur (asse x)
        \fill[blue!5, opacity=0.8] (Q_start) -- (3,0,-1.5) -- (3,0,3) -- (0,0,3) -- cycle;
        \node[blue!40!black, rotate=-20, font=\small] at (2,0,2.5) {Piano generato da $\bm{v}$ e $\bm{u}_r$};

        % --- 3. LINEA DI MOTO E CARICA ---
        \draw[dashed, gray] (Q_start) -- (Line_end) node[above] {Direzione moto};
        \draw[->, ultra thick, red] (Q_pos) -- (V_end) node[midway, left] {$\bm{v}$};
        \shade[ball color=red] (Q_pos) circle (0.12) node[left=0.2cm] {$q$};

        % --- 4. VETTORI POSIZIONE ---
        \draw[->, thick, black] (Q_pos) -- (P) node[midway, below] {$\bm{r}$};
        
        % Versore ur (piccolo)
        \draw[->, thick, blue!80!black] (Q_pos) -- (1,0,0) node[below] {$\bm{u}_r$};

        % --- 5. GEOMETRIA DEL CAMPO B ---
        % Circonferenza (Linea di campo) sul piano XY (perpendicolare a v)
        \draw[blue, thick] (0,0,0) circle (3);
        \node[blue, below] at (2,-2.5,0) {Linea di campo};

        % Punto P
        \fill (P) circle (1.5pt) node[below right] {$P$};
        
        % Proiezione (Coincide con Q in questo istante per semplicità del disegno)
        %\fill (Proj) circle (1pt); 
        %\draw[dotted] (P) -- (Proj);

        % Vettore B
        % Deve essere perpendicolare al piano (quindi lungo Y) e tangente al cerchio
        \draw[->, ultra thick, blue] (P) -- (3, 1.5, 0) node[right] {$\bm{B}$};

        % --- 6. SIMBOLI DI ORTOGONALITÀ ---
        % B perpendicolare al piano (v, ur)
        \draw[thin] (3, 0.3, 0) -- (2.7, 0.3, 0) -- (2.7, 0, 0); 
        
    \end{tikzpicture}
    \caption{Relazione geometrica tra velocità e campo magnetico. Il vettore $\bm{B}$ è perpendicolare al piano individuato dalla velocità $\bm{v}$ e dalla posizione $\bm{r}$, risultando tangente alla linea di campo circolare.}
    \label{fig:geometria_B_carica}
\end{figure}
\noindent
La costante di proporzionalità che completa la legge è $k = \frac{\mu_0}{4\pi}$, con $\mu_0 = 4\pi \cdot 10^{-7} \, \si{\tesla\cdot\meter\per\ampere}$ la \textit{permeabilità magnetica del vuoto}
\begin{equation}
\bm{B}(P) =  \frac{\mu_0}{4\pi} \frac{q\bm{v} \times \bm{u}_r}{|\bm{r}-\bm{r}'|^2}.
\label{eq:B_carica}
\end{equation}
Il vettore campo magnetico è dunque perpendicolare al piano generato da $\bm{u}_r$ e $\bm{v}$: nel caso di una carica puntiforme in moto, quindi, $\bm{B}$ è tangente alla circonferenza che ha per centro la proiezione di $P$ sulla direzione istantanea di moto di $q$ e raggio la distanza fra $P$ e la sua proiezione.


\subsection{Campo magnetico e corrente: legge di Ampere-Laplace}
Si vuole a questo punto estendere la \eqref{eq:B_carica} nel caso di un filo di lunghezza $l$ e sezione $S$ ($\sqrt{S} \ll l$) percorso da corrente. In un elemento di volume $dV = Sdl$, si ha una densità volumetrica $n$ di portatori di carica in moto alla velocità di deriva $v_d$. Il campo magnetico prodotto dalla corrente in un punto $P$ esterno al filo è dato dalla sovrapposizione dei contributi dovuti a ciascuna carica che, nell'approssimazione di un flusso di portatori stazionario e laminare, devono essere paralleli:
\begin{equation}
d\bm{B} = \frac{\mu_0}{4\pi} Sdl'  \frac{nq\bm{v} \times \bm{u}_r}{|\bm{r}-\bm{r}'|^2}.
\end{equation}
Assumendo che ogni carica si muove tangenzialmente al filo e ricordando la \eqref{eq:j} si ha
\begin{equation}
d\bm{B} = \frac{\mu_0}{4\pi} Sdl  \frac{\bm{J} \times \bm{u}_r}{|\bm{r}-\bm{r}'|^2} = \frac{\mu_0}{4\pi} \frac{Sdl(J\bm{u}_t) \times \bm{u}_r}{|\bm{r}-\bm{r}'|^2} = \frac{\mu_0}{4\pi} \frac{JS(dl\bm{u}_t) \times \bm{u}_r}{|\bm{r}-\bm{r}'|^2} = \frac{\mu_0I}{4\pi} \frac{d\bm{l}' \times \bm{u}_r}{|\bm{r}-\bm{r}'|^2},
\end{equation}
dove si è \textit{vettorializzata} la lunghezza infinitesima di filo $dl$ (il verso di $d\bm{l}'$ è concorde al moto dei portatori\footnote{Si è usato l'apostrofo per il vettore $d\bm{l}'$ per specificare il fatto che la sua posizione, appartenente al filo, è individuata da $\bm{r}'$.}). Integrando ambo i membri sull'intero filo conduttore $\gamma$ si ottiene la \textit{legge di Ampere-Laplace in forma integrale}:
\begin{equation}
\bm{B} = \int_\gamma \frac{\mu_0I}{4\pi} \frac{d\bm{l}' \times \bm{u}_r}{|\bm{r}-\bm{r}'|^2}.
\end{equation}

% Imposta l'angolazione della vista 3D (theta=60, phi=120)
\tdplotsetmaincoords{60}{120}

\begin{figure}[htbp]
\begin{center}
\begin{tikzpicture}[tdplot_main_coords, scale=1, >=stealth]

    % --- 1. Coordinate di base ---
    \coordinate (O) at (0,0,0);
    \coordinate (Sorgente) at (2,1,0);  % Posizione di dl' (r') al centro del filo
    \coordinate (PuntoP) at (2,3.5,3);  % Posizione del campo (r)

    % --- 2. Assi Cartesiani ---
    \draw[->, thick, gray!60] (0,0,0) -- (4,0,0) node[anchor=north east, text=black]{$x$};
    \draw[->, thick, gray!60] (0,0,0) -- (0,5,0) node[anchor=north west, text=black]{$y$};
    \draw[->, thick, gray!60] (0,0,0) -- (0,0,4) node[anchor=south, text=black]{$z$};

    % --- 3. Il Filo (Cilindretto curvo) ---
    % Usiamo 'double' per creare l'effetto tubo.
    % draw=... definisce il colore dei bordi esterni
    % double=... definisce il colore di riempimento interno (grigio chiaro metallico)
    % double distance=... definisce lo spessore del cilindro
    \draw[smooth, draw=black!70, thick, double=gray!25, double distance=7pt]
        plot coordinates {(0.5, -0.5, 0) (Sorgente) (3.5, 2, 0)};

    % Etichette del filo
    \node at (4.2, 1.8, 0) [right, xshift=3pt] {$\gamma$};
    \node at (1, -0.5, 0.2) [below, yshift=-3pt] {$I$};

    % --- 4. Vettori Posizione ---
    % Disegnati "sotto" gli elementi fisici per chiarezza
    \draw[dashed, ->, gray!80] (O) -- (Sorgente) node[midway, above right, text=black] {$\mathbf{r}'$};
    \draw[dashed, ->, gray!80] (O) -- (PuntoP) node[midway, above left, text=black] {$\mathbf{r}$};

    % --- 5. Elementi Fisici (Biot-Savart) ---

    % A) Elemento di corrente dl' (Rosso)
    % Parte dal centro del cilindretto
    \draw[very thick, red, ->] (Sorgente) -- ++(0.6, 0.45, 0) node[anchor= west] {$d\bm{l}' \parallel \bm{J}$};

    % B) Vettore Distanza (Blu)
    \draw[thick, blue, ->] (Sorgente) -- (PuntoP) node[midway, right] {$\mathbf{r} - \mathbf{r}'$};

    % C) Campo Magnetico elementare dB (Verde scuro)
    \draw[very thick, green!40!black, ->] (PuntoP) -- ++(-1, -0.5, 0) node[anchor=south] {$d\mathbf{B}$};

    % --- 6. Marcatore Punto P ---
    \fill (PuntoP) circle (1.5pt) node[anchor=south west] {$P$};

\end{tikzpicture}
\caption{Rappresentazione della situazione fisica descritta nella valutazione della legge di Ampere-Laplace.}
    \label{fig:geometria_B_carica}
    \end{center}
\end{figure}



\section{Il flusso di \texorpdfstring{$\bm{B}$}{campo magnetico} attraverso una superficie chiusa}
\label{sec:nablaB}
Detta $\Sigma$ una superficie chiusa, si vuole valutare il flusso del campo magnetico prodotto da un filo conduttore $\gamma$ attraverso $\Sigma$. Applicando il Teorema della divergenza alla legge di Ampere-Laplace si ha:

\begin{equation}
    \nabla \cdot \bm{B} = \nabla \cdot \frac{\mu_0 I}{4\pi} \int_\gamma \frac{d\bm{l}' \times (\bm{r}-\bm{r}')}{|\bm{r}-\bm{r}'|^3},
\end{equation}
dove si è riscritto $\frac{\bm{u}_r}{|\bm{r}-\bm{r}'|^2}$ come $\frac{\bm{r}-\bm{r}'}{|\bm{r}-\bm{r}'|^3}$ per rendere più evidenti considerazioni successive. Si noti che l'operatore $\nabla$ agisce sulle coordinate del punto in cui si valuta il campo $\bm{r}$, mentre l'integrale opera sulle coordinate della sorgente $\bm{r}'$: le due variabili sono indipendenti, gli operatori pertanto commutano:
\begin{equation}
    \nabla \cdot \bm{B} = \frac{\mu_0 I}{4\pi} \int_\gamma \nabla \cdot \left( d\bm{l}' \times \frac{\bm{r}-\bm{r}'}{|\bm{r}-\bm{r}'|^3} \right).
    \label{eq:div_integral}
\end{equation}
Si utilizza ora l'identità vettoriale per la divergenza di un prodotto vettoriale:
\begin{equation}
    \nabla \cdot(\bm{A} \times \bm{B}) = \bm{B}\cdot(\nabla\times\bm{A}) -\bm{A}\cdot(\nabla\times\bm{B});
\end{equation}
ponendo $\bm{A} = d\bm{l}'$, $\bm{B} = \frac{\bm{r}-\bm{r}'}{|\bm{r}-\bm{r}'|^3}$ e sostituendo nella \eqref{eq:div_integral}, si ottiene:
\begin{equation}
    \nabla \cdot \bm{B} = \frac{\mu_0 I}{4\pi} \int_\gamma \left[ \frac{\bm{r}-\bm{r}'}{|\bm{r}-\bm{r}'|^3} \cdot (\nabla \times d\bm{l}') - d\bm{l}' \cdot \left(\nabla \times \frac{\bm{r}-\bm{r}'}{|\bm{r}-\bm{r}'|^3} \right)\right].
\end{equation}
Analizzando i due termini dell'integranda:
\begin{itemize}
    \item $\nabla \times d\bm{l}' = \bm{0}$: $d\bm{l}'$ dipende unicamente dalla posizione $\bm{r}'$ lungo il filo. Poiché $\nabla$ opera derivate parziali rispetto a $\bm{r}$, $d\bm{l}$ si comporta come una costante vettoriale;
    \item $\nabla \times \frac{\bm{r}-\bm{r}'}{|\bm{r}-\bm{r}'|^3} = \bm{0}$: il termine vettoriale corrisponde al campo generato da una sorgente puntiforme (da un \textit{monopolo} come, ad esempio, una carica puntiforme per il campo elettrostatico), esprimibile come il gradiente (cambiato di segno) del potenziale scalare $1/|\bm{r}-\bm{r}'|$, ovvero 
    \begin{equation}
 \frac{\bm{r}-\bm{r}'}{|\bm{r}-\bm{r}'|^3} = -\nabla \left(\frac{1}{|\bm{r}-\bm{r}'|}\right).
 \label{eq:help}
 \end{equation}
  Poiché il rotore di un gradiente è sempre identicamente nullo\footnote{Questo risultato può anche essere ricavato in maniera differente: a meno di costanti, campi centrali \textit{conservativi} come quello elettrico e quello gravitazionale si presentano con la stessa dipendenza dalla distanza $\sim r^{-2}$. In entrambi i casi, data la conservatività, è garantita la condizione di rotore nullo.} ($\nabla \times \nabla \phi \equiv 0$), il termine si annulla.
\end{itemize}
L'integranda è dunque nulla ovunque: si perviene così alla \textit{seconda equazione di Maxwell in forma differenziale} (legge di Gauss per il magnetismo):
\begin{equation}
{
    \nabla \cdot \bm{B} = 0}.
    \label{eq:2_max}
\end{equation}
Per questioni di nomenclatura il campo magnetico, come tutti i campi vettoriali a divergenza unicamente nulla, è detto \textit{solenoidale}.

\subsection{Formulazione integrale}
La formulazione integrale della seconda equazione di Maxwell discende direttamente dall'applicazione del Teorema della divergenza alla \eqref{eq:2_max}. Detto $V$ il volume racchiuso dalla superficie $\Sigma$, 
\begin{equation}
{
    \oint_{\Sigma = \partial V} \bm{B} \cdot d\bm{S} = \int_V (\nabla \cdot \bm{B}) \, dV = 0}.
\end{equation}

\subsubsection{Significato fisico}
Serve ricordare il significato della divergenza di un campo vettoriale: l'operatore $(\nabla \cdot )$ suggerisce informazioni riguardo la tendenza che ha un campo a sgorgare o convergere in un "volumetto" infinitesimo di spazio. Essendo la divergenza di $\bm{B}$ sempre unicamente nulla, si conclude che il campo magnetico non tende a convergere \textit{e neppure a sgorgare} da nessun punto dello spazio: non esistono \textit{monopoli magnetici}\footnote{A differenza del caso elettrico in cui i monopoli esistono eccome e sono, ad esempio, le cariche puntiformi. Si ragioni sul fatto che $\nabla \cdot \bm{E} = \frac{\rho}{\varepsilon_0}$.}, con ciò a implicare che, spezzando un magnete in due parti, non si riesce a isolare il polo negativo da quello positivo. Ci si limita infatti a creare due magneti più piccoli.

\section{La circuitazione di \texorpdfstring{$\bm{B}$}{campo magnetico} lungo una curva chiusa}
Per comprendere a dovere il fondamento che risiede dietro il calcolo esplicito della circuitazione del campo magnetico lungo una curva chiusa $\gamma$, è necessario osservare $\bm{B}$ \textit{da un'altra prospettiva}. Ecco che arriva in soccorso il \textit{Teorema di Helmholtz} e la definizione di \textit{potenziale vettore}.




\subsection{Il potenziale vettore \texorpdfstring{$\bm{A}$}{A}}
Il Teorema di Helmholtz garantisce che un qualsiasi campo vettoriale, dovutamente continuo e con opportune condizioni di decadimento all'infinito, è scomponibile nella somma di una componente puramente solenoidale (rotore di un campo vettoriale) e di una puramente irrotazionale (gradiente di un campo scalare). È quindi possibile analizzare la struttura del campo magnetico statico alla luce della sua natura intrinsecamente solenoidale:
\begin{equation}
    \nabla \cdot \bm{B} = 0.
\end{equation}
Per il Teorema di Helmholtz deve esistere dunque un campo vettoriale $\bm{A}(\bm{r})$, detto \textit{potenziale vettore}, tale che:
\begin{equation}
    \bm{B} = \nabla \times \bm{A}.
    \label{eq:def_potenziale_vettore}
\end{equation} 

\begin{approfondimento}[colback= white, colframe=green!35!black, no shadow, title={Il Teorema di Helmholtz}]
Il \textit{Teorema fondamentale del calcolo vettoriale}, o Teorema di Helmholtz, fornisce una base teorica rigorosa per la scomposizione dei campi vettoriali. Esso afferma che un qualsiasi campo vettoriale $\bm{F} : \mathbb{R}^m \supseteq D \rightarrow \mathbb{R}^m$, $\bm{F} \in \mathbb{C}^1(D, \mathbb{R}^m)$ con componenti $f_i(r)$ tali che $f_i \in o(\frac{1}{r})$ per $r \rightarrow \infty$, è univocamente determinato se sono note la sua divergenza (le sue sorgenti scalari) e il suo rotore (le sue sorgenti vettoriali o vorticità) in tutto lo spazio.
Matematicamente, il Teorema asserisce che ogni campo vettoriale può essere scomposto nella somma di due componenti:
\begin{equation}
    \bm{F} = \bm{F}_{\text{irr}} + \bm{F}_{\text{sol}},
\end{equation}
dove:
\begin{itemize}
    \item $\bm{F}_{\text{irr}}$ è la componente \textit{irrotazionale} (o longitudinale), tale che $\nabla \times \bm{F}_{\text{irr}} = \bm{0}$. Per le proprietà del gradiente\footnote{$\nabla \times (\nabla \phi) = 0.$}, questa componente può essere espressa come il gradiente di un campo scalare $\phi$:
    \begin{equation}
        \bm{F}_{\text{irr}} = -\nabla \phi;
    \end{equation}
    \item $\bm{F}_{\text{sol}}$ è la componente \textit{solenoidale} (o trasversale), tale che $\nabla \cdot \bm{F}_{\text{sol}} = 0$. Per le proprietà del rotore\footnote{$\nabla \cdot (\nabla \times \bm{A}) = 0.$}, questa componente può essere espressa come il rotore di un campo vettoriale $\bm{A}$:
    \begin{equation}
        \bm{F}_{\text{sol}} = \nabla \times \bm{A}.
    \end{equation}
\end{itemize}
In definitiva, il Teorema di Helmholtz garantisce che un generico campo $\bm{F}$ possa essere decomposto come:
\begin{equation}
    \bm{F} = -\nabla \phi + \nabla \times \bm{A}.
    \end{equation}
\end{approfondimento}

\subsubsection{Calcolo esplicito}
Si richiama innanzitutto la legge di Ampere-Laplace e la si riscrive sfruttando la \eqref{eq:help}:
\begin{equation}
    \bm{B} = -\frac{\mu_0 I}{4\pi} \int_{\gamma} d\bm{l}' \times \nabla \left( \frac{1}{|\bm{r}-\bm{r}'|} \right).
\end{equation}
Sfruttando una nota identità vettoriale che lega il prodotto vettoriale fra un vettore ($\bm{u} = d\bm{l}'$) e il gradiente di un campo scalare $\phi(\bm{r}) = \frac{1}{|\bm{r}-\bm{r}'|}$,
\begin{equation}
    \bm{u} \times \nabla \phi = \phi \cdot (\nabla \times \bm{u}) - \nabla \times (\phi \bm{u})
\end{equation}
l'equazione diventa:
\begin{equation}
    \bm{B} = -\frac{\mu_0 I}{4\pi} \int_{\gamma}\left[ \frac{\nabla \times d\bm{l}'}{|\bm{r}-\bm{r}'|}  - \nabla \times \left( \frac{d\bm{l}'}{|\bm{r}-\bm{r}'|}\right)\right].
\end{equation}
L'operatore $(\nabla \times)$ deriva rispetto alle coordinate del punto in cui si valuta il campo: la posizione dell'infinitesimo di filo è dunque una costante. Se ne deriva che $\nabla \times d\bm{l}' = \bm{0}$:
\begin{equation}
    \bm{B} = \frac{\mu_0 I}{4\pi} \int_{\gamma} \nabla \times \left( \frac{d\bm{l}'}{|\bm{r}-\bm{r}'|}\right).
\end{equation}
Come già discusso, l'integrale e l'operatore di rotore in questo caso specifico commutano:
\begin{equation}
    \bm{B} = \nabla \times \frac{\mu_0 I}{4\pi} \int_{\gamma}  \frac{d\bm{l}'}{|\bm{r}-\bm{r}'|} = \nabla \times \frac{\mu_0}{4\pi} \int_{V'} \frac{\bm{J}(\bm{r}')}{|\bm{r} - \bm{r}'|} dV',
    \label{eq:soluzione_A}
\end{equation}
dove si è sfruttato $Id\bm{l}' = JSd\bm{l} = \bm{J}dV'$ e si è cambiato il dominio di integrazione originale con il volume entro cui scorre corrente. 

Confrontando  quest'ultima espressione ottenuta con la \eqref{eq:def_potenziale_vettore}, si evince immediatamente che:
\begin{equation}
    \bm{A} = \frac{\mu_0 I}{4\pi} \int_{\gamma}  \frac{d\bm{l}'}{|\bm{r}-\bm{r}'|} = \frac{\mu_0}{4\pi} \int_{V'} \frac{\bm{J}(\bm{r}')}{|\bm{r} - \bm{r}'|} dV'.
    \end{equation}
    
    \subsubsection{L'Invarianza di Gauge e la Gauge di Coulomb}
La \eqref{eq:def_potenziale_vettore} non determina $\bm{A}$ univocamente. Infatti, se si trasforma $\bm{A}$ sommandolo al gradiente di una funzione scalare arbitraria $\phi(\bm{r})$ si ottiene
\begin{equation}
    \bm{A}' = \bm{A} + \nabla \phi
\end{equation}
e il campo magnetico associato ad $\bm{A}'$ rimane invariato:
\begin{equation}
    \nabla \times \bm{A}' = \nabla \times (\bm{A} + \nabla \phi) = \nabla \times \bm{A} + \nabla \times (\nabla \phi) = \bm{B} + \bm{0} = \bm{B}
\end{equation}
essendo nullo il rotore del gradiente. Questa proprietà è detta \textit{invarianza di gauge} e si può sfruttare per imporre una condizione comoda sulla divergenza di $\bm{A}$. In magnetostatica, la scelta standard è detta \textit{Gauge di Coulomb}:
\begin{equation}
    \nabla \cdot \bm{A}' = \nabla \cdot (\bm{A} + \nabla \phi) = 0.
    \label{eq:inv_gauge}
\end{equation}
Si cerca quindi un potenziale scalare che verifichi la \eqref{eq:inv_gauge}. Supponendo che $\nabla \cdot \bm{A}$ sia non nullo\footnote{Se fosse nullo, il potenziale scalare che rende vera la Gauge di Coulomb sarebbe banalmente la funzionane identicamente nulla.}, si ha: 
\begin{equation}
    \nabla \cdot \bm{A} = -\nabla^2   \phi,
\end{equation}
che è l'\textit{equazione di Poisson} e ammette soluzione per ogni $\bm{A}$. Esiste sempre un campo scalare $\phi$ che rende vera la \eqref{eq:inv_gauge} e quindi permette la scelta di Coulomb.

\subsection{L'equazione di Poisson}
Per calcolare la circuitazione di $\bm{B}$ lungo $\gamma$ chiusa, ci si avvale del Teorema di Stokes. Detta $\Omega$ una qualsiasi superficie \textit{aperta} che poggia su $\gamma$, vale:
\begin{equation}
    \oint_\gamma \bm{B}\cdot d\bm{l} = \int_{\Omega} (\nabla \times \bm{B}) \cdot d\bm{S}.
    \label{eq:circ_B}
\end{equation}
Dalla definizione di potenziale vettore, poi:
\begin{equation}
    \int_{\Omega} (\nabla \times \bm{B}) \cdot d\bm{S} = \int_{\Omega} [\nabla \times( \nabla \times\bm{A})] \cdot d\bm{S} = \int_{\Omega} [\nabla(\nabla \cdot \bm{A}) - \nabla^2\bm{A}] \cdot d\bm{S} = - \int_{\Omega} \nabla^2\bm{A} \cdot d\bm{S},
\end{equation}
dove si è usata la relazione $\nabla \times( \nabla \times\bm{A}) = \nabla(\nabla \cdot \bm{A}) - \nabla^2\bm{A}$ e la Gauge di Coulomb. Dalla definizione del potenziale vettore poi, sfruttando la commutatività dell'integrale con l'operatore di Laplace, quindi:
\begin{equation}
    \nabla^2\bm{A} = \nabla^2 \left(\frac{\mu_0}{4\pi} \int_{V'} \frac{\bm{J}(\bm{r}')}{|\bm{r} - \bm{r}'|} dV' \right)= \frac{\mu_0}{4\pi} \int_{V'} \nabla^2 \left(\frac{\bm{J}(\bm{r}')}{|\bm{r} - \bm{r}'|} \right) dV'.
\end{equation}
Forti del fatto che $\bm{J}(\bm{r}')$ si comporta come una costante se le viene applicato l'operatore di Laplace, si continua: 
\begin{equation}
    \nabla^2 \left(\frac{\bm{J}(\bm{r}')}{|\bm{r} - \bm{r}'|} \right) =  \bm{J}(\bm{r}') \nabla^2\left(\frac{1}{|\bm{r} - \bm{r}'|} \right) = - \bm{J}(\bm{r}') 4\pi \delta(\bm{r} - \bm{r}')
    \label{eq:lapl_pot_vett}
\end{equation}
e sostituendo nell'integrale si ottiene:
\begin{equation}
    \nabla^2\bm{A} =- {\mu_0} \int_{V'} \bm{J}(\bm{r}')  \delta(\bm{r} - \bm{r}') dV' = -\mu_0 \bm{J}(\bm{r})
\end{equation}
per le proprietà della \textit{delta di Dirac}. Si ottiene così l'\textit{equazione di Poisson per il magnetismo}\footnote{Questa equazione è formalmente identica all'equazione di Poisson dell'elettrostatica ($\nabla^2 V = -\rho/\epsilon_0$), con la differenza che qui si opera su vettori.}:
\begin{equation}
    \nabla^2\bm{A} = -\mu_0 \bm{J}(\bm{r}).
    \label{eq:poisson_m}
\end{equation}


\begin{approfondimento}[colback= white, colframe=green!35!black, no shadow, title={Calcolo del laplaciano del potenziale vettore}]
Si vuole ora esplicare a dovere come è stata valutata la \eqref{eq:lapl_pot_vett}. Sia ${R} = |\bm{r}-\bm{r}'|$, il campo scalare di cui valutare il laplaciano diventa quindi 
\begin{equation}
\chi = \frac{1}{{R}},
\end{equation}
singolare per ${R} = {0}$. Si applica quindi l'operatore di Laplace per ${R} \neq {0}$ e nella singolarità.

\subsubsection{$\bm{R} \neq \bm{0}$}
Fuori dalla singolarità, $\chi \in \mathbb{C}^{\infty}(\mathbb{R}^m)$. Non si hanno quindi particolari problemi nell'applicare il laplaciano. Data la simmetria del sistema, sfruttando le coordinate sferiche si ottiene:
\begin{equation}
\nabla^2 \chi  = \frac{1}{R^2} \frac{d}{dR}\left( R^2 \frac{d\chi}{dR} \right) = 0.
\end{equation}

\subsubsection{$\bm{R} = \bm{0}$}
Dato il comportamento singolare per ${R} = {0}$, per il calcolo del laplaciano si cambia approccio: integrando su un volume $V$ che circonda l'origine si ottiene
\begin{equation}
\int_V \nabla^2\left( \frac{1}{{R}}\right) dV = \int_V \nabla \cdot \nabla\left( \frac{1}{{R}}\right) dV.
\end{equation}
Vale il Teorema della divergenza, per cui:
\begin{equation}
\int_V \nabla \cdot \nabla\left( \frac{1}{{R}}\right) dV = \oint_{\partial V} \nabla\left( \frac{1}{{R}}\right) \cdot d\bm{S} = -\oint_{\partial V} \frac{1}{{R^2}}\bm{u}_R \cdot d\bm{S}.
\end{equation}
Essendo la scelta del volume di integrazione arbitraria, si impone $V$ come una sfera di raggio $R$ centrata nell'origine (e dunque nella singolarità del campo). Il versore radiale $\bm{u}_R$ risulta quindi parallelo al vettore superficie $d\bm{S}$:
\begin{equation}
\int_V \nabla \cdot \nabla\left( \frac{1}{{R}}\right) dV = -\oint_{\partial V} \frac{1}{{R^2}} d{S} = -4\pi,
\end{equation}
come dimostrato nel calcolo del flusso del campo elettrostatico. 

Il laplaciano del campo $\chi$ è dunque identicamente nullo $\forall R \neq 0$ mentre il suo integrale su un volume contenente l'origine è uguale a $-4\pi$. Questa descrizione risponde alle caratteristiche della funzione di distribuzione \textit{delta di Dirac}, infatti: 
\[
\delta(x - x_0) = 
\begin{cases}
0 &  x \neq x_0 \\
+\infty & x = x_0 
\end{cases}
\implies \int_{-\infty}^{\infty} \delta(x-x_0)f(x) dx = f(x_0).
\]
Questa definizione, declinata al caso in esame, suggerisce che l'applicazione dell'operatore ($\nabla^2$) a $\chi$ restituisce una delta di Dirac centrata nella singolarità del campo e moltiplicata per la funzione costante $f = -4\pi$:
\begin{equation}
\nabla^2\left(\frac{1}{R}\right) = -4\pi \delta(R) \implies \nabla^2\left(\frac{1}{|\bm{r}-\bm{r}'|}\right) = -4\pi\delta(\bm{r}-\bm{r}').
\end{equation}


\end{approfondimento}


\subsection{La legge di Ampere: forma differenziale}
Ricordando che $-\nabla^2\bm{A} = \nabla \times \bm{B}$, si ottiene la \textit{Legge di Ampere in forma locale}:
\begin{equation}
   \nabla \times \bm{B} = \mu_0 \bm{J}(\bm{r}).
   \label{eq:ampere_diff}
\end{equation}
Questo risultato è anche noto come \textit{quarta equazione di Maxwell in forma differenziale} (nel caso statico).

\subsection{La legge di Ampere: forma integrale}
Sostituendo la \eqref{eq:ampere_diff} nella \eqref{eq:circ_B} si ha:
\begin{equation}
    \oint_\gamma \bm{B}\cdot d\bm{l} = \int_{\Omega} (\mu_0 \bm{J}) \cdot d\bm{S} = \mu_0 \int_{\Omega} \bm{J} \cdot d\bm{S}.
\end{equation}
L'integrale a secondo membro rappresenta il flusso della densità di corrente attraverso la superficie $\Omega$ che ha come bordo la curva $\gamma$. Per definizione, questa è la corrente totale \textit{concatenata} ($I_{\text{conc}}$) con la curva.
Si ottiene così la \textit{forma integrale della Legge di Ampere} (o \textit{quarta equazione di Maxwell in forma integrale}):
\begin{equation}
    {
    \oint_\gamma \bm{B}\cdot d\bm{l} = \mu_0 I_{\text{conc}}
    }
\end{equation}
In conclusione, la circuitazione di $\bm{B}$ non dipende dalla forma della forma del filo conduttore, ma esclusivamente dalle correnti che in esso fluisce. 

\subsubsection{Significato fisico}
La circuitazione del campo magnetico rivela una proprietà fondamentale: mentre il campo elettrostatico è conservativo (circuitazione nulla), il campo magnetico non lo è. La legge di Ampere afferma che la circuitazione di B lungo una qualsiasi linea chiusa $\gamma$ è direttamente proporzionale alla corrente netta di conduzione che attraversa una superficie avente $\gamma$ come bordo. Questo significa che le sorgenti del campo magnetico sono le correnti elettriche e che le linee di campo di $\bm{B}$ formano curve chiusi attorno ad esse. 


\section{Campo magnetico, elettrico: un paradosso}

Finora si sono trattatati il campo elettrico $\bm{E}$ e il campo magnetico $\bm{B}$ come due entità distinte: il primo generato da cariche (statiche o in moto), il secondo solo da correnti (cariche in moto).
Tuttavia, l'analisi della Forza di Lorentz suggerisce un apparente paradosso che può essere risolto solo ammettendo che $\bm{E}$ e $\bm{B}$ siano due manifestazioni diverse di un'unica entità fisica: il \textit{Campo Elettromagnetico}.

\subsubsection{Il paradosso della carica in moto}
Si consideri una carica puntiforme $q$ che si muove con velocità costante $\bm{v}$ in una regione di spazio in cui (nel sistema di laboratorio $S$) esiste un campo magnetico uniforme $\bm{B}$ ed è assente il campo elettrico ($\bm{E}=0$).

Le forze agenti sulla carica appaiono differenti se analizzate da due punti di vista (sistemi di riferimento) diversi:

\begin{enumerate}
    \item {sistema di riferimento del Laboratorio ($S$):}
    L'osservatore vede la carica muoversi con velocità $\bm{v}$. Misura una forza magnetica (Forza di Lorentz) che agisce sulla carica:
    \begin{equation}
        \bm{F} = q \bm{v} \times \bm{B}.
    \end{equation}
    Questa forza è reale e misurabile: se $\bm{v} \perp \bm{B}$, l'osservatore vede la carica curvare (accelerazione centripeta);

    \item {sistema di riferimento solidale alla carica ($S'$):}
    Si immagini un osservatore che si muove insieme alla particella. In questo sistema $S'$, la particella è \textbf{ferma} ($\bm{v}' = 0$).
    Applicando ciecamente la formula della forza di Lorentz, si otterrebbe:
    \begin{equation}
        \bm{F}'_{S'} = q \bm{v}' \times \bm{B}' = q \cdot \bm{0} \times \bm{B}' = 0.
    \end{equation}
    Se la forza fosse nulla, la particella non dovrebbe subire accelerazione. Ma l'accelerazione è un evento fisico reale: se il percorso curva in $S$, deve curvare anche per l'osservatore in $S'$ (seppur descritto diversamente). Non può esistere un sistema inerziale in cui la particella curva e uno in cui va dritta.
\end{enumerate}
\textbf{La soluzione del paradosso:}
Affinché in $S'$ la particella subisca una forza (necessaria per giustificare l'accelerazione osservata) pur avendo velocità nulla, l'unica possibilità è che in $S'$ esista un \textbf{campo elettrico} $\bm{E}'$ che non esisteva in $S$.
La forza osservata in $S'$ è di natura puramente elettrica:
\begin{equation}
    \bm{F}' = q \bm{E}'.
\end{equation}

\subsubsection{Trasformazione dei campi}
Questo esperimento ideale dimostra qualcosa che verrà meglio trattato successivamente: la distinzione tra campo elettrico e magnetico non è assoluta, ma \textit{relativa al sistema di riferimento}.
Ciò che un osservatore percepisce come puro campo magnetico, un altro osservatore in moto rispetto al primo può percepirlo come una combinazione di campo magnetico e campo elettrico.

La Teoria della Relatività Ristretta fornisce le leggi di trasformazione corrette. Per velocità $v \ll c$ (limite galileiano), la relazione che lega i campi tra un sistema $S$ (dove c'è solo $\bm{B}$) e un sistema $S'$ in moto con velocità $\bm{v}$ è approssimativamente:
\begin{equation}
    \bm{E}' \approx \bm{v} \times \bm{B}.
\end{equation}
Sostituendo questa espressione nella forza in $S'$:
\begin{equation}
    \bm{F}' = q \bm{E}' = q (\bm{v} \times \bm{B}).
\end{equation}
Si ritrovano esattamente lo stesso valore della forza misurata in $S$. La fisica è salva: l'interazione è la stessa, cambia solo il "nome" che si da al campo responsabile (magnetico in $S$, elettrico in $S'$). In generale, per ovviare a problemi inerenti il sistema di riferimento, si fa riferimento alla \textit{forza di Lorentz} come quella che si presenta come la somma dei contributi magnetico e elettrico:
\begin{equation}
    \bm{F}_L = q (\bm{E} +\bm{v} \times \bm{B}).
\end{equation}

\begin{center} % Usa center, NON figure
    % Imposta l'angolo di visuale
    \tdplotsetmaincoords{70}{120} 
    
    \begin{tikzpicture}[scale=1.3, tdplot_main_coords, >=Stealth]
        
        % --- SISTEMA S (Laboratorio) ---
        \begin{scope}[xshift=-3cm]
            \node[font=\bfseries] at (-2.6, -0.2) {Sistema $S$ (Lab)};
            \node at (3, 2) {$\bm{v} \neq 0, \bm{B} \neq 0, \bm{E}=0$};
            
            % Campo B (uscente)
            \foreach \x in {-1, 0, 1} \foreach \y in {0, 1, 2} 
                \fill[blue!50] (\x, \y) circle (1.5pt);
            \node[blue, right] at (1.2, 2) {$\bm{B}$};

            % Carica e Velocità
            \shade[ball color=red] (0, 1) circle (0.2) node[above right] {$q$};
            \draw[->, thick, black] (0, 1) -- (1.5, 1) node[right] {$\bm{v}$};
            
            % Forza Magnetica (Verso il basso)
            \draw[->, ultra thick, red] (0, 1) -- (0, 0) node[above right] {$\bm{F}_m$};
            
        \end{scope}

      

        % --- SISTEMA S' (In moto) ---
        \begin{scope}[xshift=3cm]
            \node[font=\bfseries] at (-2.5, -0.5) {Sistema $S'$ (Solidale)};
            \node at (2, 3) {$\bm{v}'=0, \bm{E}' \neq 0$};
            
            % Campo B' (ancora presente, leggermente diverso in RG ma qui semplifichiamo)
            \foreach \x in {-1, 0, 1} \foreach \y in {0, 1, 2} 
                \fill[blue!20] (\x, \y) circle (1.5pt);

            % Campo Elettrico INDOTTO E'
            \foreach \x in {-1, 0, 1} 
                \draw[->, orange, thick] (\x, 2.2) -- (\x, -0.2);
            \node[orange, right] at (1.2, 1.5) {$\bm{E}'$};

            % Carica ferma
            \shade[ball color=red] (0, 1) circle (0.2) node[above right] {$q$};
     
            
            % Forza Elettrica (Verso il basso)
            \draw[->, ultra thick, red] (0, 1) -- (0, 0) node[above right] {$\bm{F}_e$};
           
        \end{scope}

    \end{tikzpicture}
    \captionof{figure}{Rappresentazione 3D del dipolo elettrico. Il potenziale in $P$ dipende dall'angolo polare $\theta$.}
    \label{fig:dipolo_3d}
\end{center}

\begin{approfondimento}[colback= white, colframe=green!35!black, no shadow, title={Unificazione Elettromagnetica}]
    Questa osservazione è la base concettuale che ha portato Einstein a formulare la Relatività Ristretta (il cui articolo originale del 1905 si intitola proprio \textit{"Sull'elettrodinamica dei corpi in movimento"}).
    Matematicamente, $\bm{E}$ e $\bm{B}$ non sono vettori indipendenti, ma componenti di un oggetto geometrico a 4 dimensioni chiamato \textbf{Tensore Elettromagnetico} $F_{\mu\nu}$. Cambiando sistema di riferimento, le componenti di questo tensore si mescolano tra loro, trasformando parte del campo magnetico in elettrico e viceversa.
\end{approfondimento}




\chapter{Sorgenti e campi magnetici notevoli}

\section{Il filo rettilineo percorso da corrente}
\label{sec:biot_savart}

Si consideri un filo rettilineo di lunghezza infinita percorso da una corrente stazionaria $I$. Si vuole determinare il campo magnetico in un punto $P$ posto a distanza $r$ dal filo.

Per ragioni di simmetria cilindrica:
\begin{itemize}
    \item il campo $\bm{B}$ non può avere componenti parallele al filo né componenti radiali (divergenti dal filo), come dedotto dall'esperimento di Oersted;
    \item il modulo del campo $B$ dipende esclusivamente dalla distanza radiale $r$;
    \item le linee di campo devono essere circonferenze concentriche al filo: i contributi al campo magnetico di ogni carica (e quindi di ciascun infinitesimo di lunghezza del filo) sono paralleli fra loro, nell'assunzione di un flusso di corrente stazionario con velocità dei portatori parallele fra loro e dunque si sommano costruttivamente\footnote{Si noti una differenza sostanziale rispetto al campo elettrostatico: per un filo carico le componenti parallele al filo generate da elementi simmetrici si elidono.}; \\
    \\
\textit{Dimostrazione}: La direzione del vettore infinitesimo $d\bm{B}$ è dettata dal prodotto vettoriale $d\bm{l}' \times \bm{u}_r$ (dove $d\bm{l}' = dl\,\bm{u}_t$).
Nel caso specifico di un filo \textbf{rettilineo}, l'intera retta su cui giace il filo e il punto di osservazione $P$ individuano un unico piano $\pi$. Ogni elemento di filo $d\bm{l}'$ e ogni vettore posizione $\bm{r}$ (che congiunge l'elemento a $P$) devono quindi giacere su tale piano $\pi$. Poiché il prodotto vettoriale di due vettori appartenenti a un piano è sempre perpendicolare al piano stesso, si deduce che \textbf{tutti} i contributi $d\bm{B}$ generati dai diversi tratti del filo sono perpendicolari a $\pi$.
Essendo tutti i vettori infinitesimi paralleli ed equiversi, il modulo del campo totale $|\bm{B}|$ è dato semplicemente dalla somma (integrale) dei moduli dei singoli contributi.
    

\end{itemize}
Applicando la legge di Ampere-Laplace e definendo $r = |\bm{r}-\bm{r}'|^2$:
\begin{equation}
    d\bm{B} = \frac{\mu_0 I}{4\pi} \frac{d\bm{l}' \times \bm{u}_r}{r^2}.
\end{equation}
In modulo, essendo $d\bm{l}' \perp \bm{u}_r$ solo nel punto di minima distanza, in generale vale $dl \sin\theta$, dove $\theta$ è l'angolo tra il filo e il vettore posizione. Definendo $\alpha = \frac{\pi}{2} - \theta$, il modulo del contributo infinitesimo è:
\begin{equation}
    dB = \frac{\mu_0 I}{4\pi} \frac{dx}{r^2} \cos\alpha.
\end{equation}
Esprimendo tutte le variabili in funzione dell'angolo $\alpha$:
\begin{itemize}
    \item distanza: $r = \frac{R}{\cos\alpha}$;
    \item posizione sul filo: $x = R \tan\alpha \implies dx = R \frac{d\alpha}{\cos^2\alpha}$.
\end{itemize}
Sostituendo nell'espressione differenziale:
\begin{equation}
    dB = \frac{\mu_0 I}{4\pi} \frac{R \frac{d\alpha}{\cos^2\alpha}}{(R/\cos\alpha)^2} \cos\alpha = \frac{\mu_0 I}{4\pi R} \cos\alpha \, d\alpha.
\end{equation}
Per un filo infinito, l'angolo $\alpha$ varia da $-\pi/2$ a $+\pi/2$. Integrando:
\begin{equation}
    B = \int_{-\pi/2}^{+\pi/2} \frac{\mu_0 I}{4\pi R} \cos\alpha \, d\alpha = \frac{\mu_0 I}{4\pi R} [\sin\alpha]_{-\pi/2}^{+\pi/2} = \frac{\mu_0 I}{4\pi R} (1 - (-1)).
\end{equation}
Si ottiene infine la \textit{legge di Biot-Savart}:
\begin{equation}
    B(R) = \frac{\mu_0 I}{2\pi R}.
    \label{eq:filo_infinito}
\end{equation}

\begin{figure}[htbp]
    \centering
    \begin{tikzpicture}[scale=1.5, >=Stealth]
        
        % --- 1. IL FILO E IL PUNTO P ---
        % Filo sull'asse X (per comodità di disegno 2D)
        \draw[ultra thick, gray] (-4,0) -- (4,0) node[right] {Filo $\infty$};
        \draw[->, ultra thick, orange] (-3.5,0) -- (-2.9,0) node[above left] {$I$};
        
        % Punto P a distanza R
        \coordinate (P) at (0, 2);
        \fill (P) circle (2pt) node[above right] {$P$};
        \draw[dashed] (0,0) -- (P) node[midway, right] {$R$};
        \draw (0, 0.2) -- (0.2, 0.2) -- (0.2, 0); % Angolo retto
        
        % --- 2. ELEMENTO dx ---
        \coordinate (dX) at (-2.5, 0);
        \draw[ultra thick, red] (-2.7, 0) -- (-2.3, 0) node[midway, below] {$d\bm{l}'$};
        
        % Vettore r (distanza)
        \draw[->, thick, blue] (dX) -- (P) node[midway, above left] {$r$};
        
        % --- 3. ANGOLI ---
        % Angolo alpha (rispetto alla normale)
        \draw[dashed, gray] (dX) -- (0,0) node[midway, below] {$x$};
        
        % Disegno l'angolo alpha in P
        \draw[thick, orange] (0, 1.5) arc (270:218:0.5);
        \node[orange] at (-0.2, 1.3) {$\alpha$};
        
        % Angolo Theta (classico tra dl e r)
        \draw[thick, gray] (-2, 0) arc (0:38:0.5);
        \node[gray] at (-1.8, 0.2) {$\theta$};

        % --- 4. CAMPO B (Risultato) ---
        \node[violet, right, align=left] at (0.5, 2) {$\bm{B}$ uscente};
        \draw[violet, thick] (0.4, 2) circle (0.15);
        \fill[violet] (0.4, 2) circle (0.05);

    \end{tikzpicture}
    \caption{Geometria per il calcolo del campo magnetico generato da un filo rettilineo. Tutte le componenti infinitesime $d\bm{B}$ generate dai tratti $dx$ sono perpendicolari al piano del foglio (uscenti) e si sommano algebricamente.}
    \label{fig:biot_savart_integrale}
\end{figure}

\subsection{La forza magnetica agente su un filo percorso da corrente}
Sia $\gamma$ un filo conduttore percorso da corrente $I$ che si trova in una regione di campo magnetico $\bm{B}$. Ogni carica in moto risente di una forza magnetica $\bm{F} = q\bm{v} \times \bm{B}$; analizzata una sezione infinitesima di filo lunga $dl$, di superficie $A$ con densità media di portatori $n$ costante, si ha:
\begin{equation}
d\bm{F} =  S dl (nq\bm{v}) \times \bm{B} = (JS) dl \bm{u}_t \times \bm{B} = I d\bm{l}\times \bm{B}.
\end{equation} 
Integrando si ottiene:
\begin{equation}
\bm{F} =  \int_\gamma I d\bm{l}\times \bm{B},
\end{equation} 
dove $d\bm{l}$ rappresenta l'infinitesimo di filo vettorializzato, orientato come $\bm{J}$.

\begin{figure}[htbp]
    \centering
    \begin{tikzpicture}[x={(1cm,0cm)}, y={(0.5cm,0.4cm)}, z={(0cm,1cm)}, scale=1.5, >=Stealth]

        % --- DEFINIZIONI ---
        \coordinate (W1_start) at (0,0,0);
        \coordinate (W1_end) at (0,0,4);
        \coordinate (W2_start) at (3,0,0);
        \coordinate (W2_end) at (3,0,4);
        
        % Punti a metà filo per i vettori
        \coordinate (Mid1) at (0,0,2);
        \coordinate (Mid2) at (3,0,2);

        % --- FILO 1 (Sorgente) ---
        \draw[ultra thick, gray] (W1_start) -- (W1_end) node[above] {Filo 1};
        \draw[->, ultra thick, orange] (0,0,0.5) -- (0,0,1.5) node[midway, left] {$I_1$};

        % --- FILO 2 (Subisce la forza) ---
        \draw[ultra thick, gray] (W2_start) -- (W2_end) node[above] {Filo 2};
        \draw[->, ultra thick, orange] (3,0,0.5) -- (3,0,1.5) node[midway, right] {$I_2$};

        % --- CAMPO B1 GENERATO DA I1 SU I2 ---
        % B1 ruota attorno a W1. A distanza W2, B1 è perpendicolare al piano W1-W2.
        % Se I1 sale, a destra (su W2) il campo entra nel foglio (o va "giù" nella prospettiva 3D scelta)
        % In questa prospettiva (x destra, z alto), y è profondità.
        % Campo B1 in Mid2 è diretto verso Y negativo (entrante)
        \draw[->, thick, blue] (Mid2) -- ++(0, -1.5, 0) node[below] {$\bm{B}_1$};
        
        % Linea di campo tratteggiata per far capire
        \draw[dashed, blue!40] (0,0,2) ellipse (3 and 1); % Ellisse prospettica

        % --- FORZA F SU FILO 2 ---
        % F = I2 (l) x B1.
        % l è verso Z (alto). B1 è verso -Y (dentro/dietro).
        % Z x (-Y) = X (verso sinistra, verso il Filo 1).
        % Attrazione!
        \draw[->, ultra thick, red] (Mid2) -- ++(-1.5, 0, 0) node[midway, above] {$\bm{F}_{1 \to 2}$};
        
        % Forza reciproca su Filo 1 (Opzionale)
        \draw[->, thick, red!70] (Mid1) -- ++(1.5, 0, 0) node[midway, above] {$\bm{F}_{2 \to 1}$};

        % --- DISTANZA R ---
        \draw[<->, thin, black] (0,0,0) -- (3,0,0) node[midway, below] {$R$};
        
        % --- NOTE ---
        \node[align=center, font=\small] at (1.5, -2, 0) {Correnti concordi $\to$ Attrazione};

    \end{tikzpicture}
    \caption{Interazione tra due fili paralleli percorsi da correnti concordi. Il campo $\bm{B}_1$ generato dal primo filo (blu) interagisce con la corrente $I_2$ generando una forza attrattiva $\bm{F}$ (rossa).}
    \label{fig:forza_fili_paralleli}
\end{figure}

\subsubsection{La forza magnetica agente fra due fili paralleli percorsi da corrente}
Nel caso in cui si abbiano due fili rettilinei $f_1$ e $f_2$ di lunghezza $L_1$ e $L_2$, posti nel vuoto a distanza $R \ll L_i$ l'uno dall'altro e percorsi rispettivamente dalle correnti $I_1$ e $I_2$, la forza indotta da $f_1$ su $f_2$\footnote{Uguale in modulo e opposta in verso rispetto a quella generata da $f_2$ su $f_1$ per il \textit{terzo principio} di Newton.} si valuta come:
\begin{equation}
\bm{F} =  \int_\gamma I_2 d\bm{l}\times \bm{B}_1.
\end{equation} 
Essendo la distanza reciproca fra i fili costante, il modulo del campo magnetico indotto su $f_2$ deve essere costante\footnote{Questa assunzione è giustificata dal fatto che, data la geometria del sistema filo, come esplorato nella sezione \ref{sec:biot_savart}, il modulo del campo prodotto varia solo in direzione radiale.}:
\begin{equation}
\bm{F} =  \int_\gamma I_2 d\bm{l}\times \bm{B}_1 = I_2 \left(\int_\gamma d\bm{l}\right)\times \bm{B}_1 = I_2 L \bm{u}_t \times \bm{B}.
\end{equation} 
Il verso del vettore forza risultante è dunque quello del prodotto vettoriale fra il vettore campo e il versore orientato come $\bm{J}$ (e quindi come il flusso di corrente). In generale, vale che:
\begin{itemize}
    \item se le correnti sono \textit{concordi}, i fili si \textit{attraggono};
    \item se le correnti sono \textit{discordi}, i fili si \textit{respingono}.
\end{itemize}

\subsubsection{La forza magnetica agente fra due fili paralleli percorsi da corrente stazionaria}
Si noti come nella trattazione precedente non sono state fatte assunzioni sulla forma matematica del campo magnetico, accontentandosi esclusivamente di informazioni di natura più geometrica. Ora, invece, si analizza un caso più specifico e semplice: due fili nelle stesse condizioni precedenti che rispondono alle ipotesi della legge di Biot-Savart.
In questa nuova configurazione (si faccia riferimento alla figura (\ref{fig:forza_fili_paralleli})), il filo $f_1$ genera un campo magnetico $\bm{B}_1$ che, alla distanza $R$ (dove si trova il filo $f_2$), ha modulo:
\begin{equation}
    B_1 = \frac{\mu_0 I_1}{2\pi R}.
\end{equation}
La direzione di $\bm{B}_1$ è perpendicolare al piano individuato dai due fili e, in generale, entrante o uscente a seconda del verso di $I_1$, in ossequio alla regola della mano destra. 
Poiché il filo $f_2$ è immerso in questo campo ed è perpendicolare alle linee di campo, così come la corrente che in esso fluisce ($\bm{u}_t \perp \bm{B}_1$), risente di una forza magnetica $\bm{F}_{1 \to 2}$:
\begin{equation}
    F = \int_\gamma I_2 \, dl \, B_1 \sin\left(\frac{\pi}{2}\right) = I_2 L B_1.
\end{equation}
Sostituendo l'espressione di $B_1$, si ottiene il modulo della forza di interazione:
\begin{equation}
    F = \frac{\mu_0 I_1 I_2 L}{2\pi R}.
    \label{eq:forza_ampere}
\end{equation}



\begin{approfondimento}[colback= white, colframe=green!35!black, no shadow, title={Definizione di Ampere e calcolo di \texorpdfstring{$\bm{\mu_0}$}{mu zero}}]

\subsubsection{Definizione di Ampere (\bm{$=[A]$})}
    Questa relazione è stata utilizzata storicamente per definire l'unità di misura della corrente nel SI:
    \textit{"L'Ampere è quella corrente costante che, mantenuta in due conduttori paralleli rettilinei di lunghezza infinita posti a un metro di distanza nel vuoto, produce tra questi conduttori una forza uguale a $2 \times 10^{-7}$ Newton per ogni metro di lunghezza."}\\
    \subsubsection{Calcolo di \bm{$\mu_0$}}
   Dalla definizione di Ampere, 
\begin{equation}
    \frac{F}{L} = \frac{\SI{2e7}{\newton}}{\SI{1}{\meter}} = \frac{\mu_0 I_1 I_2}{2\pi R} = \frac{\mu_0 \cdot \SI{1}{\ampere} \cdot \SI{1}{\ampere}}{2\pi \cdot \SI{1}{\meter}}.
\end{equation}
\begin{equation}
   \implies \mu_0 = 4\pi \times 10^{-7} \, \si{\newton\per\ampere\squared}.
\end{equation}
\end{approfondimento}

\section{Il solenoide infinito}

\section{Il toroide}

\chapter{Campi magnetici nella materia}

Lo studio dell'elettromagnetismo condotto finora si è basato sull'ipotesi semplificativa che le sorgenti di campo (correnti e cariche) fossero poste nel vuoto.
Tuttavia, la materia non è inerte rispetto ai campi magnetici. Gli atomi e le molecole sono costituiti da cariche elettriche in moto (elettroni orbitali e spin) che generano campi magnetici microscopici e risentono di azioni meccaniche se immersi in un campo esterno.

In questo capitolo si analizzerà come la materia reagisce alla presenza di un campo magnetico esterno (\textit{magnetizzazione}) e come questa risposta modifichi il campo stesso, estendendo la validità delle equazioni di Maxwell ai mezzi materiali.

\section{Il momento magnetico}

Il punto di partenza per la descrizione microscopica del magnetismo non può essere la carica magnetica (o monopolo), la cui esistenza non è mai stata verificata sperimentalmente e la cui assenza è sancita dalla seconda equazione di Maxwell ($\nabla \cdot \bm{B} = 0$).
L'ente fondamentale del magnetismo è invece il \textit{dipolo magnetico}, idealizzabile fisicamente come una spira percorsa da corrente.

\subsection{Definizione formale}
Si consideri una spira piana di forma arbitraria, racchiudente una superficie di area $S$, percorsa da una corrente stazionaria $I$.
Si definisce \textit{momento di dipolo magnetico} (o semplicemente momento magnetico) il vettore $\bm{m}$ tale che:

\begin{equation}
    \bm{m} = I\bm{S} = I S \bm{u}_n \quad [\text{A} \cdot \text{m}^2],
    \label{eq:momento_magnetico}
\end{equation}

dove:
\begin{itemize}
    \item $I$ è l'intensità della corrente che circola nella spira;
    \item $S$ è l'area della superficie piana delimitata dalla spira;
    \item $\bm{u}_n$ è il versore normale alla superficie, il cui verso è determinato dalla \textit{regola della mano destra} rispetto al verso di percorrenza della corrente.
\end{itemize}
\begin{figure}[htbp]
    \centering
    \begin{tikzpicture}[scale=1.2, >=Stealth]
        % Coordinate 3D simulate
        \coordinate (O) at (0,0);
        
        % Disegno della spira (ellisse prospettica)
        \draw[thick, gray] (0,0) ellipse (2 and 0.8);
        
        % Freccia corrente (davanti)
        \draw[->, ultra thick, orange] (0, -0.8) -- (0.1, -0.8) node[below] {$I$};
        % Freccia corrente (dietro - opzionale per chiarezza)
        %\draw[->, thick, orange!50] (0, 0.8) -- (-0.1, 0.8);
        
        % Superficie S
        \fill[blue!5, opacity=0.5] (0,0) ellipse (2 and 0.8);
        \node[blue!60] at (-1.2, 0) {$S$};
        
        % Vettore normale n / momento m
        \draw[->, ultra thick, blue] (O) -- (0, 2) node[above] {$\bm{m} = I S \bm{u}_n$};
        
        % Regola mano destra (piccolo arco)
        \draw[->, thin, gray] (0.4, 0.2) arc (10:170:0.4 and 0.1);
        
    \end{tikzpicture}
    \caption{Definizione geometrica del momento magnetico $\bm{m}$ associato a una spira piana percorsa da corrente $I$.}
    \label{fig:def_momento_m}
\end{figure}
Nel caso limite in cui le dimensioni della spira tendono a zero ($S \to 0$) ma il prodotto $IS$ rimane finito, si parla di \textit{dipolo magnetico puntiforme}. Questa astrazione è particolarmente utile per modellizzare gli elettroni atomici, che possiedono un momento magnetico intrinseco (spin) e orbitale pur essendo particelle elementari.

\subsection{Azioni meccaniche in un campo uniforme}
Per comprendere l'origine della magnetizzazione nei materiali, è necessario analizzare le forze agenti su un dipolo immerso in un campo magnetico esterno $\bm{B}$.
Si supponga che il campo $\bm{B}$ sia \textit{uniforme} su tutta l'estensione della spira $\gamma$. In tal caso:
\begin{enumerate}
    \item \textit{la forza risultante è nulla:}
    \begin{equation}
        \bm{F} = I \oint_{\gamma} d\bm{l} \times \bm{B} = I \left( \oint_{\gamma} d\bm{l} \right) \times \bm{B} = \bm{0} \times \bm{B} = \bm{0}.
    \end{equation}
    Essendo il campo costante, esce dall'integrale; l'integrale di linea di un percorso chiuso è vettorialmente nullo. Il dipolo non trasla;

    \item \textit{il momento torcente non è nullo:}
    le forze agenti sui vari tratti della spira generano una coppia. Detta $dq$ una quantità infinitesima di carica in moto nella spira, il momento torcendo generato dalla forza magnetica su di essa agente è pari a:
        \begin{equation}
        d\bm{\tau} = \bm{r} \times d\bm{F} = \bm{r} \times (dq \,\bm{v} \times \bm{B}) = \bm{r} \times (dq \,\frac{d\bm{l}}{dt} \times \bm{B}) = \bm{r} \times (I \,d\bm{l} \times \bm{B}),
    \end{equation}
    con $\bm{r}$ il braccio della forza rispetto al centro di massa della spira. Integrando su $\gamma$ si ottiene:
            \begin{equation}
        \bm{\tau} = I\oint_{\gamma} \bm{r} \times (I \,d\bm{l} \times \bm{B}) = I \oint [d\bm{l}(\bm{r}\cdot \bm{B}) - \bm{B}(\bm{r}\cdot d\bm{l})] = I\left[ \oint_{\gamma}d\bm{l}(\bm{r}\cdot \bm{B}) - \oint_{\gamma} \bm{B}(\bm{r}\cdot d\bm{l})\right].
    \end{equation}
    Si può dimostrare che il primo integrale equivale a $\bm{S} \times \bm{B}$, con $\bm{S}$ il vettore superficie associato alla spira. Per risolvere il secondo si sfrutta il fatto che $\bm{B}$ è costante e si applica il Teorema di Stokes:
                \begin{equation}
         \oint_{\gamma} \bm{B}(\bm{r}\cdot d\bm{l}) = \bm{B}\oint_{\gamma} (\nabla \times \bm{r})\cdot d\bm{S} = \bm{0},
    \end{equation}
    essendo nullo il rotore di un qualsiasi vettore. Si ottiene pertanto
    \begin{equation}
        \bm{\tau} = I \bm{S} \times \bm{B} = \bm{m} \times \bm{B}.
    \end{equation}
    L'effetto di tale momento è quello di far ruotare il dipolo fino ad allineare il vettore $\bm{m}$ parallelamente al vettore campo $\bm{B}$.
\end{enumerate}

\subsection{Energia potenziale magnetica}
Il lavoro compiuto dalle forze del campo per ruotare il dipolo è associato a una variazione di energia potenziale. L'energia potenziale $U_m$ di un dipolo in un campo esterno è data da:
\begin{equation}
    U_m = - \bm{m} \cdot \bm{B} = - m B \cos\theta,
\end{equation}
dove $\theta$ è l'angolo compreso tra $\bm{m}$ e $\bm{B}$.
Dall'analisi energetica si deduce che:
\begin{itemize}
    \item \textit{Equilibrio stabile ($\theta = 0$):} L'energia è minima ($U_{min} = -mB$) quando il dipolo è parallelo ed equiverso al campo.
    \item \textit{Equilibrio instabile ($\theta = \pi$):} L'energia è massima ($U_{max} = +mB$) quando il dipolo è antiparallelo al campo.
\end{itemize}

\section{La magnetizzazione della materia}

In un campione macroscopico di materiale contenente multipli del numero di Avogadro di atomi, l'effetto cumulativo dei momenti magnetici microscopici diventa rilevante.

L'agitazione termica, in assenza di campi esterni, fa sì che i momenti magnetici atomici $\bm{m}_i$ siano orientati casualmente in tutte le direzioni spaziali. La somma vettoriale su un volume macroscopico è pertanto \textit{statisticamente nulla}: il materiale non manifesta proprietà magnetiche.
Quando però il campione viene immerso in un campo magnetico esterno $\bm{B}$, il momento torcente $\bm{\tau}_i = \bm{m}_i \times \bm{B}$ agente su ogni dipolo atomico, tende ad allinearne il momento magnetico parallelamente alle linee di campo, mirando alla minima energia potenziale.
Sebbene l'agitazione termica contrasti questo ordinamento, si ottiene un risultato netto diverso da zero: il materiale si è \textit{magnetizzato}.

\subsection{Il vettore Magnetizzazione \texorpdfstring{$\bm{M}$}{M}}
Per descrivere quantitativamente questo stato di polarizzazione magnetica, si introduce un nuovo campo vettoriale macroscopico, il vettore \textbf{Magnetizzazione} $\bm{M}(\bm{r})$, definito come il momento di dipolo magnetico netto per unità di volume\footnote{Simmetricamente al vettore di polarizzazione elettrica.}.

Considerato un volume elementare $\Delta V$ (sufficientemente piccolo a livello macroscopico, ma contenente un numero statisticamente rilevante di atomi) centrato in $\bm{r}$:
\begin{equation}
    \bm{M}(\bm{r}) = \lim_{\Delta V \to 0} \frac{\sum_{i} \bm{m}_i}{\Delta V} \quad \left[ \frac{\text{A}}{\text{m}} \right].
\end{equation}
Il vettore $\bm{M}$ riveste nel magnetismo lo stesso ruolo che il vettore polarizzazione $\bm{P}$ gioca nell'elettrostatica dei dielettrici.

\subsection{Le correnti di magnetizzazione (o Amperiane)}
La presenza di una magnetizzazione $\bm{M}$ nel materiale ha un effetto fisico tangibile: la comparsa di correnti elettriche macroscopiche, dette \textit{correnti di magnetizzazione} (o correnti legate/amperiane).
È cruciale comprendere che queste correnti non sono dovute a un trasporto di cariche libere (come negli elettroni di conduzione in un filo), ma sono l'effetto collettivo delle micro-correnti atomiche stazionarie.

\subsubsection{Origine fisica: l'elisione interna}
Si immagini di modellizzare gli atomi del materiale come piccole spire percorse da corrente (\textit{spire amperiane}), tutte allineate (o parzialmente allineate) al campo esterno:
\begin{figure}[htbp]
    \centering
    \begin{tikzpicture}[scale=1.2, >=Stealth]
        % Blocco di materiale
        \draw[thick, fill=gray!5] (0,0) rectangle (5,3);
        \node[above, gray] at (2.5, 3) {Materiale Magnetizzato Uniformemente};

        % Griglia di spire microscopiche
        \foreach \x in {0.8, 1.9, 3.0, 4.1} {
            \foreach \y in {0.8, 1.9} { % Due righe
                % Cerchio spira
                \draw[blue, thin] (\x, \y) circle (0.45);
                % Frecce sulla spira (senso antiorario)
                \draw[->, blue, thick] (\x+0.45, \y) -- (\x+0.45, \y+0.1); 
                \draw[->, blue, thick] (\x-0.45, \y) -- (\x-0.45, \y-0.1);
                \draw[->, blue, thick] (\x, \y+0.45) -- (\x-0.1, \y+0.45);
                \draw[->, blue, thick] (\x, \y-0.45) -- (\x+0.1, \y-0.45);
            }
        }

        % Focus sull'elisione interna (tra due spire)
        \draw[red, dashed] (2.45, 1.4) rectangle (3.55, 2.4);
        \node[red, font=\footnotesize, align=center] at (2.5, 1.36) {$\bm{J}_{net} = \bm{0}$};
        \node[red, font=\tiny, align=center] at (0.85, 1.36) {(Elisione)};

        % Frecce risultanti sulla superficie
        \draw[ultra thick, orange, ->] (0.2, 0.2) -- (4.8, 0.2);
        \draw[ultra thick, orange, ->] (4.8, 0.2) -- (4.8, 2.8);
        \draw[ultra thick, orange, ->] (4.8, 2.8) -- (0.2, 2.8);
        \draw[ultra thick, orange, ->] (0.2, 2.8) -- (0.2, 0.2);

        \node[orange, right] at (5, 1.5) {$\bm{K}_m$, Corrente superficiale};

    \end{tikzpicture}
    \caption{Sezione del conduttore, modello delle correnti amperiane. All'interno del volume, nel punto di contatto tra due spire adiacenti, le correnti microscopiche scorrono in verso opposto elidendosi a vicenda. Sulla superficie esterna non vi è elisione, risultando in una corrente netta $\bm{K}_m$ che circonda il materiale.}
    \label{fig:correnti_amperiane}
\end{figure}
\begin{itemize}
    \item {Nel volume:} nel punto di contatto tra due spire atomiche adiacenti, le correnti scorrono in verso opposto. Se la magnetizzazione è uniforme, le intensità sono identiche e si cancellano perfettamente. La corrente netta interna è nulla $\bm{J}_{net} = \bm{0}$;
    \item {sulla superficie:} le spire che si trovano sul bordo del materiale non hanno "vicini" con cui cancellarsi verso l'esterno. La somma di tutti i tratti di corrente superficiali non elisi costituisce una corrente macroscopica che fluisce sulle pareti esterne del materiale.
\end{itemize}

\subsubsection{Relazioni differenziali}
Formalmente, le densità di corrente di magnetizzazione sono legate alle derivate spaziali del vettore $\bm{M}$.

\begin{enumerate}
    \item \textbf{Densità di corrente volumetrica ($\bm{J}_m$):}
    se la magnetizzazione $\bm{M}$ \textit{non è uniforme} (ovvero varia da punto a punto all'interno del materiale), l'elisione delle correnti interne non è perfetta. Si genera una densità di corrente nel volume del materiale data dal rotore della magnetizzazione:
    \begin{equation}
        \bm{J}_m = \nabla \times \bm{M};
        \label{eq:J_mag}
    \end{equation}

    \item \textbf{densità di corrente superficiale ($\bm{K}_m$):}
    sulla superficie di discontinuità tra il materiale e il vuoto (o tra due materiali diversi), compare una densità di corrente superficiale data da:
    \begin{equation}
        \bm{K}_m = \bm{M} \times \bm{u}_n,
        \label{eq:K_mag}
    \end{equation}
    dove $\bm{u}_n$ è il versore normale alla superficie uscente dal materiale.
\end{enumerate}
Queste correnti $\bm{J}_m$ e $\bm{K}_m$ (dove  $\bm{K}_m$ non è altro che $\bm{J}_m$ vista dal bordo del materiale) sono sorgenti reali di campo magnetico $\bm{B}$ esattamente come le correnti di conduzione nei fili. Il campo totale sarà dunque la sovrapposizione del campo esterno e del campo prodotto da queste correnti legate.


\begin{approfondimento}[colback= white, colframe=green!35!black, no shadow, title={Derivazione formale di \texorpdfstring{$\bm{K}_m$}{Km}}]
    Le correnti superficiali $\bm{K}_m$ non sono entità distinte da quelle di volume, ma ne rappresentano il caso limite sulla superficie.
    
    Matematicamente, si consideri la densità di corrente di magnetizzazione generale $\bm{J}_m = \nabla \times \bm{M}$.
    Sulla superficie di separazione tra il materiale (dove $\bm{M} \neq 0$) e il vuoto (dove $\bm{M} = 0$), il vettore $\bm{M}$ subisce una discontinuità "a gradino".
    
    Applicando il teorema del rotore a un percorso rettangolare infinitesimo $\gamma$ a cavallo della superficie:
    \begin{equation}
        \oint_\gamma \bm{M} \cdot d\bm{l} = \int_S (\nabla \times \bm{M}) \cdot d\bm{S} = I_{m, \text{concatenata}}.
    \end{equation}
    Se l'altezza del rettangolo tende a zero, il flusso del rotore (che è infinito sulla discontinuità) diventa finito e corrisponde alla corrente superficiale che attraversa il rettangolo.
    Si ottiene così che la densità di corrente superficiale è data dalla componente tangenziale di $\bm{M}$ ruotata di $90^\circ$:
    \begin{equation}
        \bm{K}_m = \bm{M} \times \bm{u}_n.
    \end{equation}
    \textbf{Fisicamente:} Un magnete uniforme ha $\nabla \times \bm{M} = 0$ all'interno (nessuna corrente di volume), ma tutto il suo campo magnetico è generato proprio dalla corrente $\bm{K}_m$ che scorre sulla sua "pelle".
\end{approfondimento}


\section{Il campo magnetico \texorpdfstring{$\bm{H}$}{H}}

Per classificare i materiali è necessario distinguere il contributo al campo magnetico dovuto alle correnti "libere" (quelle che scorrono nei circuiti e che possiamo controllare) da quello dovuto alla risposta del materiale (correnti di magnetizzazione). Dalla legge di Ampere, sapendo che le sorgenti di $\bm{B}$ sono sia le correnti libere ($\bm{J}_{lib}$) che quelle di magnetizzazione ($\bm{J}_m = \nabla \times \bm{M}$), si ha:
\begin{equation}
    \nabla \times \bm{B} = \mu_0 (\bm{J}_{lib} + \bm{J}_m) = \mu_0 \bm{J}_{lib} + \mu_0 (\nabla \times \bm{M}).
\end{equation}
Portando i termini contenenti i rotori a primo membro:
\begin{equation}
    \nabla \times \left( \frac{\bm{B}}{\mu_0} - \bm{M} \right) = \bm{J}_{lib}.
\end{equation}
Si definisce così il vettore \textbf{Campo Magnetizzante} (o campo magnetico ausiliario) $\bm{H}$:
\begin{equation}
    \bm{H} = \frac{\bm{B}}{\mu_0} - \bm{M} \quad \left[ \frac{\text{A}}{\text{m}} \right].
    \label{eq:def_H}
\end{equation}
La proprietà fondamentale di $\bm{H}$ è che la sua circuitazione dipende \textit{esclusivamente} dalle correnti libere concatenate, indipendentemente dal materiale presente:
\begin{equation}
    \oint_\gamma \bm{H} \cdot d\bm{l} = I_{lib, conc}.
\end{equation}



\subsection{La suscettività magnetica}
Per una vasta classe di materiali (detti \textit{lineari} e \textit{isotropi}), la magnetizzazione indotta $\bm{M}$ è parallela e proporzionale al campo magnetizzante $\bm{H}$, così come al campo magnetico esterno $\bm{B}$ e vale che:
\begin{equation}
    \bm{M} = \chi_m \bm{H}.
\end{equation}
Il coefficiente scalare adimensionale $\chi_m$ è detto \textit{suscettività magnetica}\footnote{Simmetrico rispetto alla suscettibilità elettrica $\chi_e$.}.
Sostituendo nella definizione \eqref{eq:def_H}, si ottiene la relazione costitutiva:
\begin{equation}
    \bm{B} = \mu_0 (\bm{H} + \chi_m \bm{H}) = \mu_0 (1 + \chi_m) \bm{H} = \mu_0 \mu_r \bm{H} = \mu \bm{H},
\end{equation}
dove $\mu_r$ è la permeabilità magnetica relativa del materiale mentre $\mu$ quella assoluta.
In base al segno e al valore di $\chi_m$, si distinguono tre comportamenti fisici fondamentali.

\section{Diamagnetismo (\texorpdfstring{$\chi_m < 0$}{chi < 0})}

Il diamagnetismo è una proprietà intrinseca di tutta la materia, legata alla variazione del moto orbitale degli elettroni indotta dal campo esterno.

\subsection{Meccanismo microscopico (Precessione di Larmor)}
Classicamente, si può immaginare un elettrone che orbita attorno al nucleo come una micro-spira. Quando si applica un campo magnetico $\bm{B}$, il flusso attraverso l'orbita varia. Per la legge di Lenz (applicata a livello atomico), si induce una variazione di velocità dell'elettrone tale da generare un momento magnetico indotto $\Delta \bm{m}$ che si \textbf{oppone} al campo esterno.

\subsection{Proprietà ed Esempi}
Poiché il momento indotto è opposto al campo ($\bm{M}$ antiparallelo a $\bm{H}$), la suscettività è negativa:
\begin{equation}
    \chi_m < 0 \quad (\text{ordine di } -10^{-5}).
\end{equation}
Un materiale diamagnetico viene debolmente \textit{respinto} dalle regioni a campo magnetico intenso.
\begin{itemize}
    \item {Esempi tipici:} Rame ($Cu$), Oro ($Au$), Argento ($Ag$), Acqua ($H_2O$), Bismuto ($Bi$).
    \item {Caso notevole:} I \textit{superconduttori} si comportano come diamagneti perfetti ($\chi_m = -1$, $\bm{B}=0$ all'interno), fenomeno noto come \textit{Effetto Meissner}.
\end{itemize}

\section{Paramagnetismo (\texorpdfstring{$\chi_m > 0$}{chi > 0})}

Il paramagnetismo si manifesta nei materiali i cui atomi o molecole possiedono un momento magnetico intrinseco \textit{permanente} (solitamente dovuto alla presenza di \textit{elettroni spaiati negli orbitali esterni}).

\subsection{Meccanismo microscopico}
In assenza di campi esterni, l'agitazione termica fa sì che i momenti $\bm{m}$ siano orientati a caso ($\bm{M}=0$).
Applicando un campo esterno, si instaura una competizione tra:
\begin{enumerate}
    \item \textbf{Energia magnetica ($-\bm{m} \cdot \bm{B}$):} tende ad allineare i dipoli paralleli al campo (ordine).
    \item \textbf{Energia termica ($\propto k_B T$):} tende a disordinare i dipoli.
\end{enumerate}
Il risultato è un allineamento parziale statistico, con $\bm{M}$ parallelo a $\bm{H}$.

\subsection{Proprietà ed Esempi}
La suscettività è positiva ma generalmente piccola:
\begin{equation}
    \chi_m > 0 \quad (\text{ordine di } 10^{-5} - 10^{-3}).
\end{equation}
Un materiale paramagnetico viene debolmente \textit{attratto} verso le regioni a campo intenso.
La dipendenza dalla temperatura è regolata dalla \textit{Legge di Curie}: la suscettività è inversamente proporzionale alla temperatura assoluta (a temperature maggiori, il disordine termico prevale):
\begin{equation}
    \chi_m \propto \frac{1}{T}.
\end{equation}
\begin{itemize}
    \item {Esempi tipici:} Alluminio ($Al$), Platino ($Pt$), Tungsteno ($W$), Ossigeno liquido.
\end{itemize}

\section{Ferromagnetismo (\texorpdfstring{$\chi_m \gg 0$}{chi >> 0})}

I materiali ferromagnetici (come il Ferro, il Nichel, il Cobalto e le loro leghe) presentano proprietà magnetiche macroscopiche estremamente intense e complesse, in generale accomunate dalla presenza di \textit{elettroni spaiati nel guscio di valenza}. In questi materiali, la relazione $\bm{M} = \chi_m \bm{H}$ perde la sua validità lineare: $\bm{M}$ non è proporzionale ad $\bm{H}$ e dipende dalla storia precedente del campione (\textit{isteresi}).

\subsection{Domini di Weiss}
La caratteristica distintiva è l'esistenza di una forte interazione quantistica tra gli spin degli elettroni vicini, detta \textit{interazione di scambio}, che li costringe ad allinearsi parallelamente tra loro anche in assenza di campo esterno.
Si formano così regioni macroscopiche ($10^{-6}-10^{-4}$ m), dette \textit{Domini di Weiss}, all'interno delle quali la magnetizzazione è massima (saturazione).

\begin{itemize}
    \item {Materiale non magnetizzato:} i domini sono orientati casualmente; la somma vettoriale è nulla;
    \item {materiale magnetizzato:} sotto l'azione di un campo esterno, i domini orientati favorevolmente si espandono a scapito degli altri, imponendo una \textit{magnetizzazione permanente}.
\end{itemize}

\subsection{Ciclo di Isteresi}
La magnetizzazione di un ferromagnete non è un processo reversibile. Riportando su un grafico l'andamento del campo magnetico indotto $B$ in funzione del campo applicato $H$ (ciclicamente), si ottiene una curva chiusa detta ciclo di isteresi.

\begin{figure}[htbp]
    \centering
    \begin{tikzpicture}[scale=1.1, >=Stealth]
        % Assi
        \draw[->] (-3,0) -- (3,0) node[right] {$H$ [A/m]};
        \draw[->] (0,-2.5) -- (0,2.5) node[above] {$B$ [T]};
        
        % Curva di prima magnetizzazione (dashed)
        \draw[dashed, thick, gray] (0,0) .. controls (0.5,0.2) and (1,1.5) .. (2.5,1.8) node[right] {\footnotesize Prima Magnet.};
        
        % Ciclo Isteresi
        \draw[ultra thick, red, smooth] plot coordinates {
            (2.5, 1.8)     % Saturazione +
            (1.0, 1.5)
            (0, 0.9)       % Br
            (-1.1, 0)      % Hc
            (-2.5, -1.8)   % Saturazione -
            (-1.0, -1.5)
            (0, -0.9)      % -Br
            (1.1, 0)       % +Hc
            (2.5, 1.8)     % Chiusura
        };
        
        % Punti notevoli
        \fill[blue] (0, 0.9) circle (2pt) node[anchor=south west] {$B_r$};
        \node[blue, right, font=\footnotesize] at (0.1, 0.9) {Remanenza};
        
        \fill[blue] (-1.1, 0) circle (2pt) node[anchor=south east] {$-H_c$};
        \node[blue, left, font=\footnotesize] at (-1.1, -0.3) {Campo Coercitivo};

        % Frecce di direzione
        \draw[->, thick, red] (1.5, 1.6) -- (1.4, 1.58);
        \draw[->, thick, red] (-1.5, -1.6) -- (-1.4, -1.58);
        
    \end{tikzpicture}
    \caption{Ciclo di isteresi tipico di un materiale ferromagnetico. Si noti che per $H=0$ il campo $B$ non è nullo ($B_r$), trasformando il materiale in un magnete permanente.}
    \label{fig:isteresi}
\end{figure}

Definizioni chiave dal grafico:
\begin{enumerate}
    \item \textbf{Remanenza ($B_r$):} È il valore del campo magnetico che residua nel materiale quando il campo esterno $H$ viene spento. È ciò che rende possibile l'esistenza dei \textit{magneti permanenti}.
    \item \textbf{Campo Coercitivo ($H_c$):} È l'intensità del campo inverso che bisogna applicare per annullare la magnetizzazione residua (smagnetizzare il materiale).
\end{enumerate}

\subsection{Temperatura di Curie}
L'ordine ferromagnetico è una competizione tra l'interazione di scambio (ordinante) e l'agitazione termica (disordinante). Esiste una temperatura critica specifica per ogni materiale, detta \textit{Temperatura di Curie} ($T_C$), al di sopra della quale l'agitazione termica distrugge l'allineamento spontaneo dei domini.
Per $T > T_C$, il materiale perde le proprietà ferromagnetiche e diventa un semplice paramagnete.
\begin{equation}
    \chi_m \propto \frac{1}{T - T_C}.
\end{equation}


\chapter{Le equazioni di Maxwell: il caso statico}

In questo capitolo si sistematizzano i risultati ottenuti finora per l'elettrostatica e la magnetostatica.
Le leggi fondamentali che governano i campi elettrico $\bm{E}$ e magnetico $\bm{B}$ in condizioni stazionarie (ovvero quando le sorgenti $\rho$ e $\bm{J}$ sono costanti nel tempo e, di conseguenza, $\frac{\partial \bm{E}}{\partial t} = 0$ e $\frac{\partial \bm{B}}{\partial t} = 0$) costituiscono il nucleo delle \textit{Equazioni di Maxwell nel caso statico}.

\section{Quadro sinottico}

In regime stazionario, si osserva un perfetto \textit{disaccoppiamento} tra i fenomeni elettrici e magnetici:
\begin{itemize}
    \item Le sorgenti del campo elettrico sono le cariche statiche ($\rho$). Il campo è conservativo.
    \item Le sorgenti del campo magnetico sono le correnti stazionarie ($\bm{J}$). Il campo è solenoidale.
\end{itemize}
Di seguito sono riportate le quattro equazioni fondamentali sia nella loro forma locale (differenziale), utile per descrivere il campo punto per punto, sia nella forma globale (integrale), che lega i campi alle sorgenti contenute in un volume o concatenate a un percorso.

\begin{table}[htbp]
    \centering
    \renewcommand{\arraystretch}{2.2} % Aumenta lo spazio verticale per le formule
    \begin{tabular}{p{4cm} c c}
        \hline
        \textbf{Legge fisica} & \textbf{Forma Differenziale} & \textbf{Forma Integrale} \\
        \hline
        
        \textbf{1. Legge di Gauss} \newline (Elettrica) & 
        $\nabla \cdot \bm{E} = \dfrac{\rho}{\varepsilon_0}$ & 
        $\displaystyle \oint_{\Sigma} \bm{E} \cdot d\bm{S} = \frac{Q_{int}}{\varepsilon_0}$ \\
        
        \textbf{2. Legge di Gauss} \newline (Magnetica) & 
        $\nabla \cdot \bm{B} = 0$ & 
        $\displaystyle \oint_{\Sigma} \bm{B} \cdot d\bm{S} = 0$ \\
        
        \textbf{3. Legge di Faraday-Neumann} \newline (Caso statico) & 
        $\nabla \times \bm{E} = \bm{0}$ & 
        $\displaystyle \oint_{\gamma} \bm{E} \cdot d\bm{l} = 0$ \\
        
        \textbf{4. Legge di Ampere} \newline (Senza spostamento) & 
        $\nabla \times \bm{B} = \mu_0 \bm{J}$ & 
        $\displaystyle \oint_{\gamma} \bm{B} \cdot d\bm{l} = \mu_0 I_{conc}$ \\
        \hline
    \end{tabular}
    \caption{Le equazioni di Maxwell nel vuoto per campi stazionari. Si noti come il campo elettrico sia irrotazionale (conservativo) e il campo magnetico sia solenoidale (linee chiuse).}
    \label{tab:maxwell_statiche}
\end{table}

\section{Significato fisico del disaccoppiamento}
L'aspetto più rilevante del caso statico è l'assenza di termini che leghino il campo elettrico a quello magnetico.
\begin{itemize}
    \item La terza equazione ($\nabla \times \bm{E} = 0$) implica che, in assenza di variazioni temporali magnetiche, il campo elettrostatico ammette un potenziale scalare ($V$).
    \item La quarta equazione ($\nabla \times \bm{B} = \mu_0 \bm{J}$) implica che il campo magnetico è generato solo da correnti di conduzione, senza il contributo della "corrente di spostamento" (che, come si vedrà, dipende dalla variazione temporale di $\bm{E}$).
\end{itemize}

\section{Appendice storica: La revisione di Heaviside}

È comune riferirsi alle quattro formule in Tabella \ref{tab:maxwell_statiche} come "Equazioni di Maxwell". Tuttavia, se James Clerk Maxwell vedesse questa pagina, probabilmente non le riconoscerebbe immediatamente.

\begin{approfondimento}[colback=white, colframe=black!70, no shadow,  title=Dai Quaterioni ai Vettori: il ruolo di Oliver Heaviside]
    Nel suo trattato fondamentale del 1873, \textit{"A Treatise on Electricity and Magnetism"}, Maxwell formulò la sua teoria utilizzando un sistema di ben \textbf{20 equazioni scalari} in \textbf{20 variabili}.
    
    Maxwell utilizzava il formalismo dei \textit{quaternion}i e poneva grande enfasi sul potenziale vettore $\bm{A}$ (che chiamava "quantità di moto elettrocinetica") piuttosto che sui campi $\bm{E}$ e $\bm{B}$.
    
    La forma elegante, simmetrica e compatta che studiamo oggi (le 4 equazioni vettoriali) è opera del fisico e matematico autodidatta britannico \textbf{Oliver Heaviside} (e parallelamente di Josiah Willard Gibbs e Heinrich Hertz) alla fine del XIX secolo.
    
    \paragraph{Il contributo di Heaviside:}
    \begin{itemize}
        \item \textbf{Vettorializzazione:} Eliminò i quaternioni, giudicandoli inutilmente complessi ("una sciagura per la fisica"), e promosse l'uso del calcolo vettoriale moderno (divergenza, gradiente, rotore).
        \item \textbf{Centralità dei Campi:} Spostò l'attenzione dai potenziali ($\phi, \bm{A}$) ai campi di forza ($\bm{E}, \bm{B}$), ritenuti le vere entità fisiche misurabili.
        \item \textbf{Sintesi:} Ridusse le 20 equazioni originali che descrivevano i campi a sole 4, evidenziando la straordinaria simmetria tra elettricità e magnetismo (rotta solo dall'assenza di monopoli magnetici).
    \end{itemize}
    
    Quelle che oggi vengono chiamate "Equazioni di Maxwell" dovrebbero, a rigore storico, essere chiamate "Equazioni di Maxwell-Heaviside".
\end{approfondimento}









\chapter{Campi tempo-varianti}
Finora sono stati analizzati il campo elettrico $\bm{E}$ e magnetico $\bm{B}$ come due entità pressoché indipendenti e esclusivamente nel caso in cui siano \textit{costanti nel tempo}. Nella trattazione ora proposta si cerca di porre le basi dell'analisi dei due campi in condizioni \textit{tempo-varianti}, dimostrando eventualmente la loro stretta correlazione. 

\section{Una relazione inattesa: l'esperimento di Faraday}

Fino ai primi decenni del XIX secolo la comunità scientifica, forte delle scoperte di Oersted e Ampere, aveva consolidato la conoscenza per cui una corrente elettrica genera un campo magnetico. Per ragioni di simmetria naturale, molti fisici iniziarono a domandarsi se fosse possibile il fenomeno inverso: può un campo magnetico generare una corrente elettrica?
I primi tentativi, condotti immergendo circuiti in intensi campi magnetici statici, diedero esito negativo. Fu Michael Faraday, nel 1831, a intuire che la chiave del fenomeno non risiedeva nell'intensità del campo, ma nella sua \textit{variazione}.

\subsection{Descrizione fenomenologica}
L'esperimento classico di Faraday consiste in un circuito semplice, composto da una bobina (solenoide) collegata in serie a un galvanometro (strumento sensibile capace di rilevare correnti anche molto deboli), in assenza di generatori di tensione (batterie).
Si dispone inoltre di una sorgente di campo magnetico, come una barretta di magnetite o un magnete permanente. L'esperimento procede osservando il comportamento dell'ago del galvanometro in diverse configurazioni:

\begin{enumerate}
    \item \textit{magnete in quiete:} se il magnete viene posizionato fermo all'interno o nelle vicinanze della bobina, il galvanometro non segna alcuna corrente ($I=0$). La semplice presenza di un campo magnetico statico $\bm{B}$, per quanto intenso, non produce effetti elettrici;
    
    \item \textit{magnete in moto (avvicinamento):} Se si muove il magnete verso la bobina (ad esempio inserendo il polo Nord nel solenoide), l'ago del galvanometro devia istantaneamente, segnalando il passaggio di una corrente $I(t) \neq 0$;
    
    \item \textit{magnete in moto (allontanamento):} Se si estrae il magnete dalla bobina, il galvanometro segna nuovamente una corrente, ma di verso opposto rispetto al caso precedente;
    
    \item \textit{variazione di velocità:} Si osserva inoltre che l'intensità della corrente registrata è proporzionale alla rapidità con cui il magnete viene mosso: movimenti più veloci generano correnti più intense.
\end{enumerate}



\begin{figure}[htbp]
    \centering
    \begin{tikzpicture}[scale=1, >=Stealth]
        
        % --- BOBINA (Solenoide) ---
        % Disegniamo una serie di ellissi parziali per simulare il solenoide
        \foreach \x in {0, 0.3, ..., 3} {
            \draw[thick, gray!70] (\x, 0.5) arc (90:270:0.3 and 0.5); % Parte dietro (meno visibile se volessimo 3D vero)
            \draw[thick, black] (\x, -0.5) arc (-90:90:0.3 and 0.5); % Parte davanti
        }
        
        % Fili verso il galvanometro
        \draw[thick, black] (0, -0.5) -- (0, -1.5) -- (1, -1.5);
        \draw[thick, black] (3, -0.5) -- (3, -1.5) -- (2, -1.5);
        
        % --- GALVANOMETRO ---
        \draw[thick, fill=white] (1.5, -1.5) circle (0.5);
        \node[font=\footnotesize] at (1.5, -1.8) {G};
        % Ago che devia
        \draw[->, thick, red] (1.5, -1.5) -- (1.8, -1.1);
        \draw[thin] (1.2, -1.2) arc (135:45:0.4); % Scala graduata
        
        % --- MAGNETE ---
        \begin{scope}[xshift=-2.5cm, yshift=0cm]
            \draw[fill=blue!30] (0, -0.3) rectangle (1, 0.3);
            \draw[fill=red!30] (1, -0.3) rectangle (2, 0.3);
            \node at (0.5, 0) {S};
            \node at (1.5, 0) {N};
            
            % Linee di campo del magnete
            \draw[->, blue, dashed] (2, 0) -- (3.5, 0);
            \draw[->, blue, dashed] (2, 0.2) .. controls (3, 0.5) .. (3.5, 0.8);
            \draw[->, blue, dashed] (2, -0.2) .. controls (3, -0.5) .. (3.5, -0.8);
            
            % Vettore Velocità
            \draw[->, ultra thick, orange] (0.5, 0.5) -- (1.5, 0.5) node[midway, above] {$\bm{v}$};
        \end{scope}
        
        % Freccia corrente indotta sulla bobina
        \node[red, font=\bfseries] at (1.5, 0.8) {$I_{ind}$};
        \draw[->, red, thick] (1.2, 0.6) -- (1.8, 0.6);

    \end{tikzpicture}
    \caption{Rappresentazione schematica dell'esperimento di Faraday. Il moto del magnete verso il solenoide induce una variazione del campo magnetico concatenato, generando una corrente misurata dal galvanometro.}
    \label{fig:faraday_exp}
\end{figure}

\subsection{Interpretazione fisica: il campo elettrico indotto}
L'osservazione sperimentale porta a una conclusione rivoluzionaria per l'epoca. Poiché nel circuito scorre una corrente, le cariche libere del conduttore devono essere soggette a una forza che compie lavoro (detta \textit{forza elettromotrice, fem = $\mathcal{E}$}).
In assenza di generatori chimici (pile), l'unica spiegazione è la presenza di un {campo elettrico} $\bm{E}$ indotto lungo il filo.
\noindent
Questo campo elettrico, tuttavia, non è generato da cariche statiche (non vi è accumulo di carica elettrostatica), ma è causato dalla variazione temporale del campo magnetico.
Si può dunque affermare che:
\begin{center}
    \textit{Un campo magnetico variabile nel tempo genera un campo elettrico.}
\end{center}
Si noti una differenza sostanziale rispetto all'elettrostatica: il campo elettrico indotto $\bm{E}_{ind}$ non è conservativo. Essendo responsabile della circolazione di corrente in un circuito chiuso, la sua circuitazione (che corrisponde al lavoro complessivo fatto lungo la curva chiusa\footnote{Si noti che l'integrale del campo elettrico lungo una lunghezza corrisponde a una differenza di potenziale che, se non nulla, implica un moto di cariche.}) non è nulla:
\begin{equation}
   \mathcal{E} = fem = \oint_{\gamma} \bm{E}_{ind} \cdot d\bm{l} \neq 0.
\end{equation}
Questo risultato infrange la terza equazione di Maxwell per il caso statico ($\nabla \times \bm{E} = 0$) e richiede una nuova formulazione che leghi la circuitazione di $\bm{E}$ alla variazione del campo magnetico $\bm{B}$.

\section{La legge di Faraday-Neumann-Lenz}
Da evidenze sperimentali di cui Faraday si è fatto pioniere, venne formulata la seguente legge:
\begin{center}
    \textit{Un campo magnetico variabile nel tempo genera un campo elettrico la cui circuitazione lungo una curva chiusa $\gamma$ è pari e opposta alla derivata temporale del flusso del campo magnetico attraverso una qualsiasi superficie avente come sostegno $\gamma$.}
\end{center}
Si ricorda che, a causa di reminiscenze storiche, il campo elettrico indotto è anche conosciuto come forza elettromotrice\footnote{Nonostante non sia una forza!}. In formule, il flusso del campo magnetico attraverso una superficie $\Sigma$ che ha $\gamma$ come sostegno vale:
\begin{equation}
\Phi_{\Sigma}(\bm{B}) = \int_{\Sigma} \bm{B} \cdot d\bm{S} \quad [\Phi_{\Sigma}(\bm{B})] = \si{\tesla\meter\squared} = \si{\weber},
\end{equation} 
con $\si{\weber}$ l'unità di misura del flusso magnetico, il Weber. Secondo la formulazione di Faraday-Neumann-Lenz, quindi:
\begin{equation}
\Gamma_{\gamma}(\bm{E}) = \mathcal{E} = \oint_\gamma \bm{E} \cdot d\bm{l} =  -\dot{\Phi}_{\Sigma}(\bm{B}) = -\frac{d \Phi_{\Sigma}(\bm{B})}{d t} = -\frac{d }{d t} \int_{\Sigma} \bm{B} \cdot d\bm{S}.
\label{eq:far_neu_len}
\end{equation} 
Nel caso tempo-invariante,
\begin{equation}
\Gamma_{\gamma}(\bm{E}) = \mathcal{E} = -\dot{\Phi}_{\Sigma}(\bm{B}) = 0,
\end{equation} 
come prescritto dalla terza equazione di Maxwell nel caso statico. 

\subsection{La terza equazione di Maxwell: caso tempo-variante}
Riscrivendo la circuitazione nella \eqref{eq:far_neu_len} sfruttando il Teorema di Stokes e commutando, a membro destro dell'equazione, l'operatore di derivata temporale con quello di integrale\footnote{Si assume che la superficie $\Sigma$ sia fissa nello spazio (non si deformi nel tempo).} si ottiene:
\begin{equation}
   \oint_\gamma \bm{E}\cdot d\bm{l} = \int_{\Sigma} (\nabla \times \bm{E}) \cdot d\bm{S} = -\frac{d }{d t} \int_{\Sigma} \bm{B} \cdot d\bm{S} = \int_{\Sigma} \left( -\frac{\partial \bm{B}}{\partial t} \right) \cdot d\bm{S}.
\end{equation}
Affinché questa uguaglianza valga per \textit{qualunque} superficie $\Sigma$ arbitraria scelta, gli argomenti degli integrali devono essere uguali punto per punto. Si ottiene così la \textit{terza equazione di Maxwell in forma differenziale}:

\begin{equation}
    { \nabla \times \bm{E} = -\frac{\partial \bm{B}}{\partial t} }.
    \label{eq:maxwell_3_diff}
\end{equation}
Il parterre di equazioni finora presentate descrive così il fenomeno dell'\textit{induzione elettromagnetica}.

\subsubsection{Discussione del risultato}
L'equazione \eqref{eq:maxwell_3_diff} rappresenta una rottura fondamentale rispetto all'elettrostatica e introduce concetti fisici nuovi:

\begin{enumerate}
    \item \textit{Natura non conservativa del campo elettrico:}
    in elettrostatica valeva $\nabla \times \bm{E} = 0$, condizione necessaria e sufficiente\footnote{Solo nel caso in cui si possa applicare il Lemma di Poincaré, dunque con un campo vettoriale definito su un dominio semplicemente connesso.} per definire il campo come conservativo ($\bm{E} = -\nabla V$).
    In regime dinamico, la presenza di un campo magnetico variabile ($\frac{\partial \bm{B}}{\partial t} \neq 0$) rende il rotore del campo elettrico diverso da zero.
    {Conseguenze:}
    \begin{itemize}
        \item il campo elettrico indotto \textit{non è conservativo};
        \item non è possibile definire un potenziale scalare $V$ globale tale che $\bm{E} = -\nabla V$ (sarà necessaria l'introduzione dei potenziali ritardati e del potenziale vettore);
        \item le linee di forza del campo elettrico indotto sono \textit{linee chiuse} su se stesse (vortici), prive di sorgenti (cariche) di inizio o fine, analogamente alle linee del campo magnetico.
    \end{itemize}

    \item \textit{Accoppiamento tra campi:}
    l'equazione stabilisce un legame locale diretto tra elettricità e magnetismo. Una variazione temporale del campo magnetico in un punto $P$ dello spazio genera istantaneamente una vorticità del campo elettrico nello stesso punto $P$. I due campi non sono più entità indipendenti.

    \item \textit{Significato del segno meno (Legge di Lenz locale):}
    il segno meno è l'espressione matematica dell'inerzia elettromagnetica. Il rotore di $\bm{E}$ (e quindi la direzione in cui il campo tende a far circolare le cariche) si oppone alla variazione di $\bm{B}$. 
    
    Si immagini di avere una spira circolare $\gamma$ immersa in una regione di campo magnetico tempo-variante: il verso con cui fluisce corrente indotta nel circuito (orario o antiorario) è dettato dal contributo di Lenz: \textit{la corrente scorre in modo tale da generare un contro-campo che si opponga alla variazione del flusso del campo esterno.}
    
    \item \textit{Lavoro dei campi:}
è fondamentale notare una sottigliezza termodinamica.
    Come dimostrato, la forza magnetica è sempre perpendicolare alla velocità della carica e, pertanto, \textit{non compie mai lavoro} su di essa ($W = \int \bm{F} \cdot d\bm{l} = 0$).
    Chi fornisce dunque l'energia necessaria a mettere in moto le cariche (generare corrente) e a dissipare calore per effetto Joule?
Nel caso dell'induzione, il campo magnetico variabile agisce come \textit{intermediario}: la sua variazione genera un {campo elettrico indotto} che esercita la forza tangenziale $\bm{F}_e = q\bm{E}$ che compie lavoro sulle cariche.
         L'energia proviene sempre dall'\textit{agente esterno} (es. l'operatore che muove il magnete nell'esperimento di Faraday). Per la legge di Lenz, infatti, la corrente indotta genera un contro-campo che si oppone al moto del magnete: l'operatore deve compiere un lavoro meccanico per vincere questa forza frenante. Tale lavoro viene convertito in energia elettrica.


    \end{enumerate}



\begin{figure}[htbp]
    \centering
    \begin{tikzpicture}[scale=1.5, >=Stealth]
        
        % --- Vettore B(t) crescente ---
        \draw[->, ultra thick, blue] (0, -1) -- (0, 1.5) node[above] {$\bm{B}(t)$ (crescente)};
        
        % Vettore dB/dt
        \draw[->, thick, blue!60] (0.3, 0.5) -- (0.3, 1.2) node[midway, right] {$\frac{\partial \bm{B}}{\partial t}$};
        
        % --- Piano ortogonale ---
        \draw[dashed, gray] (0,0) ellipse (2 and 0.8);
        
        % --- Campo Elettrico Indotto (Vortice) ---
        % B cresce verso l'alto -> dB/dt verso l'alto.
        % Rot E = - dB/dt (verso il basso).
        % Regola mano destra col pollice in giù -> E gira orario (visto da sopra).
        
        \draw[->, thick, red] (1.5, 0.1) arc (10:-20:1.5 and 0.6);
        \draw[->, thick, red] (1.0, -0.6) arc (-50:-130:1.5 and 0.6);
        \draw[->, thick, red] (-1.5, -0.1) arc (190:170:1.5 and 0.6);
        \draw[thick, red] (-1.5, 0.1) arc (170:120:1.5 and 0.6); % Tratto dietro parziale
        \draw[->, thick, red] (-0.5, 0.75) arc (100:60:1.5 and 0.6); % Tratto dietro
        
        \node[red] at (1.8, 0) {$\bm{E}_{ind}$};
        
        % --- Annotazione Rotore ---
        \node[align=center, font=\small] at (3, 1) {$\nabla \times \bm{E}$ opposto\\a $\frac{\partial \bm{B}}{\partial t}$};

    \end{tikzpicture}
    \caption{Rappresentazione locale della terza equazione. Se il campo magnetico $\bm{B}$ aumenta verso l'alto, il campo elettrico indotto $\bm{E}$ forma vortici in senso orario (opponendosi alla variazione), dimostrando la natura non conservativa del campo.}
    \label{fig:rotore_E}
\end{figure}


\begin{approfondimento}[colback=white, colframe=blue!60, no shadow, title=Validità generale della legge di Faraday-Neumann-Lenz]

È fondamentale chiarire un dubbio comune: poiché per ricavare la forma locale $\nabla \times \bm{E} = -\frac{\partial \bm{B}}{\partial t}$ si è ipotizzata una spira fissa, questa equazione cessa di valere se il circuito si muove o si deforma?

\textbf{La risposta è no. La forma differenziale è sempre valida.}

Le equazioni di Maxwell in forma differenziale sono leggi \textit{locali}: descrivono il comportamento dei campi elettromagnetici in un punto geometrico dello spazio $(x,y,z)$ in un istante $t$. In un punto geometrico non esiste "forma" o "circuito". L'equazione comunica semplicemente che:
\begin{center}
    \textit{"Se in un punto P il campo magnetico cambia nel tempo, allora in quel punto P il campo elettrico ha un rotore."}
\end{center}
Questo accade nel vuoto, dentro una stella o vicino a un filo, indipendentemente dalla presenza di circuiti fisici.

\subsubsection{La Regola del Flusso e la Forza di Lorentz}
Considerando un circuito fisico reale che si muove o si deforma (ad esempio una sbarra mobile su binari), la variazione totale del flusso $\frac{d\Phi}{dt}$ è dovuta a due contributi.
Applicando l'analisi vettoriale alla derivata del flusso (Teorema del trasporto di Reynolds), si ottiene la scomposizione della forza elettromotrice totale:

\begin{equation}
   \mathcal{E} = \underbrace{\oint_{\gamma} \bm{E}_{ind} \cdot d\bm{l}}_{\text{Termine di Campo}} + \underbrace{\oint_{\gamma} (\bm{v} \times \bm{B}) \cdot d\bm{l}}_{\text{Termine Mozionale}} = - \frac{d\Phi}{dt}.
\end{equation}

\begin{enumerate}
    \item \textit{Il termine di Campo ($\bm{E}_{ind}$):} Dipende da $\frac{\partial \bm{B}}{\partial t}$. È l'unico termine descritto dalla terza equazione di Maxwell ($\nabla \times \bm{E} = -\frac{\partial \bm{B}}{\partial t}$). Agisce anche se il circuito è fermo.
    \item \textit{Il termine Mozionale ($\bm{v} \times \bm{B}$):} Dipende dalla velocità $\bm{v}$ del conduttore. Non deriva dalle equazioni di Maxwell per il campo, ma dalla \textit{Forza di Lorentz}.
\end{enumerate}

\subsubsection{Sintesi}
    La legge integrale "del flusso" ($\mathcal{E} = -d\Phi/dt$) è una "super-legge" che racchiude in un'unica formula due fenomeni fisici distinti: l'induzione di Maxwell (campo elettrico vorticoso) e la forza magnetica di Lorentz sulle cariche in moto.
    La forma differenziale di Maxwell descrive solo la prima parte (la natura del campo), ed è per questo che è universalmente valida in ogni punto.
\end{approfondimento}

\section{Conseguenze dell'induzione}
La corrente che fluisce in un circuito genera un campo magnetico nello spazio circostante. Quali sono gli effetti della sua variazione sul circuito stesso o su altri circuiti soggetti al campo da questo prodotto? 

\subsection{L'autoinduzione}
La corrente $I(t)$, eventualmente variabile nel tempo, che fluisce in un circuito \textit{chiuso} $\gamma$ produce un campo magnetico nello spazio:
\begin{equation}
\bm{B} = \frac{\mu_0 I(t)}{4\pi} \oint_\gamma \frac{d\bm{l}' \times \bm{u}_r}{|\bm{r}-\bm{r}'|^2}.
\end{equation}
Il flusso di $\bm{B}$ attraverso una qualsiasi superficie $\Sigma$ che poggia su $\gamma$ vale:
\begin{equation}
\Phi_\Sigma(\bm{B}) = \int_\Sigma \bm{B}(\bm{r}) \cdot d\bm{S} =I(t)\int_\Sigma \left(\frac{\mu_0 }{4\pi} \oint_\gamma \frac{d\bm{l}' \times \bm{u}_r}{|\bm{r}-\bm{r}'|^2}\right) \cdot d\bm{S},
\end{equation}
 dove si è trattata $I(t)$ come una costante lungo il filo (\textit{approssimazione quasi-stazionaria}) in quanto unicamente dipendente dal tempo. Il fattore moltiplicativo della corrente è una costante (dimensionata) caratteristica del circuito e viene detto \textit{coefficiente di autoinduzione} o \textit{(auto)induttanza}, indicato con $L$:
 \begin{equation}
L = \int_\Sigma \left(\frac{\mu_0 }{4\pi} \oint_\gamma \frac{d\bm{l}' \times \bm{u}_r}{|\bm{r}-\bm{r}'|^2}\right) \cdot d\bm{S} \quad [L] =  \si{\weber\per\ampere} = \si{\henry}.
\end{equation}
La natura intrinsecamente geometrica di $L$ emerge applicando il Teorema di Stokes:
 \begin{equation}
L = \frac{1}{I(t)}\int_\Sigma \bm{B} \cdot d\bm{S} = \frac{1}{I(t)}\int_\gamma \bm{A} \cdot d\bm{l}'
\end{equation}
e sostituendo l'espressione del potenziale vettore:
\begin{equation}
L = \frac{1}{I(t)} \oint_\gamma \bm{A} \cdot d\bm{l} = \oint_\gamma \left( \frac{\mu_0 }{4\pi} \oint_\gamma \frac{d\bm{l}'}{|\bm{r}-\bm{r}'|} \right) \cdot d\bm{l}.
\end{equation}
Poiché $I(t)$ è costante lungo tutto il filo, può essere portata fuori da entrambi gli integrali. Si ottiene così la \textit{Formula di Neumann} per l'autoinduzione:
\begin{equation}
L = \frac{\mu_0}{4\pi} \oint_\gamma \oint_\gamma \frac{d\bm{l} \cdot d\bm{l}'}{|\bm{r}-\bm{r}'|},
\end{equation}
evidentemente dipendente dalle caratteristiche morfologiche del circuito ($|\bm{r}-\bm{r}'|$ corrisponde, infatti, alla distanza fra due punti generici di $\gamma$) e dalla permeabilità magnetica del mezzo (in questo caso il vuoto, $\mu_0$). Se il circuito è immerso in un materiale magnetico lineare con permeabilità $\mu$, $L$ dipenderà linearmente da $\mu$.

\subsubsection{$\bm{fem}$ autoindotta}
Dalla definizione di coefficiente di autoinduzione, l'autoflusso del campo magnetico risulta:
\begin{equation}
\Phi_\Sigma(\bm{B}) = LI(t).
\end{equation}
Nel caso di una corrente tempo-variante, la forza elettromotrice autoindotta è:
\begin{equation}
\mathcal{E} = -\dot{\Phi}_\Sigma(\bm{B}) = -L\dot{I}(t).
\end{equation}
L'equazione esprime analiticamente la legge di Lenz per i circuiti auto-induttivi. Il segno meno indica che la forza elettromotrice autoindotta si oppone sempre alla \textit{variazione} della corrente che la ha generata:
\begin{itemize}
    \item se la corrente aumenta ($\dot{I} > 0$), la $fem$ autoindotta ha verso opposto alla corrente, comportandosi come un generatore che tenta di frenarne la crescita;
    \item se la corrente diminuisce ($\dot{I} < 0$), la $fem$ autoindotta ha lo stesso verso della corrente, tentando di sostenerla e ritardarne l'estinzione.
\end{itemize}
In tal senso, il coefficiente di autoinduzione $L$ rappresenta una misura dell'\textit{inerzia elettromagnetica} del circuito. Analogamente a come la massa in meccanica si oppone alle variazioni di velocità ($\bm{F} = m\bm{a}$), l'induttanza $L$ si oppone alle variazioni di corrente elettrica. Un circuito con elevata $L$ (ad esempio un solenoide con nucleo ferromagnetico) rende difficile variare rapidamente l'intensità di corrente che vi scorre.

\noindent
Il valore numerico corrispondente al coefficiente di autoinduzione di un circuito è in generale abbastanza piccolo da poter considerare l'autoinduzione \textit{trascurabile}.


\subsection{Energia immagazzinata nel campo magnetico}
\subsubsection{Lavoro della forza elettromotrice auotindotta}
Sia $\iota$ un solenoide attraverso cui fluisce una corrente $I$ variabile nel tempo (si immagini, ad esempio, di accendere un generatore di tensione collegato in serie a $\iota$ e che la corrente, di conseguenza, vada incontro a un transitorio da $I = 0$ a $I = I_0$). Per quanto discusso precedentemente, l'autoflusso variabile nel tempo attraverso la sezione di $\iota$ induce una $fem$ nel circuito che fa lavoro sulle cariche. Detta $dq$ una carica infinitesima:
\begin{equation}
\delta W = \mathcal{E}dq = -L\frac{dI}{dt}dq = -L\frac{dq}{dt}dI = -LIdI
\end{equation}
\begin{equation}
\implies W = -\int_0^{I_0} LIdI = -\frac{1}{2}LI_0^2.
\label{eq:lavoro_fem}
\end{equation}
Il significato del segno - risiede nel fatto che il campo elettrico autoindotto lavora \textit{contro} il transitorio della corrente: nel caso in esame, il transitorio porta la corrente ad aumentare $\rightarrow$ la $fem$ lavora contro l'aumento. 

\subsubsection{Dal particolare al generale: il solenoide ideale}
L'autoflusso complessivo attraverso un solenoide ideale (tenendo in considerazione il contributo di ogni singolo avvolgimento e il fatto che il vettore campo magnetico, data la geometria del sistema, è sempre parallelo al vettore superficie delle singole spire), detti $N$ il numero di spire, $l$ la lunghezza e $S$ la sezione, vale:
\begin{equation}
\Phi(\bm{B}) = LI = \bm{B} S N= \frac{\mu_0 I N}{l} SN = \underbrace{\frac{\mu_0 N^2 S}{l}}_{L} I.
\end{equation}
Sfruttando la \eqref{eq:lavoro_fem}, il lavoro della $fem$:
\begin{equation}
W = -\frac{1}{2} \frac{\mu_0 N^2 S}{l} I^2 = -\frac{1}{2} \frac{B^2}{\mu_0} Sl,
\end{equation}
dove $Sl$ corrisponde al volume di $\iota$. Si definisce \textit{densità volumetrica di energia magnetica} il lavoro compiuto contro la forza elettromotrice autoindotta, per unità di volume:
\begin{equation}
w_m = -\frac{W}{V} =\frac{1}{2} \frac{B^2}{\mu_0}.
\end{equation}
Come nel caso della densità volumetrica di energia del campo elettrico, il risultato ottenuto è generalizzabile per ogni regione di spazio immersa in un campo magnetico $\bm{B}$.

\begin{approfondimento}[colback=white, colframe=blue!60, no shadow, title=Generalizzazione: l'energia del campo magnetico]

Il risultato ottenuto per il solenoide, ovvero che la densità di energia è proporzionale al quadrato del campo magnetico, possiede una validità generale. Si dimostra di seguito che l'energia magnetica totale non è localizzata sul circuito conduttore, ma è distribuita in tutto lo spazio sede del campo.

Si consideri un sistema di correnti stazionarie descritte da una densità di corrente $\bm{J}$. L'energia magnetica totale $U_m$ può essere espressa, generalizzando la formula del circuito filiforme ($U = \frac{1}{2}I\Phi$), tramite l'integrale di volume:
\begin{equation}
U_m = \frac{1}{2} \int_V \bm{J} \cdot \bm{A} \, dV,
\end{equation}
dove $\bm{A}$ è il potenziale vettore e $V$ è il volume contenente le correnti.
Sfruttando la legge di Ampère in forma locale ($\nabla \times \bm{B} = \mu_0 \bm{J}$), si sostituisce $\bm{J}$:
\begin{equation}
U_m = \frac{1}{2\mu_0} \int_V (\nabla \times \bm{B}) \cdot \bm{A} \, dV.
\end{equation}
Si fa ora uso della seguente identità vettoriale relativa alla divergenza di un prodotto vettoriale:
\begin{equation}
\nabla \cdot (\bm{A} \times \bm{B}) = \bm{B} \cdot (\nabla \times \bm{A}) - \bm{A} \cdot (\nabla \times \bm{B}).
\end{equation}
Da cui si ricava il termine integrando:
\begin{equation}
\bm{A} \cdot (\nabla \times \bm{B}) = \bm{B} \cdot (\nabla \times \bm{A}) - \nabla \cdot (\bm{A} \times \bm{B}) = \bm{B} \cdot \bm{B} - \nabla \cdot (\bm{A} \times \bm{B}).
\end{equation}
Sostituendo nell'espressione dell'energia:
\begin{equation}
U_m = \frac{1}{2\mu_0} \int_V B^2 \, dV - \frac{1}{2\mu_0} \int_V \nabla \cdot (\bm{A} \times \bm{B}) \, dV.
\end{equation}
Applicando il \textit{Teorema della Divergenza} al secondo integrale, lo si trasforma in un flusso attraverso la superficie $\Sigma$ che delimita il volume $V$:
\begin{equation}
\int_V \nabla \cdot (\bm{A} \times \bm{B}) \, dV = \oint_\Sigma (\bm{A} \times \bm{B}) \cdot d\bm{S}.
\end{equation}
Se si estende il dominio di integrazione a tutto lo spazio ($V \to \infty$), la superficie $\Sigma$ si sposta all'infinito. Poiché per sorgenti localizzate i campi decadono con la distanza (asintoticamente $\bm{A} \sim 1/r$ e $\bm{B} \sim 1/r^2$), il termine integrando decade come $1/r^3$, mentre la superficie cresce come $r^2$. Il flusso all'infinito tende dunque a zero.

Rimane quindi solo il primo termine:
\begin{equation}
U_m = \int_{\text{spazio}} \frac{B^2}{2\mu_0} \, dV.
\end{equation}
Questo conferma che l'energia può essere considerata distribuita nello spazio con una \textit{densità volumetrica di energia magnetica} $w_m$ data da:
\begin{equation}
w_m = \frac{B^2}{2\mu_0}.
\end{equation}
Nel caso di materiali magnetici lineari ($\bm{B} = \mu \bm{H}$), l'espressione generalizzata diviene:
\begin{equation}
w_m = \frac{1}{2} \bm{B} \cdot \bm{H}.
\end{equation}
\end{approfondimento}

\subsubsection{Un esempio: il circuito RL}
Un circuito con autoinduzione non trascurabile può essere modellizzato inserendo nella trattazione elettrotecnica un solenoide. L'esempio più semplice è quello di un circuito elementare costituito da un generatore ideale di forza elettromotrice costante $\mathcal{E}_0$, un resistore di resistenza $R$ e un induttore ideale di coefficiente di autoinduzione $L$, tutti collegati in serie.
Si supponga che all'istante $t=0$ il circuito venga chiuso tramite un interruttore. In assenza dell'induttore, la corrente raggiungerebbe istantaneamente il valore di regime $I_0 = \mathcal{E}_0/R$; la presenza di $L$, tuttavia, introduce un'inerzia elettromagnetica che modifica la dinamica del sistema.

\begin{center}
\begin{tikzpicture}[
    circuit ee IEC,           % Abilita la libreria IEC
    x=1.5cm, y=1.5cm,         % Scala delle coordinate
    thick,                    % Spessore delle linee
    every info/.style={font=\footnotesize}, % Stile delle etichette dei componenti
    small circuit symbols     % Dimensioni dei simboli
    ]

  % Disegno del circuito in un unico percorso chiuso
  % Si parte dall'angolo in basso a sinistra (0,0) e si procede in senso orario
  \draw (0,0)
    % Ramo sinistro: Generatore di tensione continua
    % 'battery' è il simbolo. 'info' è l'etichetta.
    % La polarità standard ha il + (linea lunga) in alto se disegnato dal basso verso l'alto.
    to [battery={info={$\mathcal{E}_0$}}] (0,2)

    % Ramo superiore: Interruttore e Resistore
    % 'make contact' è l'interruttore normalmente aperto.
    to [make contact={info={$S$}}] (1.5,2)
    % Aggiunta freccia per la corrente i(t)
    to [current direction'={info={$I(t)$}, color=red!80!black, >=Latex}] (2.5,2)
    % 'resistor' è il resistore (simbolo zigzag IEC)
    to [resistor={info={$R$}}] (4,2)

    % Ramo destro: Induttore
    % 'inductor' è l'induttore (simbolo bobina)
    to [inductor={info={$L$}}] (4,0)

    % Ramo inferiore: Filo di chiusura
    -- (0,0);
\end{tikzpicture}

\end{center}

\paragraph{Equazione del circuito}
Applicando la seconda legge di Kirchhoff (legge delle maglie) all'unica maglia presente, si impone che la somma algebrica delle tensioni sia nulla. Considerando la caduta di potenziale ohmica e la forza elettromotrice autoindotta (che si oppone all'aumento di corrente), si ottiene l'equazione differenziale del primo ordine:
\begin{equation}
\mathcal{E}_0 - R I(t) - L \frac{dI}{dt} = 0.
\end{equation}
Riordinando i termini per evidenziare la derivata della corrente:
\begin{equation}
L \frac{dI}{dt} = \mathcal{E}_0 - R I(t),
\end{equation}
che ammette come soluzione:
\begin{equation}
I(t) = \frac{\mathcal{E}_0}{R} \left( 1 - e^{-\frac{t}{\tau}} \right),
\end{equation}
dove si è introdotta la \textit{costante di tempo caratteristica} del circuito:
\begin{equation}
\tau = \frac{L}{R} \quad [\si{\second}].
\end{equation}

\paragraph{Discussione fisica}
Dall'analisi della soluzione $I(t)$ emergono le proprietà fondamentali del comportamento induttivo:
\begin{itemize}
    \item \textit{Istante iniziale ($t=0$):} si ha $I(0) = 0$. L'induttore si comporta come un \textit{circuito aperto}. L'elevata derivata della corrente genera una $fem$ autoindotta che eguaglia quasi perfettamente la tensione del generatore ($\mathcal{E}_L \approx -\mathcal{E}_0$), impedendo il passaggio di cariche.
    \item \textit{Regime stazionario ($t \to \infty$):} l'esponenziale tende a zero e la corrente tende asintoticamente al valore ohmico $I_{\text{regime}} = \mathcal{E}_0/R$. La derivata temporale si annulla, la $fem$ autoindotta svanisce e l'induttore si comporta come un \textit{corto circuito} (un semplice filo ideale).
    \item \textit{Significato di $\tau$:} la costante $\tau$ rappresenta il tempo necessario affinché la corrente raggiunga circa il $63\%$ del suo valore finale. Un valore elevato di $L$ (grande inerzia) implica un $\tau$ grande, e quindi una salita della corrente molto lenta.
\end{itemize}

\paragraph{Bilancio energetico}
Riprendendo l'equazione della maglia e moltiplicando ogni termine per la corrente istantanea $i(t)$, si ottiene un'equazione di bilancio per la potenza:
\begin{equation}
\underbrace{\mathcal{E}_0 I}_{\mathcal{P} \text{erogata}} = \underbrace{R I^2}_{\text{P dissipata}} + \underbrace{L I \frac{di}{dt}}_{\text{P immagazzinata}}.
\end{equation}
Questa relazione conferma quanto discusso nelle sezioni precedenti: la potenza fornita dal generatore non viene interamente dissipata in calore sulla resistenza (come avverrebbe in regime stazionario), ma una parte viene costantemente prelevata per accrescere l'energia del campo magnetico ($dW_m/dt$) durante tutta la fase del transitorio. Solo a regime ($di/dt=0$), tutta la potenza erogata viene dissipata per effetto Joule.



\subsubsection{La risposta del circuito}
Le equazioni che governano un circuito $RL$ (resistenza-induttanza) sono matematicamente isomorfe a quelle che descrivono il moto di un corpo dotato di massa in presenza di attrito viscoso. In questa analogia, l'induttanza $L$ gioca il ruolo della \textit{massa inerziale}. Nella Tabella (\ref{tab:analogia_elettro_mecc}) sono riportate le corrispondenze dirette tra le variabili elettriche e quelle meccaniche.
\begin{table}[htbp]
    \centering
    \renewcommand{\arraystretch}{1.5}
    \begin{tabular}{l c | l c}
        \hline
        \multicolumn{2}{c|}{\textit{Dominio Elettrico}} & \multicolumn{2}{c}{\textit{Dominio Meccanico}} \\
        \hline
        Carica & $q$ & Posizione & $x$ \\
        Corrente & $I = \dot{q}$ & Velocità & $v = \dot{x}$ \\
        Variazione di corrente & $dI/dt = \ddot{q}$ & Accelerazione & $a = \dot{v} = \ddot{x}$ \\
        \hline
        \textit{Induttanza} & $L$ & \textit{Massa} & $m$ \\
        Resistenza & $R$ & Coeff. di attrito & $\beta$ \\
        Generatori ($fem$) & $\mathcal{E}$ & Forza esterna & $F_{ext}$ \\
        \hline
    \end{tabular}
    \caption{Quadro sinottico dell'analogia elettromeccanica. L'induttanza rappresenta l'inerzia del sistema elettrico.}
    \label{tab:analogia_elettro_mecc}
\end{table}
Sulla base di queste corrispondenze, è possibile reinterpretare i processi energetici discussi precedentemente:

\begin{enumerate}
    \item \textit{Inerzia e Autoinduzione:}
    così come la massa si oppone alle variazioni di velocità generando una forza d'inerzia ($F_{in} = -ma$), l'induttanza si oppone alle variazioni di corrente generando una forza elettromotrice autoindotta ($\mathcal{E}_L = -L \dot{I}$). In entrambi i casi, il sistema reagisce per mantenere il proprio stato di moto (o di corrente) invariato.

    \item \textit{Accumulo di Energia:}
    il lavoro compiuto dal generatore per vincere la $fem$ autoindotta durante un transitorio di accensione non viene dissipato (come avviene sulla resistenza), ma viene immagazzinato nel sistema.
    \begin{itemize}
        \item In meccanica, il lavoro speso per accelerare la massa diventa \textit{energia cinetica}:
        \[ K = \frac{1}{2} m v^2. \]
        \item In elettromagnetismo, il lavoro speso per aumentare la corrente diventa \textit{energia magnetica}:
        \[ U_m = \frac{1}{2} L I^2. \]
    \end{itemize}
    
    \item \textit{Restituzione dell'energia:}
    se viene rimossa la spinta esterna (generatore spento / forza cessata), l'energia accumulata impedisce l'arresto istantaneo. La massa continua a muoversi per inerzia fino a che l'attrito non dissipa l'energia cinetica; analogamente, la corrente continua a scorrere (extra-corrente di apertura) spinta dall'energia magnetica, fino a che la resistenza non la dissipa completamente per effetto Joule.
\end{enumerate}
    Lavorare contro la $fem$ autoindotta significa spendere energia per "costruire" il campo magnetico, ovvero per mettere in moto i "volani elettromagnetici" del sistema. Tale energia è recuperabile ed è distinta da quella dissipata irreversibilmente sulla componente resistiva.

\subsection{La mutua induzione}
Siano $\gamma_1$ e $\gamma_2$ due circuiti chiusi e $I_1(t)$ e $I_2(t)$ le correnti\footnote{Si ritiene sempre valida l'approssimazione di quasi stazionarietà.} che in essi fluiscono. Il flusso del campo magnetico prodotto da $\gamma_1$ in $\gamma_2$ è:
\begin{equation}
\Phi_{1\rightarrow 2}(\bm{B}_1) = \oint_{\gamma_2} \bm{B}_1 \cdot d\bm{S}_2.
\end{equation}
Con considerazioni assolutamente affini a quelle discusse nel caso dell'autoinduzione, si dimostra che il flusso non è altro che il prodotto della corrente che scorre nel circuito che produce il campo magnetico per un coefficiente numerico \textit{geometrico} detto, in questo contesto, \textit{coefficiente di mutua induzione} $M$. I flussi concatenati ai circuiti, quindi:
\begin{equation}
\begin{cases}
\Phi_{1\rightarrow 2}(\bm{B}_1) = M_{1\rightarrow 2}I_1(t) \\
\Phi_{2\rightarrow 1}(\bm{B}_2) = M_{2\rightarrow 1}I_2(t)
\end{cases}
,
\end{equation}
dove con $M_{i\rightarrow j}$ si sono indicati rispettivamente i coefficienti di mutua induzione di un circuito rispetto all'altro.

\subsubsection{La relazione fra i coefficienti di mutua induzione}
Per stabilire un legame tra i coefficienti di mutua induzione si analizza il bilancio energetico di un sistema composto da due circuiti filiformi rigidi, $\gamma_1$ e $\gamma_2$. Poiché l'energia magnetica è una funzione di stato, il lavoro totale speso per portare il sistema a una configurazione finale di correnti $(I_1, I_2)$ deve essere indipendente dall'ordine cronologico di accensione dei generatori.
Si consideri inizialmente l'accensione sequenziale $1 \rightarrow 2$. Si porta la corrente in $\gamma_1$ al valore $I_1$ (con $\gamma_2$ aperto), spendendo un lavoro di autoinduzione $W_1 = \frac{1}{2}L_1I_1^2$. Successivamente, mantenendo $I_1$ costante, si porta la corrente in $\gamma_2$ al valore $I_2$. Durante questo secondo transitorio vengono compiuti due lavori: quello di autoinduzione sul circuito 2 ($W_2 = \frac{1}{2}L_2I_2^2$) e quello compiuto dal generatore nel circuito 1 per contrastare la $fem$ indotta dal flusso variabile proveniente da 2.
Tale $fem$ vale $\mathcal{E}_1 = -M_{2\rightarrow 1}\dot{I}_2(t)$. Il lavoro infinitesimo compiuto dal generatore 1 per mantenere la corrente costante vale:
\begin{equation}
\delta W_{mut} = -\mathcal{E}_1 dq_1 = M_{2\rightarrow 1}\frac{dI_2}{dt} (I_1 dt) = M_{2\rightarrow 1}I_1 dI_2.
\end{equation}
Integrando per $I_2$ che va da $0$ a $I_2$, si ottiene l'energia di interazione: $W_{mut} = M_{2\rightarrow 1}I_1I_2$.
L'energia totale immagazzinata seguendo questo percorso risulta dunque:
\begin{equation}
U_{tot}^{(A)} = \frac{1}{2}L_1I_1^2 + \frac{1}{2}L_2I_2^2 + M_{2\rightarrow 1}I_1I_2.
\end{equation}
Si ipotizzi ora di invertire il processo (accensione $2 \rightarrow 1$), raggiungendo la medesima configurazione finale. Si accende prima $\gamma_2$ (lavoro $\frac{1}{2}L_2I_2^2$) e successivamente $\gamma_1$. In questa seconda fase, sarà il generatore in $\gamma_2$ a dover compiere lavoro extra per mantenere $I_2$ costante contro la $fem$ indotta dalla variazione di corrente in $\gamma_1$ ($\mathcal{E}_2 = -M_{1\rightarrow 2}\dot{I}_1$). Con un calcolo speculare al precedente, il termine di lavoro mutuo risulta $M_{1\rightarrow 2}I_1I_2$.
L'energia totale per questo secondo percorso è:
\begin{equation}
U_{tot}^{(B)} = \frac{1}{2}L_2I_2^2 + \frac{1}{2}L_1I_1^2 + M_{1\rightarrow 2}I_1I_2.
\end{equation}
Dal momento che la configurazione finale è la medesima in entrambi i casi, si procede a uguagliare le due espressioni dell'energia ($U_{tot}^{(A)} = U_{tot}^{(B)}$) e, semplificando i termini di autoinduzione, si ottiene:
\begin{equation}
M_{2\rightarrow 1} = M_{1\rightarrow 2} = M.
\end{equation}
Si conclude che l'accoppiamento magnetico è descritto da un unico parametro geometrico $M$, detto \textit{coefficiente di mutua induzione} permettendo di scrivere l'energia magnetica totale del sistema come:
\begin{equation}
U_m = \frac{1}{2}L_1 I_1^2 + \frac{1}{2}L_2 I_2^2 + M I_1 I_2.
\end{equation}

\section{Il contributo di Maxwell: la Legge di Ampere-Maxwell}
Dall'analisi della Legge di Ampere per la circuitazione del campo magnetico, Maxwell si rese conto di un incongruenza di carattere \textit{puramente matematico}: in certe condizioni (non stazionarie), il principio di conservazione della carica veniva infatti violato. 

\subsection{La corrente di spostamento} 
Sia $\Sigma$ una superficie chiusa e stazionaria, immersa in una regione in cui è presenta una densità di corrente $\bm{J}$. Il flusso netto di carica nell'unità di tempo uscente da $\Sigma$ è:
\begin{equation}
\Phi _\Sigma = \oint_\Sigma \bm{J} \cdot d\bm{S}.
\end{equation}
Sia $q$ la carica contenuta in $\Sigma$: per il principio di conservazione, un flusso netto positivo verso l'esterno comporta una diminuzione della carica interna. Segue che:
\begin{equation}
-\frac{dq}{dt} = \oint_\Sigma \bm{J} \cdot d\bm{S}.
\label{eq:cons_carica}
\end{equation}
Vale la Legge di Gauss, per cui $q =\varepsilon_0 \oint_\Sigma \bm{E}\cdot d\bm{S}$. Sostituendo nella \eqref{eq:cons_carica} si ottiene:
\begin{equation}
\oint_\Sigma \bm{J} \cdot d\bm{S} = -\frac{d}{dt} \oint_\Sigma \varepsilon_0\bm{E}\cdot d\bm{S} = -\oint_\Sigma \varepsilon_0\frac{\partial\bm{E}}{\partial t} \cdot d\bm{S} \implies \oint_\Sigma \left(\bm{J} + \varepsilon_0\frac{\partial\bm{E}}{\partial t} \right)\cdot d\bm{S}  = 0,
\label{eq:dens_corr}
\end{equation}
dove il termine $\varepsilon_0\frac{\partial\bm{E}}{\partial t} $ è detto \textit{densità di corrente di spostamento}. Perché la \eqref{eq:dens_corr} valga per qualsiasi superficie $\Sigma$, serve che l'integranda sia identicamente nulla: il flusso netto della \textit{densità di corrente generalizzata} $\bm{J}_{tot} = \bm{J} + \varepsilon_0\frac{\partial\bm{E}}{\partial t}$, attraverso una superficie stazionaria, deve quindi essere nullo.   

\subsection{La quarta equazione di Maxwell}

\subsubsection{Il paradosso del condensatore}
Si vuole ora cercare di trovare un parallelismo tangibile della discrepanza trovata da Maxwell nella Legge di Ampere tramite l'analisi di un circuito composto da un generatore di tensione e un condensatore. Detta $\gamma$ una curva chiusa che cinge il circuito in prossimità di una delle facce del capacitore, la circuitazione del campo magnetico su $\gamma$, secondo Ampere:
\begin{equation}
\oint_\gamma \bm{B} \cdot d\bm{l} = \mu_0 I_{conc}
\end{equation}
e la corrente concatenata a $\gamma$ deve essere il flusso della densità di corrente attraverso una superficie $\Sigma$ che poggia su $\gamma$:
\begin{equation}
\mu_0 I_{conc} = \mu_0 \oint_\Sigma \bm{J} \cdot d\bm{S}
\end{equation}
e quindi:
\begin{equation}
\oint_\gamma \bm{B} \cdot d\bm{l} = \mu_0 \oint_\Sigma \bm{J} \cdot d\bm{S}.
\end{equation}
Nel limite per cui $\gamma$ viene \textit{stretta} fino a far chiudere la superficie $\Sigma$, la circuitazione:
\begin{equation}
\lim_{\gamma \rightarrow 0} \oint_\gamma \bm{B} \cdot d\bm{l} = 0 = \mu_0 \oint_\Sigma \bm{J} \cdot d\bm{S},
\end{equation}
\begin{figure}[htbp]
    \centering
    \begin{tikzpicture}[
        scale=1,
        >=Latex,
        % Stili
        current/.style={thick, ->, color=red!80!black},
        surface/.style={fill=green!20, opacity=0.6, draw=green!40!black},
        plate/.style={fill=gray!30, draw=black, thick},
        midarrow/.style={postaction={decorate, decoration={markings,mark=at position 0.6 with {\arrow{>}}}}}
    ]

        % --- COORDINATE ---
        \coordinate (WireStart) at (-5, 0);
        \coordinate (Plate1) at (-1, 0);
        \coordinate (Neck) at (-3, 0); % Dove si trova la curva gamma

        % --- 1. FILO E PIATTI (Sfondo) ---
        \draw[thick] (WireStart) -- (Plate1); % Filo
        \draw[plate] (-1, -1.5) rectangle (-0.8, 1.5); % Piatto sx (dentro la superficie)
        \draw[plate] (0.8, -1.5) rectangle (1, 1.5);   % Piatto dx (fuori)

        % --- 2. SUPERFICIE SIGMA (Sacco che avvolge il piatto) ---
        % Parte posteriore della "bocca"
        \draw[green!40!black, dashed] (Neck) ellipse (0.3 and 0.6);

        % Contorno del sacco
        % Parte superiore
        \draw[surface] (-3, 0.6) 
            .. controls (-2, 2) and (0, 2) .. (0.2, 0) % Curva sopra
            .. controls (0, -2) and (-2, -2) .. (-3, -0.6) % Curva sotto
            -- cycle; % Chiude sul retro (approssimazione visiva)
        
        % Ridisegniamo la parte davanti per l'effetto 3D
        \begin{scope}
            \clip (-3, -2) rectangle (0.5, 2);
             \fill[surface] (-3, 0.6) 
                .. controls (-2, 2) and (0, 2) .. (0.2, 0)
                .. controls (0, -2) and (-2, -2) .. (-3, -0.6)
                -- cycle;
        \end{scope}

        % --- 3. CURVA GAMMA (Il collo che si stringe) ---
        % Disegniamo l'ellisse anteriore
        \draw[red, thick, dashed] (-3, -0.6) arc (270:90:0.3 and 0.6);
        \draw[red, thick] (-3, 0.6) arc (90:-90:0.3 and 0.6); % Parte davanti solida
        
        % Frecce di restringimento (limite gamma -> 0)
        \draw[->, black, thick] (-3, 0.9) -- (-3, 0.65);
        \draw[->, black, thick] (-3, -0.9) -- (-3, -0.65);
        \node[black, font=\footnotesize] at (-3, 1.1) {$\gamma \to 0$};

        % --- 4. DETTAGLI VETTORIALI ---
        
        % Corrente che entra
        \draw[current] (-4.5, 0.2) -- (-3.5, 0.2) node[midway, above] {$I$};
        
        % Normale alla superficie (dall'altra parte rispetto a I, o uscente)
        % Nota: per il teorema della divergenza su superficie chiusa, dS è verso l'esterno.
      
        
        % Carica q interna
        \node at (-1.3, -1.8) {Carica $q(t)$};
        \draw[->, thin] (-1.3, -1.6) -- (-0.9, -1.0);

        % Etichette
        \node[red] at (-3.5, -0.5) {$\gamma$};
        \node[green!30!black] at (-0.5, 0.5) {$\Sigma$};
        \node[gray] at (0.9, 1.8) {Piatto esterno};

    \end{tikzpicture}
    \caption{La curva $\gamma$ (in rosso) si stringe attorno al filo fino a chiudersi. In questo limite, la superficie $\Sigma$ (in verde) diventa una superficie chiusa che ingloba l'armatura e la carica $q$, ma attraverso la quale passa la corrente $I$ (entrante dal "buco" che si sta chiudendo).}
    \label{fig:ragionamento_limite}
\end{figure}
evidentemente falso nel caso di una corrente non stazionaria. Nel caso analizzato, infatti, attraverso $\Sigma$ "entra" corrente ma non esce: il flusso netto non può quindi essere nullo. Maxwell comprese dunque che la Legge di Ampere mancava di un termine che correggesse l'apparente paradosso del condensatore: la soluzione fu trovata nella \textit{corrente di spostamento}, dipendente dalla variazione temporale del campo elettrico.


\subsubsection{Derivazione analitica - forma differenziale}
Dalla \eqref{eq:dens_corr} segue che il campo vettoriale della densità di corrente generalizzata è solenoidale e, dal momento che soddisfa le ipotesi del Teorema di Helmholtz\footnote{Si noti infatti che, a parte condizioni di continuità che si danno per scontate, il decadimento all'infinito è garantito sia per il campo elettrico $\dot{E} \propto r^{-2}$ sia per la densità di corrente, per la quale si assume un numero finito di sorgenti (e quindi un valore nullo nel limite per uno spazio infinito).}, deve esistere un campo $\bm{F}$ il cui rotore è proprio $\bm{J}_{tot}$:
\begin{equation}
\nabla \times \bm{F} = \bm{J}_{tot} = \bm{J} + \varepsilon_0\frac{\partial\bm{E}}{\partial t}.
\end{equation}
In condizioni stazionarie (tempo-invarianti):
\begin{equation}
\nabla \times \bm{F} = \bm{J}_{tot} = \bm{J},
\end{equation}
da cui si ricava immediatamente, dalla Legge di Ampere per la magnetostatica, $\bm{F} = \bm{B}/\mu_0$:  
\begin{equation}
\nabla \times \bm{B} = \mu_0\bm{J}_{tot} = \mu_0\bm{J} +\mu_0\varepsilon_0\frac{\partial\bm{E}}{\partial t},
\label{eq:4_max}
\end{equation}
dove $\mu_0\varepsilon_0\frac{\partial\bm{E}}{\partial t}$ è detta \textit{corrente di spostamento}. Questo risultato è noto come \textit{quarta equazione di Maxwell} o \textit{Legge di Ampere-Maxwell}. 

\subsubsection{Equazione di continuità}
La correzione di Maxwell permette di scrivere l'equazione di continuità per il caso elettromagnetico. Applicando l'operatore di divergenza ad ambo i membri della \eqref{eq:4_max} e ricordando che la divergenza di un rotore è identicamente nulla, si ottiene:
\begin{equation}
\nabla \cdot \bm{J}_{tot} = \nabla \cdot \bm{J} + \varepsilon_0 \nabla \cdot \frac{\partial\bm{E}}{\partial t} = \nabla \cdot \bm{J} + \frac{\partial\rho}{\partial t} = 0.
\end{equation}
Si osservi che, nel passaggio intermedio, è stato lecito invertire l'ordine degli operatori di derivazione spaziale e temporale (essendo le grandezze funzioni continue) e si è fatto ricorso alla prima equazione di Maxwell (Legge di Gauss, $\nabla \cdot \bm{E} = \rho/\varepsilon_0$).
Questo risultato dimostra la perfetta coerenza interna della teoria elettromagnetica: la correzione introdotta da Maxwell fa sì che il principio di conservazione della carica locale non debba essere imposto come condizione esterna, ma scaturisca come identità matematica direttamente dalle equazioni di campo. La validità della quarta equazione implica, dunque, la conservazione della carica elettrica.

\subsubsection{Forma integrale e considerazioni fisiche}
Per riottenere la formulazione globale, si integra la relazione differenziale su una superficie aperta $\Sigma$ avente come bordo la curva chiusa $\gamma$. Applicando il Teorema del Rotore (Stokes), si giunge alla \textit{forma integrale della quarta equazione di Maxwell}:
\begin{equation}
\oint_\gamma \bm{B} \cdot d\bm{l} = \mu_0 \left( I_{conc} + \varepsilon_0 \frac{d}{dt}\int_\Sigma \bm{E} \cdot d\bm{S} \right) = \mu_0 (I_{conc} + I_{d}),
\end{equation}
dove $I_d$ è la corrente di spostamento.
È importante sottolineare una differenza fisica sostanziale: mentre la corrente di conduzione $I_c$ è associata al moto ordinato di portatori di carica e, nei conduttori reali, a dissipazione di energia (effetto Joule), la corrente di spostamento $I_d$ rappresenta una variazione temporale del campo elettrico e non implica trasporto di materia né dissipazione ohmica nel vuoto. Tuttavia, entrambe le grandezze sono sorgenti equivalenti del campo magnetico. 

\chapter{Le equazioni di Maxwell: il caso generale}

In questo capitolo si giunge alla sintesi finale dell'elettromagnetismo classico. Abbandonando l'ipotesi di stazionarietà (ovvero ammettendo che $\rho$ e $\bm{J}$ possano variare nel tempo), si osserva che i campi elettrico e magnetico cessano di essere entità distinte e indipendenti. Essi si fondono in un'unica entità fisica, il \textit{campo elettromagnetico}, in cui la variazione temporale di uno agisce come sorgente per l'altro.

\section{Quadro sinottico tempo-variante}

Il sistema completo delle quattro equazioni differenziali lineari accoppiate è riportato nella Tabella (\ref{tab:maxwell_dinamiche}). Si noti come le prime due equazioni (quelle di Gauss) rimangano formalmente identiche al caso statico, mentre le leggi di circuitazione (rotori) acquisiscano termini dipendenti dal tempo che rompono la simmetria statica.

\begin{table}[htbp]
    \centering
    \renewcommand{\arraystretch}{2.4}
    \begin{tabular}{p{4cm} c c}
        \hline
        \textit{Legge fisica} & \textit{Forma Differenziale} & \textit{Forma Integrale} \\
        \hline
        
        \textit{1. Legge di Gauss} \newline (Elettrica) & 
        $\nabla \cdot \bm{E} = \dfrac{\rho}{\varepsilon_0}$ & 
        $\displaystyle \oint_{\Sigma} \bm{E} \cdot d\bm{S} = \frac{Q_{int}}{\varepsilon_0}$ \\
        
        \textit{2. Legge di Gauss} \newline (Magnetica) & 
        $\nabla \cdot \bm{B} = 0$ & 
        $\displaystyle \oint_{\Sigma} \bm{B} \cdot d\bm{S} = 0$ \\
        
        \textit{3. Legge di Faraday} \newline (Induzione) & 
        $\nabla \times \bm{E} = -\dfrac{\partial \bm{B}}{\partial t}$ & 
        $\displaystyle \oint_{\gamma} \bm{E} \cdot d\bm{l} = - \frac{d\Phi(\bm{B})}{dt}$ \\
        
        \textit{4. Legge di Ampère-Maxwell} & 
        $\nabla \times \bm{B} = \mu_0 \bm{J} + \mu_0\varepsilon_0 \dfrac{\partial \bm{E}}{\partial t}$ & 
        $\displaystyle \oint_{\gamma} \bm{B} \cdot d\bm{l} = \mu_0 (I_{conc} + I_{s})$ \\
        \hline
    \end{tabular}
    \caption{Le equazioni di Maxwell nel vuoto in forma completa. Il termine $I_s = \varepsilon_0 \frac{d\Phi_E}{dt}$ rappresenta la corrente di spostamento.}
    \label{tab:maxwell_dinamiche}
\end{table}

\section{Le implicazioni della dinamica}
L'introduzione delle derivate temporali porta a conseguenze fisiche di portata rivoluzionaria:
\begin{itemize}
    \item \textit{La propagazione ondosa:} manipolando le equazioni 3 e 4 nel vuoto ($\rho=0, \bm{J}=0$), si ottengono equazioni delle onde per $\bm{E}$ e $\bm{B}$. Ciò implica che i campi possono staccarsi dalle sorgenti e propagarsi nello spazio alla velocità della luce $c = 1/\sqrt{\varepsilon_0\mu_0}$\footnote{Come verrà successivamente discusso.}.
    \item \textit{La conservazione della carica:} la quarta equazione corretta garantisce la validità dell'equazione di continuità $\nabla \cdot \bm{J} + \partial \rho / \partial t = 0$ anche in regime non stazionario, cosa impossibile con la sola legge di Ampère originale.
\end{itemize}

\subsection{Appendice: L'audacia matematica della corrente di spostamento}

La quarta equazione, nota oggi come legge di Ampère-Maxwell, contiene il contributo cruciale di James Clerk Maxwell: la \textit{corrente di spostamento}:
\begin{equation}
I_s = \mu_0\bm{J}_s = \mu_0\varepsilon_0 \frac{\partial \bm{E}}{\partial t}.
\end{equation}
È interessante notare come l'introduzione di questo termine non nacque da un'osservazione sperimentale (come fu per Faraday o Ampère), ma da una necessità puramente teorica e matematica, generando non poco scetticismo nella comunità scientifica dell'epoca.

\begin{tcolorbox}[colback=gray!10, colframe=black!70, title=Il "paradosso" e la simmetria mancante]
    Maxwell si accorse di una falla logica nella legge di Ampère originale ($\nabla \times \bm{B} = \mu_0 \bm{J}$). Applicando l'operatore divergenza a entrambi i membri, si ottiene:
    \[ \nabla \cdot (\nabla \times \bm{B}) = \mu_0 (\nabla \cdot \bm{J}). \]
    Poiché la divergenza di un rotore è identicamente nulla, ne risulterebbe $\nabla \cdot \bm{J} = 0$, condizione vera solo per correnti stazionarie. Per processi variabili (es. la carica di un condensatore), dove $\nabla \cdot \bm{J} = -\partial \rho/\partial t \neq 0$, la legge di Ampère portava a un assurdo matematico.
    
    Maxwell aggiunse il termine $\mu_0\varepsilon_0 \frac{\partial \bm{E}}{\partial t}$ "a tavolino" per ripristinare la coerenza con l'equazione di continuità, guidato da un forte senso di simmetria ed eleganza formale.
\end{tcolorbox}



\paragraph{Perché nessuno l'aveva vista?}
L'audacia di Maxwell risiede nel fatto che, in condizioni di laboratorio standard per l'epoca, il termine di spostamento è praticamente invisibile.
Considerando che $\mu_0\varepsilon_0 = 1/c^2$, il coefficiente che moltiplica la variazione del campo elettrico è piccolissimo ($\approx 10^{-17}$).
Affinché la corrente di spostamento ($\bm{J}_s$) sia paragonabile alla corrente di conduzione ($\bm{J}$), sono necessarie frequenze di oscillazione elevatissime, irraggiungibili con gli strumenti dell'Ottocento.
Fu solo vent'anni dopo la morte di Maxwell, con gli esperimenti di Hertz sulle onde elettromagnetiche (frequenze molto alte), che l'esistenza fisica di questo termine fu confermata, dimostrando che la "correzione matematica" descriveva una realtà fisica fondamentale: la luce stessa è un fenomeno elettromagnetico.





\chapter{Le onde elettromagnetiche}
Evidentissimo eppur nascosta agli occhi di tutti è la fondamentale conseguenza delle equazioni di Maxwell: il campo elettrico e quello magnetico non sono altro che la manifestazione di un unica entità di carattere intrinsecamente \textit{ondulatorio}. In questo capitolo, quindi, si cerca di evidenziare la stretta correlazione fra la formulazione di Maxwell e l'equazione delle onde e le principali conseguenze fisiche.

\section{Da Maxwell a d'Alembert}
Sufficientemente lontano dalle sorgenti dei campi, in regioni in cui si può assumere $\rho = 0$ e $\bm{J} = \bm{0}$, la terza e la quarta equazione di Maxwell assumono una forma altamente simmetrica:
\begin{equation}
\begin{dcases}
\nabla \times \bm{E} = -\frac{\partial\bm{B}}{\partial t} \\
\nabla \times \bm{B} = \mu_0\varepsilon_0\frac{\partial\bm{E}}{\partial t}
\end{dcases}
.
\label{eq:max_ver}
\end{equation}
Si applica a questo punto l'operatore rotore ad ambo i membri delle due equazioni. Ricordando l'identità vettoriale $\nabla \times (\nabla \times \bm{A}) = \nabla (\nabla \cdot \bm{A}) - \nabla^2 \bm{A}$ e notando che nel vuoto le divergenze sono nulle ($\nabla \cdot \bm{E} = 0$, $\nabla \cdot \bm{B} = 0$), si ottiene:
\begin{equation}
\begin{dcases}
- \nabla^2 \bm{E} = - \frac{\partial }{\partial t} (\nabla \times\bm{B}) \\
- \nabla^2 \bm{B} = \mu_0\varepsilon_0\frac{\partial}{\partial t}(\nabla \times\bm{E})
\end{dcases}
.
\end{equation}
Sostituendo i rotori presenti al membro destro tramite le \eqref{eq:max_ver}, si giunge a:
\begin{equation}
\begin{dcases}
\nabla^2 \bm{E} = \mu_0\varepsilon_0\frac{\partial^2 \bm{E}}{\partial t^2} \\
\nabla^2 \bm{B} = \mu_0\varepsilon_0\frac{\partial^2 \bm{B}}{\partial t^2}
\end{dcases}
.
\label{eq:wave_eq}
\end{equation}
Si è dunque ottenuto il disaccoppiamento dei campi: ogni componente vettoriale di $\bm{E}$ e $\bm{B}$ soddisfa autonomamente l'\textit{equazione di d'Alembert} per le onde. 

\begin{approfondimento}[
    colback=white,
    colframe=blue!50!black, % O il colore che preferisci (es. green!35!black)
    boxed title style={
        colback=blue!50!black,
        sharp corners,
        frame hidden
    },
    title=Caratteristiche di un'onda]
Una qualsiasi onda, ovvero una perturbazione che si propaga nello spazio e nel tempo, è contraddistinta da una direzione di propagazione e da una velocità di fase $v$ con cui avanza nel in un generico mezzo. Oltre a queste caratteristiche cinematiche, una perturbazione periodica (in particolare un'onda armonica) è univocamente determinata dai seguenti parametri fisici:
\begin{itemize}
    \item \textit{Periodo temporale} $T$: rappresenta l'intervallo di tempo necessario affinché un punto dello spazio compia un'oscillazione completa. È la periodicità nel tempo della funzione d'onda.
    
    \item \textit{Frequenza} $\nu$: definita come l'inverso del periodo ($\nu = 1/T$), rappresenta il numero di oscillazioni compiute nell'unità di tempo. Si misura in Hertz ($\si{\hertz}$).
    
    \item \textit{Pulsazione} $\omega$: è la frequenza angolare, legata alla frequenza dalla relazione $\omega = 2\pi\nu$. È il parametro naturale per descrivere l'evoluzione temporale nelle funzioni trigonometriche.
    
    \item \textit{Lunghezza d'onda} $\lambda$: rappresenta la distanza minima tra due punti dello spazio che vibrano in fase (ad esempio due massimi consecutivi). È la periodicità spaziale dell'onda.
    
    \item \textit{Numero d'onda} $k$ (o vettore d'onda $\bm{k}$): analogamente alla pulsazione per il tempo, il numero d'onda rappresenta la frequenza spaziale, definita come $k = 2\pi/\lambda$.
    In forma vettoriale, il \textit{vettore d'onda} $\bm{k}$ ha modulo $k$ e direzione e verso concordi alla direzione di propagazione.
    
    \item \textit{Ampiezza} $A$: corrisponde al massimo spostamento della grandezza oscillante rispetto alla posizione di equilibrio. Nel caso elettromagnetico, corrisponde al modulo massimo dei campi $\bm{E}_0$ e $\bm{B}_0$.
    
    \item \textit{Fase} $\varphi$: rappresenta lo stato dell'oscillazione all'istante iniziale ($t=0$) e nell'origine spaziale ($\bm{r}=0$).
\end{itemize}
    Per le onde vettoriali, come quelle elettromagnetiche, è necessario definire un ulteriore parametro: la \textit{polarizzazione}. Essa descrive la direzione di oscillazione del vettore campo ($\bm{E}$) rispetto alla direzione di propagazione $\bm{\kappa}$. Le onde elettromagnetiche nel vuoto sono trasversali, ovvero $\bm{E}$ è sempre ortogonale a $\bm{\kappa}$.
\end{approfondimento}


\subsubsection{La velocità della luce}
L'equazione delle onde classica per una funzione $\Psi$ scalare di variabile vettore è $\nabla^2 \Psi = \frac{1}{v^2}\frac{\partial^2 \Psi}{\partial t^2}$. Confrontando questa struttura con la \eqref{eq:wave_eq}, si identifica immediatamente la velocità di propagazione $v$ delle onde elettromagnetiche:
\begin{equation}
v = \frac{1}{\sqrt{\varepsilon_0 \mu_0}}.
\end{equation}
Inserendo i valori numerici delle costanti del vuoto ($\varepsilon_0 \approx \SI{8.85e-12}{\farad\per\meter}$ e $\mu_0 = 4\pi \cdot 10^{-7} \, \si{\henry\per\meter}$):
\begin{equation}
c = \frac{1}{\sqrt{\varepsilon_0 \mu_0}} \approx \SI{2.998e8}{\meter\per\second}.
\end{equation}
Questo risultato, che coincide con la velocità della luce misurata sperimentalmente nell'ottica, portò Maxwell a concludere che la luce stessa non è altro che un'onda elettromagnetica.

\subsection{La relazione fra $\bm{E}$ e $\bm{B}$}
Si consideri un'onda piana che si propaga nella direzione individuata dal versore $\bm{u}_n$ con velocità $c$. Il campo elettrico, ad esempio, può essere descritto senza perdita di generalità da una funzione vettoriale arbitraria $\bm{E}$ dipendente da un'unica variabile scalare $\xi$ (fase totale), combinazione lineare di spazio e tempo:
\begin{equation}
\bm{E}(\bm{r}, t) = \bm{E}(\xi), \quad \text{con } \xi = \bm{r} \cdot \bm{u}_n - ct.
\label{eq:onde3d}
\end{equation}
Per determinare il campo magnetico associato, si applica la Legge di Faraday:
\begin{equation}
\nabla \times \bm{E} = -\frac{\partial \bm{B}}{\partial t}.
\end{equation}
Si calcola innanzitutto il rotore del campo elettrico utilizzando la regola della catena per gli operatori differenziali. Indicando con l'apice la derivata rispetto alla variabile $\xi$ ($\bm{E}' = d\bm{E}/d\xi$):
\begin{equation}
\nabla \times \bm{E}(\xi) = (\nabla \xi) \times \frac{d\bm{E}}{d\xi} = \bm{u}_n \times \bm{E}'.
\end{equation}
Successivamente, si esprime la derivata temporale del campo elettrico in funzione della stessa derivata $\bm{E}'$:
\begin{equation}
\frac{\partial \bm{E}}{\partial t} = \frac{d\bm{E}}{d\xi} \frac{\partial \xi}{\partial t} = \bm{E}' (-c) \implies \bm{E}' = -\frac{1}{c} \frac{\partial \bm{E}}{\partial t}.
\end{equation}
Sostituendo l'espressione di $\bm{E}'$ nella formula del rotore precedentemente calcolata, si ottiene:
\begin{equation}
\nabla \times \bm{E} = \bm{u}_n \times \left( -\frac{1}{c} \frac{\partial \bm{E}}{\partial t} \right) = -\frac{1}{c} \frac{\partial}{\partial t} (\bm{u}_n \times \bm{E}).
\end{equation}
Confrontando questo risultato con la Legge di Faraday ($\nabla \times \bm{E} = -\partial_t \bm{B}$), si eguagliano i termini temporali:
\begin{equation}
-\frac{\partial \bm{B}}{\partial t} = -\frac{1}{c} \frac{\partial}{\partial t} (\bm{u}_n \times \bm{E}).
\end{equation}
Integrando rispetto al tempo (e trascurando eventuali campi statici sovrapposti, non pertinenti alla propagazione ondosa), si giunge alla relazione vettoriale fondamentale:
\begin{equation}
\bm{B} = \frac{1}{c} (\bm{u}_n \times \bm{E}).
\label{eq:onda_em_B}
\end{equation}
Questa equazione porta con se un nuovo modo di intendere campi elettrici e magnetici: $\bm{E}$ e $\bm{B}$ sono infatti ovunque perpendicolari tra loro ($\bm{B} \perp \bm{E}$). Ancora, valutando la divergenza della generica funzione che descrive il campo elettrico e imponendo il risultato di Gauss nel vuoto $\nabla \cdot \bm{E} = 0$, si ottiene:
\begin{equation}
\nabla \cdot \bm{E} = 0 = \nabla(\xi) \cdot \frac{d\bm{E}}{d\xi} = \bm{u}_n \cdot \bm{E}',
\end{equation}
dove $\bm{E}'$ rappresenta la variazione dell'onda. La Legge di Gauss garantisce quindi che la variazione del campo elettrico può avvenire solo perpendicolarmente (prodotto scalare nullo) alla direzione di propagazione dell'onda, con ciò intendendo che la componente di campo parallela alla direzione di propagazione è costante. Dal momento che un'onda è definita come la perturbazione di una grandezza, si deduce che anche la componente elettrica dell'onda elettromagnetica deve essere perpendicolare a $\bm{u}_n$. 

\subsubsection{Campi mediati dalla velocità della luce}
Essendo $\bm{E} \perp \bm{B} \perp \bm{u}_n$, i tre vettori generano una terna destrorsa:
\begin{equation}
\{\bm{E}, \bm{B}, \bm{u}_n\}.
\end{equation}
Sfruttando questa relazione, la \eqref{eq:onda_em_B} e valutandone il modulo:
\begin{equation}
|\bm{B}| = \frac{1}{c} |\bm{u}_n \times \bm{E}| \implies B = \frac{1}{c}E.
\end{equation}
Il rapporto dell'ampiezza massima della perturbazione della componente elettrica e magnetica dell'onda è dunque equivalente alla velocità dell'onda ($=c$ nel vuoto).

\section{Le onde armoniche piane}
\subsection{Onde armoniche}
Si introduce una prima semplificazione alla trattazione delle onde che non nuoce alla generalità: si analizzeranno d'ora in poi fenomeni perturbatori che si propagano lungo un asse privilegiato, quello delle ascisse. In questo contesto, la \eqref{eq:onde3d} si semplifica e le funzioni che descrivono l'evoluzione spazio temporale dei campi elettrico e magnetico sono scalari:
\begin{equation}
\begin{cases}
E(x,t) = f_E (x \pm ct) \\
B(x,t) = f_B (x \pm ct) 
\end{cases}
,
\end{equation}
dove $f_E$ e$f_B$ sono funzioni generiche che hanno come argomento la \textit{fase} $x \pm ct$ dell'onda.
Ecco una seconda approssimazione: le funzioni d'onda sono sinusoidi \textit{progressive}\footnote{La progressività di un'onda corrisponde ad un evoluzione nel verso positivo dell'asse di propagazione con l'avanzare del tempo. Caratteristica denotabile nel segno $-$ davanti al fattore $ct$.}. Con il fine di rendere adimensionato il loro argomento, si introduce il fattore $\kappa$, detto \textit{numero d'onda}\footnote{In una trattazione multidimensionale, il numero d'onda è in realtà un vettore detto \textit{vettore d'onda} $\bm{\kappa}$.}:
\begin{equation}
\begin{cases}
E(x,t) = E_0 \sin [\kappa(x-ct)] \\
B(x,t) = B_0 \sin[\kappa(x-ct)] 
\label{sis:EB}
\end{cases}
,
\end{equation}
con $E_0$ e $B_0$ le ampiezze di massima oscillazione\footnote{Si ricorda che vale $E_0/B_0 = c$.}. Le sinusoidi sono contraddistinte da una \textit{lunghezza d'onda} caratteristica, detta $\lambda$, che corrisponde allo spazio percorso dalla perturbazione dopo un'oscillazione di $2\pi$ e per la quale vale che $\sin(x) = \sin(x+\lambda)$:
\begin{equation}
\sin [\kappa(x+\lambda-ct)] = \sin [\kappa(x-ct)] 
\end{equation}
\begin{equation}
\implies \kappa(x+\lambda-ct) =  \kappa(x-ct) + 2\pi
\end{equation}
da cui si ottiene: 
\begin{equation}
 \kappa =  \frac{2\pi}{\lambda}.
\end{equation}
La \eqref{sis:EB} così diventa: 
\begin{equation}
\begin{cases}
E(x,t) = E_0 \sin (\kappa x-\kappa ct) =E_0 \sin (\kappa x- \frac{2\pi}{\lambda} ct) = E_0\sin (\kappa x- \omega t) \\
B(x,t) = B_0 \sin (\kappa x-\kappa ct) =B_0 \sin (\kappa x- \omega t)
\label{sis:EB_k}
\end{cases}
,
\end{equation}
dove si è definito $c/\lambda = \nu$, con $\nu = 1/T$ la \textit{frequenza di oscillazione}, e $2\pi \nu = \omega$ la \textit{pulsazione angolare}.

\begin{approfondimento}[
    colback=white,
    colframe=blue!50!black, % O il colore che preferisci (es. green!35!black)
    boxed title style={
        colback=blue!50!black,
        sharp corners,
        frame hidden
    },
    title={Perdita di generalità? L'analisi di Fourier e l'isotropia}
]
La scelta di limitare lo studio a un'onda \textit{armonica} (sinusoidale) che si propaga lungo un \textit{singolo asse} (monodimensionale) potrebbe apparire riduttiva. Tuttavia, due argomentazioni fondamentali assicurano che i risultati ottenuti siano applicabili a qualsiasi perturbazione fisica reale.

\begin{enumerate}
    \item \textit{Teorema di Fourier e Linearità:}
    in natura non esistono onde perfettamente sinusoidali ed estese all'infinito. Tuttavia, l'analisi armonica (serie e trasformata di Fourier) garantisce che qualsiasi funzione d'onda "fisica" (limitata e regolare) possa essere scomposta in una somma (o integrale) di infinite componenti sinusoidali, ciascuna con la propria ampiezza e frequenza.
    Poiché le equazioni di Maxwell (nel vuoto o in mezzi lineari) sono \textit{lineari}, vale il \textit{Principio di Sovrapposizione}: il comportamento di un impulso complesso è semplicemente la somma dei comportamenti delle sue componenti elementari. Risolvere per la sinusoide significa dunque risolvere per qualunque onda.

    \item \textit{Isotropia dello spazio:}
    un'onda piana è definita da un vettore d'onda $\bm{\kappa}$ che ne indica la direzione di propagazione. Lo spazio vuoto è \textit{isotropo} (le leggi fisiche non dipendono dall'orientamento). È pertanto sempre possibile operare una rotazione del sistema di riferimento cartesiano tale da allineare l'asse $x$ con la direzione di $\bm{\kappa}$. In questo nuovo riferimento ruotato, il problema tridimensionale si riduce formalmente a uno monodimensionale ($\bm{\kappa} \cdot \bm{r} = \kappa x$) senza alterare la fisica sottostante.
\end{enumerate}

\end{approfondimento}

\subsection{Onde piane}
Dalla \eqref{sis:EB_k} si ha che, fissato un generico punto $x_0$ nello spazio e un generico istante $t$ nel tempo, il valore di $E(x,t)$ è costante per ogni punto di coordinate $[x_0, y, z]$, con $y,z \in \mathbb{R}^3$, quindi sul piano di equazione $x = x_0$. Da qui la dicitura \textit{onda piana}. Si segnala che, essendo il piano un'entità geometrica illimitata, un'onda piana ideale costituisce un'\textit{astrazione matematica} fisicamente irrealizzabile nella sua interezza. Se il campo elettrico mantenesse ampiezza costante su una superficie infinita, la densità di energia totale trasportata dall'onda risulterebbe divergente (infinita), condizione inaccettabile per qualsiasi sorgente reale.
Tuttavia, il modello dell'onda piana rappresenta un'approssimazione locale eccellente e indispensabile per la descrizione dei fenomeni fisici reali.
Le sorgenti fisiche (come un'antenna o un atomo) sono limitate nello spazio ed emettono tipicamente \textit{onde sferiche}, i cui fronti d'onda sono superfici sferiche concentriche che si espandono radialmente.
Quando l'osservatore si pone a una distanza $r$ molto grande rispetto alle dimensioni della sorgente e alla lunghezza d'onda ($\lambda$), il raggio di curvatura del fronte d'onda diviene così elevato da rendere la superficie localmente indistinguibile da un piano.
In sintesi, un'onda sferica osservata a grande distanza (o in una regione di spazio limitata) degenera asintoticamente in un'onda piana.
Dal punto di vista analitico, questo corrisponde a trascurare i termini di ordine superiore nello sviluppo della curvatura, permettendo di utilizzare la semplice formulazione matematica sinusoidale della \eqref{sis:EB_k}.


\section{Energia e quantità di moto di un'onda}
In un sistema isolato, un'onda elettromagnetica esercita una forza sulle particelle cariche che incontra nel suo propagarsi. In figura (\ref{fig:press_rad}) è riportata una perturbazione che si propaga lungo l'asse delle ascisse: nell'incidere sul piano $x = 0$, incontra una particella carica $q$. 
\begin{figure}[h]
    \centering
\begin{tikzpicture}[scale=1, >=Stealth, x={(-0.7cm,-0.4cm)}, y={(1cm,0cm)}, z={(0cm,1cm)}]

    % --- Definizioni Parametri ---
    \def\waveStart{2.5*pi} % Punto di inizio dell'onda (a destra) % Posizione Y della particella
    \def\amplitude{1.5}    % Ampiezza dell'onda

    % --- Assi Cartesiani ---
    % Disegno l'asse X negativo (direzione di provenienza)
    \draw[->] (9,0,0) -- (-1,0,0) node[above left] {$x$};
    % Asse Y
    \draw[->] (0,-2,0) -- (0,2.5,0) node[above] {$z$};
    % Asse Z (necessario per il campo B)
    \draw[->] (0,0,-2) -- (0,0,2.5) node[above right] {$y$};
    \node[below right] at (0,0,0) {O};

    % --- La Particella Carica ---
    % Posizionata sull'asse Y
    \coordinate (Particle) at (0, 0, 0);
    \shade[ball color=orange] (Particle) circle (0.15cm);
\node[right=0.2cm, orange, font=\bfseries] at (Particle)  {$q$};

    % --- Onda Elettromagnetica ---
    % Disegniamo l'onda che arriva da x positivo e colpisce la particella a x=0

    % 1. Campo Magnetico (B) - Piano XZ (Rosso)
    % Disegno la curva sinusoidale
    \draw[red!80, thick, smooth, samples=200, domain=0:\waveStart]
        plot (\x, 0, {\amplitude*sin(\x r + 90)}); % +90 per farla arrivare col picco a x=0

    % Disegno i vettori del campo B lungo l'onda
    \foreach \x in {0, 2*pi} {
         \draw[->, red, thick] (\x, 0, 0) -- (\x, 0, \amplitude);
    }
    \foreach \x in {0.5*pi, 2.5*pi} {
         \draw[->, red, thick] (\x, 0, 0) -- (\x, 0, -\amplitude);
    }
    % Etichetta campo B
    \node[red, above] at (2*pi, 0, \amplitude) {$\bm{E}$};


    % 2. Campo Elettrico (E) - Piano XY (Blu)
    % Disegno la curva sinusoidale
    \draw[blue!80, very thick, smooth, samples=200, domain=0:\waveStart]
        plot (\x, {\amplitude*sin(\x r + 90)}, 0);

    % Disegno i vettori del campo E lungo l'onda
    % Vettori positivi
    \foreach \x in {0, 2*pi} {
        \draw[->, blue, very thick] (\x, 0, 0) -- (\x, \amplitude, 0);
    }
    % Vettori negativi
    \foreach \x in {0.5*pi, 2.5*pi} {
        \draw[->, blue, very thick] (\x, 0, 0) -- (\x, -\amplitude, 0);
    }
    % Etichetta campo E
    \node[blue, below] at (2*pi, \amplitude, 0) {$\bm{B}$};

\end{tikzpicture}
\label{fig:press_rad}
\caption{Un'onda elettromagnetica piana sinusoidale che incide sull'asse $x = 0$ dove si trova un particella carica $q$.}
\end{figure}
Dall'interazione, si registra una forza di Lorentz del tipo
\begin{equation}
\bm{F} = \bm{F}_m + \bm{F}_e = q(\bm{v} \times \bm{B} + \bm{E}).
\end{equation}
La forza elettrostatica $\bm{F}_e = qE \bm{u}_y$, dal Teorema dell'Energia Cinetica, modifica l'energia $U$ della particella mentre la forza magnetica $\bm{F}_m = qvB (\bm{u}_x)$ induce una variazione della quantità di moto\footnote{Si ricorda che la forza magnetica non compie lavoro.} $p_x$:
\begin{equation}
\begin{dcases}
\frac{\partial U}{\partial t} = qEv \\
\frac{\partial p_x}{\partial t} = qvB
\end{dcases}
.
\end{equation}
Dal momento che $B = E/c$:
\begin{equation}
\frac{\partial p_x}{\partial t} = qvB = qE \frac{v}{c} = \frac{1}{c} \frac{\partial U}{\partial t} \implies U = p_xc.
\label{eq:E_vs_p}
\end{equation}
La derivazione presentata, data l'isotropia dello spazio vuoto, è generalizzabile per ogni direzione di propagazione. La particella, quindi, assorbe energia e vede variare la sua quantità di moto. Essendo il sistema isolato per ipotesi e poiché valgono i teoremi di conservazione dell'energia e della quantità di moto, si desume che l'onda, per cedere simultaneamente energia e quantità di moto secondo la \eqref{eq:E_vs_p}, deve \textit{possedere essa stessa una quantità di moto intrinseca} pari a $U/c$. 

\subsection{Densità di energia di un'onda elettromagnetica}
Come discusso nei capitoli precedenti, ai campi elettrico e magnetico è associata una densità di energia, indicata rispettivamente con $w_E$ e $w_B$.
Nello spazio libero permeato da un'onda elettromagnetica, la densità di energia totale $w$ è data dalla sovrapposizione dei contributi dei due campi:
\begin{equation}
w = w_e + w_m = \frac{1}{2} \varepsilon_0 E^2 + \frac{1}{2} \frac{B^2}{\mu_0}.
\end{equation}
Ricordando la relazione tra i moduli dei campi nel vuoto $B = E/c$ e la definizione della velocità della luce $c = 1/\sqrt{\varepsilon_0\mu_0}$, si osserva che i contributi energetici elettrico e magnetico sono identici:
\begin{equation}
w_m = \frac{1}{2\mu_0} \left( \frac{E}{c} \right)^2 = \frac{1}{2} \frac{E^2}{\mu_0} (\varepsilon_0 \mu_0) = \frac{1}{2} \varepsilon_0 E^2 = w_e.
\end{equation}
La densità totale risulta dunque il doppio della singola componente elettrica (o magnetica):
\begin{equation}
w = \varepsilon_0 E^2 = \frac{B^2}{\mu_0}.
\label{eq:ene_dens}
\end{equation}

\subsection{La pressione di radiazione}
La quantità di moto è, nella meccanica classica, una grandezza fisica intrinsecamente legata alla massa ($\bm{p} = m\bm{v}$). Un'onda elettromagnetica, pur essendo priva di massa, trasporta quantità di moto e può trasferirla alla materia interagente.
Si consideri un'onda piana che incide perpendicolarmente su una superficie assorbente. In un intervallo di tempo $\Delta t$, il fronte d'onda avanza di una distanza $h = c \Delta t$. Considerando una sezione di area $A$, l'onda occupa un volume cilindrico ideale $\Omega$ di base $A$ e altezza $h$.
L'energia totale $U$ contenuta in tale volume è il prodotto della densità volumetrica $w$ per il volume stesso:
\begin{equation}
U = w \cdot (A h) = w A c \Delta t.
\end{equation}
Sfruttando la relazione precedentemente dimostrata $p = U/c$, la quantità di moto contenuta nel volume è:
\begin{equation}
p = \frac{U}{c} = w A \Delta t.
\end{equation}
Per il Teorema dell'Impulso, la forza media $F$ esercitata dall'onda sulla superficie (assumendo che la quantità di moto venga ceduta completamente nell'intervallo $\Delta t$) è pari alla variazione della quantità di moto nel tempo:
\begin{equation}
F = \frac{\Delta p}{\Delta t} = \frac{w A \Delta t}{\Delta t} = wA.
\end{equation}
La \textit{pressione di radiazione} $P_{rad}$, definita come il rapporto tra la forza normale e la superficie d'incidenza, risulta infine:
\begin{equation}
P_{rad} = \frac{F}{A} = w.
\end{equation}
Si nota quindi che la pressione di radiazione esercitata da un'onda su una superficie perfettamente assorbente (corpo nero) coincide numericamente con la sua densità di energia.

\begin{approfondimento}[
    colback=white,
    colframe=gray!50!black,
    boxed title style={colback=gray!50!black, frame hidden, sharp corners},
    title={Nota su assorbimento e riflessione}
]
Il risultato $P_{rad} = w$ è valido nel caso di \textit{assorbimento completo} (la quantità di moto finale dell'onda è nulla).
Se la superficie è perfettamente \textit{riflettente}, l'onda rimbalza invertendo il verso di propagazione. La variazione di quantità di moto risulta doppia ($\Delta p = p_{in} - (-p_{in}) = 2p_{in}$), e di conseguenza la pressione di radiazione raddoppia:
\[ P_{rad, rifl} = 2w. \]
\end{approfondimento}

\section{Bilancio energetico del campo elettromagnetico}
Dopo aver caratterizzato la densità di energia associata ai campi elettrico e magnetico ed aver introdotto il concetto di trasporto energetico per mezzo della radiazione, risulta necessario generalizzare la trattazione per formulare un principio di conservazione globale valido in condizioni non stazionarie.
In un sistema fisico generico, i campi non sono entità statiche: essi evolvono nel tempo, interagiscono con la materia esercitando forze sulle cariche (compiendo lavoro meccanico) e propagano energia attraverso lo spazio. Affinché il principio di conservazione dell'energia sia rispettato, deve sussistere un preciso bilancio tra la variazione dell'energia immagazzinata nei campi, il lavoro compiuto sulle sorgenti di corrente e il flusso di energia che attraversa la superficie di confine del sistema.
La formalizzazione matematica di tale bilancio, che assume la struttura di un'equazione di continuità, è nota come \textit{Teorema di Poynting}.

\subsection{Il teorema di Poynting}
Si consideri una regione di spazio di volume $V$ delimitata da una superficie chiusa $\Sigma$, al cui interno sia presente una distribuzione di carica con densità $n$. Quando un campo elettromagnetico attraversa tale regione, esso interagisce con i portatori di carica esercitando su di essi una forza.
La potenza erogata dal campo alla singola carica risulta:
\begin{equation}
\mathcal{P} = \frac{dW}{dt} = \bm{F} \cdot \bm{v} = q(\bm{E} + \bm{v} \times \bm{B}) \cdot \bm{v}.
\end{equation}
Il termine magnetico non compie lavoro. La potenza trasferita alla singola carica si riduce a:
\begin{equation}
\mathcal{P} = q\bm{E} \cdot \bm{v}.
\end{equation}
La densità di potenza erogata per unità di volume deve quindi valere $n q\bm{E} \cdot \bm{v} = \bm{E} \cdot \bm{J}$. Integrando sull'intero volume $V$, la potenza totale meccanica trasferita dal campo alla materia risulta:
\begin{equation}
\mathcal{P}_{mec} = \int_V \bm{E} \cdot \bm{J} \, dV.
\label{eq:potenza_joule}
\end{equation}
Per esprimere questa grandezza in funzione dei soli campi, si ricorre alla legge di Ampere-Maxwell ($\bm{J} = \frac{1}{\mu_0}\nabla \times \bm{B} - \varepsilon_0 \frac{\partial \bm{E}}{\partial t}$):
\begin{equation}
\mathcal{P}_{mec} = \int_V \left( \frac{1}{\mu_0} \bm{E} \cdot (\nabla \times \bm{B}) - \varepsilon_0 \bm{E} \cdot \frac{\partial \bm{E}}{\partial t} \right) dV.
\end{equation}
Si ricordi l'identità vettoriale per la divergenza di un prodotto vettoriale:
\begin{equation}
\nabla \cdot (\bm{E} \times \bm{B}) = \bm{B} \cdot (\nabla \times \bm{E}) - \bm{E} \cdot (\nabla \times \bm{B}) \implies \bm{E} \cdot (\nabla \times \bm{B}) = \bm{B} \cdot (\nabla \times \bm{E}) - \nabla \cdot (\bm{E} \times \bm{B}).
\end{equation}
Sostituendo nell'integrale e sfruttando la legge di Faraday ($\nabla \times \bm{E} = -\frac{\partial \bm{B}}{\partial t}$):
\begin{equation}
\mathcal{P}_{mec} = \int_V \left[ -\frac{1}{\mu_0} \bm{B} \cdot \frac{\partial \bm{B}}{\partial t} - \frac{1}{\mu_0} \nabla \cdot (\bm{E} \times \bm{B}) - \varepsilon_0 \bm{E} \cdot \frac{\partial \bm{E}}{\partial t} \right] dV.
\end{equation}
Raggruppando i termini e notando che $\bm{A} \cdot \partial_t \bm{A} = \frac{1}{2}\partial_t A^2$:
\begin{equation}
\int_V \bm{E} \cdot \bm{J} \, dV = - \frac{\partial}{\partial t} \int_V \underbrace{\left( \frac{1}{2}\varepsilon_0 E^2 + \frac{B^2}{2\mu_0} \right)}_{w} dV - \frac{1}{\mu_0} \int_V \nabla \cdot (\bm{E} \times \bm{B}) \, dV.
\end{equation}
Applicando il Teorema della Divergenza (o di Gauss) all'ultimo termine, si giunge alla formulazione finale del \textit{Teorema di Poynting}:
\begin{equation}
\int_V \bm{E} \cdot \bm{J} \, dV + \oint_{\Sigma} \frac{\bm{E} \times \bm{B}}{\mu_0} \cdot d\bm{S} = - \frac{\partial}{\partial t} \int_V w \, dV.
\label{eq:poynting_finale}
\end{equation}

\subsubsection{Significato fisico e vettore di Poynting}
L'equazione \eqref{eq:poynting_finale} rappresenta la legge generale di \textit{conservazione dell'energia} per il campo elettromagnetico. Essa stabilisce che la variazione temporale dell'energia elettromagnetica immagazzinata in un volume $V$ deve eguagliare la somma dell'energia dissipata (o convertita) all'interno del volume e dell'energia che fluisce attraverso la superficie di confine.
Per rendere esplicito tale bilancio, si definisce il \textit{Vettore di Poynting} $\bm{S}$:
\begin{equation}
\bm{S} = \frac{\bm{E} \times \bm{B}}{\mu_0} \quad \left[ \si{\watt\per\square\meter} \right].
\end{equation}
Riscrivendo l'equazione di bilancio spostando i termini:
\begin{equation}
-\frac{\partial U_{em}}{\partial t} = \mathcal{P}_{mec} + \oint_{\Sigma} \bm{S} \cdot d\bm{S}.
\end{equation}
I termini assumono il seguente significato fisico:
\begin{itemize}
    \item \textbf{primo membro ($-\partial_t U_{em}$):} rappresenta il tasso di diminuzione dell'energia immagazzinata nei campi elettrico e magnetico all'interno del volume. L'energia "spesa" dal sistema;
    \item \textbf{termine di sorgente ($\mathcal{P}_{mec}$):} rappresenta il lavoro compiuto dai campi sulle cariche in moto (effetto Joule $\bm{J} \cdot \bm{E}$). È l'energia convertita da elettromagnetica a meccanica (o termica);
    \item \textbf{termine di flusso ($\Phi_{\Sigma}(\bm{S})$):} Rappresenta il flusso del vettore di Poynting attraverso la superficie chiusa. Se positivo, indica che l'energia sta "uscendo" dal volume sotto forma di radiazione.
\end{itemize}
Il vettore $\bm{S}$ possiede dunque una duplice valenza: la sua direzione indica la direzione di propagazione dell'energia (e dunque dell'onda: si noti che $\bm{S} \parallel \bm{\kappa}$), mentre il suo modulo rappresenta l'intensità istantanea $I$ dell'onda (energia trasportata per unità di tempo e di superficie).

\subsection{Intensità dell'onda elettromagnetica}

Si è visto come il vettore di Poynting $\bm{S}$ rappresenti la potenza trasportata per unità di superficie. Il suo modulo istantaneo, esplicitando la dipendenza temporale dei campi (e quindi di $\bm{S}$), $S(t)$, nel caso di un'onda piana nel vuoto in cui $\bm{E} \perp \bm{B}$ e $E = cB$, si calcola come:
\begin{equation}
S(t) = \frac{|\bm{E} \times \bm{B}|}{\mu_0} = \frac{E(t)B(t)}{\mu_0} = \frac{E(t)^2}{c\mu_0}.
\end{equation}
Ricordando che $c = 1/\sqrt{\varepsilon_0\mu_0}$, è possibile riscrivere l'espressione in termini di densità di energia istantanea $w(t)$:
\begin{equation}
S(t) = \frac{E(t)^2}{\mu_0} \sqrt{\varepsilon_0\mu_0} = c \varepsilon_0 E(t)^2 = c w(t).
\end{equation}
Tuttavia, poiché i campi oscillano con frequenze elevatissime, la grandezza di interesse fisico non è il valore istantaneo bensì il valore medio nel tempo. Si definisce \textit{intensità} $I$ dell'onda la media temporale del modulo del vettore di Poynting su un periodo $T$:
\begin{equation}
I = \langle S \rangle_T = \frac{1}{T} \int_0^T S(t) \, dt.
\end{equation}
Considerando un'onda armonica del tipo $E(t) = E_0 \cos(kx - \omega t)$, il calcolo diviene:
\begin{equation}
I = c \varepsilon_0 \langle E_0^2 \cos^2(kx - \omega t) \rangle_T.
\end{equation}
Poiché il valore medio della funzione $\cos^2$ su un periodo è pari a $1/2$, si ottiene l'espressione finale dell'intensità in funzione dell'ampiezza massima del campo elettrico $E_0$:
\begin{equation}
I = \frac{1}{2} c \varepsilon_0 E_0^2 = \frac{1}{2} \frac{E_0^2}{c\mu_0}.
\end{equation}
In termini di valore efficace del campo ($E_{eff} = E_0/\sqrt{2}$), l'intensità assume la forma $I = c \varepsilon_0 E_{eff}^2$.



\section{La polarizzazione delle onde elettromagnetiche}
Una proprietà fondamentale che distingue le onde trasversali (come quelle elettromagnetiche) dalle onde longitudinali (come quelle sonore nei fluidi) è la \textit{polarizzazione}. Poiché i vettori campo elettrico $\bm{E}$ e campo magnetico $\bm{B}$ giacciono sul piano ortogonale alla direzione di propagazione $\bm{\kappa}$, essi possiedono un grado di libertà relativo al loro orientamento su tale piano.
Per convenzione, la direzione della polarizzazione è definita dalla direzione di oscillazione del vettore campo elettrico $\bm{E}$.

\subsection{Tipologie di polarizzazione}

Lo stato di polarizzazione di un'onda è determinato dalla relazione di fase e di ampiezza tra le due componenti ortogonali del campo elettrico (ad esempio $E_x$ ed $E_y$). Si distinguono tre casi principali:
\begin{enumerate}
\item \textit{Polarizzazione lineare:}
si verifica quando il vettore $\bm{E}$ oscilla mantenendo una direzione costante nel tempo. Ciò accade se esiste una sola componente del campo (es. solo lungo $x$) oppure se le due componenti ortogonali sono in fase (sfasamento $\delta = 0$ o $\delta = \pi$). In questo caso, la risultante vettoriale traccia un segmento di retta nel piano trasversale.

\item \textit{Polarizzazione circolare:}
si ottiene quando le due componenti ortogonali hanno \textit{eguale ampiezza} ma sono sfasate di $\pi/2$ ($90^{\circ}$). In tale configurazione, il modulo del vettore campo elettrico $|\bm{E}|$ rimane costante, ma la sua direzione ruota uniformemente nel tempo attorno all'asse di propagazione (con velocità angolare $\omega$). La punta del vettore descrive una circonferenza.

\item \textit{Polarizzazione ellittica:}
rappresenta il caso più generale. Si manifesta quando le componenti hanno ampiezze diverse e/o uno sfasamento arbitrario (diverso da $0, \pi, \pi/2$). Il vettore risultante cambia sia in modulo che in direzione, tracciando un'ellisse.
\end{enumerate}
La luce naturale (es. solare) è definita \textit{non polarizzata} o stocastica: essa è composta dalla sovrapposizione incoerente di un numero enorme di pacchetti d'onda emessi da singoli atomi, ciascuno con una polarizzazione casuale. In media, il campo elettrico non predilige alcuna direzione specifica.

\subsection{Filtri polarizzatori e Legge di Malus}

È possibile selezionare una specifica direzione di oscillazione del campo elettrico interponendo sul cammino ottico un \textit{filtro polarizzatore} (o polaroide). Tali dispositivi sono costituiti da materiali anisotropi (spesso catene polimeriche stirate) che assorbono selettivamente la componente del campo elettrico parallela a un asse molecolare e trasmettono quella perpendicolare.
Si definisce \textit{asse di trasmissione} del filtro la direzione lungo la quale il campo elettrico viene trasmesso senza attenuazione.
Quando un fascio di luce non polarizzata di intensità $I_0$ attraversa un polarizzatore ideale, l'intensità trasmessa $I_1$ risulta dimezzata:
\begin{equation}
I_1 = \frac{1}{2} I_0.
\end{equation}
Questo avviene perché, decomponendo il campo stocastico su due assi ortogonali, in media l'energia è equipartita tra le due componenti; il filtro ne elimina una e trasmette l'altra.
\subsubsection{La legge di Malus}
L'intensità di un'onda elettromagnetica è proporzionale al quadrato dell'ampiezza del campo elettrico ($I \propto E^2$).
Si consideri il vettore campo elettrico incidente $\bm{E}_0$. Esso può essere scomposto vettorialmente in due componenti rispetto al sistema di riferimento del filtro:
\begin{itemize}
    \item una componente parallela all'asse di trasmissione: $E_{\parallel} = E_0 \cos \theta$;
    \item una componente perpendicolare all'asse di trasmissione: $E_{\perp} = E_0 \sin \theta$.
\end{itemize}
Il filtro ideale è trasparente per la componente parallela e totalmente opaco per quella perpendicolare. Di conseguenza, il campo elettrico emergente $\bm{E}_{out}$ ha modulo pari alla sola componente parallela:
\begin{equation}
E_{out} = E_{\parallel} = E_0 \cos \theta.
\end{equation}
Calcolando l'intensità trasmessa $I_{out}$ come grandezza proporzionale al quadrato del campo:
\begin{equation}
I_{out} \propto (E_{out})^2 = (E_0 \cos \theta)^2 = E_0^2 \cos^2 \theta.
\end{equation}
Essendo $I_0 \propto E_0^2$, si ottiene la relazione finale:
\begin{equation}
I_{out} = I_0 \cos^2 \theta.
\end{equation}
Si nota che l'intensità è massima ($I=I_0$) quando asse e polarizzazione sono allineati ($\theta = 0$) e nulla quando sono ortogonali ($\theta = \pi/2$, condizione di estinzione).
Se un'onda \textit{già linearmente polarizzata} incide su un polarizzatore, l'intensità trasmessa dipende dall'angolo relativo tra la direzione di polarizzazione dell'onda incidente e l'asse di trasmissione del filtro. 


\chapter{L'interferenza di onde elettromagnetiche}
\label{chap:interferenza}

Il fenomeno dell'interferenza rappresenta la manifestazione più diretta ed evidente della natura ondulatoria della radiazione elettromagnetica. Mentre l'ottica geometrica approssima la propagazione della luce tramite raggi rettilinei indipendenti l'uno dall'altro, l'ottica fisica studia gli effetti derivanti dalla sovrapposizione dei campi.
In questo capitolo viene analizzata la ridistribuzione spaziale dell'energia che si verifica quando due o più onde elettromagnetiche occupano simultaneamente la stessa regione di spazio, dando luogo a zone di amplificazione (interferenza costruttiva) e zone di attenuazione (interferenza distruttiva).

\section{Il principio di sovrapposizione: richiami}
Le equazioni di Maxwell nel vuoto sono equazioni differenziali lineari. Da ciò discende direttamente, come già discusso, il \textit{Principio di Sovrapposizione}: il campo elettrico risultante in un punto $P$ dello spazio è dato dalla somma vettoriale dei singoli campi generati dalle diverse sorgenti.
Siano $\bm{E}_1(\bm{r}, t)$ ed $\bm{E}_2(\bm{r}, t)$ i campi elettrici associati a due onde distinte. Il campo totale è:
\begin{equation}
    \bm{E}_{\text{tot}} = \bm{E}_1 + \bm{E}_2.
\end{equation}
Tuttavia, l'occhio umano e i comuni rilevatori ottici (fotodiodi, lastre fotografiche) non sono in grado di seguire l'oscillazione istantanea del campo (che avviene a frequenze dell'ordine di $10^{14}$ Hz), ma rispondono all'\textit{intensità} $I$, definita come la media temporale del modulo del vettore di Poynting (o, a meno di costanti, del quadrato dell'ampiezza del campo elettrico).
L'intensità risultante è dunque proporzionale alla media temporale del quadrato del campo totale:
\begin{equation}
    I \propto \langle |\bm{E}_{\text{tot}}|^2 \rangle = \langle (\bm{E}_1 + \bm{E}_2) \cdot (\bm{E}_1 + \bm{E}_2) \rangle
\end{equation}
Sviluppando il prodotto scalare:
\begin{equation}
    I \propto \underbrace{\langle E_1^2 \rangle}_{I_1} + \underbrace{\langle E_2^2 \rangle}_{I_2} + \underbrace{2 \langle \bm{E}_1 \cdot \bm{E}_2 \rangle}_{I_{12}}
\end{equation}
Si osserva che l'intensità totale non è semplicemente la somma delle intensità delle singole onde ($I \neq I_1 + I_2$), ma contiene un terzo termine $I_{12}$, detto \textit{termine di interferenza}, che dipende intrinsecamente (stante la sua definizione da prodotto scalare) dalla differenza di fase dei campi elettrici.

\subsection{Significato fisico del termine di interferenza}
È opportuno analizzare il caso in cui il termine misto si annulli ($I_{12} = 0$). In tale circostanza, l'intensità totale si riduce alla semplice somma aritmetica delle singole intensità:
\begin{equation}
    I_{\text{tot}} = I_1 + I_2
\end{equation}
Questa situazione descrive il comportamento di \textit{sorgenti incoerenti} (come due comuni lampadine). In assenza del termine $I_{12}$, non si osserva alcuna modulazione spaziale dell'energia (frange chiare e scure), ma semplicemente un'illuminazione diffusa e uniforme.
Il fenomeno dell'interferenza risiede dunque interamente nell'esistenza di $I_{12}$, il quale è responsabile della \textit{ridistribuzione spaziale dell'energia}: l'energia che "manca" nei minimi di interferenza (distruttiva) è esattamente quella che si ritrova "in eccesso" nei massimi (costruttiva), nel pieno rispetto del principio di conservazione dell'energia globale.

\subsection{Condizioni di interferenza}
Affinché il termine $I_{12}$ sia diverso da zero e dia luogo a una figura di interferenza stazionaria (osservabile), devono essere soddisfatte specifiche condizioni fisiche:
\begin{enumerate}
    \item \textit{stessa frequenza:} le sorgenti devono essere monocromatiche e avere la medesima pulsazione $\omega$. Se le frequenze fossero diverse, la media temporale del loro prodotto si annullerebbe (battimenti rapidissimi non rilevabili);
    \item \textit{coerenza di fase:} la differenza di fase $\delta = \phi_1 - \phi_2$ tra le due onde deve mantenersi costante nel tempo. Sorgenti naturali indipendenti emettono treni d'onda con fasi casuali che cambiano ogni $10^{-8}$ s, rendendo la media temporale $\langle \vec{E}_1 \cdot \vec{E}_2 \rangle$ nulla (sorgenti incoerenti);
    \item \textit{polarizzazione parallela:} il prodotto scalare $\vec{E}_1 \cdot \vec{E}_2$ implica che onde con polarizzazioni ortogonali non interferiscono. Si deve verificare inoltre l'ampiezza massima dei campi elettrici sia la medesima.
\end{enumerate}
Verificate queste condizioni, si parla di \textit{onde coerenti.}


\subsection{Trattazione complessa}
Sfruttando l'ipotesi di coerenza, che garantisce uno sfasamento costante tra i campi interagenti, è vantaggioso adottare una rappresentazione nel piano complesso.
Date due onde monocromatiche di uguale ampiezza:
\begin{equation}
    E_1(x_1,t) = E_0\sin(kx_1 -\omega t) \quad E_2(x_2,t) = E_0\sin(kx_2 -\omega t),
\end{equation}
la loro differenza di fase relativa è data da:
\begin{equation}
    \delta = (kx_1 -\omega t) - (kx_2 -\omega t) = k(x_1-x_2).
\end{equation}
A tali grandezze si associano, sul piano di Gauss, due vettori rotanti:
\begin{equation}
    \bm{E}_1 = E_0 e^{i(kx_1 -\omega t)} \quad \bm{E}_2 = E_0 e^{i(kx_2 -\omega t)},
\end{equation}
nella chiara assunzione che le onde originali siano:
\begin{equation}
    E_1(x_1,t) = \mathbb{I}\text{m}(\bm{E}_1) \quad E_2(x_2,t) = \mathbb{I}\text{m}(\bm{E}_2).
\end{equation}
Si osserva che, sebbene entrambi i vettori ruotino con velocità angolare $\omega$, la loro \textit{posizione reciproca risulta costante nel tempo}.
Infatti, l'angolo compreso tra i due vettori corrisponde esattamente alla differenza di fase $\delta$: essendo quest'ultima costante (indipendente dal tempo), i due vettori mantengono invariata la loro separazione angolare durante la rotazione.

\section{Il Principio di Huygens-Fresnel}

La comprensione dei fenomeni di interferenza e diffrazione richiede un modello che descriva come un'onda si propaghi nello spazio e come interagisca con gli ostacoli. Tale modello è fornito dal \textit{Principio di Huygens}, successivamente perfezionato da Fresnel.

\subsection{Enunciato e costruzione geometrica}
Il principio afferma che la propagazione di un fronte d'onda può essere descritta geometricamente attraverso la seguente costruzione:
\begin{itemize}
    \item \textit{sorgenti secondarie:} ogni punto di un fronte d'onda primario $\Sigma(t)$ agisce come una sorgente puntiforme virtuale di onde sferiche secondarie (dette \textit{wavelets}), aventi la stessa frequenza dell'onda primaria e velocità di propagazione $v$ caratteristica del mezzo;
    \item \textit{inviluppo:} la posizione del fronte d'onda all'istante successivo $t + \Delta t$, indicato con $\Sigma(t+\Delta t)$, è data dall'\textit{inviluppo tangenziale} (o superficie tangente) di tutte le onde sferiche secondarie emesse dai punti di $\Sigma(t)$ e propagatesi per un raggio $r = v \Delta t$.
\end{itemize}


\subsection{Discussione fisica e limiti}
Il principio di Huygens rappresenta uno strumento potente per spiegare le leggi della riflessione e della rifrazione (tramite la variazione della velocità $v$ nel mezzo), nonché la propagazione rettilinea della luce in assenza di ostacoli.
Tuttavia, nella sua formulazione originaria, presentava un limite fisico: la costruzione geometrica ammetterebbe anche un inviluppo "all'indietro" (onda retrograda), mai osservato sperimentalmente nella propagazione libera.
Tale ambiguità fu risolta da Augustin-Jean Fresnel, il quale integrò il principio geometrico con il concetto di \textit{interferenza}. Secondo il principio rivisitato (Huygens-Fresnel):
\begin{itemize}
    \item l'ampiezza del campo risultante in un punto P è data dalla \textit{sovrapposizione coerente} (somma vettoriale) di tutte le onde secondarie che giungono in P;
    \item si introduce un \textit{fattore di obliquità} $K(\theta)$, il quale annulla l'ampiezza delle onde secondarie nella direzione opposta alla propagazione ($\theta = 180^\circ$) e la massimizza nella direzione frontale.
\end{itemize}
È proprio grazie a questo principio che è possibile spiegare perché la luce, incontrando una fenditura di dimensioni comparabili alla lunghezza d'onda, non proietta un'ombra netta geometrica, ma "aggira" l'ostacolo diffondendosi in tutte le direzioni (fenomeno della diffrazione).


\section{Interferenza alla Young}

Si consideri un sistema ottico costituito da due sorgenti puntiformi $S_1$ ed $S_2$, separate da una distanza $d$ e immerse in un mezzo omogeneo (vuoto o aria, con indice di rifrazione $n \approx 1$).
Le sorgenti emettono onde elettromagnetiche monocromatiche di lunghezza d'onda $\lambda$, aventi la medesima ampiezza $E_0$ e la medesima polarizzazione.

\subsection{Setup geometrico e ipotesi}
La figura di interferenza viene osservata su uno schermo posto parallelamente alla congiungente delle sorgenti, a una distanza $L$. Al fine di semplificare la trattazione analitica, si introducono le seguenti ipotesi operative:

\begin{enumerate}
    \item \textit{condizione di campo lontano (o di Fraunhofer):} si assume che la distanza dello schermo sia molto maggiore della separazione tra le sorgenti ($L \gg d$). Questa approssimazione implica che i raggi che partono da $S_1$ e $S_2$ e giungono in un generico punto $P$ sullo schermo possano essere considerati tra loro \textit{paralleli} in prossimità delle sorgenti;
    \item \textit{sorgenti in fase (sincrone):} si ipotizza che le due sorgenti oscillino in perfetta fase tra loro ($\phi_{01} = \phi_{02}$). Di conseguenza, la differenza di fase totale $\delta$ in un punto $P$ dipenderà esclusivamente dalla differenza di cammino geometrico percorso dalle onde e non da sfasamenti intrinseci alla sorgente.
\end{enumerate}

\subsection{Differenza di cammino ottico}
Sia $P$ un punto sullo schermo, individuato dalla coordinata verticale $y$ (misurata rispetto all'asse del segmento congiungente le due sorgenti) o dall'angolo $\theta$ sotto cui il punto è visto dal centro del sistema di fenditure.
Siano $r_1$ ed $r_2$ le distanze percorse rispettivamente dall'onda emessa da $S_1$ e da $S_2$ per raggiungere $P$.
Nell'approssimazione $L \gg d$ (e per piccoli angoli, $y \ll L$), la differenza di cammino ottico $\Delta$:
\begin{equation}
    \Delta = r_2 - r_1 
\end{equation}
\begin{figure}[htbp]
    \centering
    \begin{tikzpicture}[scale=0.8, >=latex]
        % --- DEFINIZIONE VARIABILI ---
        \def\d{2.5}      % Distanza tra le fenditure (esagerata per chiarezza)
        \def\L{7}        % Distanza schermo
        \def\y{2.5}      % Altezza punto P
        \def\slitgap{0.3}% Apertura grafica fenditura

        % --- COORDINATE ---
        \coordinate (S1) at (0, \d/2);
        \coordinate (S2) at (0, -\d/2);
        \coordinate (O) at (0,0);       % Centro fenditure
        \coordinate (C_screen) at (\L,0); % Centro schermo
        \coordinate (P) at (\L, \y);    % Punto di osservazione

        % --- 1. SCHERMO SORGENTI (SX) ---
        \draw[thick, fill=gray!20] (-0.2, \d/2+\slitgap) -- (-0.2, \d+0.5) -- (0.2, \d+0.5) -- (0.2, \d/2+\slitgap) -- cycle;
        \draw[thick, fill=gray!20] (-0.2, -\d/2-\slitgap) -- (-0.2, -\d-0.5) -- (0.2, -\d-0.5) -- (0.2, -\d/2-\slitgap) -- cycle;
        \draw[thick, fill=gray!20] (-0.2, -\d/2+\slitgap) -- (-0.2, \d/2-\slitgap) -- (0.2, \d/2-\slitgap) -- (0.2, -\d/2+\slitgap) -- cycle;

        \node[left] at (-0.3, \d/2) {$S_1$};
        \node[left] at (-0.3, -\d/2) {$S_2$};

        % --- 2. SCHERMO DI OSSERVAZIONE (DX) ---
        \draw[thick] (\L, -3.5) -- (\L, 3.5) node[above] {Schermo};
        \draw[->, dashed, gray] (C_screen) -- (\L, \y) node[midway, right] {$y$};
        \fill (P) circle (2pt) node[right] {$P$};

        % --- 3. ASSE OTTICO E RAGGI ---
        \draw[dashed, gray] (-1,0) -- (\L+1,0);
        
        % Raggi r1 e r2
        \draw[blue, thick] (S1) -- (P) node[midway, above=-2pt] {$r_1$};
        \draw[blue, thick] (S2) -- (P) node[midway, below right=-2pt] {$r_2$};

        % Raggio dal centro (per definire theta)
        \draw[dashed, thin, black!70] (O) -- (P);

        % --- 4. GEOMETRIA DIFFERENZA DI CAMMINO (ZOOM) ---
        % Piede della perpendicolare H da S1 su r2
        % Per approssimazione di campo lontano, consideriamo i raggi paralleli vicino all'origine
        % Angolo theta calcolato
        \pgfmathsetmacro{\thetaangle}{atan(\y/\L)}
        
        % Proiezione ortogonale H approssimata (ruotata di theta)
        \coordinate (H) at ($(S2) + ({cos(90-\thetaangle)*\d}, {sin(90-\thetaangle)*\d})$); 
        % In realtà, geometricamente H sta sul raggio S2->P. 
        % Usiamo la proiezione geometrica vera per precisione TikZ:
        \coordinate (H_true) at ($(S2)!(S1)!(P)$);
        
        \draw[dashed, red] (S1) -- (H_true);
        \draw[red, very thick] (S2) -- (H_true) node[midway, below, font=\small] {$\Delta \approx d\sin\theta$};
        
        % Angolo retto in H
        \draw[red] ($(H_true)!0.2cm!(S2)$) -- ($(H_true)!0.2cm!(S2)!0.2cm!90:(S2)$) -- ($(H_true)!0.2cm!(S1)$);

        % --- 5. ANGOLI E ETICHETTE ---
        % Angolo Theta all'origine
        \draw (1.5,0) arc (0:\thetaangle:1.5);
        \node at (1.8, 0.3) {$\theta$};

        % Angolo Theta nel triangolino (è uguale a quello all'origine)
        \draw[red] ($(S1)+(0,-0.6)$) arc (270:{270+\thetaangle}:0.6);
        \node[red, right, font=\footnotesize] at ($(S1)+(0.1,-0.8)$) {$\theta$};

        % Etichetta d
        \draw[<->] (-0.8, -\d/2) -- (-0.8, \d/2) node[midway, fill=white] {$d$};

        % Etichetta L
        \draw[<->] (0, -3.8) -- (\L, -3.8) node[midway, fill=white] {$L$};

    \end{tikzpicture}
    \caption{Schema geometrico dell'esperimento di Young. Le fenditure $S_1$ e $S_2$ distano $d$. Il punto $P$ sullo schermo (distante $L$) è individuato dall'angolo $\theta$. In rosso è evidenziata la differenza di cammino ottico $\Delta = r_2 - r_1$, che nell'approssimazione di campo lontano ($L \gg d$) vale $d\sin\theta$.}
    \label{fig:young_scheme}
\end{figure}

\noindent
Nell'assunzione che i cammini ottici siano pressoché paralleli, la loro differenza all'angolo $\theta$ formato tra le onde e l'orizzontale:
\begin{equation}
    \Delta = d \sin\theta = r_2 - r_1.
\end{equation}
Tale differenza di percorso si traduce in uno sfasamento $\delta$ proporzionale al numero d'onda $k$:
\begin{equation}
    \delta = k(r_2 - r_1) = \frac{2\pi}{\lambda} d \sin\theta
\end{equation}
È proprio la variazione di $\theta$ (e quindi di $\Delta$) muovendosi lungo lo schermo a generare l'alternanza di massimi e minimi di intensità.

\subsection{Intensità dell'onda risultante}
Essendo le onde coerenti, ci si può avvalere della loro rappresentazione complessa per semplificare notevolmente i conti trigonometrici. 

\begin{figure}[htbp]
    \centering
    \begin{tikzpicture}[scale=1.2, >=latex]
        % --- Librerie necessarie (da mettere nel preambolo se mancano) ---
        % \usetikzlibrary{calc, angles, quotes}

        % --- Parametri ---
        \def\R{3}       % Modulo dei vettori (E0)
        \def\angA{25}   % Angolo di fase E1 (omega*t + phi1)
        \def\angB{75}   % Angolo di fase E2 (omega*t + phi2)
        % Delta = 75 - 25 = 50 gradi

        % --- Coordinate ---
        \coordinate (O) at (0,0);
        \coordinate (X) at (5,0); % Asse Reale
        \coordinate (Y) at (0,5); % Asse Immaginario
        \coordinate (E1) at (3,0);
        \coordinate (Ex) at (4,0);
        \coordinate (E2) at (5,1.8);
        \coordinate (R) at (1.8,4.5);
        \coordinate (Etot) at ($ (E2) $); % Somma vettoriale

        % --- Assi del Piano di Gauss ---
        \draw[->, gray, thick] (-1,0) -- (X) node[below, black] {Re};
        \draw[->, gray, thick] (0,-1) -- (Y) node[left, black] {Im};

        % --- Vettori Componenti (E1, E2) ---
        \draw[->, thick, blue] (O) -- (E1) node[midway, above right] {$\mathbf{E}_1$};
        \draw[->, thick, blue] (E1) -- (E2) node[midway,  left] {$\mathbf{E}_2$};

        % --- Costruzione del Parallelogramma ---
        \draw[dashed, blue!50] (E1) -- (R) node[midway, below, black] {$R$};
        \draw[dashed, blue!50] (O) -- (R) node[midway, below, black] {$R$};
        \draw[dashed, blue!50] (E2) -- (R) node[midway, below, black] {$R$};

        % --- Vettore Risultante (Etot) ---
        \draw[->, very thick, red] (O) -- (Etot) node[midway, above] {$\mathbf{E}_{\text{tot}}$};

        % --- Angolo Delta (Sfasamento) ---
        \pic [draw, <->, "$\delta$", angle radius=0.6cm, angle eccentricity=1.3] {angle = Ex--E1--E2};
        \pic [draw, <->, "$\alpha$", angle radius=0.6cm, angle eccentricity=1.3, green] {angle = R--E1--O};
        \pic [draw, <->, "$\alpha$", angle radius=0.6cm, angle eccentricity=1.3, green] {angle = R--E2--E1};
        \pic [draw, <->, "$\alpha$", angle radius=0.6cm, angle eccentricity=1.3, green] {angle = E2--E1--R};
        \pic [draw, <->, "$\alpha$", angle radius=0.6cm, angle eccentricity=1.3, green] {angle = E1--O--R};
        \pic [draw, <->, "$\delta$", angle radius=0.6cm, angle eccentricity=1.3, orange] {angle = O--R--E1};
        \pic [draw, <->, "$\delta$", angle radius=0.6cm, angle eccentricity=1.3, orange] {angle = E1--R--E2};

        % --- Indicazione della Rotazione (Omega) ---
        % Un arco esterno che indica che tutto il sistema ruota

        % --- Proiezione (Opzionale: Campo fisico) ---
        % Decommentare se si vuole mostrare che il campo reale è la proiezione
        % \draw[dotted, red] (Etot) -- (Etot |- O) node[below, font=\footnotesize] {$E(t)$};

    \end{tikzpicture}
    \caption{Rappresentazione dei campi nel piano complesso (metodo dei fasori). I vettori $\mathbf{E}_1$ e $\mathbf{E}_2$, di modulo $E_0$, sono sfasati di un angolo costante $\delta$. Il vettore risultante $\mathbf{E}_{\text{tot}}$ (in rosso) è la somma vettoriale dei due. L'intero sistema ruota con velocità angolare $\omega$ attorno all'origine, mantenendo invariata la geometria relativa del parallelogramma.}
    \label{fig:fasori_interferenza}
\end{figure}

\noindent
Dalla Figura si deriva facilmente che:
\begin{equation}
    E_1 = E_2 = E_0 = 2R \sin\frac{\delta}{2}
\end{equation}
e $E_{\text{tot}}$:
\begin{equation}
    E_{\text{tot}} = 2R\sin\delta.
\end{equation}
da cui:
\begin{equation}
    \frac{E_{\text{tot}}}{E_{0}} = \frac{\sin\delta}{\sin\frac{\delta}{2}} = \frac{2\sin\frac{\delta}{2}\cos\frac{\delta}{2}}{\sin\frac{\delta}{2}} = 2\cos \left(\frac{\pi}{\lambda} d \sin\theta\right).
\end{equation}
Le intensità delle singole onde e della risultante sono: 
\begin{equation}
    I_{\text{tot}} \propto E_{\text{tot}}^2 = 4E_0^2\cos \left(\frac{\pi}{\lambda} d \sin\theta\right)^2.
\end{equation}
La funzione $I(\theta)$ può assumere valori compresi nell'intervallo $[0, 4E_0^2]$. I minimi si trovano quindi in occorrenza degli zeri del del coseno, i massimi quando $\cos(\cdot)^2 = 1$.
\begin{equation}
\begin{aligned}
\text{massimi}& \\
\\
\text{minimi} &
\end{aligned}
    \begin{dcases}
    \cos \left(\frac{\pi}{\lambda} d \sin\theta\right)^2 = 1 \implies \frac{\pi}{\lambda} d \sin\theta = n\pi, \quad n \in \mathbb{Z} \\
    \cos \left(\frac{\pi}{\lambda} d \sin\theta\right)^2 = 0 \implies \frac{\pi}{\lambda} d \sin\theta = (2m+1)\frac{\pi}{2}, \quad m \in \mathbb{Z}.
    \end{dcases}
\end{equation}
da cui i valori di $\sin\theta$ per cui si verificano massimi e minimi di interferenza
\begin{equation}
\begin{aligned}
\text{massimi}& \\
\text{minimi} &
\end{aligned}
    \begin{cases}
      \sin\theta = n\frac{\lambda}{d}, \quad n \in \mathbb{Z} \\
      \sin\theta = (2m+1)\frac{\lambda}{2d}, \quad m \in \mathbb{Z}.
    \end{cases}
\end{equation}
Si noti come si verifica interferenza costruttiva (massimi) quando la differenza di cammino ottico $\Delta$ è un numero intero di lunghezze d'onda, distruttiva (minimi) se $\Delta$ è un numero semi intro di lunghezze d'onda.

\subsubsection*{Osservazioni fisiche sui risultati}
Dall'analisi dell'equazione dell'intensità e delle condizioni di massimo e minimo, emergono tre considerazioni fondamentali:

\begin{enumerate}
    \item \textit{ridistribuzione dell'energia:} nei punti di interferenza costruttiva, l'intensità totale è $I_{\text{max}} = 4I_0$ (dove $I_0 \propto E_0^2$ è l'intensità della singola fenditura), mentre nei minimi è nulla.
    Si noti che l'intensità media integrata su un periodo spaziale è esattamente pari a $2I_0$. Ciò conferma che il fenomeno dell'interferenza non viola il principio di conservazione dell'energia: l'energia non viene né creata né distrutta, ma semplicemente \textit{ridistribuita} nello spazio, concentrandosi nei massimi a scapito dei minimi.

    \item \textit{interfrangia e linearità:} nell'approssimazione di piccoli angoli, valida per $L \gg y$, si può assumere $\sin\theta \approx \tan\theta \approx y/L$. La posizione dei massimi sullo schermo diventa lineare rispetto all'ordine $n$:
    \begin{equation}
        y_n = n \frac{\lambda L}{d}
    \end{equation}
    Si definisce \textit{interfrangia} $i$ la distanza tra due massimi consecutivi, la quale risulta costante su tutto lo schermo (nella zona parassiale):
    \begin{equation}
        i = y_{n+1} - y_n = \frac{\lambda L}{d}
    \end{equation}
    Questa relazione permette di misurare la lunghezza d'onda $\lambda$ (microscopica) tramite grandezze macroscopiche facilmente misurabili ($L, d, i$).

    \item \textit{dispersione cromatica:} poiché la posizione dei massimi dipende da $\lambda$, se si utilizza una sorgente policromatica (es. luce bianca), le frange dei diversi colori non coincideranno.
    Il massimo centrale ($n=0$, $\Delta=0$) sarà l'unico a rimanere bianco (tutte le lunghezze d'onda interferiscono costruttivamente nello stesso punto). Per gli ordini superiori ($n \neq 0$), si osserverà uno spettro con la componente blu ($\lambda$ minore) più vicina al centro e quella rossa ($\lambda$ maggiore) più esterna.
\end{enumerate}

\subsection{La genesi storica: l'esperimento del 1801}
La realizzazione pratica dell'esperimento, compiuta da Thomas Young nel 1801, rappresentò una sfida tecnica notevole per l'epoca, data l'assenza di sorgenti di luce intrinsecamente coerenti (come i moderni laser).
Le sorgenti disponibili, come la luce solare o le fiamme delle candele, sono infatti sorgenti termiche estese e incoerenti: i diversi punti della sorgente emettono treni d'onda con fasi casuali e rapidamente variabili, impedendo l'osservazione di una figura di interferenza stabile.

\subsubsection{Il metodo della divisione del fronte d'onda}
Young risolse brillantemente il problema della coerenza spaziale e temporale ideando un setup a tre stadi, basato sul principio della \textit{divisione del fronte d'onda}.
Il dispositivo sperimentale consisteva in:
\begin{enumerate}
    \item una sorgente primaria intensa (luce solare filtrata o una candela);
    \item un primo schermo opaco dotato di un singolo foro puntiforme (pinhole) $S_0$;
    \item un secondo schermo, posto a valle del primo, dotato di due piccoli fori (o fenditure) $S_1$ e $S_2$ molto ravvicinati.
\end{enumerate}
La funzione del primo foro $S_0$ è quella di selezionare una porzione spazialmente limitata del fronte d'onda emesso dalla sorgente estesa, agendo come una sorgente puntiforme virtuale. Le onde sferiche in uscita da $S_0$ si propagano fino a investire il secondo schermo\footnote{Condizione necessaria perché $S_0$ potesse essere sorgente è che la sua dimensione fosse nell'ordine della lunghezza d'onda della radiazione, $\sqrt{S_0} \approx \lambda$.}.
In accordo con il \textit{Principio di Huygens-Fresnel}, i due fori $S_1$ e $S_2$, venendo investiti simultaneamente dallo stesso fronte d'onda primario generato da $S_0$, agiscono come sorgenti secondarie di onde sferiche.
Poiché l'emissione in $S_1$ e $S_2$ è pilotata dalla medesima perturbazione incidente, le due sorgenti secondarie risultano essere \textit{intrinsecamente coerenti} (o sincrone): la loro differenza di fase iniziale è nulla (o costante geometricamente), condizione necessaria per l'interferenza stazionaria.

\subsection{Osservazioni e rottura con la fisica classica}
L'osservazione sullo schermo di rivelazione mostrò una serie di frange luminose alternate a frange scure. Tale risultato era in netto contrasto con le previsioni dell'ottica corpuscolare newtoniana allora dominante, secondo la quale le particelle di luce passanti per i due fori avrebbero dovuto generare semplicemente due macchie luminose distinte (o una sovrapposizione additiva delle intensità).
La presenza di zone di buio laddove arrivava luce da entrambe le fenditure (interferenza distruttiva) costrinse la comunità scientifica ad accettare la natura ondulatoria della luce, unico modello capace di spiegare come la somma di due luminosità potesse dare luogo all'oscurità.


\section{Interferenza di N sorgenti coerenti}
\subsection{Il caso a 3 sorgenti}
Mantenendo il medesimo setup sperimentale sfruttato da Young (schermo di rilevazione molto lontano e parallelo alla congiungente delle sorgenti), si definiscono $S_1$, $S_2$ e $S_3$ tre sorgenti coerenti e allineate, distanti reciprocamente $d$. Le onde emesse hanno campo elettrico in notazione complessa del tipo:
\begin{equation}
\bm{E}_1 = E_0e^{i\overbrace{(kr_1-\omega t)}^{\phi_1}} \quad \bm{E}_2 = E_0e^{i\overbrace{(kr_2-\omega t)}^{\phi_2}}  \quad \bm{E}_3 = E_0e^{i\overbrace{(kr_3-\omega t)}^{\phi_3}} , 
\end{equation}
con differenza di fase:
\begin{equation}
\phi_2-\phi_1 = k(r_2-r_1) = \delta = \phi_3 -\phi_2 = k(r_3-r_2).
\end{equation}
Dall'analisi della Figura (\ref{fig:fasori_tre_onde}), in modo assolutamente analogo all'esperimento delle due fenditure, si ottiene che il campo elettrico risultante e $E_0$ si possono esprimere come:
\begin{equation}
E_{\text{tot}} = 2R\sin\left(\frac{3}{2}\delta\right) \quad E_{\text{0}} = 2R\sin\left(\frac{\delta}{2}\right), 
\end{equation}
da cui:
\begin{equation}
\frac{E_{\text{tot}}}{E_0}= \frac{\sin\left(\frac{3}{2}\delta\right)}{\sin\left(\frac{\delta}{2}\right)}.
\end{equation}
L'intensità registrata sullo schermo di rivelazione, quindi:
\begin{equation}
I_{\text{tot}} \propto E_{\text{tot}}^2 = E_0^2 \frac{\sin\left(\frac{3}{2}\delta\right)^2}{\sin\left(\frac{\delta}{2}\right)^2}.
\end{equation}
\begin{figure}[htbp]
    \centering
    \begin{tikzpicture}[scale=1.2, >=latex]
        % --- Parametri ---
        \def\Ezero{2.5}      % Modulo dei vettori (E0)
        \def\angDelta{40}    % Sfasamento delta tra un vettore e il successivo

        % --- Coordinate ---
        \coordinate (O) at (0,0);
        \coordinate (X) at (6,0); % Asse Reale esteso
        \coordinate (Y) at (0,5.5); % Asse Immaginario
        
        % Calcolo dei vettori "testa-coda"
        % E1 parte dall'origine, angolo 0
        \draw[->, thick, blue] (O) -- ++(0:\Ezero) coordinate (E1) node[midway, below] {$\mathbf{E}_1$};
        
        % E2 parte da E1, ruotato di delta
        \draw[->, thick, blue] (E1) -- ++(\angDelta:\Ezero) coordinate (E2) node[midway, left] {$\mathbf{E}_2$};
        
        % E3 parte da E2, ruotato di 2*delta
        \draw[->, thick, blue] (E2) -- ++(2*\angDelta:\Ezero) coordinate (E3) node[midway, above left] {$\mathbf{E}_3$};

        % --- Assi del Piano di Gauss ---
        \draw[->, gray, thick] (-1,0) -- (X) node[below, black] {Re};
        \draw[->, gray, thick] (0,-1) -- (Y) node[left, black] {Im};

        % --- Vettore Risultante ---
        \draw[->, very thick, red] (O) -- (E3) node[midway, below right] {$\mathbf{E}_{\text{tot}}$};

        % --- Costruzione Geometrica (Centro del Poligono) ---
        % Calcolo del centro C del cerchio circoscritto
        % Il raggio R = E0 / (2*sin(delta/2))
        % L'angolo alla base del triangolo isoscele O-C-E1 è (180-delta)/2 = 90 - delta/2
        \pgfmathsetmacro{\Radius}{\Ezero / (2*sin(\angDelta/2))}
        \coordinate (C) at ($(O) + ({90-\angDelta/2}:\Radius)$);

        % Raggi tratteggiati verso il centro
        \draw[dashed, blue!40] (O) -- (C) node[midway, left, font=\footnotesize, black] {$R$};
        \draw[dashed, blue!40] (E1) -- (C);
        \draw[dashed, blue!40] (E2) -- (C);
        \draw[dashed, blue!40] (E3) -- (C) node[midway, right, font=\footnotesize, black] {$R$};

        % --- Angoli ---
        % Prolungamento per mostrare delta tra E1 e E2
        \coordinate (E1_ext) at ($(E1) + (1.5,0)$);
        \draw[dotted, gray] (E1) -- (E1_ext);
        \pic [draw, ->, "$\delta$", angle radius=0.8cm, angle eccentricity=1.3] {angle = E1_ext--E1--E2};

        % Prolungamento per mostrare delta tra E2 e E3
        % Calcolo un punto sulla retta di E2 prolungata
        \coordinate (E2_ext) at ($(E2) + ({\angDelta}:1.5)$);
        \draw[dotted, gray] (E2) -- (E2_ext);
        \pic [draw, ->, "$\delta$", angle radius=0.8cm, angle eccentricity=1.3] {angle = E2_ext--E2--E3};

        % Angolo al centro (opzionale, per mostrare la geometria)
        \pic [draw, "$\delta$", angle radius=0.6cm, angle eccentricity=1.5, orange] {angle = E1--C--O};
        \pic [draw, "$\delta$", angle radius=0.6cm, angle eccentricity=1.5, orange] {angle = E2--C--E1};

    \end{tikzpicture}
    \caption{Somma vettoriale di tre campi elettrici coerenti nel piano complesso. I tre fasori $\mathbf{E}_1, \mathbf{E}_2, \mathbf{E}_3$, di eguale modulo $E_0$, sono disposti a catena, ciascuno ruotato di un angolo $\delta$ rispetto al precedente. Il vettore risultante $\mathbf{E}_{\text{tot}}$ (in rosso) è la corda del poligono inscritto in una circonferenza di raggio $R$.}
    \label{fig:fasori_tre_onde}
\end{figure}

\noindent
Assumendo che le differenze di cammino ottico siano uguali per coppie di onde consecutive, stante l'approssimativo parallelismo:
\begin{equation}
\Delta_1 = r_2 - r_1 = d\sin\theta = \Delta_2 = r_3 - r_2,
\end{equation}
si ottiene:
\begin{equation}
I_{\text{tot}} \propto E_{\text{tot}}^2 = E_0^2 \frac{\sin\left(\frac{3}{2}kd\sin\theta\right)^2}{\sin\left(\frac{kd\sin\theta}{2}\right)^2} = E_0^2 \frac{\sin\left(\frac{d}{\lambda}3\pi\sin\theta\right)^2}{\sin\left(\frac{\pi d\sin\theta}{\lambda}\right)^2}.
\label{eq:3_sorg}
\end{equation}

\subsection{Generalizzazione a N sorgenti}
Generalizzando per induzione la \eqref{eq:3_sorg} per un numero $N$ di sorgenti (che emettono radiazione con campi elettrici di ampiezza $E_0$) coerenti e equidistanti $d$, si ottiene la seguente relazione:
\begin{equation}
I_{\text{tot}} \propto E_{\text{tot}}^2 = E_0^2 \frac{\sin\left(\frac{\delta}{2}N\right)^2}{\sin\left(\frac{\delta}{2}\right)^2}.
\end{equation}
La differenza di fase, poi, $\delta = kd\sin\theta = 2\pi\sin\theta/\lambda$. Quindi:
\begin{equation}
I_{\text{tot}} \propto E_0^2 \frac{\sin\left(\frac{N}{2} kd\sin\theta \right)^2}{\sin\left(\frac{kd\sin\theta}{2}\right)^2} = E_0^2 \frac{\sin\left(\frac{d}{\lambda} N\pi\sin\theta \right)^2}{\sin\left(\frac{ d}{\lambda}\pi\sin\theta\right)^2}.
\end{equation}
L'espressione dell'intensità appena ricavata presenta una struttura periodica modulata dal rapporto di due funzioni seno. Si procede ora all'analisi dei punti stazionari della funzione.

\subsubsection{Massimi principali}
I massimi di intensità più rilevanti, detti \textit{massimi principali}, si verificano quando il denominatore della frazione tende a zero. Ciò accade quando l'argomento del seno è un multiplo intero di $\pi$:
\begin{equation}
    \frac{\delta}{2} = m\pi \implies \frac{\pi d \sin\theta}{\lambda} = m\pi, \quad \text{con } m \in \mathbb{Z}.
\end{equation}
Da cui si ritrova la classica condizione di interferenza costruttiva (identica al caso di due fenditure):
\begin{equation}
    \sin\theta = m\frac{\lambda}{d}.
\end{equation}
In tali punti, anche il numeratore si annulla ($\sin(Nm\pi) = 0$), generando una forma indeterminata del tipo $0/0$. Espandendo in serie di McLaurin numeratore e denominatore, si ottiene:
\begin{equation}
    \lim_{\delta/2 \to m\pi} \frac{\sin(N\delta/2)}{\sin(\delta/2)} = \lim_{\delta/2 \to m\pi} \frac{N\delta/2}{\delta/2} = N.
\end{equation}
L'intensità dei massimi principali scala dunque con il quadrato del numero di fenditure:
\begin{equation}
    I_{\text{max}} \propto (N E_0)^2 = N^2 I_0.
\end{equation}
Questo risultato è fisicamente coerente: in fase, i campi si sommano scalarmente ($N E_0$) e l'intensità è il quadrato del campo totale.

\subsubsection{Minimi di intensità}
L'intensità nulla si ottiene quando si annulla il numeratore, a patto che il denominatore sia diverso da zero (per non ricadere nei massimi principali).
\begin{equation}
    \sin\left(\frac{N\delta}{2}\right) = 0 \implies \frac{N\delta}{2} = n\pi, \quad n \in \mathbb{Z}.
\end{equation}
La condizione per i minimi è dunque:
\begin{equation}
    \delta = \frac{2\pi}{N} n \implies \sin\theta = \frac{n}{N}\frac{\lambda}{d}, \quad \text{con } n \neq mN.
\end{equation}
La clausola $n \neq mN$ esclude i valori di $n$ che porterebbero a $m\lambda$ (che corrispondono ai massimi principali).
Si osserva che tra due massimi principali consecutivi (es. ordine $m_0=0$ e $m_1=1$), $n$ può assumere tutti i valori tali per cui:
\begin{equation}
    m_0 < \frac{n}{N} < m_1 \implies Nm_0 < n < Nm_1.
\end{equation}
Essendo $m_1$ e $m_0$ consecutivi, 
\begin{equation}
    0 < n < N(m_1-m_0) = N,
\end{equation}
da cui si evince che $n$ può assumere $N-1$ valori: vi sono dunque esattamente \textit{$N-1$ minimi} fra due massimi consecutivi.

\subsubsection{Massimi secondari}
Tra gli $N-1$ minimi presenti tra due picchi principali, la funzione deve necessariamente risalire. Dal momento che ogni massimo principale deve essere necessariamente anticipato (e seguito) da un minimo secondario e che tra due minimi consecutivi si trova un massimo, si ha che per ognuno degli $N-1$ minimi a parte l'ultimo (dopo cui si ha massimo primario)\footnote{O il primo, con il ragionamento in viceversa.} si trova un \textit{massimo secondario}, in numero complessivo di $N-2$.
Questi picchi minori si trovano approssimativamente a metà strada tra i minimi:
\begin{equation}
    \frac{N\delta}{2} \approx (2n+1)\frac{\pi}{2}.
\end{equation}
Sostituendo tale condizione nell'equazione dell'intensità, si nota che il numeratore vale 1, mentre il denominatore cresce allontanandosi dal massimo principale. Di conseguenza, l'intensità dei lobi secondari decresce rapidamente.
Per grandi $N$, il primo massimo secondario ha un'intensità pari a circa il $4.5\%$ del massimo principale.

\subsubsection{Considerazione conclusiva}
All'aumentare di $N$, si verificano due fenomeni congiunti:
\begin{enumerate}
    \item L'intensità dei massimi principali cresce come $N^2$.
    \item La larghezza dei massimi principali diminuisce come $1/N$ (poiché il primo zero si trova a $\Delta \theta \propto \lambda/Nd$).
\end{enumerate}
Per $N$ molto grande, la figura di interferenza diventa una serie di picchi estremamente intensi e stretti (delta di Dirac approssimate) su uno sfondo quasi buio. Questo è il principio di funzionamento del \textit{reticolo di diffrazione}.


\section{Il Potere Risolutivo e il Criterio di Rayleigh}
Nell'analisi di qualsiasi sistema ottico, sia esso un telescopio deputato alla risoluzione spaziale o un reticolo di diffrazione per l'analisi spettrale, è fondamentale distinguere tra \textit{ingrandimento} e \textit{risoluzione}. Mentre il primo è un fattore puramente geometrico, limitato solo dalle aberrazioni o dalla costruzione tecnica, la risoluzione rappresenta un limite fisico fondamentale imposto dalla natura ondulatoria della luce.\\

\noindent
A causa della diffrazione, l'immagine di una sorgente puntiforme (o di una riga spettrale monocromatica ideale) non è mai un punto geometrico, bensì una distribuzione estesa di irradianza (ad esempio, un disco di Airy o un picco principale circondato da lobi secondari).
Quando si analizza la radiazione emessa da due sorgenti distinte, o composta da due lunghezze d'onda vicine $\lambda$ e $\lambda + \Delta\lambda$, le rispettive figure di diffrazione si sovrappongono parzialmente. Esiste dunque una distanza critica al di sotto della quale le due figure si fondono in un unico inviluppo, rendendo impossibile distinguere le due componenti originali: in tale condizione, l'informazione spettrale o spaziale è irrimediabilmente persa.
Per quantificare tale limite, si adotta universalmente il \textit{Criterio di Rayleigh}. Tale criterio, di natura euristica ma estremamente efficace, stabilisce una condizione operativa di "appena risolvibilità":
\begin{quote}
    \textit{Due immagini di diffrazione si considerano risolte quando il massimo principale dell'una cade in corrispondenza del primo minimo di intensità dell'altra.}
\end{quote}
Sulla base di questo principio, è possibile definire il \textit{Potere Risolutivo} $R$ di uno strumento come la capacità di discriminare variazioni relative di lunghezza d'onda:
\begin{equation}
    R = \frac{\lambda}{\Delta\lambda_{\min}}.
\end{equation}

\subsection{Approccio analitico}
Siano due onde elettromagnetiche di lunghezza d'onda rispettivamente $\lambda$ e $\lambda' = \lambda + \Delta \lambda$ incidenti su un reticolo di diffrazione con $N$ fenditure equidistanti $d$. Perché in occorrenza del massimo di ordine $m$ della prima onda la seconda sia risolvibile, per il Criterio di Rayleigh, deve avere il suo massimo \textit{almeno} nel primo zero successivo al massimo in questione. I minimi si trovano alle posizioni angolari dettate dalla seguente relazione:
\begin{equation}
    d\sin\theta = \frac{n}{N}\lambda, \quad n \in \mathbb{Z}.
\end{equation}
L'$m$-esimo massimo si trova per $n = mN$: il minimo successivo, dunque con $n = mN +1$:
\begin{equation}
    d\sin\theta = \frac{mN+1}{N}\lambda = \lambda m + \frac{\lambda}{N}, \quad m \in \mathbb{Z}.
\end{equation}
Imponendo che la posizione angolare appena individuata descriva il massimo di ordine $m$ per l'onda contraddistinta da $\lambda'$ si ottiene:
\begin{equation}
    m\lambda' = \lambda m + \frac{\lambda}{N}.
\end{equation}
Da cui, la minima variazione di lunghezza d'onda corrisponde a:
\begin{equation}
    \Delta\lambda = \frac{\lambda}{mN}.
\end{equation}
Il potere risolutivo, infine:
\begin{equation}
    R = \frac{\lambda}{\Delta\lambda} = {mN}.
\end{equation}

\section{Diffrazione alla Fraunhofer}
L'analisi condotta sui sistemi discreti di emettitori (doppia fenditura e reticoli) costituisce il fondamento per comprendere il comportamento della luce quando interagisce con ostacoli o aperture di estensione continua.
Il fenomeno che emerge quando la luce incontra un'apertura di dimensioni comparabili alla sua lunghezza d'onda è noto come \textit{diffrazione}. Sebbene storicamente distinto dall'interferenza, fisicamente i due fenomeni sono identici: la diffrazione non è altro che l'\textit{auto-interferenza} delle infinite sorgenti secondarie continue che costituiscono il fronte d'onda.

\subsection{Dal discreto al continuo}
Si consideri una fenditura rettangolare di larghezza $a$. In accordo con il principio di Huygens-Fresnel, ogni punto infinitesimo $dx$ all'interno dell'apertura agisce come una sorgente secondaria di onde sferiche.
Il passaggio dalla trattazione precedente a quella attuale è un'operazione di limite: si immagini di avere $N$ sorgenti discrete distribuite sulla larghezza $a$ e di far tendere $N \to \infty$ mantenendo costante l'ampiezza totale del campo emesso.
Matematicamente, la sovrapposizione dei campi non sarà più descritta da una sommatoria discreta ($\sum$), bensì da un integrale definito esteso a tutta l'apertura:
\begin{equation}
    E_{\text{tot}} = \sum_{i=1}^N E_i \xrightarrow{N\to\infty} \int_{-a/2}^{a/2} dE(x).
\end{equation}

\subsection{L'approssimazione di Fraunhofer (Campo lontano)}
La trattazione matematica del problema si semplifica notevolmente se ci si pone nel regime già sfruttato di \textit{Fraunhofer}. Tale regime è caratterizzato da due condizioni geometriche fondamentali:
\begin{enumerate}
    \item la sorgente primaria è posta all'infinito, in modo che l'onda incidente sulla fenditura sia \textit{piana};
    \item lo schermo di osservazione è posto a una distanza $L$ sufficientemente grande rispetto alle dimensioni dell'apertura ($L \gg a$).
\end{enumerate}
Sotto queste ipotesi, i raggi che partono dai diversi punti dell'apertura e convergono in un punto $P$ sullo schermo possono essere considerati paralleli.

\subsection{Analisi quantitativa}
Passando al campo complesso, si hanno dunque infinite sorgenti $S_i$ \textit{coerenti} (la coerenza è garantita dal fatto che discendono dalla medesima onda), sfasata a due a due di un infinitesimo $d\delta$. Dalla poligonale sfruttata per analizzare il caso di Young e con 3 sorgenti, si passa così al continuo dell'\textit{arco di circonferenza} di raggio $R$, $\Lambda$. L'angolo al centro sotteso da $\Lambda$ corrisponde alla differenza di fase totale fra i fronti secondari che si trovano in corrispondenza dei bordi dell'apertura (di ampiezza $b$). Con le solite convenzioni sui simboli, si ottiene $\delta = kb\sin\theta$. Il campo risultante dalla sovrapposizione di tutti gli infinitesimi contributi deve essere la corda che sottende $\Lambda$:
\begin{equation}
    E_{\text{tot}} = 2R\sin\left(\frac{\delta}{2}\right) = 2R\sin\left(\frac{kd\sin\theta}{2}\right) = 2R\sin\left(\frac{b}{\lambda}\pi \sin\theta\right).
\end{equation}
La somma di tutti i contributi di campo elettrico (ovvero il campo che si otterrebbe se tutti i contributi fossero in fase) restituisce invece:
\begin{equation}
    E_{R} = R\delta = R \frac{2\pi}{\lambda}d\sin \theta,
\end{equation}
da cui:
\begin{equation}
    \frac{E_{\text{tot}}}{E_R} = \frac{2R\sin\left(\frac{\delta}{2}\right)}{R\delta} = \frac{\sin\left(\frac{\delta}{2}\right)}{\frac{\delta}{2}} = \frac{\sin\left(\frac{b}{\lambda}\pi \sin\theta\right)}{ \frac{b}{\lambda}\pi\sin \theta}.
\end{equation}
L'intensità, dunque:
\begin{equation}
   I_{\text{tot}} \propto E_{\text{tot}}^2 = E_R^2 \frac{\sin\left(\frac{b}{\lambda}\pi \sin\theta\right)^2}{\left( \frac{b}{\lambda}\pi\sin \theta\right)^2}.
\end{equation}

\subsection{Studio della figura di diffrazione}
Per analizzare la distribuzione dell'intensità sullo schermo, è conveniente introdurre il parametro adimensionale $\beta$, che racchiude la dipendenza angolare:
\begin{equation}
    \beta = \frac{\pi b \sin\theta}{\lambda}.
\end{equation}
L'espressione dell'intensità assume così la forma compatta:
\begin{equation}
    I(\theta) = I_0 \left( \frac{\sin\beta}{\beta} \right)^2,
\end{equation}
dove $I_0$ rappresenta l'intensità massima ideale ($E_R^2$). Si procede ora all'individuazione dei punti stazionari.

\subsubsection{Minimi di intensità (Frange scure)}
L'intensità si annulla quando il numeratore della funzione è nullo, a condizione che il denominatore non lo sia contemporaneamente.
\begin{equation}
    \sin\beta = 0 \quad \wedge \quad \beta \neq 0.
\end{equation}
Ciò si verifica per valori interi di $\pi$:
\begin{equation}
    \beta = m\pi, \quad \text{con } m \in \mathbb{Z} \setminus \{0\}.
\end{equation}
Sostituendo l'espressione di $\beta$, si ottiene la condizione fisica per i minimi di diffrazione:
\begin{equation}
    \frac{\pi b \sin\theta}{\lambda} = m\pi \implies b\sin\theta = m\lambda, \quad m = \pm 1, \pm 2, \dots
\end{equation}
Si nota che per ogni ordine $m$ (escluso lo zero) esiste una coppia di frange scure simmetriche rispetto all'asse ottico.

\subsubsection{Massimo centrale}
Per $\theta \to 0$, e quindi $\beta \to 0$, si incontra una forma indeterminata $0/0$. Il limite è noto:
\begin{equation}
    \lim_{\beta \to 0} \frac{\sin\beta}{\beta} = 1.
\end{equation}
Di conseguenza $I(0) = I_0$. Questo corrisponde al \textit{massimo centrale}, che risulta essere il picco di intensità dominante dell'intera figura.
La larghezza angolare del massimo centrale è definita dalla distanza tra i due primi minimi ($m = +1$ e $m = -1$):
\begin{equation}
    \Delta\theta \approx 2\sin\theta_1 = 2\frac{\lambda}{b}.
\end{equation}
Tale estensione è esattamente doppia rispetto alla larghezza dei lobi successivi (compresi tra $m$ e $m+1$).

\subsubsection{Massimi secondari}
Oltre al massimo centrale, la funzione presenta una serie di massimi locali di intensità decrescente. La loro posizione si individua annullando la derivata prima dell'intensità rispetto a $\beta$:
\begin{equation}
    \frac{dI}{d\beta} = 2 \frac{\sin\beta}{\beta} \left( \frac{\beta\cos\beta - \sin\beta}{\beta^2} \right) = 0.
\end{equation}
Escludendo il caso del massimo centrale ($\beta \to 0$), la condizione di stazionarietà richiede che:
\begin{equation}
    \beta\cos\beta = \sin\beta \implies \tan\beta = \beta.
\end{equation}
Le soluzioni di questa equazione trascendente si trovano graficamente intersecando la retta $y=\beta$ con i rami della funzione tangente. Le radici non si trovano esattamente a metà tra i minimi ($\beta \neq (m+1/2)\pi$), ma tendono asintoticamente a tale valore per $m$ grandi.
Le prime soluzioni approssimate sono:
\begin{itemize}
    \item $1^\circ$ ordine: $\beta \approx 1.43\pi \implies I_1 \approx I_0 (0.0472)$ (circa il $4.7\%$ del picco centrale).
    \item $2^\circ$ ordine: $\beta \approx 2.46\pi \implies I_2 \approx I_0 (0.0165)$ (circa l'$1.6\%$ del picco centrale).
\end{itemize}
La rapida decrescita dell'intensità conferma che la maggior parte dell'energia luminosa è concentrata nel lobo principale.

\subsubsection{Influenza della larghezza della fenditura}
La posizione del primo minimo di diffrazione ($m = \pm 1$) determina l'apertura angolare del fascio diffratto e dipende criticamente dal rapporto $\lambda/b$:
\begin{equation}
    \sin\theta_1 = \frac{\lambda}{b}.
\end{equation}
Si possono distinguere tre regimi fisici di interesse:

\begin{enumerate}
    \item \textit{regime di massima diffrazione ($b \approx \lambda$):}
    nel caso limite in cui la larghezza della fenditura è comparabile alla lunghezza d'onda ($b \to \lambda$), si ha $\sin \theta_1 \to 1$, ovvero $\theta_1 \to \pi/2$.
    Il massimo centrale si allarga fino a coprire l'intero semispazio oltre lo schermo ($\Delta\theta \approx 180^\circ$). In questa condizione, la luce non proietta più un fascio direzionale ma si diffonde in tutte le direzioni;

    \item \textit{regime sub-lunghezza d'onda ($b < \lambda$):}
    se la fenditura è più piccola della lunghezza d'onda, il rapporto $\lambda/b$ risulta maggiore di 1. Poiché il seno non può eccedere l'unità, l'equazione per i minimi non ammette soluzioni reali.
    Fisicamente, ciò significa che \textit{non esistono minimi di diffrazione}: l'intensità decresce monotonamente senza mai annullarsi e la fenditura si comporta a tutti gli effetti come una sorgente puntiforme isotropa di onde sferiche;

    \item \textit{limite dell'ottica geometrica ($b \gg \lambda$):}
    se la fenditura è molto larga rispetto alla lunghezza d'onda (situazione macroscopica), allora $\lambda/b \to 0$. Di conseguenza, $\sin\theta_1 \approx \theta_1 \to 0$.
    Il massimo centrale diventa estremamente stretto e l'intera figura di diffrazione si contrae fino a diventare indistinguibile dalla proiezione geometrica della fenditura. Solo in questo limite è lecito parlare di propagazione rettilinea della luce e trascurare gli effetti ondulatori.
\end{enumerate}










\chapter{La crisi della fisica classica}
Verso la fine del XIX secolo, il panorama della fisica teorica appariva agli occhi della comunità scientifica come un edificio maestoso e sostanzialmente completo. La sintesi operata dalla Meccanica di Newton e dalla Teoria Elettromagnetica di Maxwell sembrava in grado di descrivere la quasi totalità dei fenomeni naturali, dal moto dei pianeti alla propagazione della luce. L'opinione diffusa era che il compito dei fisici futuri si sarebbe limitato a misurare le costanti naturali con maggiore precisione, piuttosto che a scoprire nuove leggi fondamentali.
Tuttavia, all'orizzonte di questo quadro apparentemente idilliaco, si addensavano quelle che Lord Kelvin definì, in una celebre conferenza del 1900, "due piccole nubi": incongruenze sperimentali che il paradigma classico non riusciva a dissipare.
La prima riguardava l'assenza del vento d'etere (esperimento di Michelson-Morley), preludio alla Relatività Ristretta; la seconda riguardava la distribuzione dell'energia nella radiazione di corpo nero, la cui analisi avrebbe portato al crollo della fisica classica.

\section{Le crepe nel modello classico}
Nonostante la potenza predittiva delle equazioni di Maxwell e della termodinamica statistica, tre fenomeni specifici misero in crisi l'approccio continuo e deterministico dell'epoca:

\begin{enumerate}
    \item \textit{la radiazione di corpo nero:} l'applicazione dell'equipartizione dell'energia al campo elettromagnetico in una cavità portava a un risultato assurdo, noto come \textit{catastrofe ultravioletta}. Secondo la teoria classica, un corpo nero avrebbe dovuto emettere una quantità infinita di energia alle alte frequenze, in palese contrasto con l'evidenza sperimentale e con il principio di conservazione dell'energia;
    
    \item \textit{l'effetto fotoelettrico:} l'emissione di elettroni da una superficie metallica colpita da radiazione mostrava una dipendenza dalla frequenza della luce incidente e non dalla sua intensità, comportamento inspiegabile assumendo la natura puramente ondulatoria della luce;
    
    \item \textit{la stabilità atomica:} con la scoperta dell'elettrone e del nucleo, il modello planetario dell'atomo risultava instabile secondo l'elettromagnetismo classico. Una carica accelerata (l'elettrone in orbita) deve irradiare energia, collassando sul nucleo in una frazione di secondo; la stabilità della materia rimaneva dunque un mistero insoluto.
\end{enumerate}
Queste discrepanze, riassumibili da \textit{6 evidenze sperimentali fondamentali}, resero evidente che le leggi di Newton e Maxwell, valide nel mondo macroscopico, fallivano nel descrivere i processi su scala atomica. La risoluzione di tali paradossi richiese l'introduzione di un concetto rivoluzionario: la \textit{quantizzazione} delle grandezze fisiche.

\section{Il problema del corpo nero}
Max Planck, figura chiave della fisica di fine Ottocento, si confrontò con un problema teorico e pratico che avrebbe segnato la nascita della meccanica quantistica. In quel periodo, anche su spinta di necessità industriali legate all'ottimizzazione delle sorgenti di illuminazione, divenne prioritario derivare un modello analitico in grado di descrivere lo spettro di emissione termica universale in funzione della temperatura.\\
\\
Il quadro teorico classico a disposizione di Planck si fondava sulle seguenti certezze:
\begin{itemize}
    \item la materia è composta da oscillatori microscopici carichi (atomi del reticolo) il cui moto è sostenuto dall'\textit{agitazione termica};
    \item il moto è di carattere oscillatorio attorno alle posizioni di equilibrio;
    \item secondo l'\textit{elettrodinamica classica}, una carica accelerata emette radiazione elettromagnetica continua, perdendo energia;
    \item deve valere il \textit{Teorema di Equipartizione dell'Energia}: all'equilibrio termico, ad ogni modo di oscillazione del campo elettromagnetico deve essere associata un'energia media pari a $k_B T$.
\end{itemize}

\subsection{Il modello sperimentale: il corpo nero}
L'oggetto di studio ideale è il cosiddetto \textit{corpo nero}, definito come un corpo capace di assorbire interamente tutta la radiazione elettromagnetica incidente (riflettività nulla). Secondo la legge di Kirchhoff, un perfetto assorbitore è anche un perfetto emettitore.
In laboratorio, il corpo nero viene approssimato da una \textit{cavità isotermica} (un contenitore chiuso con pareti mantenute a temperatura $T$) dotata di un piccolissimo foro verso l'esterno. La radiazione che entra nel foro subisce innumerevoli riflessioni interne e assorbimenti, raggiungendo l'equilibrio termodinamico con le pareti.
La radiazione che riesce a fuoriuscire dal foro è dunque rappresentativa dell'equilibrio termico e il suo spettro (distribuzione dell'intensità in funzione della frequenza) dipende \textit{esclusivamente} dalla temperatura $T$, indipendentemente dal materiale e dalla forma della cavità. \\
\\
Per descrivere analiticamente lo spettro, si assume che la radiazione all'interno della cavità sia in equilibrio termodinamico con gli emettitori sulle pareti, modellizzati come un insieme di oscillatori armonici carichi (dipoli oscillanti).
La grandezza fisica di interesse è la \textit{densità spettrale di energia} $u(\nu, T)$, definita in modo tale che $u(\nu, T)d\nu$ rappresenti l'energia per unità di volume contenuta nell'intervallo di frequenze $[\nu, \nu + d\nu]$.
L'approccio teorico scompone il calcolo di $u(\nu, T)$ nel prodotto di due fattori distinti:
\begin{equation}
    u(\nu, T) = g(\nu) \cdot \langle E \rangle,
    \label{eq:fattorizzazione_planck}
\end{equation}
dove:
\begin{itemize}
    \item $g(\nu)$ rappresenta la \textit{densità dei modi normali} (o densità degli stati), ovvero il numero di onde stazionarie che possono esistere geometricamente all'interno della cavità per unità di volume e per intervallo di frequenza;
    \item $\langle E \rangle$ rappresenta l'\textit{energia media} associata a ciascun modo di oscillazione alla temperatura $T$, determinata dallo scambio energetico con gli oscillatori delle pareti.
\end{itemize}
Mentre il calcolo di $N(\nu)$ è un problema di elettromagnetismo classico (conteggio delle onde stazionarie in una scatola), la determinazione di $\langle E \rangle$ è il passaggio critico dove la fisica classica fallisce e dove Planck introdurrà la sua ipotesi rivoluzionaria.

\subsubsection{Analisi elettromagnetica}
Per determinare il fattore $g(\nu)$, si consideri una cavità ideale di forma cubica e spigolo $L$ con pareti perfettamente conduttrici.
Affinché un'onda elettromagnetica possa esistere in regime stazionario all'interno della cavità, il campo elettrico deve annullarsi sulle pareti. Tale condizione al contorno ($\bm{E}_{\parallel} = 0$) discende dalla natura di conduttore ideale delle pareti e garantisce l'assenza di dissipazione energetica e la riflessione totale dell'onda incidente.
Considerando un sistema di riferimento cartesiano con origine in un vertice del cubo e assi allineati agli spigoli, si analizzi inizialmente il problema in termini monodimensionali lungo l'asse delle ascisse. Il risultato verrà successivamente esteso alle tre dimensioni spaziali sfruttando l'indipendenza delle componenti cartesiane.
Si ipotizzi un'onda piana progressiva che si propaga lungo l'asse $x$, descritta dalla funzione:
\begin{equation}
E_1(x,t) = E_0 \cos(kx-\omega t).
\end{equation}
In seguito alla riflessione sulla parete in $x=L$, si genera un'onda regressiva:
\begin{equation}
E_2(x,t) = E'_0 \cos(kx+\omega t).
\end{equation}
Per il principio di sovrapposizione, il campo elettrico risultante all'interno della cavità è dato dalla somma dei due contributi: $E(x,t) = E_1(x,t) + E_2(x,t)$. Si impongono ora le condizioni al contorno.
In $x=0$, il campo deve essere nullo per ogni istante $t$:
\begin{equation}
E(0,t) = E_0 \cos(-\omega t) + E'_0 \cos(\omega t) = (E_0 + E'_0)\cos(\omega t) = 0 \implies E'_0 = -E_0.
\end{equation}
Sostituendo questa relazione nell'espressione del campo totale, si ottiene:
\begin{equation}
E(x,t) = E_0 [\cos(kx-\omega t) - \cos(kx+\omega t)].
\end{equation}
Applicando le formule di prostaferesi, l'espressione si trasforma nel prodotto di una parte spaziale e una temporale:
\begin{equation}
E(x,t) = 2E_0 \sin(kx)\sin(\omega t).
\label{eq:onda_stazionaria_1D}
\end{equation}
Imponendo infine la condizione di nullità sulla parete opposta ($x=L$):
\begin{equation}
E(L,t) = 2E_0 \sin(kL)\sin(\omega t) = 0 \quad \forall t.
\end{equation}
Affinché l'equazione sia soddisfatta per ogni tempo $t$, deve annullarsi il termine spaziale, ovvero l'argomento del seno deve essere un multiplo intero di $\pi$:
\begin{equation}
kL = n\pi \quad \implies \quad k = \frac{n\pi}{L},
\end{equation}
con $n \in \mathbb{N}$ (numero intero positivo).
Espandendo il ragionamento alle tre dimenioni, si evince che una qualsiasi onda esistente nella scatola si può descrivere in funzione di un vettore d'onda $\bm{k} = (k_x, k_y, k_z)$, le cui componenti devono quindi soddisfare la condizione di \textit{quantizzazione}:
\begin{equation}
    k_x = \frac{n_x \pi}{L}, \quad k_y = \frac{n_y \pi}{L}, \quad k_z = \frac{n_z \pi}{L},
\end{equation}
dove $n_x, n_y, n_z$ sono numeri interi positivi ($n_i \in \mathbb{N}^+$), essendo il numero d'onda $k$ sempre positivo per definizione. Si deduce allora che $k_{i}-k_{i-1} = \pi/L$. In un sistema di riferimento tridimensionale, il generico vettore $\bm{k}$ può descrivere un ottante di sfera, di volume:
\begin{equation}
    V_O = \frac{\pi |\bm{k}|^3}{6}.
\end{equation}
Un "cubetto" $\mathcal{C}$ di lato $\pi/L$ con i lati paralleli al sistema tridimensionale, descrive in linea di principio tante terne di numeri d'onda, e dunque onde, quanti sono i suoi vertici (8): ogni vertice è però condiviso con 8 cubi. Ogni $\mathcal{C}$ descrive in media dunque un'onda. Il volume di $\mathcal{C}$:
\begin{equation}
V_{\mathcal{C}} = \left(\frac{\pi}{L}\right)^3,
\end{equation}
da cui il nuemero totale di cubetti, e quindi di onde, nell'ottante risulta:
\begin{equation}
\mathcal{N}_{\text{onde}}(\bm{k}) = \frac{V_O}{V_{\mathcal{C}}} = \frac{\pi |\bm{k}|^3}{6} \left(\frac{L}{\pi}\right)^3 = \frac{L^3|\bm{k}|^3}{6\pi^2}.
\end{equation}
Ricordando che $|\bm{k}| = 2\pi/\lambda = 2\pi\nu/c$, si ottiene:
\begin{equation}
\mathcal{N}_{\text{onde}}(\nu) = \frac{4}{3}\pi \frac{L^3}{c^3}\nu^3.
\end{equation}
Per completare la trattazione è necessario rendersi conto che una qualsiasi onda nella scatola è scomponibile come combinazione lineare di un'onda con campo elettrico parallelo al generico $E_0$ e una con $E \perp E_0$: manca pertanto un fattore 2:
\begin{equation}
\mathcal{N}_{\text{tot}}(\nu) = 2\mathcal{N}_{\text{onde}}(\nu) = \frac{8}{3}\pi \frac{L^3}{c^3}\nu^3.
\end{equation}
Dal momento che interessa esclusivamente il numero di onde alla frequenza $\nu$ comprese nell'intervallo $[\nu, \nu+d\nu]$, si sfrutta l'espansione di Taylor:
\begin{equation}
d\mathcal{N}_{\text{tot}}(\nu) =  \mathcal{N}_{\text{tot}}(\nu+d\nu) - \mathcal{N}_{\text{tot}}(\nu) \approx \frac{\partial \mathcal{N}_{\text{tot}}(\nu)}{\partial \nu}d\nu = \frac{8\pi L^3}{c^3}\nu^2 d\nu,
\end{equation}
da cui:
\begin{equation}
\frac{d\mathcal{N}_{\text{tot}}(\nu)}{d\nu} = \frac{8\pi L^3}{c^3}\nu^2.
\label{eq:densita_modi}
\end{equation}
La densità di modi, essendo il numero di onde per unità di volume e intervallo di frequenza, risulta dunque:
\begin{equation}
g(\nu) = \frac{8\pi}{c^3}\nu^2.
\end{equation}



\subsubsection{Analisi meccanico-statistica}
Una volta determinato il numero di modi per unità di volume, il problema si sposta sulla determinazione dell'energia media $\langle E \rangle$ posseduta da ciascun modo di oscillazione alla temperatura $T$. È in questo frangente che emerge la discrepanza fondamentale tra la previsione classica e i dati sperimentali.
Secondo la fisica classica, basata sulla statistica di Maxwell-Boltzmann applicata a variabili continue, vale il \textit{Teorema di Equipartizione dell'Energia}. Esso afferma che, all'equilibrio termico, ad ogni grado di libertà quadratico dell'Hamiltoniana è associata un'energia media pari a $\frac{1}{2}k_B T$.
Poiché un oscillatore armonico possiede due gradi di libertà quadratici (energia cinetica $K \propto v^2$ ed energia potenziale $U \propto x^2$), l'energia media prevista è:
\begin{equation}
    \langle E \rangle_{classica} = k_B T,
\end{equation}
dove $k_B$ è la costante di Boltzmann.
Sostituendo tale valore nella formula della densità spettrale \eqref{eq:fattorizzazione_planck} insieme alla densità dei modi \eqref{eq:densita_modi}, si ottiene la legge di \textit{Rayleigh-Jeans}:
\begin{equation}
    u(\nu, T) = \frac{8\pi \nu^2}{c^3} k_B T.
\end{equation}
Tale relazione riproduce bene i dati sperimentali alle basse frequenze, ma diverge per $\nu \to \infty$ ($u \to \infty$). Questa divergenza, che implica un'energia totale infinita nella cavità, è nota storicamente come \textit{catastrofe ultravioletta}.

\subsubsection{L'ipotesi di Planck e la quantizzazione}
Per risolvere l'incongruenza, Planck introdusse un'ipotesi \textit{ad hoc}, priva di giustificazione nella meccanica classica: l'energia scambiata dagli oscillatori con il campo elettromagnetico non è una variabile continua, ma è discretizzata in "pacchetti" (quanti) proporzionali alla frequenza.
L'energia $E_n$ di un oscillatore a frequenza $\nu$ può assumere solo valori discreti:
\begin{equation}
    E_n = n h \nu, \quad \text{con } n = 0, 1, 2, \dots
\end{equation}
dove $h$ è una costante fondamentale (costante di Planck).
Si ricalcola dunque l'energia media $\langle E \rangle$ utilizzando la distribuzione di Boltzmann per livelli discreti. La probabilità $P(n)$ che un oscillatore si trovi nello stato energetico $n$-esimo è data dal quoziente della sua energia ($\propto \exp{(-\frac{E_n}{k_B T}})$) per la somma di tutte le energie disponibili:
\begin{equation}
    P(n) = \frac{e^{-\frac{E_n}{k_B T}}}{\sum_{n=0}^{\infty} e^{-\frac{E_n}{k_B T}}}.
\end{equation}
Il valore medio è quindi la somma di tutte le energia moltiplicate per la loro probabilità associata:
\begin{equation}
    \langle E \rangle = \frac{\sum_{n=0}^{\infty} E_n e^{-\frac{E_n}{k_B T}}}{\sum_{n=0}^{\infty} e^{-\frac{E_n}{k_B T}}} = \frac{\sum_{n=0}^{\infty} n h \nu \, e^{-\frac{n h \nu}{k_B T}}}{\sum_{n=0}^{\infty} e^{-\frac{n h \nu}{k_B T}}}.
\end{equation}
Per semplificare il calcolo, si ponga $x = e^{-\frac{h \nu}{k_B T}}$. Poiché l'esponente è negativo, si ha $|x| < 1$.
Il denominatore è una serie geometrica convergente:
\begin{equation}
    \sum_{n=0}^{\infty} x^n = \frac{1}{1-x}.
\end{equation}
Il numeratore può essere riscritto sfruttando la derivata della serie geometrica:
\begin{equation}
    \sum_{n=0}^{\infty} n x^n = x \frac{d}{dx} \left( \sum_{n=0}^{\infty} x^n \right) = x \frac{d}{dx} \left( \frac{1}{1-x} \right) = \frac{x}{(1-x)^2}.
\end{equation}
Sostituendo queste espressioni nella formula dell'energia media:
\begin{equation}
    \langle E \rangle = h\nu \frac{\frac{x}{(1-x)^2}}{\frac{1}{1-x}} = h\nu \frac{x}{1-x} = \frac{h\nu}{\frac{1}{x} - 1}.
\end{equation}
Reinserendo la definizione $x = e^{-\frac{h \nu}{k_B T}}$, si ottiene l'espressione finale dell'energia media di Planck:
\begin{equation}
    \langle E \rangle_{Planck} = \frac{h \nu}{e^{\frac{h \nu}{k_B T}} - 1}.
\end{equation}
La relazione ottenuta risulta un modello teorico incredibilmente aderente ai risultati sperimentali: la costante $h$ venne infatti ricavata per interpolazione. Un'ipotesi tanto bella da sembrare sbagliata a Planck...

\subsubsection{La legge di Planck}
Combinando la densità dei modi \eqref{eq:densita_modi} con l'energia media quantistica appena derivata, si giunge alla formulazione completa della densità spettrale di energia per il corpo nero:
\begin{equation}
    u(\nu, T) = \frac{8\pi \nu^2}{c^3} \cdot \frac{h \nu}{e^{\frac{h \nu}{k_B T}} - 1} = \frac{8\pi h \nu^3}{c^3} \frac{1}{e^{\frac{h \nu}{k_B T}} - 1}.
    \label{eq:legge_planck}
\end{equation}
Questa funzione descrive perfettamente lo spettro di corpo nero:
\begin{itemize}
    \item per basse frequenze ($h\nu \ll k_B T$), espandendo l'esponenziale al primo ordine ($e^x \approx 1+x$), si ritrova la legge classica di Rayleigh-Jeans ($\langle E \rangle \to k_B T$);
    \item per alte frequenze ($h\nu \gg k_B T$), il termine esponenziale al denominatore domina, portando la funzione a zero e risolvendo la catastrofe ultravioletta.
\end{itemize}

\section{L'effetto fotoelettrico}
Nel 1887 Heinrich Hertz, durante i suoi celebri esperimenti sulle onde elettromagnetiche, notò incidentalmente che la scarica elettrica tra due elettrodi metallici era facilitata se questi venivano illuminati da luce ultravioletta.
Questo fenomeno, noto come \textit{effetto fotoelettrico}, fu oggetto di indagini quantitative approfondite nei primi anni del '900 (in particolare da P. Lenard e R. Millikan).
L'esperimento consisteva nell'illuminare una superficie metallica (catodo) con radiazione monocromatica e misurare le caratteristiche degli elettroni emessi (fotocorrente) al variare della frequenza $\nu$ e dell'intensità $\mathcal{I}$ dell'onda incidente.

\subsection{Le aspettative della fisica classica}
Secondo la teoria ondulatoria classica (Maxwell), l'energia è distribuita uniformemente sul fronte d'onda e l'intensità è proporzionale al quadrato dell'ampiezza del campo elettrico ($\mathcal{I} \propto E^2$). L'elettrone, interagendo con l'onda, dovrebbe accumulare energia fino a superare il lavoro di estrazione del metallo.
Le previsioni classiche erano dunque:
\begin{enumerate}
    \item \textit{dipendenza dall'intensità:} l'energia cinetica $K_e$ degli elettroni emessi dovrebbe dipendere dall'intensità della luce. Un'onda più intensa (campo elettrico $E$ maggiore) trasferisce più energia e forza ($F = qE$) alle cariche: ci si attende $K_e \propto \mathcal{I}$;
    \item \textit{indipendenza dalla frequenza:} la frequenza $\nu$ determina il colore della luce (o l'oscillazione del campo), ma non la quantità di energia trasportata, che dipende solo dall'ampiezza. Non ci si attende alcuna correlazione tra $K_e$ e $\nu$;
    \item \textit{tempo di latenza:} poiché l'energia è distribuita sull'intero fronte d'onda, un singolo elettrone (che occupa una sezione d'urto piccolissima) necessiterebbe di un tempo finito e calcolabile per accumulare l'energia sufficiente a fuggire dal reticolo. Per luci di bassa intensità, questo ritardo dovrebbe essere misurabile (anche dell'ordine dei minuti);
    \item \textit{emissione a qualsiasi frequenza:} con un'intensità sufficientemente alta o un tempo di esposizione sufficientemente lungo, l'emissione dovrebbe avvenire a qualsiasi frequenza.
\end{enumerate}

\subsection{Le evidenze sperimentali}
I risultati ottenuti da Lenard (1902) furono in netto contrasto con le previsioni classiche, delineando una profonda crisi del modello ondulatorio:

\begin{enumerate}
    \item \textit{energia cinetica e frequenza:} l'energia cinetica massima degli elettroni emessi \textit{non dipende dall'intensità} della luce, bensì dipende linearmente dalla sua frequenza $\nu$;
    
    \item \textit{frequenza di soglia:} esiste una frequenza minima $\nu_{\text{min}}$ (o di soglia), caratteristica del materiale, al di sotto della quale \textit{non si registra alcuna emissione}, indipendentemente dall'intensità della luce o dalla durata dell'esposizione;
    
    \item \textit{intensità e numero di elettroni:} l'intensità della luce incidente ($\mathcal{I}$) influenza esclusivamente il \textit{numero} di elettroni emessi nell'unità di tempo (l'intensità della fotocorrente), ma non la loro energia;
    
    \item \textit{assenza di ritardo:} l'emissione di elettroni è istantanea ($\Delta t < \SI{e-9}{\second}$) anche per intensità luminose bassissime, contraddicendo l'ipotesi dell'accumulo graduale di energia.
\end{enumerate}
Si riporta di seguito l'interpretazione fisica dei dati sperimentali, la cui validazione definitiva fu condotta da R. Millikan.
L'analisi mostrò che la pendenza $m$ della retta di interpolazione coincideva numericamente con la costante $h$ introdotta da Planck nello studio del corpo nero. \\
\begin{figure}[h]
\centering
\begin{tikzpicture}[>=Stealth, scale=1.5]

    % --- Definizioni ---
    \def\nuZero{2}       % Frequenza di soglia (x-intercept)
    \def\workFunc{1.5}   % Lavoro di estrazione (valore assoluto y-intercept)
    \def\slope{0.75}     % Pendenza (h)
    \def\maxNu{6}        % Estensione asse x

    % --- Assi ---
    % Asse X (Frequenza)
    \draw[->, thick] (-0.5, 0) -- (\maxNu, 0) node[below left] {$\nu$};
    % Asse Y (Energia Cinetica)
    \draw[->, thick] (0, -2) -- (0, 3.5) node[above left] {$K_{\max}$};
    \node[below left] at (0,0) {O};

    % --- Il Grafico ---
    % Coordinate chiave
    \coordinate (Y_intercept) at (0, -\workFunc);
    \coordinate (Threshold) at (\nuZero, 0);
    \coordinate (End_point) at (5.5, {(\slope * (5.5 - \nuZero))});

    % 1. Parte estrapolata (non fisica, K < 0) - Tratteggiata
    \draw[dashed, thick, red!70!black] (Y_intercept) -- (Threshold);

    % 2. Parte fisica (K > 0) - Linea solida
    \draw[very thick, blue!80!black] (Threshold) -- (End_point) node[right, font=\small] {$K_{\max} = h\nu - \phi$};

    % --- Elementi significativi ---
    
    % Punto Frequenza di soglia
    \filldraw[black] (Threshold) circle (1.5pt);
    \node[below right] at (Threshold) {$\nu_0$};

    % Punto Lavoro di estrazione (-phi)
    \filldraw[red!70!black] (Y_intercept) circle (1.5pt);
    \node[left, red!70!black] at (Y_intercept) {$-\phi$};

    % Triangolo della pendenza (h)
    \coordinate (A) at (3.5, {\slope * (3.5 - \nuZero)});
    \coordinate (B) at (4.5, {\slope * (3.5 - \nuZero)});
    \coordinate (C) at (4.5, {\slope * (4.5 - \nuZero)});

    \draw[thin, gray] (A) -- (B) node[midway, below, font=\scriptsize] {$\Delta \nu$};
    \draw[thin, gray] (B) -- (C) node[midway, right, font=\scriptsize] {$\Delta K$};
    \draw[thin, gray] (C) -- (A);
    
    % Etichetta pendenza
    \node[font=\small, align=center] at (3.8, 2) {pendenza\\$m = h$};

    % --- Annotazioni extra (opzionale) ---
    % Linea orizzontale tratteggiata per phi (opzionale, per chiarezza geometrica)
    \draw[dotted, gray] (0,0) -- (Threshold); 

\end{tikzpicture}
\end{figure}
\\
Albert Einstein, nel 1905, fornì la spiegazione teorica di tale evidenza: il fatto che la stessa costante $h$ governasse due fenomeni apparentemente distinti suggeriva una natura intrinseca della radiazione stessa.
Il tedesco postulò che l'energia dell'onda non fosse distribuita in modo continuo sul fronte d'onda, ma concentrata in pacchetti discreti (quanti di luce), successivamente denominati \textit{fotoni}, ciascuno con energia $E = h\nu$.
L'interazione radiazione-materia viene dunque descritta come un urto (o meglio, un assorbimento) puntuale tra un singolo fotone e un singolo elettrone. Applicando il \textit{principio di conservazione dell'energia}, l'energia del fotone incidente viene interamente trasferita all'elettrone per vincere il potenziale di estrazione $\phi$ e conferire energia cinetica.
L'equazione di Einstein per l'effetto fotoelettrico è dunque:
\begin{equation}
K_{\max} = h\nu - \phi.
\end{equation}
Si specifica $K_{\max}$ poiché questa relazione descrive gli elettroni più superficiali, che non subiscono perdite energetiche per collisioni interne al reticolo prima di essere emessi.
Da questa relazione discende naturalmente l'esistenza di una frequenza di soglia $\nu_{\text{min}}$, al di sotto della quale il fotone non possiede energia sufficiente per estrarre l'elettrone ($h\nu < \phi$):
\begin{equation}
\nu_{\text{min}} = \frac{\phi}{h}.
\end{equation}
\\
Se Young aveva dimostrato che la radiazione elettromagnetica deve essere la manifestazione di un fenomeno ondulatorio, Einstein discute delle sue caratteristiche corpuscolari: un'ambivalenza classicamente inammissibile.

\section{Il modello atomico a righe spettrali}
A seguito dell'esperimento di Rutherford che comprovò l'impossibilità del \textit{modello a panettone} per l'atomo di Thompson, Bohr si trovò ad avanzare ipotesi che supponevano le seguenti conoscenze: 
\begin{itemize}
    \item \textit{la validità del modello nucleare:} in accordo con i risultati di Rutherford (scattering delle particelle $\alpha$), l'atomo deve possedere un nucleo centrale carico positivamente, in cui è concentrata la quasi totalità della massa, attorno al quale orbitano gli elettroni;
    
    \item \textit{l'instabilità elettrodinamica classica:} secondo le equazioni di Maxwell e la formula di Larmor, una carica elettrica accelerata (come un elettrone in moto circolare uniforme) deve irradiare energia con continuità\footnote{Una carica accelerata produce, infatti, un campo magnetico tempo variante.}. Perdendo energia, l'elettrone dovrebbe collassare sul nucleo in tempi dell'ordine di $10^{-10}$ secondi, rendendo la materia instabile;
    
    \item \textit{la legge empirica di Rydberg-Balmer:} dagli studi di spettroscopia era nota una precisa relazione matematica ($1/\lambda = R (1/n_f^2 - 1/n_i^2)$) in grado di descrivere le lunghezze d'onda delle righe spettrali dell'idrogeno. Tale formula, tuttavia, era puramente fenomenologica: descriveva \textit{come} erano distribuite le righe, ma non spiegava il \textit{perché} fisico di quella regolarità basata su numeri interi;
    
    \item \textit{la quantizzazione dell'energia:} la recente introduzione del concetto di \textit{quanto di azione} da parte di Planck ($E=h\nu$) e la sua applicazione alla natura della luce da parte di Einstein suggerivano che la discretizzazione fosse una proprietà fondamentale dei processi microscopici.
\end{itemize}

\subsection{L'ipotesi di Bohr}
Bohr, che analizzò l'atomo di idrogeno sperando di riuscire ad estendere la trattazione per atomi più complessi, assunse il modello orbitale newtoniano per il moto dell'elettrone attorno al nucleo (un protone), a meno di un \textit{vincolo aggiunto ad hoc sul momento angolare}:
\begin{equation}
L = n\frac{h}{2\pi} = n\hbar,
\end{equation}
dove si è introdotta la \textit{costante di Planck ridotta} $\hbar$ e $n \in \mathbb{N}$. Fisicamente, l'introduzione di una quantizzazione del momento angolare suggerisce implicazioni fondamentali. 

\subsubsection{Il modello atomico di Bohr}
L'elettrone è soggetto all'attrazione coulombiana esercitata dal protone attorno al quale orbita. La risultante delle forze sulla particella si scrive:
\begin{equation}
m\frac{v^2}{r} = \frac{e^2}{4 \pi \varepsilon_0 r^2}.
\label{eq:new_H}
\end{equation}
La condizione di quantizzazione implica:
\begin{equation}
L = mvr = n\hbar \implies r = \frac{n\hbar}{mv}
\end{equation}
e sostituendo nella \eqref{eq:new_H} si ottiene:
\begin{equation}
mv^2 = \frac{e^2}{4\pi \varepsilon_0} \frac{mv}{n\hbar} \implies v(n) = \frac{e^2}{2\varepsilon_0 h}\frac{1}{n}.
\end{equation}
La velocità orbitale deve variare nel discreto, fatto classicamente impensabile. L'espressione per il raggio orbitale $r$, allora:
\begin{equation}
r(n) = \frac{nh}{2m\pi}\frac{2\varepsilon_0 h n}{e^2} = \frac{h^2\varepsilon_0}{m\pi e^2}n^2,
\end{equation}
da cui anche $r$ deve essere quantizzato. L'energia meccanica dell'elettrone nel sistema orbitale, proponendo una similitudine con il caso gravitazionale\footnote{$E = -\frac{mMG}{2R}$.} ancora
\begin{equation}
E(n) = -\frac{e^2}{4\pi\varepsilon_0}\frac{1}{2r} = -\frac{e^2}{8\pi\varepsilon_0} \frac{m\pi e^2}{n^2 h^2 \varepsilon_0} = -\frac{e^4 m}{8\varepsilon_0^2 h^2}\frac{1}{n^2},
\label{eq:E_bohr}
\end{equation}
con $n > 0$. \textit{Anche i livelli energetici concessi all'elettrone devono essere discreti}. 

\subsection{Conseguenze dell'ipotesi di Bohr}
Il modello di Bohr si sposa perfettamente con diversi risultati sperimentali.

\subsubsection{Raggio atomico}
Nel caso dell'idrogeno, l'energia potenziale minore (e dunque la configurazione all'equilibrio) si ottiene per $n = 1$. Il valore di raggio orbitale corrispondente
\begin{equation}
r(n = 1) = \frac{h^2\varepsilon_0}{m\pi e^2} \approx \SI{0.529e-10}{\meter}
\end{equation}
è assolutamente consistente con le misurazioni sperimentalmente effettuate.
 
\subsubsection{Energia di prima ionizzazione}
Il lavoro necessario per estrarre l'elettrone all'atomo di idrogeno è l'inverso della sua energia meccanica a riposo (quindi con $n = 1$):
\begin{equation}
W = -E(n = 1) = \frac{e^4 m}{8\varepsilon_0^2 h^2} \approx \SI{2e-18}{\joule},
\end{equation}
valore coerente con le previsioni derivanti dalle analisi chimiche.

\subsubsection{Legge di Rydberg-Balmer}
Quando un elettrone si sposta da un livello energetico $n$ ad uno "più basso", contraddistinto da $m < n$, perde energia. Bohr, alla luce dell'interpretazione di Einstein dell'effetto fotoelettrico, ipotizzò che la diminuzione di energia meccanica corrispondesse all'emissione di un quanto di energia (un fotone $\implies$ emissione di radiazione) $h\nu$:
\begin{equation}
\Delta E = E(n) - E(m) = h\nu = \frac{hc}{\lambda}.
\end{equation}
Dalla \eqref{eq:E_bohr}, inoltre:
\begin{equation}
\begin{aligned}
E(n) - E(m) &= \frac{e^4 m}{8\varepsilon_0^2 h^2}\left(\frac{1}{m^2}-\frac{1}{n^2} \right) = \frac{hc}{\lambda} \\
\implies \frac{1}{\lambda} &= \underbrace{\frac{e^4 m}{8\varepsilon_0^2 h^3c}}_{=R}\left(\frac{1}{m^2}-\frac{1}{n^2} \right)
\end{aligned}
\end{equation}
con $R$ la costante di Rydberg, calcolata sperimentalmente.

\section{Calori specifici}
Il calore specifico associato ad una trasformazione di un gas, in termodinamica classica, corrisponde al calore scambiato da una mole di particelle per unità di temperatura: 
\begin{equation}
c = \frac{\delta Q}{dT} = \frac{dU + \delta W}{dT}.
\end{equation}

\subsubsection{Teorema di equipartizione dell'energia}
Il \textit{Teorema di equipartizione dell'energia}, risultato classico che deriva dall'assunzione che l'energia si distrbuisca nel ontinuo, asserisce che a ogni grado di libertà quadratico di una particella è associato un termine energetico del tipo:
\begin{equation}
E = \frac{1}{2}k_BT,
\end{equation}
con $k_B$ la costante di Boltzmann. Nel caso di molecole biatomiche, si hanno $l = 7$ gradi di libertà derivanti dall'analisi dei moti indipendenti che la molecola può compiere nello spazio.
L'Hamiltoniana classica del sistema può essere scomposta nella somma di tre contributi indipendenti: traslazionale, rotazionale e vibrazionale.

\begin{enumerate}
    \item \textit{Gradi di libertà traslazionali (3):}
    Il centro di massa della molecola può muoversi liberamente nello spazio tridimensionale. L'energia cinetica associata è:
    \begin{equation}
        K_{\text{trasl}} = \frac{1}{2}M(v_x^2 + v_y^2 + v_z^2),
    \end{equation}
    dove $M$ è la massa totale della molecola. Poiché le tre componenti della velocità ($v_x, v_y, v_z$) sono indipendenti e compaiono al quadrato, esse contribuiscono con \textit{3 termini} $k_B T/2$ all'energia media.
    
    \item \textit{Gradi di libertà rotazionali (2):}
    Una molecola biatomica è lineare. Definendo l'asse $z$ coincidente con l'asse di legame interatomico, il momento d'inerzia rispetto a tale asse è trascurabile ($I_z \approx 0$) trattandosi di masse puntiformi. La molecola può ruotare con energia significativa solo attorno ai due assi ortogonali al legame ($x$ e $y$). L'energia cinetica rotazionale è:
    \begin{equation}
        K_{\text{rot}} = \frac{1}{2}I(\omega_x^2 + \omega_y^2).
    \end{equation}
    Le due velocità angolari indipendenti ($\omega_x, \omega_y$) generano \textit{2 termini} quadratici, contribuendo per $2 \times k_B T/2 = k_B T$.
    
    \item \textit{Gradi di libertà vibrazionali (2):}
    Gli atomi possono oscillare lungo l'asse di legame avvicinandosi e allontanandosi l'uno dall'altro. Questo moto è modellizzabile come un oscillatore armonico. L'energia totale di un oscillatore è somma di un termine cinetico (velocità di vibrazione) e di un termine potenziale (elasticità del legame):
    \begin{equation}
        E_{\text{vib}} = K_{\text{vib}} + U_{\text{vib}} = \frac{1}{2}\mu \dot{r}^2 + \frac{1}{2}k (r-r_{eq})^2,
    \end{equation}
    dove $\mu$ è la massa ridotta e $k$ la costante elastica.
    È fondamentale notare che, sebbene vi sia una sola coordinata di vibrazione ($r$), essa introduce nell'Hamiltoniana \textit{2 termini quadratici} (uno cinetico e uno potenziale). Per il teorema di equipartizione, il contributo energetico è dunque $2 \times k_B T/2 = k_B T$.
\end{enumerate}
Sommando i contributi si ottiene il totale classico di $3 + 2 + 2 = 7$ termini quadratici, corrispondenti a un'energia totale di: 
    \begin{equation}
        E = \frac{7}{2}k_BT.
    \end{equation}
    L'energia interna di una mole di particelle, quindi:
        \begin{equation}
        E = \frac{7}{2}N_Ak_BT = \frac{7}{2}RT,
    \end{equation}
dove si è sfruttato $N_Ak_B = R$. Il calore specifico a volume costante ($\delta W = 0$) per una mole di gas biatomico, di conseguenza risulta:
        \begin{equation}
        c = \frac{dU}{dT} = \frac{7}{2}R\frac{dT}{dT} = \frac{7}{2}R,
    \end{equation}
    costante al variare della temperatura.  
   

    
    
    \subsection{Il "congelamento" dei gradi di libertà e la natura discreta dell'energia}
Il confronto con i dati sperimentali evidenzia una netta discrepanza con la previsione classica di un calore specifico costante ($c_V = \frac{7}{2}R$). Analizzando l'andamento di $c_V$ in funzione della temperatura per un gas biatomico (come $H_2$), si osserva una curva "a gradini" che smentisce l'equiparietà continua dell'energia:
\begin{itemize}
    \item a basse temperature ($T < \SI{100}{\kelvin}$), il gas si comporta come se fosse monoatomico ($c_V = \frac{3}{2}R$), possedendo solo energia traslazionale;
    \item a temperatura ambiente, il calore specifico sale a $\frac{5}{2}R$, segnalando l'attivazione dei gradi di libertà rotazionali;
    \item solo ad alte temperature si raggiunge il valore asintotico classico $\frac{7}{2}R$, con l'attivazione dei moti vibrazionali.
\end{itemize}
Tale comportamento si spiega ammettendo che l'energia associata ai moti interni non vari con continuità, ma sia quantizzata. Affinché un grado di libertà possa assorbire calore (e quindi contribuire a $c_V$), l'energia termica disponibile $k_B T$ deve essere comparabile con la differenza di energia tra il livello fondamentale e il primo stato eccitato ($\Delta E$). Se $k_B T \ll \Delta E$, il grado di libertà rimane "congelato" nello stato fondamentale.
La giustificazione teorica di tali soglie energetiche si fonda direttamente sulle ipotesi di Planck e Bohr precedentemente discusse:
\begin{enumerate}
    \item \textit{Quantizzazione rotazionale (Bohr):} Analogamente a quanto postulato da Bohr per l'atomo di idrogeno, il momento angolare $L$ della molecola è quantizzato ($L = n\hbar$). Di conseguenza, l'energia rotazionale assume valori discreti:
    \begin{equation}
        E(n) = \frac{L^2}{2I} = \frac{n^2\hbar^2}{2I}.
    \end{equation}
    Essendo il momento d'inerzia $I$ molto piccolo, il "gap" energetico è significativo e richiede temperature dell'ordine delle centinaia di Kelvin per essere superato.
    
    \item \textit{Quantizzazione vibrazionale (Planck):} Il moto oscillatorio degli atomi lungo l'asse di legame ricade perfettamente nel modello dell'oscillatore armonico quantistico introdotto da Planck. L'energia può assumere solo valori:
    \begin{equation}
        E(n) = nh\nu = n\hbar\omega.
    \end{equation}
    Le frequenze di vibrazione molecolare $\omega$ sono tipicamente molto elevate, rendendo il "salto" $\hbar\omega$ molto grande: per questo motivo i moti vibrazionali si attivano solo ad alte temperature.
\end{enumerate}


 \section{La quantità di moto del fotone e l'ipotesi di De Broglie}
Per un oggetto dotato di massa a riposo nulla (come il fotone ipotizzato da Einstein), la meccanica newtoniana ($p=mv$) perde di validità. Si deve invece ricorrere alla \textit{Relatività Ristretta}, la quale stabilisce una relazione generale tra energia totale $E$, quantità di moto $p$ e massa a riposo $m_0$:
\begin{equation}
E^2 = (pc)^2 + (m_0c^2)^2.
\end{equation}
Per un fotone ($m_0 = 0$), l'equazione si semplifica in $E = pc$. Combinando questo risultato con la relazione di Einstein per l'effetto fotoelettrico ($E = h\nu$), si ottiene l'espressione della quantità di moto associata al quanto di luce:
\begin{equation}
p = \frac{E}{c} = \frac{h\nu}{c} = \frac{h}{\lambda}.
\label{eq:momento_fotone}
\end{equation}
Si delinea così un dualismo per la radiazione: essa è descritta da grandezze ondulatorie ($\lambda, \nu$) che sono indissolubilmente legate a grandezze corpuscolari ($p, E$) tramite la costante di Planck $h$.

\subsubsection{L'ipotesi di De Broglie}
Nel 1924, Louis de Broglie spinse questo dualismo oltre, guidato da considerazioni di simmetria della natura. Se un ente fisico tradizionalmente ondulatorio (la luce) può manifestare proprietà corpuscolari, è ragionevole ipotizzare che un ente tradizionalmente corpuscolare (come l'elettrone, dotato di massa $m$) possa manifestare proprietà ondulatorie. \\
De Broglie propose di estendere la validità della relazione \eqref{eq:momento_fotone} anche alla materia. A una particella che si muove con quantità di moto $p$, deve essere associata un'onda di materia con lunghezza d'onda $\lambda_{dB}$:
\begin{equation}
\lambda_{dB} = \frac{h}{p} = \frac{h}{mv}.
\end{equation}
Questa audace ipotesi implica che anche gli elettroni possano dare luogo a fenomeni di interferenza e diffrazione, fatto che fu brillantemente confermato pochi anni dopo dall'esperimento di Davisson e Germer sulla diffrazione elettronica attraverso un reticolo cristallino.


\section{Diffrazione di elettroni da cristalli}

Prova della \textit{natura ondulatoria della materia} arriva da un esperimento condotto dai fisici Davisson e Germer nel 1927. Preliminarmente, è opportuno però affrontare il fenomeno della diffrazione di onde elettromagnetiche da parte di reticoli cristallini, storicamente osservato con i raggi X.

\subsection{Diffrazione di onde elettromagnetiche da cristalli}
Affinché si verifichi diffrazione, è necessario che la lunghezza d'onda $\lambda$ della radiazione incidente sia comparabile con le dimensioni caratteristiche degli ostacoli che essa incontra. Nel caso dei solidi cristallini, gli atomi sono disposti in strutture periodiche regolari con spaziature interatomiche $d$ (passo reticolare) dell'ordine di \SI{1}{\angstrom} ($10^{-10}$ m).
Di conseguenza, i raggi X, possedendo lunghezze d'onda in questo intervallo spettrale, sono ottimi "tester" per provare il fenomeno di diffrazione. Siano i campi elettrici di due onde coerenti, che condividono lo stesso fronte, inclinate di un angolo $\theta$ rispetto alla superficie del cristallo descritte dalle seguenti equazioni scalari:
\begin{equation}
E_1(x_1,t) = E_0e^{i(kx_1-\omega t)} \quad E_2(x_2,t) = E_0e^{i(kx_2-\omega t)}.
\end{equation}
Nell'imbattersi su due piani reticolari successivi, come mostrato in Figura (\ref{fig:bragg_diffraction}), accumulano una differenza di cammino ottico pari a:
\begin{equation}
x_1 - x_2 = 2d\sin\theta.
\end{equation}

\begin{figure}[htbp]
    \centering
    \begin{tikzpicture}[scale=1.2, font=\footnotesize]
        % --- Definizione Parametri ---
        \def\d{2.0}       % Spaziatura interplanare (d)
        \def\a{1.5}       % Spaziatura atomica orizzontale
        \def\ang{30}      % Angolo di incidenza (theta)
        \def\raylen{4}    % Lunghezza dei raggi disegnati
        \def\nplanes{3}   % Numero di piani
        \def\natoms{6}    % Atomi per piano

        % --- 1. Disegno del Reticolo Cristallino ---
        \foreach \i in {0,...,\numexpr\nplanes-1} {
            % Piani reticolari (linee tratteggiate)
            \draw[dashed, gray!60] (-\a/2, -\i*\d) -- (\natoms*\a-\a/2, -\i*\d);
            % Atomi (sfere ombreggiate)
            \foreach \j in {0,...,\numexpr\natoms-1} {
                \shade[ball color=blue!20!gray] (\j*\a, -\i*\d) circle (0.2);
            }
        }
        % Etichetta per i piani
        \node[gray, right] at (\natoms*\a-\a/2, 0) {Piani reticolari};

        % Indicazione della spaziatura 'd'
        \draw[<->, >=stealth, thick] (\natoms*\a+0.2, 0) -- node[right] {$d$} (\natoms*\a+0.2, -\d);

        % --- 2. Definizione Punti Chiave per i Raggi ---
        % Scegliamo due atomi su piani adiacenti per l'interazione
        \coordinate (A1) at (3*\a, 0);   % Atomo sul piano superiore
        \coordinate (A2) at (3*\a, -\d); % Atomo sul piano inferiore

        % Punti di partenza e arrivo dei raggi (calcolati con trigonometria)
        \coordinate (Ray1InStart) at ($(A1) + ({-\raylen*cos(\ang)}, {\raylen*sin(\ang)})$);
        \coordinate (Ray1OutEnd)  at ($(A1) + ({\raylen*cos(\ang)}, {\raylen*sin(\ang)})$);
        \coordinate (Ray2InStart) at ($(A2) + ({-\raylen*cos(\ang)}, {\raylen*sin(\ang)})$);
        \coordinate (Ray2OutEnd)  at ($(A2) + ({\raylen*cos(\ang)}, {\raylen*sin(\ang)})$);

        % --- 3. Disegno dei Raggi (Incidente e Diffratto) ---
        % Raggio 1 (Superiore)
        \draw[->, red!80!black, thick, >=latex] (Ray1InStart) -- (A1) node[midway, above, sloped] {Onda incidente};
        \draw[->, red!80!black, thick, >=latex] (A1) -- (Ray1OutEnd) node[midway, above, sloped] {Onda diffratta};

        % Raggio 2 (Inferiore)
        \draw[->, red!80!black, thick, >=latex] (Ray2InStart) -- (A2);
        \draw[->, red!80!black, thick, >=latex] (A2) -- (Ray2OutEnd);

        % --- 4. Geometria della Differenza di Cammino ---
        % Proiezioni ortogonali da A1 sui raggi 2 per trovare i punti M e N
        \coordinate (M) at ($(Ray2InStart)!(A1)!(A2)$); % Piede della perpendicolare su raggio incidente
        \coordinate (N) at ($(Ray2OutEnd)!(A1)!(A2)$);  % Piede della perpendicolare su raggio diffratto

        % Disegno delle perpendicolari (linee di costruzione)
        \draw[dashed, blue, thick] (A1) -- (M);
        \draw[dashed, blue, thick] (A1) -- (N);

        % Evidenziare il cammino aggiuntivo
        \draw[blue, very thick] (M) -- (A2) -- (N);

        % --- 5. Angoli e Etichette ---
        % Linea orizzontale di riferimento per l'angolo theta su A1
        \draw[gray] (A1) -- +(-1.5,0) coordinate (Href);
        \pic [draw, "$\theta$", angle radius=0.7cm, angle eccentricity=1.3] {angle = Ray1InStart--A1--Href};

        % Angoli theta nel triangolo della differenza di cammino (geometricamente uguali a theta di incidenza)
        \pic [draw, "$\theta$", angle radius=0.5cm, angle eccentricity=1.4, blue] {angle = M--A1--A2};
        \pic [draw, "$\theta$", angle radius=0.5cm, angle eccentricity=1.4, blue] {angle = A2--A1--N};

        % Etichette per i segmenti di differenza di cammino usando parentesi graffe
        \draw[decorate, decoration={brace, amplitude=4pt, raise=2pt}, blue] (A2) -- (M)
            node [midway, below left=3pt, sloped, black] {$d\sin\theta$};
        \draw[decorate, decoration={brace, amplitude=4pt, mirror, raise=2pt}, blue] (A2) -- (N)
            node [midway, below right=3pt, sloped, black] {$d\sin\theta$};

        % Fronte d'onda incidente (opzionale, per chiarezza)
        \draw[gray, dotted, thick] (M) -- ++($(A1)-(M)$) -- ++($(Ray1InStart)-(A1)$);
        
     

    \end{tikzpicture}
    \caption{Schema della diffrazione di Bragg da un reticolo cristallino. Due raggi paralleli incidono con angolo $\theta$ sui piani reticolari distanziati di $d$. Il raggio inferiore percorre un tragitto aggiuntivo pari a $2d\sin\theta$ (tratti blu spessi) rispetto al raggio superiore. L'interferenza costruttiva si ha quando questa differenza è un multiplo intero della lunghezza d'onda $\lambda$.}
    \label{fig:bragg_diffraction}
\end{figure}
\noindent
La differenza di fase è dunque data da:
\begin{equation}
    \phi = k(x_1 - x_2) = \frac{4\pi d}{\lambda}\sin\theta
\end{equation}
Si consideri uno schermo posto a grande distanza dal punto di incidenza. L'intensità dell'onda risultante dall'interferenza è proporzionale al quadrato dell'ampiezza del campo. Utilizzando la formula generale per l'interferenza di $N$ sorgenti coerenti:
\begin{equation}
    I \propto \left[ \frac{\sin\left(\frac{N\phi}{2}\right)}{\sin\left( \frac{\phi}{2}\right)} \right]^2
\end{equation}
Nel caso in esame, considerando l'interferenza tra due piani adiacenti ($N = 2$), l'espressione si semplifica notevolmente ricordando che $\sin\phi = 2\sin(\phi/2)\cos(\phi/2)$:
\begin{equation}
    I \propto \left[ \frac{\sin\phi}{\sin(\phi/2)} \right]^2 = \left[ \frac{2\sin(\phi/2)\cos(\phi/2)}{\sin(\phi/2)} \right]^2 = 4\cos^2\left(\frac{\phi}{2}\right)
\end{equation}
Sostituendo il valore di $\phi$ calcolato precedentemente:
\begin{equation}
    I \propto \cos^2\left( \frac{2\pi d}{\lambda}\sin\theta \right)
\end{equation}
La condizione di interferenza costruttiva (massimi di intensità) si verifica quando l'argomento del coseno è un multiplo intero di $\pi$:
\begin{equation}
    \frac{2\pi d}{\lambda}\sin\theta = m\pi \quad \implies \quad \sin\theta = \frac{m\lambda}{2d} \quad (m \in \mathbb{Z})
    \label{eq:condizione_bragg}
\end{equation}
L'interferenza distruttiva (minimi d'intensità) si verifica invece quando l'argomento del coseno è un multiplo dispari di $\pi/2$. La condizione di fase è:
\begin{equation}
    \frac{2\pi d}{\lambda}\sin\theta = (2m+1)\frac{\pi}{2}
\end{equation}
Semplificando, si ottiene la condizione sulla differenza di cammino ottico:
\begin{equation}
    \sin\theta = \left(m + \frac{1}{2}\right) \frac{\lambda}{2d} \quad (m \in \mathbb{Z})
    \label{eq:condizione_bragg_minimi}
\end{equation}
In questo caso, le onde riflesse dai due piani adiacenti risultano in opposizione di fase, annullandosi reciprocamente.
Si ritrova così la celebre \textit{Legge di Bragg}.

\subsubsection{Conseguenze dell'ipotesi di De Broglie}

Se si ammette la validità dell'ipotesi di De Broglie, secondo cui a una particella materiale di quantità di moto $p$ è associata una lunghezza d'onda $\lambda = h/p$, si deve prevedere che anche un fascio di elettroni possa manifestare fenomeni diffrattivi quando interagisce con un cristallo.
Tuttavia, ciò è osservabile solo se l'impulso degli elettroni è tale da rendere la loro lunghezza d'onda di De Broglie comparabile con il passo reticolare $d$. Questa previsione teorica trova la sua conferma sperimentale nell'esperimento di Davisson e Germer (1927).

\subsection{L'esperimento di Davisson e Germer}
La conferma definitiva dell'ipotesi di De Broglie giunse nel 1927 grazie agli esperimenti condotti da C. Davisson e L.H. Germer presso i Bell Labs. Utilizzando un cristallo di nichel come reticolo di diffrazione, essi bombardarono la superficie del metallo con un fascio di elettroni accelerati da un potenziale $V$ variabile. La previsione classica voleva che gli elettroni, essendo \textit{particelle} a tutti gli effetti, assumessero una traiettoria variabile solo dall'angolo d'incidenza sul metallo. \\
\\
\noindent
L'apparato sperimentale permetteva di misurare l'intensità degli elettroni retrodiffusi a vari angoli. Si osservò che, per specifici valori della tensione acceleratrice, l'intensità diffusa presentava dei massimi pronunciati in direzioni ben precise, analogamente a quanto accade nella diffrazione dei raggi X.
In particolare, noto il passo reticolare $d$ del nichel (determinato precedentemente via cristallografia a raggi X) e misurato l'angolo $\theta$ di picco, è possibile calcolare la lunghezza d'onda associata all'elettrone tramite la legge di Bragg ($n=1$):
\begin{equation}
    \lambda_{\text{sper}} = 2d \sin\theta
\end{equation}
Parallelamente, trattando l'elettrone come una particella classica accelerata dal potenziale $V$, la sua quantità di moto $p$ è determinata dalla conservazione dell'energia $eV = p^2 / 2m$:
\begin{equation}
    p_{\text{cin}} = \sqrt{2m_e eV}
\end{equation}
Riportando in un grafico i valori della quantità di moto $p$ calcolati cineticamente in funzione del reciproco della lunghezza d'onda misurata ($1/\lambda_{\text{sper}}$), si ottiene una distribuzione di punti sperimentali che si dispongono lungo una retta passante per l'origine.
La pendenza di tale retta di interpolazione coincide, entro gli errori sperimentali, con la costante di Planck $h$.
\begin{equation}
    p = h \cdot \frac{1}{\lambda}
\end{equation}
Questo risultato sancisce la validità della relazione di De Broglie anche per particelle massive. \\

% --- Grafico P vs 1/Lambda ---
\begin{figure}[htbp]
    \centering
    \begin{tikzpicture}
        \begin{axis}[
            width=0.8\textwidth,
            height=6cm,
            xlabel={Numero d'onda spettroscopico $1/\lambda$ [\SI{e9}{\per\meter}]},
            ylabel={Quantità di moto $p$ [\SI{e-24}{\kilogram\meter\per\second}]},
            grid=major,
            title={Verifica della relazione $p = h/\lambda$},
            legend pos=north west,
            xmin=0, xmax=6,
            ymin=0, ymax=40,
            axis lines=left,
            scaled x ticks=false,
            scaled y ticks=false
        ]

        % Retta teorica di De Broglie (p = h * 1/lambda)
        % h = 6.626e-34.
        % Se x è in 10^9 m^-1 e y in 10^-24, la slope è:
        % (6.626e-34 * 1e9) / 1e-24 = 6.626 * 10^-1 = 0.6626 (fattore di scala assi)
        % Usiamo una slope visuale corretta per i dati simulati:
        % Supponiamo x=1 -> 1e9 m^-1. p = 6.6e-34 * 1e9 = 6.6e-25 = 0.66 * 10^-24.
        % Per rendere il grafico leggibile usiamo unità arbitrarie scalate o dati simulati
        % Slope approssimata per visualizzazione: y = 6.626 * x
        \addplot[
            domain=0:6,
            samples=100,
            color=blue,
            thick
        ]
        {6.626*x};
        \addlegendentry{Previsione teorica ($p=h/\lambda$)}

        % Dati simulati (con leggero rumore sperimentale)
        \addplot[
            only marks,
            mark=*,
            color=red,
            mark size=2.5pt
        ]
        coordinates {
            (1.2, 7.9)
            (2.1, 14.2)
            (3.0, 19.5)
            (3.8, 25.5)
            (4.5, 29.3)
            (5.2, 34.8)
        };
        \addlegendentry{Dati sperimentali}

        % Annotazione della pendenza
        \node at (axis cs: 4, 15) [anchor=north west, text width=4cm, align=left, font=\footnotesize] {
            La pendenza della retta\\
            corrisponde alla costante\\
            di Planck $h$.
        };

        \end{axis}
    \end{tikzpicture}
    \caption{Grafico di interpolazione dei dati sperimentali. Sull'asse delle ascisse è riportato l'inverso della lunghezza d'onda determinata per diffrazione, sulle ordinate la quantità di moto calcolata dal potenziale accelerante. La linearità conferma l'ipotesi di De Broglie.}
    \label{fig:davisson_germer_plot}
\end{figure}

\noindent
Appare dunque necessario un nuovo approccio alla descrizione dinamica dei \textit{corpi massivi}. Le evidenze sperimentali mostrano che non solo la radiazione, ma anche la materia è soggetta a fenomeni di interferenza e diffrazione. L'esperimento di Davisson e Germer ha infatti confermato che anche gli elettroni, tradizionalmente intesi come corpuscoli, necessitano di una descrizione tramite una funzione d'onda $\psi(\vec{r}, t)$ in tutti quei contesti in cui le dimensioni in gioco sono comparabili con la loro lunghezza d'onda di De Broglie.

\subsubsection{Dualità onda-corpuscolo e principio di complementarietà}
Se l'effetto fotoelettrico ha reso evidenti le proprietà corpuscolari della radiazione elettromagnetica, l'esperienza di diffrazione elettronica impone di attribuire agli enti particellari un comportamento ondulatorio.
Si giunge così a una visione unificata, seppur apparentemente contraddittoria, della realtà fisica microscopica.

Tale dualismo non implica che l'elettrone (o il fotone) sia \textit{simultaneamente} un'onda e una particella nel senso classico dei termini, bensì che esso manifesti l'una o l'altra natura a seconda del tipo di interazione a cui è sottoposto. Niels Bohr formalizzò questo concetto nel \textit{Principio di Complementarietà}: gli aspetti corpuscolare e ondulatorio sono descrizioni complementari di una stessa realtà fisica; un esperimento volto a misurare le proprietà dell'uno preclude la manifestazione delle proprietà dell'altro.

\section{Il principio di De Broglie}

Le relazioni fondamentali che connettono in maniera biunivoca le grandezze dinamiche corpuscolari (Energia $E$, quantità di moto ${p}$ e posizione $\bm{r}$) con le caratteristiche spettrali dell'onda associata (pulsazione $\omega$, vettore d'onda $\bm{\kappa}$, funzione d'onda $\psi$) sono riassunte nelle \textit{relazioni di Einstein-De Broglie}.
Utilizzando la costante di Planck ridotta $\hbar = h/2\pi$, tali relazioni assumono la forma simmetrica:
\begin{equation}
    \begin{cases}
        E \iff \hbar \omega \\
        \bm{p} \iff \hbar \bm{\kappa} \\
        \bm{r} \iff \psi(\bm{r},t) 
        
    \end{cases}
    \label{eq:einstein_debroglie}
\end{equation}
La prima equazione lega l'energia alla periodicità temporale dell'onda, mentre la seconda collega la quantità di moto alla periodicità spaziale. La terza, infine, comunica la necessità di analizzare la traiettoria spazio-temporale di una particella con una funzione d'onda. È attraverso queste relazioni che diviene possibile costruire una meccanica ondulatoria capace di descrivere l'evoluzione dei sistemi microscopici. \\

\noindent
Appaiono evidenti a questo punto almeno due domande fondamentali:
\begin{itemize}
\item da che equazione differenziale discende la funzione d'onda $\psi$?
\item un'onda è la propagazione della perturbazione di uno stato fisico: cosa si propaga? cosa è perturbato?
\end{itemize}
    
    
  \chapter{Fisica moderna: la meccanica quantistica}
L'introduzione della funzione d'onda come strumento descrittivo impone una revisione radicale dei principi della dinamica. Nella meccanica classica newtoniana, determinare lo stato di un sistema significa conoscere istante per istante la posizione $\bm{r}(t)$ e la quantità di moto $\bm{p}(t)$ di ogni particella; l'evoluzione di tali variabili è governata deterministicamente dalla legge $\dot{\bm{p}} = m\bm{a}$.
Stante il principio di De Broglie, occorre trovare un'equazione che non governi direttamente le variabili cinematiche, inutili in termini ondulatori, bensì l'evoluzione temporale della funzione d'onda $\Psi(\vec{r},t)$ stessa.
Tale legge fondamentale è l'\textit{Equazione di Schrödinger}. Essa rappresenta per la meccanica ondulatoria ciò che la seconda legge di Newton rappresenta per la meccanica classica. \\

\noindent
Si presente preliminarmente un'introduzione all'\textit{operatore Hamiltoniano}, fondamentale per la comprensione della meccanica quantistica.

\subsection{L'Hamiltoniana}
Al fine di formalizzare l'equazione d'onda, è necessario introdurre l'operatore che rappresenta l'energia totale del sistema: l'\textit{Hamiltoniana}.
In meccanica classica, la funzione di Hamilton $H(\bm{r}, \bm{p})$ per una particella di massa $m$ soggetta a un campo di forze conservativo è data dalla somma dell'energia cinetica $T(\bm{p})$ e dell'energia potenziale $V(\bm{r})$:
\begin{equation}
    H(\bm{r}, \bm{p}, t) = \frac{p^2}{2m} + V(\bm{r}, t).
\end{equation}

\section{L'equazione di Schrödinger}
\subsubsection{Considerazioni preliminari}
Partendo dai risultati di Davisson e Germer, Schrödinger ipotizza che gli elettroni, nel contesto atomico e subatomico, possano essere trattati come onde di materia.
Si discutono dunque le proprietà spaziali e temporali di un'onda piana \textit{progressiva} per risalire all'equazione generale\footnote{La funzione d'onda $\psi$ prende in argomento posizione e tempo, che qui si omettono per semplicità di trattazione.}, richiamando formalmente la struttura dell'equazione delle onde (o di d'Alembert):
\begin{equation}
    \nabla^2 \psi = \frac{1}{v_f^2}\frac{\partial^2 \psi}{\partial t^2}
\end{equation}
dove $v_f$ è la \textit{velocità di fase}, in generale diversa dalla velocità della luce $c$.
Per un elettrone in moto con velocità corpuscolare $v_p$, la quantità di moto è $p = m v_p$, mentre per l'onda associata vale $v_f = \lambda \nu$.
Dall'ipotesi di De Broglie si ottiene la relazione fondamentale:
\begin{equation}
    p = \frac{h}{\lambda} = \frac{h\nu}{v_f} \implies \frac{\nu}{v_f} = \frac{1}{\lambda}
    \label{eq:schrodi_1}
\end{equation}

\subsection{Analisi ondulatoria}
Assumendo per il generico elettrone una funzione d'onda progressiva:
\begin{equation}
    \psi(x,t) = A e^{i(kx-\omega t)}
\end{equation}
con $A$ ampiezza massima e $k$ numero d'onda, la derivata seconda temporale risulta:
\begin{equation}
    \pdv[2]{\psi}{t} = \pdv{t}\left( \pdv{\psi}{t} \right) = \pdv{t}\left( -i\omega \psi \right) = (-i\omega)(-i\omega)\psi = -\omega^2\psi
\end{equation}
Sostituendo tale risultato nella struttura di d'Alembert:
\begin{equation}
    \nabla^2 \psi = - \frac{1}{v_f^2} \omega^2\psi = -\left( \frac{\omega}{v_f} \right)^2\psi = -\left( \frac{2\pi\nu}{v_f} \right)^2\psi
\end{equation}
Utilizzando la relazione \eqref{eq:schrodi_1} e l'ipotesi di De Broglie ($p = h/\lambda = \hbar k$), si ottiene:
\begin{equation}
    \nabla^2 \psi = -\left(2\pi \frac{1}{\lambda} \right)^2\psi = -\left( \frac{p}{\hbar} \right)^2\psi
\end{equation}

\subsection{Analisi corpuscolare}
Nell'assunzione che l'onda descriva un corpo di massa $m$, si può riscrivere l'equazione in termini di energia cinetica. Moltiplicando per $-\hbar^2$:
\begin{equation}
    -\hbar^2 \nabla^2 \psi = p^2 \psi = (mv_p)^2 \psi
\end{equation}
Essendo l'energia cinetica classica definita come\footnote{Questa discussione assume validità solo fuori dal limite relativistico ($v_p \ll c$).} $T = \frac{p^2}{2m} = \frac{1}{2}mv_p^2$, dividendo per $2m$ si ottiene:
\begin{equation}
    -\frac{\hbar^2}{2m}\nabla^2 \psi = T\psi
\end{equation}
Dalla conservazione dell'energia meccanica, l'Hamiltoniana del sistema è $H = T + V$. Sostituendo l'operatore cinetico appena ricavato:
\begin{equation}
    -\frac{\hbar^2}{2m} \nabla^2 \psi + V\psi = H\psi
    \label{eq:schrodi_2}
\end{equation}

\subsubsection{Formulazione finale}
Si vuole ora esprimere l'energia totale $H$ in termini di operatori temporali. La derivata prima della funzione d'onda è:
\begin{equation}
    \pdv{\psi}{t} = -i\omega \psi
\end{equation}
Sfruttando la relazione quantistica di Planck-Einstein $E = \hbar \omega$ (dove l'energia $E$ corrisponde all'autovalore dell'Hamiltoniana $H$), si ha:
\begin{equation}
    \pdv{\psi}{t} = -i\frac{H}{\hbar}\psi \implies H \psi = i\hbar \pdv{\psi}{t}
\end{equation}
Sostituendo questa espressione per $H\psi$ nella \eqref{eq:schrodi_2}, si ottiene la formulazione completa dell'equazione di Schrödinger dipendente dal tempo:
\begin{equation}
    -\frac{\hbar^2}{2m} \nabla^2 \psi + V\psi = i\hbar \pdv{\psi}{t}
\end{equation}

\subsection{Commenti e interpretazione fisica}

\subsubsection{Il parallelo con la meccanica newtoniana}
È fondamentale osservare l'analogia formale tra l'equazione di Schrödinger e la seconda legge della dinamica ($ \bm{F} = m\bm{a} $).
In meccanica classica, la conoscenza delle forze agenti sulla particella (o equivalentemente del potenziale $V(\vec{r})$, da cui $\bm{F} = -\nabla V$) e delle condizioni iniziali permette di determinare univocamente la traiettoria $\bm{r}(t)$.
Analogamente, nell'equazione di Schrödinger, l'Hamiltoniana è interamente definita dalla conoscenza della massa $m$ della particella e del potenziale $V$ in cui essa è immersa.
Risolvendo questa equazione differenziale (nota la funzione d'onda iniziale $\psi_0$), si ottiene la funzione d'onda a ogni istante successivo.
La differenza sostanziale risiede nell'oggetto della determinazione: non più la posizione puntuale della particella lungo una traiettoria, bensì la sua distribuzione spazio temporale in termini ondulatori.

\subsubsection{La generalizzazione induttiva (l'ipotesi di Schrödinger)}
È doveroso sottolineare che il percorso logico seguito per "ricavare" l'equazione non costituisce una dimostrazione matematica rigorosa, ma piuttosto una giustificazione euristica.
Schrödinger compì un vero e proprio "atto di fede" scientifico (o, più formalmente, un'induzione empirica): partendo dalle evidenze sperimentali specifiche di Davisson e Germer sulla diffrazione elettronica e dalle intuizioni di De Broglie, postulò che la relazione trovata valesse \textit{universalmente}.
Egli assunse che \textit{qualsiasi} sistema microscopico (non solo elettroni liberi, ma anche particelle legate in un atomo o interagenti) fosse governato da tale equazione. La validità dell'equazione di Schrödinger non risiede nella sua derivazione, ma nello straordinario accordo delle sue soluzioni (come i livelli energetici dell'atomo di idrogeno) con i dati sperimentali.

\subsubsection{Nota sulle forze non conservative}
Nell'Hamiltoniana ${H} = {T} + {V}$ compare esclusivamente l'energia potenziale associata a forze conservative. Ci si potrebbe chiedere perché non vengano incluse forze dissipative (come l'attrito), comuni nella meccanica macroscopica.
La ragione è profonda: a livello microscopico fondamentale, \textit{tutte le forze in natura sono conservative}.
Le forze non conservative, come l'attrito o la resistenza viscosa, sono concetti emergenti che derivano dalla media statistica di un numero enorme di interazioni microscopiche conservative.
In un sistema quantistico elementare (un elettrone, un atomo), non esiste "dissipazione" nel senso classico; l'energia si conserva sempre se si considera il sistema completo. Quella che macroscopicamente appare come perdita di energia (calore) è, a livello microscopico, un trasferimento di energia verso altri gradi di libertà (es. vibrazioni reticolari) governati anch'essi da potenziali conservativi, ma troppo complessi per essere descritti singolarmente nell'Hamiltoniana di una singola particella.

\subsection{Il principio di corrispondenza}
Sorge spontaneo un quesito fondamentale: se la teoria di Schrödinger prevede una natura intrinsecamente probabilistica e ondulatoria, perché tali effetti non sono osservabili nel mondo macroscopico? \\

\noindent
In meccanica quantistica vige il \textit{principio di corrispondenza}, secondo il quale, nel limite di grandi numeri quantici o di grandi masse, le previsioni quantistiche devono convergere asintoticamente alla meccanica classica di Newton.
Si consideri ad esempio la lunghezza d'onda di De Broglie, $\lambda = h/mv$. Per un corpo macroscopico, la massa $m$ è talmente elevata (rispetto alla costante di Planck $h \approx 10^{-34}$ J$\cdot$s) che la lunghezza d'onda associata risulta infinitesima, vari ordini di grandezza inferiore a qualsiasi dimensione fisica misurabile.
Di conseguenza, i fenomeni tipicamente ondulatori come diffrazione e interferenza diventano trascurabili ("invisibili") e la funzione d'onda si localizza spazialmente, degenerando in quella che macroscopicamente appare come una traiettoria deterministica. Formalmente, questo raccordo è garantito dal \textit{Teorema di Ehrenfest}, il quale dimostra che i valori di aspettazione di posizione e momento evolvono nel tempo seguendo esattamente le leggi della dinamica newtoniana:
\begin{equation}
    m \frac{d\langle x \rangle}{dt} = \langle p \rangle, \quad \frac{d\langle p \rangle}{dt} = \langle -\nabla V \rangle.
\end{equation}
 

\section{Il significato fisico della funzione d'onda}
Provare a dedurre classicamente quale possa essere il significato fisico che risiede dietro la funzione d'onda associata a una particella massiva è sostanzialmente \textit{impossibile}: compare infatti nella formulazione di Schrödinger l'unità immaginaria $i$, da cui si deduce che $\psi \in \mathbb{C}$. Che senso ha? Come può una funzione che descrive un'entità \textit{reale} appartenere al mondo \textit{immaginario} dei numeri complessi?

\subsection{L'interpretazione probabilistica: l'ipotesi di Born}
\subsubsection{L'analogia ottica}
Si consideri l'esperimento di Young per l'interferenza di due onde luminose coerenti. Classicamente, l'intensità $I(P)$ registrata in un punto $P$ di uno schermo posto a grande distanza è proporzionale al quadrato del modulo del campo elettrico risultante $\bm{E}_{\text{tot}}$. In un'ottica corpuscolare, tale intensità è correlata al flusso di energia e, conseguentemente, al numero di fotoni $\mathcal{N}_f$ che incidono nell'intorno di $P$ nell'unità di tempo.
Normalizzando rispetto al numero totale di fotoni emessi $\mathcal{N}_{\text{tot}}$, si può stabilire una relazione di proporzionalità tra la distribuzione spaziale dei fotoni e il campo ondulatorio:
\begin{equation}
    \frac{\mathcal{N}_f(P)}{\mathcal{N}_{\text{tot}}} = \frac{I(P)}{I_{\text{tot}}} \propto |\bm{E}_{\text{tot}}(P)|^2.
\end{equation}
Tale rapporto rappresenta la frequenza relativa di arrivo dei fotoni in $P$, dunque la \textit{probabilità} che un fotone emesso colpisca $P$.

\subsubsection{L'ipotesi di Born}
Dall'esperimento di Davisson e Germer emerge che gli elettroni, analogamente ai fotoni, manifestano fenomeni di diffrazione e interferenza descrivibili tramite una funzione d'onda $\psi$.
Max Born estese l'analogia ottica proponendo che la natura ondulatoria della materia non risieda nella distribuzione di carica o massa, bensì nella probabilità.
Nel setup interferenziale, l'intensità rilevata in $P$ non corrisponde all'energia di un'onda classica, ma è proporzionale al numero di elettroni $\mathcal{N}_e(P)$ rilevati. La funzione d'onda risultante $\psi_{\text{tot}}$ (ottenuta per sovrapposizione lineare delle funzioni d'onda parziali) assume quindi il seguente significato fisico:
\begin{equation}
    \frac{\mathcal{N}_e(P)}{\mathcal{N}_e}  \propto |\psi_{\text{tot}}(P)|.
\end{equation}
In termini formali, il modulo quadro $|\psi|^2$ viene interpretato come \textit{densità di probabilità} di trovare la particella in un dato volume elementare dello spazio. Nell'analogia specifica discussa, $|\psi_{\text{tot}}(P)|^2$ rappresenta la probabilità per unità di volume che un elettrone impatti in $P$.

\subsection{L'interpretazione di Copenaghen}
L'ipotesi di Born, tutt'oggi ancora valida e non ancora smentita, permette di assumere che, nel suo moto, un generico elettrone non sia contraddistinto da una traiettoria specifica, bensì da un'\textit{onda di probabilità} che descrive la sua \textit{assumibile} posizione spazio-temporale. Se ne deriva che, se si confina la particella in un volume $\Omega$, la probabilità di trovarla in $\Omega$ deve essere unitaria:
\begin{equation}
    \int_\Omega |\psi(\bm{r}, t)|^2 d\Omega = 1.
\end{equation}
Sfruttando una notazione vastamente usata in meccanica quantistica per cui il differenziale di volume viene scritto come $d^3\bm{r}$ la struttura complessa della funzione d'onda e indicando con $\psi^*$ il complesso coniugato di $\psi$, l'integrale in questione può essere riscritto come:
\begin{equation}
    \int_\Omega \psi \psi^* d^3\bm{r} = 1.
\end{equation}
La funzione $\psi$ deve dunque avere la caratteristica di poter essere integrata in forma quadrata. Tutte le funzioni che soddisfano a questa proprietà sono dette \textit{a quadrato sommabile} su $V$ e appartengono allo spazio funzionale di Hilbert $\mathcal{L}^2(\Omega)$. Interessante caratteristica di $\mathcal{L}^2(\Omega)$ è il suo essere uno \textit{spazio vettoriale}: qualsiasi combinazione lineare di due funzioni d'onda definite in una regione $\Omega$ rimane quindi una funzione d'onda $\in \mathcal{L}^2(\Omega)$. Vale dunque il principio di sovrapposizione: essendo lo spazio funzionale in questione infinitamente generato, esistono infinite basi (infinite $n$-uple) sotto cui è possibile rappresentare uno stato generico di una particella. \\

\noindent
Le considerazioni interpretative appena discusse sono conosciute sotto il nome di \textit{interpretazione di Copenaghen}.

\section{Come estrarre informazioni dalla funzione d'onda}
Sino ad ora si è discussa esclusivamente la proprietà probabilistica della funzione d'onda associata a una particella di massa $m$. Se nella meccanica newtoniana la conoscenza della massa della particella e delle forze agenti su di essa era sufficiente per estrapolare \textit{qualsiasi} informazione fisica dal sistema, è necessario stabilire come fare lo stesso processo a partire dalle soluzioni dell'equazione di Schrödinger. 

\subsection{L'ipotesi di Heisenberg-Dirac}
\subsubsection{La posizione media}
Si vuole cercare in primis di trovare un modo per conoscere informazioni sulla posizione del corpuscolo. Condotte $N$ misure della sua posizione $\bm{r}$, il suo valore vero è dato dalla media aritmetica di tutte le $\bm{r}_i$:
\begin{equation}
    \langle \bm{r} \rangle = \sum_{i = 1}^N \frac{\bm{r}_i}{N}.
    \label{eq:sum_HD}
\end{equation}
Detta $\bm{r}_j$ una delle possibili $M$ posizioni in cui si può trovare la particella, la probabilità che $\langle \bm{r} \rangle = \bm{r}_j$ è frutto dell'ipotesi di Born, $ \mathcal{P}(\bm{r}_j) = |\psi(\bm{r}_j)|^2$. Il numero di volte in cui si registra $\bm{r}_j$ come posizione della particella deve essere data dal prodotto della probabilità per il numero di misure (fra loro indipendenti) $ N|\psi(\bm{r}_j)|^2$. Il fattore $\bm{r}_i$ della \eqref{eq:sum_HD} può allora essere riscritto come il prodotto del. numero di volte che si trova il corpuscolo nella generica posizione $\bm{r}_j$ per la posizione stessa: 
\begin{equation}
    \langle \bm{r} \rangle = \sum_{j = 1}^M \frac{\bm{r}_jN|\psi(\bm{r}_j)|^2}{N} = \sum_{j = 1}^M\bm{r}_j|\psi(\bm{r}_j)|^2.
\end{equation}
Nel limite in cui si considera uno spazio continuo $\Omega$, la sommatoria diventa un integrale:
\begin{equation}
    \langle \bm{r} \rangle = \int_\Omega \bm{r}|\psi(\bm{r})|^2 d^3\bm{r} = \int_\Omega  \psi(\bm{r})\bm{r}\psi^*(\bm{r}) d^3\bm{r}.
    \label{eq:r_q}
\end{equation}

\subsubsection{La velocità media}
Dalla definizione di velocità media, $\langle \dot{\bm{r}} \rangle$, si ottiene:
\begin{equation}
    \langle \dot{\bm{r}} \rangle =\frac{d}{dt} \int_\Omega \bm{r} \psi(\bm{r})\psi^*(\bm{r}) d^3\bm{r}.
\end{equation}
L'integrale agisce sullo spazio, la derivata sul tempo: i due operatori di conseguenza commutano. Si nota inoltre come la generica posizione $\bm{r}$ non può dipendere dal tempo, mentre la funzione d'onda ammette come parametro $t$:
\begin{equation}
    \langle {\bm{v}} \rangle = \int_\Omega \bm{r} \frac{\partial}{\partial t}[\psi(\bm{r})\psi^*(\bm{r})] d^3\bm{r}.
    \label{eq_v_HD}
\end{equation}
D'ora in avanti si analizza il problema in termini monodimensionali, sfruttando l'isotropia dello spazio. Dall'equazione di Schrödinger, la derivata della funzione d'onda rispetto al tempo:
\begin{equation}
    i\hbar \frac{\partial\psi}{\partial t} = -\frac{\hbar^2}{2m}\frac{\partial^2\psi}{\partial x^2} + V\psi.
\end{equation}
Moltiplicando ambo i membri per $\psi^*$ e trovando il coniugato dell'espressione derivante si ottiene.
\begin{equation}
\begin{dcases}
    \psi^*i\hbar \frac{\partial\psi}{\partial t} = -\frac{\hbar^2}{2m}\psi^*\frac{\partial^2\psi}{\partial x^2} + V\psi\psi^* \\
    -\psi i\hbar \frac{\partial\psi^*}{\partial t} = -\frac{\hbar^2}{2m}\psi\frac{\partial^2\psi^*}{\partial x^2} + V\psi^*\psi, \quad \text{coniugato},
    \end{dcases}
\end{equation}
dove il meno davanti alla seconda equazione deriva dal complesso coniugato dell'unità immaginaria. Sottraendo membro a membro le espressioni ottenute:
\begin{equation}
i\hbar\underbrace{\left[ \psi^* \frac{\partial\psi}{\partial t} + \psi \frac{\partial\psi^*}{\partial t}\right ]}_{\frac{\partial}{\partial t}(\psi\psi^*)} = -\frac{\hbar^2}{2m}\left[ \psi^*\frac{\partial^2\psi}{\partial x^2}- \psi\frac{\partial^2\psi^*}{\partial x^2} \right],
\end{equation}
da cui:
\begin{equation}
\frac{\partial}{\partial t}(\psi\psi^*) = -\frac{\hbar}{2mi}\left[ \psi^*\frac{\partial^2\psi}{\partial x^2}- \psi\frac{\partial^2\psi^*}{\partial x^2} \right].
\end{equation}
Sostituendo nella \eqref{eq_v_HD} monodimensionale (direzione $x$; si assume che gli estremi della regione sull'asse delle ascisse siano $a$ e $b$), poi:
\begin{equation}
\langle {{v}_x} \rangle = -\frac{\hbar}{2mi}\int_a^{b} x\left[ \psi^*\frac{\partial^2\psi}{\partial x^2}- \psi\frac{\partial^2\psi^*}{\partial x^2} \right] dx.
\end{equation}
Integrando per parti si ottiene:
\begin{equation}
\langle {{v}_x} \rangle = -\frac{\hbar}{2mi} \left[x\left( \psi^*\frac{\partial \psi}{\partial x}- \psi\frac{\partial\psi^*}{\partial x} \right)_{a}^{b} - \int_a^b \left( \psi^*\frac{\partial \psi}{\partial x}- \psi\frac{\partial\psi^*}{\partial x} \right)dx\right ].
\end{equation}
Serve una considerazione di carattere matematico adesso. $\psi$ deve soddisfare l'equazione di Schrödinger, dunque serve che nello spazio sia almeno derivabile $\psi \in \mathbb{C}^1(\Omega)$. Questa condizione implica continuità, anche sui bordi del dominio $a$ e $b$. Se la probabilità di trovare la particella fuori da $\Omega$ è nulla, serve che $\psi(\bar{\Omega}) = 0$. Da cui, $\psi(a) = \psi(b) = 0$. 
\begin{equation}
\langle {{v}_x} \rangle = \frac{\hbar}{2mi} \int_a^b \left( \psi^*\frac{\partial \psi}{\partial x}- \psi\frac{\partial\psi^*}{\partial x} \right)dx.
\end{equation}
Portando il termine $1/i$ all'interno e ricordando che $1/i = -i$, si ottiene:
\begin{equation}
    \langle v_x \rangle = \frac{1}{2m} \int_a^b \psi^* \left( -i\hbar \frac{\partial}{\partial x} \right) \psi \, dx + \frac{1}{2m} \int_a^b \psi \left( i\hbar \frac{\partial}{\partial x} \right) \psi^* \, dx.
\end{equation}
Il secondo integrale non è altro che il complesso coniugato del primo. Poiché la somma di un numero complesso $z$ e del suo coniugato $z^*$ è pari a $2\text{Re}(z)$ e dal momento che la velocità deve necessariamente essere un numero reale, si arriva alla conclusione per cui:
\begin{equation}
    \langle v_x \rangle = \frac{1}{m} \int_a^b \psi^* \left( -i\hbar \frac{\partial}{\partial x} \right) \psi.
\end{equation}
Moltiplicando ambo i membri per la massa $m$ e estendendo ad una generica trattazione multidimensionale, si ottiene:
\begin{equation}
    m\langle \bm{v} \rangle =\langle \bm{p} \rangle =  \int_\Omega [\psi^* ( -i\hbar \nabla) \cdot \psi] d^3\bm{r}.
    \label{eq_p_q}
\end{equation}

\subsubsection{Interpretazione quantistica}
Classicamente, note posizione $\bm{r}$ e quantità di moto $\bm{p}$ di una particella, esiste una funzione $F(\bm{r}, \bm{p})$ tramite cui si riesce a estrarre qualsiasi informazione di natura fisica dal sistema analizzato. In meccanica quantistica, invece, i valori medi delle grandezze sono univocamente determinati da \textit{operatori lineari} che vengono applicati (\textit{applicazione lineare}) alla funzione d'onda. Nell'analogia con la meccanica classica, deve quindi esistere un operatore $\hat{F}$ che ha in argomento gli operatori quantità di moto e posizione tramite cui è possibile estrarre informazioni fisiche dal sistema in esame. \\

\noindent 
Dalla simmetria delle \eqref{eq:r_q} e \eqref{eq_p_q}, si evince che gli operatori associati a posizione e quantità di moto sono:
\begin{equation}
    \hat{r} = \bm{r} \quad \hat{p} = -i\hbar \nabla.
    \end{equation}
Il valore medio di una generica grandezza ${A}$, in virtù di quanto detto in precedenza, deve poter essere scritto come: 
\begin{equation}
    \langle {A} \rangle = \int_\Omega \psi^* [ \hat{A}(\bm{r}, -i\hbar \nabla)] \psi \, d^3\bm{r}.
\end{equation}

\subsection{Operatori Hermitiani e Grandezze Osservabili}
In meccanica quantistica, sebbene la funzione d'onda $\psi$ sia intrinsecamente complessa, il risultato di una misura fisica, e di conseguenza il valore di aspettazione $\langle A \rangle$ di una qualsiasi grandezza $A$, deve essere necessariamente un numero reale.
Affinché il valore di aspettazione $\langle A \rangle$ sia reale per ogni possibile stato $\psi$, deve valere la condizione $\langle A \rangle = \langle A \rangle^*$. Esplicitando tale uguaglianza in termini integrali per un generico operatore $\hat{A}$, si richiede che:
\begin{equation}
    \int \psi^* (\hat{A} \psi) \, dx = \left( \int \psi^* (\hat{A} \psi) \, dx \right)^* = \int \psi (\hat{A} \psi)^* \, dx.
\end{equation}
Questa identità definisce la classe degli \textit{operatori hermitiani} (o autoaggiunti).
La condizione di hermitianità garantisce che la media delle misure restituisca un valore reale. Si può dimostrare, inoltre, che tale proprietà assicura non solo la realtà dei valori di aspettazione, ma anche quella dei singoli autovalori dell'operatore, i quali rappresentano i possibili esiti di una misura singola.
Pertanto, si formula il seguente postulato generale: \textit{ad ogni osservabile fisica è associato un operatore lineare ed hermitiano agente nello spazio delle funzioni d'onda.}


\subsection{Riscrittura in termini operatoriali dell'equazione di Schrödinger}
L'Hamiltoniana di una particella di massa $m$ soggetta a un campo di forze conservative $V(\bm{r}, t)$ è:
\begin{equation}
    H(\bm{r}, \bm{p}, t) = \frac{p^2}{2m} + V(\bm{r}, t).
\end{equation}
In termini operatoriali, quindi:
\begin{equation}
    \hat{H}(\hat{r}, \hat{p}, t) = \frac{(-i\hbar \nabla)(-i\hbar \nabla)}{2m} + V(\bm{r}, t) = -\frac{\hbar^2}{2m}\nabla^2 + V(\bm{r}, t).
\end{equation}
Si segnala che, descrivendo l'operatore Hamiltoniano l'energia del sistema, deve essere hermitiano. L'equazione di Schödinger allora si scrive:
\begin{equation}
    i\hbar \frac{\partial \psi}{\partial t} = -\frac{\hbar^2}{2m}\nabla^2\psi + V\psi = \left(-\frac{\hbar^2}{2m}\nabla^2 + V\right)\psi = \hat{H}\psi.
\end{equation}

\section{L'equazione di Schrödinger stazionaria}
Si vuole analizzare una generica particella di massa $m$ in una regione $\Omega$ sottoposta a un campo di forze conservative e tempo-invariante. Nell'equazione di Schrödinger,
\begin{equation}
    i\hbar \frac{\partial \psi}{\partial t} = \hat{H}\psi,
\end{equation}
l'operatore Hamiltoniano perde la dipendenza dal tempo, il primo membro dell'equazione la conserva. Nel tentare di semplificare la soluzione dell'equazione, si sfrutta il fatto che si può dimostrare come una soluzione particolare $\psi(\bm{r},t)$ di questa è data dal prodotto di un termine dipendente dal tempo $\mathcal{T}(t)$ e uno dalla posizione $\mathcal{R}(\bm{r})$:
\begin{equation}
   \psi(\bm{r},t) = \mathcal{T}(t)\cdot\mathcal{R}(\bm{r}).
\end{equation}
Sostituendo $\psi$ nell'equazione di partenza si ottiene:
\begin{equation}
   i\hbar \frac{\partial (\mathcal{T}(t)\mathcal{R}(\bm{r}))}{\partial t} =  i\hbar \mathcal{R}(\bm{r}) \frac{\partial \mathcal{T}(t)}{\partial t} = \hat{H}(\mathcal{T}(t)\mathcal{R}(\bm{r})) = \mathcal{T}(t)\hat{H}(\mathcal{R}(\bm{r})) 
\end{equation}
dove si sono sfruttate le proprietà della derivazione parziale rispetto alle costanti moltiplicative. L'equazione appena ottenuta si riscrive come
\begin{equation}
    \frac{i\hbar}{\mathcal{T}(t)}  \frac{\partial \mathcal{T}(t)}{\partial t} = \frac{\hat{H}(\mathcal{R}(\bm{r}))}{\mathcal{R}}.
\end{equation}
Risulta ora evidente come a membro sinistro si abbia una funzione esclusivamente dipendente dal tempo, a quello destro dalla posizione. Perché esista una soluzione all'equazione serve allora che entrambi i membri siano autonomamente uguali alla stessa costante $\beta$: 
\begin{equation}
\begin{dcases}
    \frac{i\hbar}{\mathcal{T}(t)}  \frac{\partial \mathcal{T}(t)}{\partial t} = \beta \\
    \frac{\hat{H}(\mathcal{R}(\bm{r}))}{\mathcal{R}} = \beta.
\end{dcases}
\end{equation}
La prima equazione è a variabili separabili e ammette come soluzione un'esponenziale
\begin{equation}
    \mathcal{T}(t) = Ae^{-\frac{i\beta}{\hbar}t},
\end{equation}
con $A$ una costante determinata dalle condizioni al contorno. La seconda, invece:
\begin{equation}
    {\hat{H}(\mathcal{R}(\bm{r}))} = \beta{\mathcal{R}(\bm{r})}
\end{equation}
non è altro che l'equazione agli autovalori dell'operatore $\hat{H}$. Espandendo l'operatore Hamiltoniano si ottiene \textit{l'equazione di Schrödinger in forma stazionaria}:
\begin{equation}
    \beta{\mathcal{R}(\bm{r})} = -\frac{\hbar^2}{2m}\nabla^2 \mathcal{R}(\bm{r}) + V(\bm{r}) \mathcal{R}(\bm{r}).
\end{equation}

\subsection{Significato fisico della costante $\bm{\beta}$}
Dal momento che una qualsiasi grandezza media può essere espressa come l'integrale sulla regione $\Omega$ del coniugato della funzione d'onda $\psi$ moltiplicato per l'applicazione del generico operatore legato all'osservabile $F$ su $\psi$, dalla riscrittura analizzata in precedenza si ottiene:
\begin{equation}
    \langle F \rangle = \int_\Omega \psi^*\hat{F}(\psi) d^3\bm{r} = \int_\Omega \mathcal{T}^*\mathcal{R}^*\hat{F}(\mathcal{T}\mathcal{R})d^3\bm{r} = \int_\Omega A^*e^{\frac{i\beta}{\hbar}t}\mathcal{R}^*\hat{F}(Ae^{-\frac{i\beta}{\hbar}t}\mathcal{R})d^3\bm{r}.
\end{equation}
Nel caso stazionario, l'operatore $\hat{F}$ non agisce sul tempo. Essendo poi $A$ una costante, risulta:
\begin{equation}
    \langle F \rangle = \int_\Omega A^*A\mathcal{R}^*\hat{F}(\mathcal{R})d^3\bm{r}.
\end{equation}
Nel caso dell'energia, l'equazione diventa:
\begin{equation}
    \langle H \rangle = \int_\Omega A^*A\mathcal{R}^*\hat{H}(\mathcal{R})d^3\bm{r} = \int_\Omega A^*A\mathcal{R}^*\beta\mathcal{R} d^3\bm{r} = \beta \int_\Omega (A\mathcal{R})^*(A\mathcal{R}) d^3\bm{r}
\end{equation}
per quanto discusso sull'equazione agli autovalori per l'operatore Hamiltoniano. Il modulo quadro della funzione d'onda, poi, lo si ottiene come:
\begin{equation}
    |\psi|^2 = \psi^*\psi = (\mathcal{T}^*\mathcal{R}^*)(\mathcal{T}\mathcal{R}) = (A^*e^{\frac{i\beta}{\hbar}t}\mathcal{R}^*)(Ae^{-\frac{i\beta}{\hbar}t}\mathcal{R}) = (A\mathcal{R})^*(A\mathcal{R}),
\end{equation}
da cui:
\begin{equation}
    \langle H \rangle = \beta \int_\Omega |\psi|^2 d^3\bm{r} = \beta.
\end{equation}
La costante che uguaglia la parte temporale e quella spaziale nell'equazione di Schrödinger non è altro, quindi, che l'energia media del sistema $E$. L'equazione di Schrödinger in forma stazionaria allora diventa:
\begin{equation}
    {\hat{H}(\mathcal{R}(\bm{r}))} = E{\mathcal{R}(\bm{r})}.
\end{equation}



\chapter{Conseguenze dell'equazione di Schrödinger}
L'equazione di Schrödinger impone la \textit{quantizzazione} di diverse grandezze fisiche associate alle particelle, con implicazioni colossali e assolutamente insensate dal punto di vista classico. In questo capitolo si discutono alcune delle principali. 

\section{Il caso della particella in una buca di potenziale}
Si immagini una particella libera\footnote{Per \textit{libera} si intende che non ci sono forze agenti sulla particella $\implies V(\bm{r}, t) = 0$.} di massa $m$ confinata in una scatola cubica di lato $a$. Come già fatto, si procede ad analizzare il caso monodimensionale estendendo in fine la trattazione sfruttando la proprietà di \textit{isotropia}. Serve però in primis trovare una tecnica risolutiva. L'equazione stazionaria trovata nel capitolo precedente permette di \textit{diagonalizzare} l'operatore Hamiltoniano\footnote{Sempre diagonalizzabile perchè hermitiano.}: si riesce così a determinare una base (tra l'altro ortonormale) dello spazio funzionale in cui vive $\psi$ formata dalle autofunzioni\footnote{A ciascuna delle quali è associata uno stato energetico preciso, l'autovalore.} di $\hat{H}$. Qualsiasi stato in cui si può trovare l'elettrone (in condizioni stazionarie e non) deve poter essere scritto come combinazione lineare della base trovata. La parte temporale, ancora, la si aggiunge "alla fine" sfruttando il metodo di separazione delle variabili discusso nel capitolo precedente\footnote{La componente temporale di soluzione $\mathcal{T}(t)$ è esprimibile come $\psi_n|_{\mathcal{T}(t)} = A\exp(-\frac{iE_n}{\hbar^2})$, dove $E_n$ è l'$n$-esimo autovalore associato all'$n$-esima autofunzione $\psi_n$.}.

\subsection{Formalizzazione matematica}
Sia allora $m$ la massa di una particella confinata nella regione dell'asse delle ascisse $[0,a]$. L'equazione di l'equazione di Schrödinger stazionaria si scrive:
\begin{equation}
   -\frac{\hbar^2}{2m}\frac{d^2\psi(x)}{d x^2} = E\psi(x) \implies \frac{d^2\psi(x)}{d x^2} = -\underbrace{\left( \frac{2mE}{\hbar^2} \right)}_{\alpha^2}\psi(x),
\end{equation}
dove si è definito $\alpha^2$ il fattore moltiplicativo di $\psi$ per semplificare i calcoli. L'equazione caratteristica associata alla differenziale analizzata ha come radici la coppia di complessi coniugati:
\begin{equation}
  \lambda^2 = -\alpha^2 \implies \lambda_{1,2} = \pm i\alpha,
\end{equation}
da cui la generica soluzione dell'ODE risulta:
\begin{equation}
\psi(x) = Ae^{i\alpha x} + Be^{-i\alpha x}
\end{equation}
con $A$ e $B$ costanti che dipendono dalle condizioni al contorno. Per determinarle, si impone innanzitutto la continuità di $\psi$ sul bordo della regione di confinamento:
\begin{equation}
\begin{aligned}
&\psi(0) = 0 \implies A+B = 0 \implies A = -B; \\
&\psi(a) = 0 \implies Ae^{i\alpha a} + Be^{-i\alpha a} = A(e^{i\alpha a} - e^{-i\alpha a}) = 0 \implies 2iA\sin(\alpha a) = 0.
\end{aligned}
\end{equation}
La soluzione $A = 0$ è inaccettabile perché condurrebbe a $\psi(x) = 0$ ovvero al fatto che la particella non si trovi nella scatola. Cercando gli zeri del seno allora:
\begin{equation}
\alpha a = n\pi \implies \frac{\sqrt{2mE}}{\hbar} = \frac{n\pi}{a}, \quad n\in \mathbb{Z}.
\label{eq:quant_a}
\end{equation}
Isolando infine l'energia, si ottiene:
\begin{equation}
E_n = \frac{n^2 \pi^2 \hbar^2}{2m a^2} = \frac{n^2h^2}{8ma^2}.
\end{equation}
Si è arrivati ad una conclusione sconcertante (dal punto di vista classico): l'energia è infatti quantizzata, può esistere solo \textit{discretamente}, non più nel continuo. 

\subsubsection{Calcolo del parametro $\bm{A}$}
Serve ancora una condizione al contorno da imporre. Seguendo le orme di Born, si impone la normalizzazione della funzione sull'intervallo di definizione $[0,a]$ :
\begin{equation}
\int_0^a |\psi(x)|^2 dx = \int_0^a |A(e^{i\alpha x} - e^{-i\alpha x})|^2 dx = 4A^2\int_0^a |\sin(\alpha x)|^2 dx = 1,
\end{equation}
da cui si ottiene:
\begin{equation}
4A^2\int_0^a \frac{1-\cos(2\alpha x)}{2} dx = 1 \implies 2A = \pm \left[ a-\frac{\sin(2\alpha a)}{2\alpha}\right]^{-\frac{1}{2}}.
\end{equation}
Sfruttando la condizione di quantizzazione precedentemente trovata \eqref{eq:quant_a} e notando che il seno risulta sempre quindi nullo, l'espressione si semplifica a:
\begin{equation}
 A = \frac{1}{\sqrt{2a}},
\end{equation}
dove si è scelto per semplicità il risultato positivo (la normalizzazione è in ogni caso verificata perché dipende dal quadrato). 

\subsubsection{La soluzione del problema}
La generica autofunzione si scrive quindi come:
\begin{equation}
 \psi_n(x) = \frac{1}{\sqrt{2a}}(e^{i\alpha x} - e^{-i\alpha x}) = i\sqrt{\frac{2}{{a}}} \sin(\alpha x) = i\sqrt{\frac{2}{{a}}} \sin\left(\frac{\sqrt{2mE_n}}{\hbar^2}x \right).
\end{equation}
Sfruttando il fatto che della costante moltiplicativa del seno interessa solo in termini di normalizzazione e che il modulo dell'unità immaginaria è 1, si può riscrivere l'ultima relazione in forma reale:
\begin{equation}
 \psi_n(x) = \sqrt{\frac{2}{{a}}} \sin(\alpha x) = \sqrt{\frac{2}{{a}}} \sin\left(\frac{\sqrt{2mE_n}}{\hbar^2}x \right).
\end{equation}
Un generico stato stazionario $\Psi(x)$ si può scrivere come combinazione lineare di queste autofunzioni di base:
\begin{equation}
 \Psi(x) = \sum_{n \geq 1} c_n\psi_n(x),
\end{equation}
aggiungendo il contributo temporale, infine:
\begin{equation}
 \Psi(x,t) = \sum_{n \geq 1} c_n\psi_n(x)e^{-\frac{iE_n}{\hbar^2}t}.
\end{equation}
Sulle altre direzioni il risultato è simmetrico.

\subsection{Influenza della quantizzazione}
Una particella di gas in una stanza è a tutti gli effetti un'ottima approssimazione del caso della buca di potenziale. Avendo questa 3 gradi di libertà, per il teorema di equipartizione dell'energia deve valere:
\begin{equation}
 E = \frac{3}{2}mv^2 = \frac{n^2h^2}{8ma^2},
\end{equation}
da cui si deduce che anche la velocità deve essere una grandezza quantizzata, infatti:
\begin{equation}
v = \frac{nh}{2ma}\sqrt{3}.
\end{equation}
Dall'ultima relazione è evidente che l'effetto a "gradino" imposto dalla quantizzazione è inversamente proporzionale alle dimensioni della regione in cui si trova, infatti detto $\Delta v = v_{n+1} - v_{n}$, con $n \geq 1$ risulta:
\begin{equation}
\Delta v = \frac{h\sqrt{3}}{2ma}.
\end{equation}
Perché l'effetto sia osservabile, serve che $a \approx 0$, essendo $h$ e la massa di una generica particella molto piccole. 



\section{Effetto tunnel}
Si vuole ora cercare di analizzare un caso classicamente \textit{banale}: una particella che affronta una barriera di potenziale di "altezza" $V_0$. Si immagini una massa puntiforme che viene sparata con energia cinetica $K$ su di un piano inclinato di altezza $H$. Si possono verificare due sole situazioni distinte, determinate rigorosamente dal pricipio di conservazione dell'energia meccanica:

\begin{itemize}
    \item \textit{caso $E < V_0$ (riflessione totale):} se l'energia totale della particella è inferiore all'energia potenziale della barriera, classicamente essa arresta il proprio moto nel punto di inversione (dove $K=0$) e viene riflessa indietro. La probabilità di trovare la particella oltre l'ostacolo è identicamente nulla; la regione al di là della barriera è una \textit{zona classicamente proibita};
    \item \textit{caso $E > V_0$ (trasmissione garantita):} se l'energia è superiore al potenziale, la particella supera l'ostacolo con certezza, subendo solo un rallentamento durante l'attraversamento della regione a potenziale non nullo.
\end{itemize}
L'analisi quantistica, tuttavia, sovverte questa logica binaria e deterministica. In virtù della natura ondulatoria della materia, descritta dalla funzione d'onda $\psi(x)$, non è possibile imporre che l'onda si arresti istantaneamente all'interfaccia della barriera. Le condizioni di continuità per $\psi$ e impongono che l'onda penetri nella regione proibita, attenuandosi ma non annullandosi.
Ne consegue un fenomeno controintuitivo: anche nel regime $E < V_0$, esiste una probabilità finita e non nulla che la particella \textit{attraversi} la barriera, "comparendo" nella regione successiva senza aver posseduto l'energia cinetica necessaria per sormontarla classicamente. Tale fenomeno, privo di analogo nella meccanica classica dei punti materiali, è noto come \textit{Effetto Tunnel}.

\subsection{Formalizzazione matematica}
Sia $V(x)$ il campo scalare potenziale in cui è immersa una particella che "vive" in $\mathbb{R}$ di massa $m$ e sia:
\begin{equation}
V(x) = 
\begin{cases}
0 \quad x < 0 \\
V_0 \quad x \geq 0.
\end{cases}
\end{equation}
Il corpuscolo, con energia $E$, si muove nella regione di potenziale nullo e "impatta" con la barriera in $x = 0$. Si procede a scrivere l'equazione di Schrödinger stazionaria\footnote{La soluzione completa dell'equazione si può trovare seguendo lo stesso procedimento sfruttato nel caso della buca di potenziale.} nelle due regioni di definizione del potenziale a gradino:
\begin{equation}
\begin{aligned}
&-\frac{\hbar^2}{2m}\frac{d^2 \psi(x)}{d x^2} = E\psi(x) \quad x < 0 \\
&-\frac{\hbar^2}{2m}\frac{d^2 \psi(x)}{d x^2} + V_0\psi(x) = E\psi(x) \quad x \geq 0,
\end{aligned}
\end{equation}
che si riscrivono come:
\begin{equation}
\begin{aligned}
&\frac{d^2 \psi(x)}{d x^2} = -\overbrace{\frac{2mE}{\hbar^2}}^{\alpha^2}\psi(x) = -\alpha^2\psi(x) \quad x < 0 \\
&\frac{d^2 \psi(x)}{d x^2}  = \underbrace{\frac{2m(V_0 - E)}{\hbar^2}}_{\gamma^2}\psi(x) = \gamma^2\psi(x) \quad x \geq 0,
\end{aligned}
\end{equation}
dove si sono definiti i parametri $\alpha$ e $\gamma$ per semplificare la trattazione. Nel caso in cui $V_0 > E$ (interessante perché, come discusso, la risposta dovrebbe essere ovvia... almeno classicamente), le due equazioni differenziali ammettono soluzioni esponenziali del tipo:
\begin{equation}
\psi(x) =
\begin{dcases}
 Ae^{i\alpha x} + Be^{-i\alpha x} \quad x < 0 \\
  Ce^{\gamma x} + De^{-\gamma x} \quad x \geq 0,
\end{dcases}
\end{equation}
Nella regione di ascisse negative si ha la combinazione lineare di due esponenziali, positivo e negativo rispettivamente. Rappresentano fisicamente un'onda progessiva e regressiva: a seguito dell'urto con la barriera di potenziale, si generano infatti due onde distinte che si muovono in verso opposto. Perché $\psi$ possa essere normalizzabile, poi, non deve divergere a $\pm \infty$, quindi serve che $C = 0$. Imponendo poi continuità e derivabilità\footnote{La funzione d'onda, essendo soluzione di un'equazione differenziale, è intrinsecamente almeno di classe $\mathbb{C}^2$ rispetto alle componenti spaziali e $\mathbb{C}^1$ nel tempo. L'equazione di Schrödinger è infatti del secondo ordine rispetto alla posizione del primo rispetto al tempo.} nell'origine (per $x = 0$, in occorrenza della barriera), si ottiene il seguente sistema:
\begin{equation}
\begin{dcases}
A + B = D  \\
i\alpha( A - B) = -\gamma D \implies A-B = \frac{i\gamma}{\alpha}D.
\end{dcases}
\end{equation}
Il sistema è evidentemente sottodeterminato: usando $A$ come variabile indipendente, si ottengono come soluzioni:
\begin{equation}
\begin{dcases}
B = \frac{\alpha - i\gamma}{\alpha + i\gamma}A \\
D = \frac{2\alpha}{\alpha + i\gamma}A.
\end{dcases}
\end{equation}
Si segnala che l'ultima incognita $A$ non si può trovare tramite il processo di normalizzazione: nella regione di ascisse negative, infatti, si ha la somma di onde piane che non sono normalizzabili su un dominio infinito (l'integrale diverge). Servirebbe introdurre concetti di riflessione che esulano dalla trattazione qui proposta. La soluzione finale, risulta quindi: 
\begin{equation}
\psi(x) =
\begin{dcases}
 A\left(e^{i\alpha x} + \frac{\alpha - i\gamma}{\alpha + i\gamma}e^{-i\alpha x}\right) \quad x < 0 \\
   \frac{2\alpha A}{\alpha + i\gamma}e^{-\gamma x} \quad x \geq 0,
\end{dcases}
\end{equation}

\subsection{Significato fisico del termine $\bm{D}$}
Dall'ipotesi di Born, il modulo quadro della funzione d'onda valutata in un punto $x_0$ dello spazio in cui è definita restituisce la densità di probabilità di trovare la particella descritta dalla funzione in $x_0$. Nel caso in esame, limitandosi alla regione che segue la barriera $x \geq 0$:
\begin{equation}
    |\psi(x)|^2 = |D e^{-\gamma x}|^2 = |A|^2 \frac{4\alpha^2}{\alpha^2 + \gamma^2} e^{-2\gamma x}.
\end{equation}
Sostituendo le espressioni esplicite dei parametri $\alpha^2 = 2mE/\hbar^2$ e $\gamma^2 = 2m(V_0-E)/\hbar^2$, si osserva che il denominatore si riduce semplicemente a $2mV_0/\hbar^2$. L'espressione finale risulta dunque:
\begin{equation}
    |\psi(x)|^2 = |A|^2 \frac{4E}{V_0} e^{-2\gamma x}.
\end{equation}
Esiste dunque una probabilità non nulla che decade esponenzialmente a zero di trovare la particella al di là della barriera.

\subsubsection{Passaggio al macroscopico}
Il principio di corrispondenza, in questo caso, si ritrova nella caratterizzazione della massa $m$. Il termine esponenziale, infatti, dipende direttamente dalla massa rapportata al quadrato della costante di Planck ridotta (valore molto piccolo). Per masse caratteristiche di realtà macroscopiche, l'esponenziale va a zero praticamente istantaneamente, rendendo la densità di probabilità analizzata pressoché nulla (e rendendo essenzialmente impossibile oltrepassare la barriera).

\begin{approfondimento}[colback= white, colframe=green!35!black, no shadow, title={Applicazione astrofisica: il modello di Gamow-Bethe}]
L'effetto tunnel rappresenta la chiave di volta per la comprensione dei meccanismi di nucleosintesi stellare, risolvendo un paradosso altrimenti inspiegabile dalla fisica classica.
All'inizio del XX secolo, le stime sulla temperatura del nucleo solare ($T \approx 1.5 \cdot 10^7$ K) ponevano un problema energetico insormontabile: l'energia cinetica media dei protoni, stimata tramite la statistica di Boltzmann ($E \approx k_B T \approx 1 \text{ keV}$), risulta essere di tre ordini di grandezza inferiore rispetto all'altezza della \textit{barriera coulombiana} ($V_c \approx 1 \text{ MeV}$) che impedisce la fusione di due nuclei di idrogeno.
Secondo la meccanica classica, la probabilità che due protoni si avvicinino a una distanza sufficiente ($\sim 10^{-15}$ m) per innescare l'interazione nucleare forte sarebbe identicamente nulla; il Sole, di conseguenza, non potrebbe brillare. \\

\noindent
Nel 1928, George Gamow applicò il formalismo della meccanica ondulatoria al nucleo atomico, intuendo che i protoni, comportandosi come pacchetti d'onda, possiedono una probabilità finita di penetrare la barriera di potenziale repulsivo $V(r) \propto 1/r$ attraverso l'effetto tunnel, anche possedendo energie $E \ll V_c$.
Successivamente, Hans Bethe (1939) integrò questa intuizione nel calcolo dei tassi di reazione per il ciclo protone-protone (catena p-p). Egli dimostrò che la fusione termonucleare avviene prevalentemente in una stretta finestra energetica (detta \textit{finestra di Gamow}), risultante dal prodotto di due funzioni in competizione: la distribuzione di Maxwell-Boltzmann (che decresce esponenzialmente con l'energia) e la probabilità di tunneling (che cresce esponenzialmente con l'energia).
È dunque l'effetto tunnel a garantire la stabilità delle stelle: se la barriera coulombiana fosse sormontabile classicamente, le reazioni divergerebbero istantaneamente; la bassa probabilità di tunneling agisce invece come un "termostato" che regola il tasso di fusione su tempi scala di miliardi di anni.


\subsubsection{Valutazione semiclassica: il fattore di Gamow}
Per quantificare la probabilità di fusione, è necessario superare l'approssimazione della barriera a gradino e considerare la forma reale del potenziale repulsivo elettrostatico tra due nuclei di carica $Z_1e$ e $Z_2e$:
\begin{equation}
    V(r) = \frac{1}{4\pi\epsilon_0} \frac{Z_1 Z_2 e^2}{r}.
\end{equation}
La particella incidente, con energia $E$, incontrerebbe classicamente un punto di inversione del moto (turning point) alla distanza $r_c$ tale per cui $V(r_c) = E$:
\begin{equation}
    r_c = \frac{Z_1 Z_2 e^2}{4\pi\epsilon_0 E}.
\end{equation}
Poiché il potenziale $V(r)$ varia spazialmente, l'attenuazione della funzione d'onda nella regione classicamente proibita non è descritta da un esponente costante, bensì dall'integrale dell'impulso immaginario lungo il percorso di tunneling (in accordo con l'approssimazione semiclassica WKB). La probabilità di penetrazione (trasmissione) $T$ assume la forma:
\begin{equation}
    T(E) \approx \exp\left( -2 \int_{R}^{r_c} \sqrt{\frac{2\mu}{\hbar^2} [V(r) - E]} \, dr \right),
\end{equation}
dove $\mu$ rappresenta la massa ridotta del sistema a due corpi e $R$ il raggio d'interazione nucleare (distanza alla quale scattano le forze attrattive forti). Assumendo $R \ll r_c$ (approssimazione puntiforme per il nucleo rispetto alla barriera coulombiana), l'integrale può essere risolto analiticamente, portando al risultato noto come \textit{Fattore di Gamow}:
\begin{equation}
    T(E) \approx e^{-2G(E)} \quad \text{con} \quad G(E) = \frac{\pi Z_1 Z_2 e^2}{\hbar} \sqrt{\frac{\mu}{2E}}.
\end{equation}
È consuetudine riscrivere l'esponente introducendo l'\textit{energia di Gamow} $E_G$, una costante caratteristica del sistema reagente:
\begin{equation}
    T(E) \propto \exp\left( -\sqrt{\frac{E_G}{E}} \right), \quad \text{dove } E_G = 2\mu c^2 (\pi \alpha Z_1 Z_2)^2.
\end{equation}
Questa espressione evidenzia la drammatica dipendenza della probabilità di tunnel dall'energia cinetica $E$: a temperature stellari tipiche, il fattore esponenziale è estremamente piccolo (dell'ordine di $10^{-20}$ per il Sole), giustificando la lentezza del processo di combustione dell'idrogeno e la longevità della sequenza principale stellare.
\end{approfondimento}

\section{Il rotore rigido}
Si analizza a questo punto il caso di una particella di massa $m$ in moto circolare attorno ad un punto fisso ad una distanza dal centro di rotazione costante e pari a $r$. Classicamente, per descrivere il sistema si farebbe uso del \textit{momento angolare} $L$ associato al corpuscolo in moto; l'analogo quantistico deve essere un operatore hermitiano e lineare: l'\textit{operatore momento angolare} $\hat{L}$. Dalla definizione, 
\begin{equation}
  \bm{L} = \bm{r} \times \bm{p} \implies \hat{L} = \hat{r} \times \hat{p}.
\end{equation}
Utilizzando una funzione ausiliaria $\psi$, si continua passando in coordinate sferiche\footnote{Scelta naturale data la simmetria sferica del problema.}:
\begin{equation}
\begin{aligned}
  \hat{L}\psi &= [\hat{r} \times \hat{p}](\psi) = \bm{r} \times[ -i\hbar \nabla (\psi)] = -i\hbar \bm{r} \times \left( \bm{u}_r \pdv{\psi}{r} + \bm{u}_{\theta} \frac{1}{r} \pdv{\psi}{\theta} + \bm{u}_{\phi} \frac{1}{r\sin\theta} \pdv{\psi}{\phi} \right) \\
  &= -i\hbar  \left[r(\bm{u}_r \times \bm{u}_r) \pdv{\psi}{r} + r(\bm{u}_r \times \bm{u}_{\theta}) \frac{1}{r} \pdv{\psi}{\theta} + r(\bm{u}_r \times \bm{u}_{\phi}) \frac{1}{r\sin\theta} \pdv{\psi}{\phi} \right]  \\
  &= -i\hbar  \left(  \bm{u}_{\phi} \pdv{\psi}{\theta} -  \bm{u}_{\theta} \frac{1}{\sin\theta} \pdv{\psi}{\phi} \right)
  \end{aligned}
\end{equation}
da cui l'operatore momento angolare in coordinate sferiche risulta:
\begin{equation}
  \hat{L} = -i\hbar  \left(  \bm{u}_{\phi} \pdv{}{\theta} -  \bm{u}_{\theta} \frac{1}{\sin\theta} \pdv{}{\phi} \right).
\end{equation}
L'energia cinetica rotazionale di un corpo è descrivibile dal quadrato del suo momento angolare $L^2$. L'operatore associato $\hat{L}^2$ si trova applicando $\hat{L}$ a se stesso:
\begin{equation}
\begin{aligned}
    \hat{L}^2 (\psi) &= -i\hbar \left( \bm{u}_{\phi} \pdv{}{\theta} - \frac{\bm{u}_{\theta}}{\sin\theta} \pdv{}{\phi} \right) \cdot \left[ -i\hbar \left( \bm{u}_{\phi} \pdv{\psi}{\theta} - \frac{\bm{u}_{\theta}}{\sin\theta} \pdv{\psi}{\phi} \right) \right] \\
    &= -\hbar^2 \left[ (\bm{u}_{\phi} \cdot \bm{u}_{\phi})\pdv[2]{\psi}{\theta} - \bm{u}_{\phi} \pdv{}{\theta}\left(\frac{\bm{u}_{\theta}}{\sin\theta} \pdv{\psi}{\phi}\right) - \frac{\bm{u}_{\theta}}{\sin\theta} \pdv{\psi}{\phi} \left(\bm{u}_{\phi} \pdv{\psi}{\theta}\right) + \frac{\bm{u}_{\theta} \cdot \bm{u}_{\theta}}{\sin^2\theta}\pdv[2]{\psi}{\phi}\right] \\
    &= -\hbar^2 \left[\pdv[2]{\psi}{\theta} - \bm{u}_{\phi} \pdv{}{\theta}\left(\frac{\bm{u}_{\theta}}{\sin\theta} \pdv{\psi}{\phi}\right) - \frac{\bm{u}_{\theta}}{\sin\theta} \pdv{\psi}{\phi} \left(\bm{u}_{\phi} \pdv{\psi}{\theta}\right) + \frac{1}{\sin^2\theta}\pdv[2]{\psi}{\phi}\right] \\
    &= -\hbar^2\left[ \frac{1}{\sin\theta}\pdv{}{\theta}\left( \sin\theta\pdv{\psi}{\theta} \right)+ \frac{1}{\sin^2\theta}\pdv[2]{\psi}{\phi} \right],
\end{aligned}
\end{equation}
da cui l'operatore $\hat{L}^2$ risulta:
\begin{equation}
\hat{L}^2 = -\hbar^2 \left[ \frac{1}{\sin\theta}\pdv{}{\theta}\left( \sin\theta\pdv{}{\theta} \right)+ \frac{1}{\sin^2\theta}\pdv[2]{}{\phi} \right].
\end{equation}
Essendo l'energia cinetica rotazionale classica definita come $E = L^2/2I$, con $I$ il \textit{momento d'inerzia} del sistema, l'operatore hamiltoniano si scrive di conseguenza:
\begin{equation}
\hat{H} = \frac{\hat{L}^2}{2I} = -\frac{\hbar^2}{2I} \left( \frac{1}{\sin\theta}\pdv{}{\theta}\left( \sin\theta\pdv{}{\theta} \right)+ \frac{1}{\sin^2\theta}\pdv[2]{}{\phi} \right)
\end{equation}
e l'equazione di Schrödinger stazionaria, quindi:
\begin{equation}
\hat{H} \psi = E\psi,
\end{equation}
con la generica autofunzione dipendente dai parametri sferici $(r,\theta,\phi)$. Si assume che la soluzione dell'equazione possa essere essere decomposta come il prodotto di una parte radiale e una parte angolare $\psi(r,\theta,\phi) = \mathcal{R}({r})\mathcal{Y}(\theta,\phi)$:
\begin{equation}
\hat{H} \mathcal{R}({r})\mathcal{Y}(\theta,\phi) = E\mathcal{R}({r})\mathcal{Y}(\theta,\phi).
\end{equation}
Dal momento che l'hamiltoniano non agisce rispetto alle variabili radiali, $\mathcal{R}({r})$ si semplifica ad ambo i membri ottenendo: 
\begin{equation}
\begin{aligned}
\hat{H}\mathcal{Y}(\theta,\phi) &= E\mathcal{Y}(\theta,\phi) \\
 -\underbrace{\frac{2IE}{\hbar^2}}_{l(l+1)}\mathcal{Y} &= \frac{1}{\sin\theta}\pdv{}{\theta}\left( \sin\theta\pdv{\mathcal{Y}}{\theta} \right)+ \frac{1}{\sin^2\theta}\pdv[2]{\mathcal{Y}}{\phi} , 
\end{aligned}
\end{equation}
dove si è uguagliato il generico autovalore dell'equazione al termine $l(l+1)$ per motivazioni che risulteranno più chiare successivamente. La prima evidente conseguenza di questa scelta è la quantizzazione dell'energia:
\begin{equation}
E = \frac{\hbar^2}{2I}l(l+1), \quad l \in \mathbb{N}
\end{equation}
con il livello energetico più basso permesso caratterizzato da $l = 0$. Continuando con l'equazione, si suppone ancora che la soluzione possa essere scritta separando le variabili come prodotto di una parte dipendente da $\theta$ e una da $\phi$: $\mathcal{Y}(\theta, \phi) = \Theta(\theta)\Phi(\phi)$. Da cui:
\begin{equation}
\begin{aligned}
-l(l+1)\Theta(\theta)\Phi(\phi) &= \frac{\Phi(\phi)}{\sin\theta}\pdv{}{\theta}\left( \sin\theta\pdv{\Theta(\theta)}{\theta} \right)+ \frac{\Theta(\theta)}{\sin^2\theta}\pdv[2]{{\Phi(\phi)}}{\phi} \\
-\frac{1}{\Phi(\phi)}\pdv[2]{{\Phi(\phi)}}{\phi} &= \frac{1}{\Theta(\theta)}\sin\theta\pdv{}{\theta}\left( \sin\theta\pdv{\Theta(\theta)}{\theta} \right) + l(l+1)\sin^2\theta.
\end{aligned}
\end{equation}
Dalla separazione delle variabili si evince come, perché l'uguaglianza di due funzioni di parametri diversi possa essere vera, i due membri dell'equazione devono essere contemporaneamente uguali alla stessa costante $m$:
\begin{equation}
\begin{dcases}
-\frac{1}{\Phi(\phi)}\pdv[2]{{\Phi(\phi)}}{\phi} = m \\
 \frac{1}{\Theta(\theta)}\sin\theta\pdv{}{\theta}\left( \sin\theta\pdv{\Theta(\theta)}{\theta} \right) + l(l+1)\sin^2\theta = m.
\end{dcases}
\end{equation}

\subsection{L'equazione nella variabile $\bm{\phi}$}
La seconda equazione in $\phi$ ammette come soluzione:
\begin{equation}
\Phi(\phi) = Ae^{-im\phi} + Be^{im\phi}.
\end{equation}
Si nota però che una $\Phi(\phi)$ così scritta non è autofunzione della componente $z$ del momento angolare $L_z$ ma lo diventa nel caso in cui si abbia $A  = 0$ o $B = 0$. Si sceglie per semplicità il caso $B = 0$. Stante la simmetria sferica, la funzione deve essere uguale a se stessa ogni rivoluzione di angolo giro, $\Phi(\phi) = \Phi(\phi + 2\pi)$, da cui:
\begin{equation}
\begin{aligned}
\Phi(\phi) = Ae^{-im\phi} &= \Phi(\phi + 2\pi) = Ae^{-im(\phi + 2\pi)} \\
Ae^{-im\phi} &= Ae^{-im\phi}e^{-im2\pi} \\
1 &= e^{-im2\pi} = \cos(2\pi m) - i\sin(2\pi m) \\
\implies 2\pi m &= 2 \pi n, \quad n \in \mathbb{Z} \\
\implies m &= n
\end{aligned}
\end{equation}
da cui si ottiene che la costante $m$ deve necessariamente essere un numero intero.

\subsection{L'equazione nella variabile $\bm{\theta}$}
La prima equazione del sistema riportato rappresenta un'equazione differenziale particolare già studiata da Legendre, che ammette come soluzione una \textit{funzione associata di Legendre}:
\begin{equation}
\Theta(\theta) = N_{ml} P_l^m(\cos\theta)
\end{equation}
con $N_{ml}$ una costante dipendente dai parametri $m$ e $l$ e 
\begin{equation}
P_l^m(x) = (1+x^2)^{\frac{|m|}{l}} \frac{d^{|m|}}{dx^{|m|}}\left[P_l(x)\right]
\end{equation}
una funzione dove $P_l(x)$ è il \textit{polinomio di Legendre} di ordine $l$. Dalla definizione, ancora
\begin{equation}
P_l(x) = \frac{1}{2^l l!} \frac{d^{l}}{dx^{l}}(x^2-1)^l
\end{equation}
si ottiene che il grado di $P_l(x)$ deve essere $l$. Osservando che l'operazione di derivazione di ordine $m$ abbassa il grado del polinomio di $m$ ordini, si ottiene che la condizione necessaria perché $P_l^m(x)$ sia non nullo è:
\begin{equation}
|m| \leq l.
\end{equation}
I valori accettabili per la costante $m$, detta \textit{numero quantico magnetico}, sono:
\begin{equation}
-l \leq m \leq l.
\end{equation}

\subsection{La soluzione generale}
Componendo la soluzione angolare generica per il rotore rigido si ottiene:
\begin{equation}
\mathcal{Y}(\theta,\phi) = N_{ml} P_l^m(\cos\theta) Ae^{-im\phi},
\end{equation}
con $\mathcal{Y}$ rappresentante le cosiddette \textit{armoniche sferiche} e autofunzione dell'operatore hamiltoniano. 


\section{L'atomo di idrogeno}


\chapter{La Teoria della Misura}
La formulazione della meccanica quantistica sviluppata nei capitoli precedenti descrive l'evoluzione temporale di un sistema isolato attraverso l'equazione di Schrödinger. Tale evoluzione è caratterizzata dall'essere \textit{deterministica}, \textit{reversibile} e, in termini matematici, \textit{unitaria}: lo stato $\Psi$ preserva la sua norma e l'informazione in esso contenuta fluisce in modo continuo nel tempo. \\

\noindent
Tuttavia, l'esperienza empirica e la verifica sperimentale della teoria richiedono necessariamente l'interazione tra il sistema quantistico microscopico e un apparato di misura macroscopico (o osservatore). È in questo frangente che emerge una singolarità logica e formale: l'atto della misura non può essere descritto dalla semplice applicazione dell'operatore di evoluzione temporale $\hat{U}(t)$. Al contrario, esso introduce un elemento di \textit{irreversibilità} e di \textit{indeterminismo intrinseco}.
La Teoria della Misura si occupa di formalizzare questo processo discontinuo, postulando che l'osservazione di una grandezza fisica causi una modifica istantanea dello stato del sistema, nota come \textit{riduzione del pacchetto d'onda} o \textit{collasso}. In questo capitolo verranno enunciati i postulati che governano tale processo, definendo il legame tra gli operatori hermitiani (osservabili), i loro autovalori (risultati di misura) e la natura probabilistica della transizione di stato, formalizzata rigorosamente nell'interpretazione di Copenaghen.

\section{Misure e incertezze associate}
Si è discusso di come, detta $F$ una caratteristica fisica osservabile di una particella in una regione $\Omega$ e $\hat{F}(\bm{r}, -i\hbar \nabla)$ l'operatore hermitiano associato, il suo valore medio \textit{teorico} può essere determinato come:
\begin{equation}
    \langle F \rangle_{th} = \int_\Omega \psi^*\hat{F}(\psi) d^3\bm{r}.
    \label{eq2}
\end{equation}
Se si procede a misurare \textit{macroscopicamente} $F$ un numero di volte $N \to \infty$, detto $f_i$ il risultato dell'$i$-esimo processo di misurazione, il valor medio \textit{sperimentale} è: 
\begin{equation}
    \langle F \rangle_{sp} = \sum_{i = 1}^{N} \frac{f_i}{N}.
    \label{eq:mis_A}
\end{equation}
Detto $f_j$ il generico risultato ottenibile dalla misura e $n_j$ il numero di volte che si ottiene $f_j$, si può riscrivere la \eqref{eq:mis_A} come:
\begin{equation}
    \langle F \rangle_{sp} = \sum_j f_j\frac{n_j}{N},
    \label{eq1}
\end{equation}
dove il termine tra parentesi $n_j/N$ rappresenta la frequenza relativa dell'evento $f_j$, la quale, per la legge dei grandi numeri, converge alla \textit{probabilità} che una singola misura restituisca il valore $f_j$.
Il postulato fondamentale della teoria della misura impone la consistenza tra le due descrizioni, ovvero l'uguaglianza $\langle F \rangle_{th} \equiv \langle F \rangle_{sp}$.

\subsection{Analisi dell'errore - sperimentale}
Una qualsiasi misura sperimentale è affetta da un'\textit{incertezza} (o \textit{errore}) causata da:
\begin{enumerate}
\item errore casuale nel processo di misura [\textit{incertezza di tipo A}];
\item errore sistematico nel processo di misura (causato, ad esempio, da strumenti mal calibrati) [\textit{incertezza di tipo B}];
\item intrinseca precisione non infinita dello strumento di misura.
\end{enumerate}
In generale l'incertezza su una grandezza, conoscendo il suo valor medio e il numero di misure condotte, si esprime tramite lo \textit{scarto quadratico medio}:
\begin{equation}
     \Delta F_{sp}^2 = \sum_{i = 1}^N  \frac{(f_i-\langle F \rangle)^2}{N} = \langle (F - \langle F \rangle) \rangle ^2.
\end{equation}
In un contesto puramente classico, tali incertezze sono considerate limiti \textit{tecnologici} e non fondamentali. Idealmente, nel limite di un numero infinito di misure ($N \to \infty$), gli errori casuali si elidono statisticamente; ipotizzando inoltre di disporre di strumentazione ideale (priva di errori sistematici e a risoluzione infinita), l'incertezza $\Delta F_{sp}$ tenderebbe asintoticamente a zero.
Ne consegue che, nella meccanica classica, lo stato fisico è determinabile con precisione arbitraria: non esiste alcun limite teorico che impedisca di conoscere simultaneamente e con esattezza tutte le proprietà del sistema.

\subsection{Analisi dell'errore - teoria}
Heisenberg cercò di istituire un profondo parallelismo formale tra la statistica classica e la nuova meccanica. In quest'ottica, le definizioni di valor medio e incertezza dovevano mantenere una struttura analoga a quella classica, adattata però alla natura operatoriale della teoria. \\

\noindent
In un certo senso, la \eqref{eq1} e la \eqref{eq2} sono simmetriche: in entrambi i casi, infatti, si combinano i valori ottenibili dalla misura con la probabilità di ottenerli. Nel tentativo di creare un parallelo anche nella definizione di incertezza, Heisenberg ipotizza che:
\begin{equation}
     \Delta F_{th}^2 = \int_\Omega \psi^* [\hat{F}-\langle F \rangle_{th}]^2(\psi) \, d^3\bm{r},
\end{equation}
dove l'operatore $[\hat{F}-\langle F \rangle_{th}]^2$ introdotto è detto di \textit{scarto quadratico medio} ed è rappresentato con $\hat{\mathcal{O}}$. Esiste quindi un nuovo certo contributo all'incertezza su una misurazione: la natura intrinsecamente probabilistica della realtà. 

\section{Il caso della misura con precisione infinita}
Si vuole cercare di scoprire se è teoricamente possibile effettuare una misura che sia affetta da incertezza nulla. 
\begin{equation}
\Delta F_{th}^2 = 0 = \int_\Omega \psi^* [\hat{F}-\langle F \rangle_{th}]^2(\psi) \, d^3\bm{r} = \int_\Omega \psi^* [\hat{F}-\langle F \rangle_{th}][\hat{F}-\langle F \rangle_{th}](\psi) \, d^3\bm{r}.
\end{equation}
Sfruttando l'hermitianità dell'operatore\footnote{L'operatore $\hat{F}$ è hermitiano perché descrive un'osservabile, $\langle F \rangle$ è un numero reale perché misura. Dunque il coniugato dell'operatore $(\hat{F} - \langle F \rangle)$ è uguale a se stesso.}, si continua:
\begin{equation}
\int_\Omega ([\hat{F}-\langle F \rangle_{th}]\psi)([\hat{F}-\langle F \rangle_{th}]\psi)^* \, d^3\bm{r},
\end{equation}
che si riscrive in modo compatto come l'integrale del modulo quadro di una funzione:
\begin{equation}
    \int_\Omega \left| (\hat{F} - \langle F \rangle_{th})\psi \right|^2 \, d^3\bm{r} = 0.
\end{equation}
Poiché il modulo quadro è una quantità definita positiva ($|f|^2 \geq 0$), l'unico modo affinché il suo integrale su tutto lo spazio sia nullo è che la funzione stessa sia identicamente nulla ovunque (a meno di un insieme di misura nulla). Ne consegue che:
\begin{equation}
    (\hat{F} - \langle F \rangle_{th})\psi = 0 \implies \hat{F}\psi = \langle F \rangle_{th}\psi.
\end{equation}
Si è giunti a una conclusione fondamentale: l'incertezza $\Delta F$ è nulla se e solo se la funzione d'onda $\psi$ soddisfa l'\textit{equazione agli autovalori} per l'operatore $\hat{F}$, con autovalore pari proprio al valor medio $\langle F \rangle_{th}$.
In termini fisici, ciò significa che una misura priva di incertezza è possibile solo se $\psi$ è autofunzione dell'operatore $\hat{F}$. In tal caso, il risultato della misura è deterministico e coinciderà con il relativo autovalore.

\subsection{La singola misura}
Si consideri il caso generale in cui lo stato del sistema $\Psi$ non è necessariamente un'autofunzione dell'operatore $\hat{F}$.
Poiché $\hat{F}$ è un'osservabile (operatore hermitiano), il teorema spettrale garantisce che le sue autofunzioni costituiscano una base ortonormale completa per lo spazio di Hilbert. È dunque possibile esprimere $\Psi$ come combinazione lineare delle autofunzioni di base:
\begin{equation}
    \Psi = \sum_i c_i\psi_i,
\end{equation}
dove $\psi_i$ è la generica autofunzione (tale che $\hat{F}\psi_i = f_i \psi_i$) e $c_i$ il relativo coefficiente di combinazione. L'equazione per il valor medio può essere scritta come:
\begin{equation}
    \langle F \rangle_{th} = \int_\Omega \left(\sum_i c_i\psi_i \right)^* \left[\hat{F}\left(\sum_i c_i\psi_i \right)\right] d^3\bm{r}.
\end{equation}
Essendo $\hat{F}$ un'operatore lineare, si continua:
\begin{equation}
        \langle F \rangle_{th} = \int_\Omega \left(\sum_i c_i\psi_i \right)^* \left[\sum_i c_i \underbrace{(\hat{F}\psi_i)}_{f_i \psi_1} \right] d^3\bm{r} = \int_\Omega \left(\sum_i c_i\psi_i \right)^* \left(\sum_i f_ic_i\psi_i \right) d^3\bm{r}.
\end{equation}
Per sviluppare il calcolo, è necessario distinguere gli indici delle due sommatorie (utilizzando ad esempio $j$ e $i$) e sfruttare la linearità dell'integrale. Portando i coefficienti costanti fuori dal segno di integrazione, si ottiene:
\begin{equation}
\begin{aligned}
    \langle F \rangle &= \sum_j c_j^* \sum_i c_i f_i \int_\Omega \psi_j^* \psi_i \, d^3\bm{r} \\
    &= \sum_j \sum_i c_j^* c_i f_i {\delta_{ji}}.
\end{aligned}
\end{equation}
L'integrale coincide con la delta di Kronecker $\delta_{ji}$ in virtù dell'ortonormalità della base. Questo permette di eliminare una delle due sommatorie (ponendo $j=i$):
\begin{equation}
    \langle F \rangle = \sum_i c_i^* c_i f_i = \sum_i |c_i|^2 f_i.
\end{equation}
Il risultato dimostra che il valore medio dell'osservabile non è altro che la media ponderata degli autovalori $f_i$, dove i pesi statistici sono forniti dai moduli quadri dei coefficienti di espansione:
\begin{equation}
    P(f_i) = |c_i|^2,
\end{equation} 
Si deduce allora che il risultato di una singola misurazione effettuata su uno stato $\Psi$ coincide con il valor medio della grandezza (e ha quindi incertezza nulla) se e solo se la distribuzione di probabilità è degenere, ovvero se un coefficiente di combinazione ha modulo quadro unitario (mentre tutti gli altri sono nulli).
Matematicamente, ciò implica che $\Psi$ deve essere un'autofunzione dell'operatore $\hat{F}$. Solo in questa specifica circostanza si verifica la condizione:
\begin{equation}
    f_{misurato} = \langle F \rangle_{teorico} = f_i.
\end{equation}
In tal caso, la natura probabilistica della meccanica quantistica lascia il posto a una predizione deterministica: ripetendo infinite volte la misura su sistemi identici (oppure effettuando una singola misura), si otterrà sempre il medesimo, \textit{unico valore deterministico} coincidente con l'autovalore. \\
Se $\Psi$ non fosse invece autofunzione di $\hat{F}$, il risultato della singola misura sarebbe uno dei suoi autovalori $f_i$, con probabilità dettata dal modulo quadro del suo coefficiente di combinazione lineare che descrive $\Psi$ stesso.

\subsection{Interpretazione fisica: Incertezza intrinseca vs Errore di misura}
Il risultato matematico appena derivato impone una profonda riconsiderazione del concetto di misura rispetto alla fisica classica.
È fondamentale osservare che l'incertezza $\Delta F$ non origina da limiti tecnologici dell'apparato sperimentale o da errori accidentali, bensì è una proprietà intrinseca dello stato $\Psi$ in relazione all'osservabile misurata.
Si possono distinguere due regimi fisici:
\begin{itemize}
    \item \textit{stato di sovrapposizione ($\Delta F \neq 0$):}
    quando lo stato $\Psi$ non è allineato con nessun autovettore della base (ovvero possiede più coefficienti $c_i \neq 0$), il sistema non possiede a priori un valore definito della grandezza $F$.
    In questo scenario, la misura agisce come una proiezione probabilistica. L'incertezza $\Delta F$ quantifica il "grado di sovrapposizione" o, geometricamente, quanto il vettore di stato sia "inclinato" rispetto agli assi definiti dall'operatore.
    La dispersione dei risultati è inevitabile, anche ipotizzando strumenti di precisione infinita;

    \item \textit{autostato ($\Delta F = 0$):}
    nel caso limite in cui il sistema sia preparato in un'autofunzione (es. $\Psi = \psi_k$), la distribuzione di probabilità collassa in una delta di Kronecker ($|c_i|^2 = \delta_{ik}$).
    Solo in questa circostanza la grandezza fisica è "ben definita" e il risultato della misura è prevedibile deterministicamente.
\end{itemize}
In conclusione, l'incertezza quantistica è la misura della \textit{incompatibilità} momentanea tra lo stato in cui il sistema è stato preparato e la domanda (l'osservabile) che gli viene posta.



\subsubsection{Interpretazione geometrica: il disallineamento vettoriale}
È utile visualizzare l'incertezza $\Delta F$ non come un semplice parametro statistico, ma come una proprietà geometrica della relazione tra \textit{Stato} e \textit{Osservabile}.
Nello spazio di Hilbert, l'incertezza emerge quando si interroga il sistema utilizzando una base rispetto alla quale lo stato $\Psi$ non è diagonale.
L'analogia vettoriale permette di fissare il concetto:
\begin{itemize}
    \item l'\textit{Operatore} $\hat{F}$ definisce un sistema di assi ortogonali di riferimento (la base degli autostati $\{\psi_i\}$);
    \item lo \textit{Stato} $\Psi$ è un vettore orientato in questo spazio astratto.
\end{itemize}
Se il vettore stato è perfettamente allineato con uno degli assi definiti dall'operatore, la proiezione è unica e il risultato certo.
Se, al contrario, il vettore è "inclinato" rispetto agli assi di riferimento, esso possiede proiezioni non nulle su molteplici direzioni di base.
In quest'ottica, $\Delta F$ rappresenta la misura quantitativa di tale \textit{disallineamento}: essa indica quanto l'informazione portata dallo stato $\Psi$ risulti "sparpagliata" (o delocalizzata) sullo spettro dei valori possibili dell'operatore $\hat{F}$.

\subsubsection{Un'analogia intuitiva: l'accordo musicale}

Per comprendere il senso fisico di quanto appena discusso, ci si può affidare a un'analogia acustica. Si immagini l'operatore $\hat{F}$ come un \textit{accordatore elettronico}, programmato per rilevare la frequenza di un suono, e lo stato $\Psi$ come l'onda sonora prodotta da uno strumento.
\begin{itemize}
    \item \textit{incertezza nulla (la nota pura):}
    si consideri il caso in cui viene emesso un fischio a una singola frequenza costante (es. 440 Hz). L'onda è "pura".
    Ogni volta che si guarda l'accordatore, questo segnerà esattamente 440 Hz. Non c'è ambiguità: la misura conferma semplicemente una proprietà che il suono possiede già in modo definito.
    Questo corrisponde al caso in cui lo stato è un'\textit{autofunzione}: il sistema \textit{ha} un valore preciso e la misura si limita a rivelarlo;

    \item \textit{incertezza non nulla (l'accordo):}
    si consideri ora il caso in cui vengono premuti tre tasti contemporaneamente su un pianoforte (un accordo, es. Do-Mi-Sol). L'onda sonora che arriva allo strumento è unica, ma è composta dalla sovrapposizione di tre frequenze distinte.
    Se l'accordatore (la misura quantistica) è costretto per sua natura a restituire un \textit{singolo} numero sul display, esso andrà in confusione: segnerà talvolta la frequenza del Do, talvolta quella del Mi, altre volte quella del Sol.
    Prima della misura, il suono conteneva tutte e tre le note simultaneamente. È l'atto di misurare (chiedere "che frequenza sei?") che costringe il sistema a "scegliere" casualmente una delle possibilità.
    L'incertezza $\Delta F$ rappresenta proprio questa ricchezza intrinseca del segnale: non è un errore dello strumento, ma la misura di quanto lo stato sia "polifonico".
\end{itemize}

\subsection{Ripetibilità statistica dell'esperimento}
In conclusione, la natura probabilistica della Meccanica Quantistica non implica che la fisica diventi arbitraria. Sebbene il risultato della singola misura sia impredicibile (se $\Delta F \neq 0$), la statistica dei risultati è rigorosamente deterministica.
Se si considera un insieme statistico (ensemble) di $N$ sistemi identici, tutti preparati nel medesimo stato $\Psi$, e si sottopongono tutti alla misura dell'osservabile $F$, la frequenza relativa $\nu_k$ con cui si osserva l'autovalore $f_k$ convergerà asintoticamente alla probabilità teorica prevista:
\begin{equation}
    \lim_{N \to \infty} \frac{N(f_k)}{N} = |c_k|^2 .
\end{equation}
Ciò significa che lo stato quantistico $\Psi$ contiene un'informazione completa e invariabile non sul destino del singolo sistema, ma sulla \textit{distribuzione di probabilità} dei risultati sperimentali. Ripetendo l'esperimento nelle stesse condizioni, si otterrà sempre, per ogni misura effettuata, la medesima distribuzione statistica\footnote{Dettata dai vari $|c_k|^2$.}.

\section{Il principio del collasso della funzione d'onda}
Fino a questo punto si è discusso della probabilità di ottenere un certo risultato \textit{prima} di effettuare la misura. È necessario ora formalizzare cosa accade allo stato del sistema \textit{immediatamente dopo} che la misura è stata eseguita e ha restituito un valore determinato.
Questo fenomeno è descritto dal postulato della riduzione del pacchetto d'onde (o collasso).

\subsection{Enunciato del postulato}
Se una misura dell'osservabile $F$ sullo stato generico $\Psi$ restituisce il risultato $f_k$ (appartenente allo spettro discreto dell'operatore $\hat{F}$), allora immediatamente dopo la misura lo stato del sistema "collassa" proiettandosi sull'autostato $\psi_k$ associato all'autovalore misurato.
In termini formali, l'atto di misurazione induce una trasformazione discontinua e non unitaria del vettore di stato:
\begin{equation}
    \Psi_{\text{pre}} = \sum_i c_i \psi_i \quad \xrightarrow{\text{misura } F = f_k} \quad \Psi_{\text{post}} = \psi_k.
\end{equation}

\subsection{Conseguenze fisiche}
Il principio del collasso comporta implicazioni fondamentali per la dinamica del sistema quantistico:

\begin{enumerate}
    \item \textit{selezione dell'informazione:}
    la misura agisce come un filtro. Tutte le componenti dello stato originale relative ad autovalori diversi da $f_k$ (i termini $c_i$ con $i \neq k$) vengono distrutte. L'informazione contenuta nella sovrapposizione iniziale viene irreversibilmente persa; dopo la misura, il sistema perde "memoria" del suo stato precedente;

    \item \textit{ripetibilità della misura:}
    una conseguenza diretta del collasso è che una seconda misurazione della stessa osservabile, effettuata immediatamente dopo la prima (prima che il sistema evolva temporalmente), fornirà lo stesso risultato $f_k$ con certezza assoluta (probabilità 1).
    Infatti, essendo ora lo stato $\Psi_{\text{post}} = \psi_k$, esso è diventato un'autofunzione dell'operatore. L'incertezza sulla grandezza $F$ è stata azzerata dall'atto di misurazione;

    \item \textit{preparazione dello stato:}
    operativamente, questo principio fornisce il metodo per "preparare" un sistema quantistico in uno stato noto. Se si desidera ottenere un sistema nello stato $\psi_k$, è sufficiente misurare la grandezza $F$ e selezionare (filtrare) solo quei sistemi che hanno restituito il valore $f_k$.
\end{enumerate}

\begin{equation}
    \text{Evoluzione (Schrödinger)} \leftrightarrow \text{Deterministica, Reversibile}
\end{equation}
\begin{equation}
    \text{Misura (Collasso)} \leftrightarrow \text{Probabilistica, Irreversibile}
\end{equation}

\subsubsection{Intepretazione fisica: la misura come forzatura irreversibile}
L'idea alla base di questo principio rivoluzionario può essere riassunta come segue: l'atto della misura non si limita a registrare una proprietà preesistente, ma \textit{forza} il sistema in maniera irreversibile a manifestarsi in uno stato definito.
Il limite intrinseco della conoscenza quantistica risiede nel fatto che misurare significa obbligare il generico vettore di stato $\Psi$ a "scegliere" un'identità precisa tra quelle disponibili, ovvero quelle dettate dalla diagonalizzazione dell'operatore utilizzato.
L'apparato di misura agisce come un setaccio che ammette solo gli autostati della base $\{\psi_i\}$:
\begin{itemize}
    \item prima della misura, il sistema contiene l'informazione completa di tutti i coefficienti $c_i$ (la sovrapposizione);
    \item a seguito del processo \textit{esterno}, l'unica informazione sopravvissuta è quella relativa allo specifico autostato $\psi_k$ su cui il sistema è precipitato (collassato).
\end{itemize}
Tutta la ricchezza informativa contenuta negli altri termini della sovrapposizione viene irrimediabilmente cancellata. La probabilità che tale collasso avvenga sullo stato $k$-esimo è dettata rigorosamente dal peso statistico $|c_k|^2$.

\begin{approfondimento}[colback= white, colframe=green!35!black, no shadow, title={Origine fisica: Il Dualismo Onda-Corpicolo}]
È legittimo interrogarsi sulla causa fisica profonda che impedisce una descrizione deterministica dei fenomeni quantistici.
La radice del comportamento probabilistico non risiede in una limitazione della conoscenza umana, ma nella natura intrinseca della materia, sintetizzata dal principio del \textit{Dualismo Onda-Corpicolo}.

\subsubsection{Conflitto tra propagazione e interazione}
I costituenti fondamentali della materia (elettroni, fotoni, atomi) esibiscono un comportamento duale:
\begin{itemize}
    \item \textit{propagazione (natura ondulatoria):}
    nell'evoluzione libera (lontano da apparati di misura), il sistema fisico si comporta come un'onda estesa nello spazio. Esso non possiede una traiettoria definita, ma evolve come un campo di probabilità delocalizzato, capace di fenomeni di interferenza e diffrazione. Lo stato $\Psi$ descrive questa realtà ondulatoria;

    \item \textit{interazione (natura corpuscolare):}
    nell'atto della misura, il sistema interagisce con l'apparato macroscopico in modo localizzato, manifestandosi come un corpuscolo discreto (un "quanto").
\end{itemize}

\subsubsection{Genesi dell'indeterminazione}
Tutto il formalismo dell'algebra lineare (autostati, sovrapposizioni, collasso) non è altro che il tentativo matematico di conciliare questi due aspetti contraddittori.
L'atto di misurare un'osservabile (come la posizione) equivale a forzare un ente che è intrinsecamente un'onda (estesa) a comportarsi come una particella (puntiforme).
\begin{itemize}
    \item l'\textit{incertezza} nasce perché concetti corpuscolari (come "posizione esatta") sono mal definiti per un'onda;
    \item la \textit{probabilità} ($\propto |\Psi|^2$) è il legame necessario tra l'intensità dell'onda (realtà pre-misura) e la frequenza degli impatti (realtà post-misura).
\end{itemize}
In sintesi, "succede tutto questo" perché si tenta di proiettare una realtà ondulatoria su una base di misura corpuscolare.

\end{approfondimento}


\section{Misura simultanea e compatibilità delle osservabili}
Dopo aver discusso la misurazione di una singola grandezza, si estende ora l'analisi al caso multivariato. Sorge un interrogativo fondamentale sulla struttura dell'informazione quantistica: è possibile determinare con precisione arbitraria e simultanea i valori di due o più grandezze fisiche distinte?
Mentre in meccanica classica lo stato di un sistema è definito dalla conoscenza congiunta di tutte le variabili dinamiche (es. posizione e impulso), nel formalismo quantistico la possibilità di accumulare informazione è subordinata alla compatibilità algebrica tra gli operatori coinvolti.
Si intende qui indagare le condizioni sotto le quali la misurazione di una grandezza non distrugge l'informazione acquisita su un'altra.

\subsection{Il caso delle due misure}
A due misure differenti corrispondono le osservabili indagate ad esse associate $F$ e $G$, i rispettivi operatori $\hat{F}$ e $\hat{G}$ e le loro equazioni agli autovalori:
\begin{equation}
\begin{dcases}
    \hat{F}\psi_i = f_i\psi_i \\
    \hat{G}\phi_j = g_j \phi_j,
    \end{dcases}
\end{equation}
dove $f_i$ e $g_j$ sono i generici autovalori associati alle autofunzioni $\psi_i$ e $\phi_j$. Si effettui a questo punto, tramite $\hat{F}$, una misura sul generico stato $\Psi$: il risultato è uno degli autovalori dell'operatore $f_i$. Lo stato $\Psi$ collassa così sull'autostato $\psi_k$. Se si sottopone ora il sistema ad un processo di misura tramite $\hat{G}$, si ottiene un risultato $g_j$ affetto, in generale, da una duplice incertezza dovuta all'incertezza sulla prima misura ($\Psi \to \psi_i$) e sulla seconda ($\psi_i \to \phi_j$). \\


\subsection{Quando si ottiene indeterminazione nulla?}
Affinché l'indeterminazione sulle misurazioni di due grandezze fisiche $F$ e $G$ sia simultaneamente \textit{nulla} ($\Delta F = 0$ e $\Delta G = 0$), è necessario che lo stato del sistema $\Psi$ sia un'autofunzione comune a entrambi gli operatori $\hat{F}$ e $\hat{G}$.
Matematicamente, deve verificarsi:
\begin{equation}
    \begin{cases}
        \hat{F}\Psi = f \Psi \\
        \hat{G}\Psi = g \Psi
    \end{cases}
\end{equation}
Questa condizione fisica garantisce che, indipendentemente dall'ordine temporale con cui vengono eseguite le misure, non vi sia perdita di informazione. Il collasso della funzione d'onda indotto dalla prima misura proietta il sistema in uno stato che è già "pronto" (autostato) anche per la seconda misura, preservando così la determinatezza di entrambe le osservabili.
Tuttavia, affinché questa proprietà non sia un caso fortuito limitato a uno specifico stato, ma una caratteristica generale delle due grandezze fisiche, è richiesta una condizione molto più forte a livello operatoriale: gli operatori devono condividere l'\textit{intero set di autostati} (ovvero ammettere una base ortonormale comune).

\subsubsection{Il Commutatore}
Per verificare se due operatori condividono la medesima base senza doverla calcolare esplicitamente, si introduce uno strumento algebrico fondamentale: il \textit{commutatore}.
Dati due operatori $\hat{F}$ e $\hat{G}$, si definisce commutatore l'operatore:
\begin{equation}
    [\hat{F}, \hat{G}] \equiv \hat{F}\hat{G} - \hat{G}\hat{F}.
\end{equation}
L'azione di tale operatore su una generica funzione d'onda di prova $\psi$ permette di valutare se l'ordine di applicazione degli operatori è rilevante:
\begin{equation}
    [\hat{F}, \hat{G}]\psi = (\hat{F}\hat{G} - \hat{G}\hat{F})\psi = \hat{F}(\hat{G}\psi) - \hat{G}(\hat{F}\psi).
\end{equation}
L'utilità del commutatore è evidente se si considerano due operatori che appunto \textit{commutano}, dunque condividono il medesimo set di autovalori:
\begin{equation}
    \begin{cases}
        \hat{F}\psi_k = f_k \psi_k \\
        \hat{G}\psi_k = g_k \psi_k.
    \end{cases}
\end{equation}
Si applichi ora $\hat{G}$ ad ambo i membri della prima uguaglianza. Sfruttando la proprietà di linearità:
\begin{equation}
    \hat{G}\hat{F}\psi_k = \hat{G}f_k \psi_k = f_k\hat{G} \psi_k = f_k g_k \psi_k.
\end{equation}
Se si applica invece $\hat{F}$ alla seconda equazione, forti ancora della linearità, si ottiene:
\begin{equation}
    \hat{F}\hat{G}\psi_k = \hat{F}g_k g_k\psi_k = \hat{F} \psi_k = f_k g_k \psi_k,
\end{equation}
da cui:
\begin{equation}
    \hat{F}\hat{G}\psi_k = \hat{G}\hat{F}\psi_k \implies (\hat{F}\hat{G} - \hat{G}\hat{F})\psi_k = 0,
\end{equation}
condizione vera per ogni $\psi_k$ se e solo se si tratta dell'applicazione nulla:
\begin{equation}
(\hat{F}\hat{G} - \hat{G}\hat{F}) = 0 \implies \hat{F}\hat{G} = \hat{G}\hat{F}.
\end{equation}
Si giunge così al teorema fondamentale della compatibilità:
\begin{itemize}
    \item se $[\hat{F}, \hat{G}] = 0$ (i due operatori commutano), allora $\hat{F}$ e $\hat{G}$ ammettono una base comune di autofunzioni. Le grandezze fisiche sono dette \textit{compatibili} e possono essere misurate simultaneamente con precisione arbitraria su qualsiasi stato\footnote{O meglio, qualsiasi sia lo stato di partenza è possibile prepararlo affinché le misure $F$ e $G$ siano eseguibili in precisione infinita.};
    \item se $[\hat{F}, \hat{G}] \neq 0$, gli operatori non ammettono una base comune completa. Le grandezze sono \textit{incompatibili} e, in generale, la misura precisa di una perturba inevitabilmente la misura dell'altra.
\end{itemize}

\subsection{Commutatività di operatori: esempi}
Si riportano a questo punto alcuni esempi di operatori che commutano e non commutanti. Preliminarmente si discute l'algebra degli operatori di commutazione.

\subsubsection{Algebra dei commutatori}
Il calcolo dei commutatori è semplificato dall'esistenza di una struttura algebrica rigorosa. L'operazione di commutazione $[\hat{A}, \hat{B}]$ gode di proprietà fondamentali che permettono di scomporre espressioni complesse in termini elementari, analogamente a quanto avviene con le regole di derivazione.
Siano $\hat{A}, \hat{B}, \hat{C}$ tre operatori generici e $\lambda \in \mathbb{C}$ uno scalare complesso. Valgono le seguenti identità:

\begin{itemize}
    \item \textit{Antisimmetria:}
    Lo scambio dell'ordine degli operatori inverte il segno del commutatore.
    \begin{equation}
        [\hat{A}, \hat{B}] = - [\hat{B}, \hat{A}].
    \end{equation}
    Da ciò discende immediatamente che ogni operatore commuta con se stesso: $[\hat{A}, \hat{A}] = 0$.

    \item \textit{Linearità:}
    Il commutatore è un'operazione lineare rispetto a entrambi gli argomenti.
    \begin{equation}
        [\hat{A}, \hat{B} + \hat{C}] = [\hat{A}, \hat{B}] + [\hat{A}, \hat{C}],
    \end{equation}
    \begin{equation}
        [\hat{A}, \lambda \hat{B}] = \lambda [\hat{A}, \hat{B}].
    \end{equation}
    Si noti che, poiché gli scalari commutano con qualsiasi operatore ($[\hat{A}, \lambda] = 0$), le costanti moltiplicative possono essere portate fuori dal segno di commutazione.

    \item \textit{Regola di Leibniz (o del prodotto):}
    Questa è l'identità più rilevante per le applicazioni fisiche, poiché regola la commutazione di prodotti di operatori. Essa distribuisce il commutatore in modo analogo alla regola di derivazione del prodotto:
    \begin{equation}
        [\hat{A}, \hat{B}\hat{C}] = [\hat{A}, \hat{B}]\hat{C} + \hat{B}[\hat{A}, \hat{C}].
    \end{equation}
    È cruciale osservare che l'ordine degli operatori non coinvolti nel commutatore parziale deve essere rigorosamente preservato: nel primo termine $\hat{C}$ rimane a destra, nel secondo termine $\hat{B}$ rimane a sinistra.
    Analogamente, per un prodotto nel primo argomento:
    \begin{equation}
        [\hat{A}\hat{B}, \hat{C}] = \hat{A}[\hat{B}, \hat{C}] + [\hat{A}, \hat{C}]\hat{B}.
    \end{equation}
\end{itemize}
L'applicazione sistematica di queste regole, combinata con i commutatori canonici fondamentali (come $[\hat{x}, \hat{p}] = i\hbar$), consente di valutare la compatibilità delle osservabili fisiche in modo puramente algebrico.

\subsubsection{Posizione e quantità di moto}
Da quanto discusso in precedenza, se gli operatori posizione $\hat{x}$ e $\hat{p}$ commutassero, dovrebbe essere $[\hat{x}, \hat{p}] = 0$. Espandendo il calcolo:
\begin{equation}
\begin{aligned}
   &[\hat{x}, \hat{p}]\psi(x) = \hat{x}[\hat{p}(\psi(x))] - \hat{p}[\hat{x}(\psi(x))] = \hat{x}\left[-i\hbar\frac{d}{dx}\psi(x)\right] - \hat{p}[x\psi(x)] \\
  &= -ix\hbar\frac{d}{dx}\psi(x) + i\hbar\frac{d}{dx}(x\psi(x)) =  -ix\hbar\frac{d}{dx}\psi(x) + ix\hbar\frac{d}{dx}\psi(x) + i\hbar  \\&= i\hbar
   \end{aligned}
\end{equation}
Essendo il commutatore diverso da zero, gli operatori quantità di moto e posizione non commutano: \textit{non si possono conoscere simultaneamente con precisione infinita}.

\subsubsection{Energia e Posizione}
Si consideri l'operatore Hamiltoniano per una particella di massa $m$ soggetta a un potenziale $V(x)$:
\begin{equation}
    \hat{H} = \frac{\hat{p}^2}{2m} + V(\hat{x}).
\end{equation}
Si calcola il commutatore tra la posizione $\hat{x}$ e l'energia totale $\hat{H}$. Sfruttando la linearità e il fatto che ogni operatore commuta con una funzione di se stesso\footnote{Nel caso in esame, ha senso pensare che se si conosce la posizione di una particella si conosce simultaneamente anche il valore del potenziale valutato nella posizione stessa.} ($[\hat{x}, V(\hat{x})] = 0$), il calcolo si riduce al termine cinetico:
\begin{equation}
\begin{aligned}
    [\hat{x}, \hat{H}] &= \left[ \hat{x}, \frac{\hat{p}^2}{2m} \right] + \underbrace{[\hat{x}, V(\hat{x})]}_{0} \\
    &= \frac{1}{2m} [\hat{x}, \hat{p}^2] = \frac{1}{2m} \left( \hat{p}[\hat{x}, \hat{p}] + [\hat{x}, \hat{p}]\hat{p} \right).
\end{aligned}
\end{equation}
Sostituendo la relazione canonica $[\hat{x}, \hat{p}] = i\hbar$:
\begin{equation}
    [\hat{x}, \hat{H}] = \frac{1}{2m} ( \hat{p}(i\hbar) + (i\hbar)\hat{p} ) = \frac{i\hbar}{m}\hat{p}.
\end{equation}
Il commutatore è diverso da zero (a meno che l'impulso non sia nullo). Ciò implica che non è possibile determinare esattamente la posizione di una particella e la sua energia totale contemporaneamente. Fisicamente, questo risultato è legato alla velocità: l'operatore velocità è infatti $\hat{v} = \frac{d\hat{x}}{dt} = \frac{1}{i\hbar}[\hat{x}, \hat{H}] = \frac{\hat{p}}{m}$.

\subsubsection{Energia e Quantità di moto}
Analogamente, si valuta la compatibilità tra l'impulso $\hat{p}$ e l'Hamiltoniano $\hat{H}$.
\begin{equation}
    [\hat{p}, \hat{H}] = \left[ \hat{p}, \frac{\hat{p}^2}{2m} \right] + [\hat{p}, V(\hat{x})].
\end{equation}
Il primo termine è nullo poiché ogni operatore commuta con le proprie potenze. Il secondo termine viene valutato applicandolo a una funzione di prova $\psi(x)$:
\begin{equation}
\begin{aligned}
    [\hat{p}, V(x)]\psi &= -i\hbar \frac{d}{dx} (V(x)\psi) - V(x) \left( -i\hbar \frac{d}{dx}\psi \right) \\
    &= -i\hbar \left( \frac{dV}{dx}\psi + V\frac{d\psi}{dx} \right) + i\hbar V \frac{d\psi}{dx} \\
    &= -i\hbar \frac{dV}{dx} \psi.
\end{aligned}
\end{equation}
Pertanto:
\begin{equation}
    [\hat{p}, \hat{H}] = -i\hbar \frac{dV}{dx}.
\end{equation}
In generale, impulso ed energia non commutano. L'unica eccezione è il caso della \textit{particella libera} ($V(x) = \text{costante}$, quindi $V'=0$).
Se è presente un potenziale variabile, l'impulso non è una costante del moto. Questo risultato è l'analogo quantistico della seconda legge di Newton, come sancito dal Teorema di Ehrenfest ($\frac{d\langle p \rangle}{dt} = \langle -\frac{dV}{dx} \rangle = \langle F \rangle$).

\subsubsection{Momento angolare e Modulo quadro}
Un caso fondamentale di grandezze compatibili riguarda il momento angolare. Si vuole verificare se è possibile conoscere simultaneamente una componente (es. lungo l'asse $z$, $\hat{L}_z$) e il modulo quadro del vettore momento angolare $\hat{L}^2 = \hat{L}_x^2 + \hat{L}_y^2 + \hat{L}_z^2$.
Si utilizzano le relazioni di commutazione fondamentali tra le componenti: $[\hat{L}_x, \hat{L}_y] = i\hbar \hat{L}_z$ (e permutazioni cicliche).

\begin{equation}
    [\hat{L}^2, \hat{L}_z] = [\hat{L}_x^2, \hat{L}_z] + [\hat{L}_y^2, \hat{L}_z] + \underbrace{[\hat{L}_z^2, \hat{L}_z]}_{0}.
\end{equation}
Applicando la regola di Leibniz ai primi due termini:
\begin{equation}
\begin{aligned}
    [\hat{L}_x^2, \hat{L}_z] &= \hat{L}_x[\hat{L}_x, \hat{L}_z] + [\hat{L}_x, \hat{L}_z]\hat{L}_x \\
    &= \hat{L}_x(-i\hbar \hat{L}_y) + (-i\hbar \hat{L}_y)\hat{L}_x = -i\hbar (\hat{L}_x \hat{L}_y + \hat{L}_y \hat{L}_x).
\end{aligned}
\end{equation}
\begin{equation}
\begin{aligned}
    [\hat{L}_y^2, \hat{L}_z] &= \hat{L}_y[\hat{L}_y, \hat{L}_z] + [\hat{L}_y, \hat{L}_z]\hat{L}_y \\
    &= \hat{L}_y(i\hbar \hat{L}_x) + (i\hbar \hat{L}_x)\hat{L}_y = +i\hbar (\hat{L}_y \hat{L}_x + \hat{L}_x \hat{L}_y).
\end{aligned}
\end{equation}
Sommando i contributi, i termini si elidono esattamente:
\begin{equation}
    [\hat{L}^2, \hat{L}_z] = 0.
\end{equation}
I due operatori commutano. È dunque possibile costruire una base di autostati comuni (le armoniche sferiche $Y_{l}^{m}$) etichettati simultaneamente dai numeri quantici $l$ (per il modulo quadro) e $m$ (per la componente $z$). Questa è la base della quantizzazione degli atomi.


\subsubsection{Classificazione: Osservabili compatibili e incompatibili}
Alla luce dell'algebra dei commutatori, è possibile suddividere le coppie (o gli insiemi) di grandezze fisiche in due categorie distinte, che definiscono il limite della conoscenza accessibile sul sistema.

\begin{enumerate}
    \item \textit{Osservabili compatibili ($[\hat{A}, \hat{B}] = 0$):}
    Due osservabili si dicono compatibili se i rispettivi operatori commutano.
    In questo scenario:
    \begin{itemize}
        \item esiste una base di autostati comuni a entrambi gli operatori.
        \item l'ordine di misurazione è irrilevante: la misura di $\hat{B}$ non distrugge l'informazione acquisita precedentemente su $\hat{A}$.
        \item è fisicamente possibile preparare il sistema in uno stato in cui entrambe le grandezze assumono valori determinati (incertezza nulla simultanea).
    \end{itemize}
    \textit{Esempio:} la posizione $\hat{x}$ e l'energia potenziale $\hat{V}(\hat{x})$; oppure il modulo del momento angolare $\hat{L}^2$ e la sua componente $\hat{L}_z$.

    \item \textit{osservabili incompatibili o complementari ($[\hat{A}, \hat{B}] \neq 0$):}
    Due osservabili sono incompatibili se i loro operatori non commutano.
    In questo caso:
    \begin{itemize}
        \item non esiste una base comune completa (sebbene possano esistere singoli autostati isolati in comune).
        \item l'ordine di misurazione è cruciale: misurare $\hat{B}$ perturba inevitabilmente lo stato del sistema, rendendo indefinita la grandezza $\hat{A}$ precedentemente misurata.
            \end{itemize}
    \textit{Esempio:} posizione $\hat{x}$ e impulso $\hat{p}$; oppure due diverse componenti del momento angolare $\hat{L}_x$ e $\hat{L}_z$.
\end{enumerate}

\subsubsection{Insieme Completo di Osservabili che Commutano (CSCO)}
L'obiettivo della descrizione quantistica è estrarre la massima quantità di informazione possibile dal sistema senza violare i principi di indeterminazione.
Si definisce \textit{Insieme Completo di Osservabili che Commutano} (CSCO) un set di operatori $\{\hat{A}, \hat{B}, \hat{C}, \dots\}$ tale che:
\begin{itemize}
    \item Tutti gli operatori commutano a coppie tra loro (sono tutti compatibili).
    \item L'insieme dei loro autovalori $\{a, b, c, \dots\}$ è sufficiente a identificare \textbf{univocamente} lo stato del sistema (senza degenerazione).
\end{itemize}
Una volta misurati simultaneamente tutti gli operatori di un CSCO, lo stato è determinato con la massima precisione consentita dalla natura e viene indicato con la notazione dei numeri quantici: $|\Psi\rangle = |a, b, c\rangle$.

\subsection{Precisione infinita negli stati stazionari}
Dalla formulazione dell'equazione di Schrödinger, emerge una proprietà cardinale per i sistemi con Hamiltoniana indipendente dal tempo (sistemi conservativi).
L'energia di uno stato stazionario $\Psi_n$ è conoscibile con precisione infinita. Per tale classe di stati, l'equazione agli autovalori assume la forma:
\begin{equation}
    \hat{H}\Psi_n = E_n\Psi_n.
\end{equation}
Questa relazione implica che l'applicazione dell'operatore energia restituisce lo stato stesso scalato di un fattore costante $E_n$. Ne consegue che:
\begin{itemize}
    \item il valor medio dell'energia coincide con l'autovalore: $\langle H \rangle = E_n$;
    \item l'incertezza statistica è identicamente nulla: $\Delta H = 0$.
\end{itemize}
Si deduce quindi che le locuzioni "stato stazionario" e "autostato dell'energia" sono fisicamente equivalenti.
Un sistema preparato in tale condizione evolve nel tempo mantenendo inalterata la distribuzione di probabilità delle sue osservabili; l'unica variazione è un fattore di fase globale $e^{-iE_n t/\hbar}$ che, scomparendo nel calcolo dei moduli quadri, garantisce la stazionarietà delle grandezze fisiche misurabili.

\section{Il Principio di Indeterminazione di Heisenberg}
La non commutatività di due operatori non rappresenta un mero dettaglio algebrico, ma costituisce il fondamento matematico di uno dei pilastri della meccanica quantistica: il Principio di Indeterminazione.
Esso stabilisce un limite teorico invalicabile alla precisione con cui è possibile preparare un sistema quantistico rispetto a due osservabili incompatibili.

\subsection{Formulazione generale (Disuguaglianza di Robertson)}
Dati due operatori hermitiani generici $\hat{A}$ e $\hat{B}$, si dimostra che il prodotto delle loro incertezze (deviazioni standard) su un qualsiasi stato $\Psi$ deve soddisfare la seguente disuguaglianza:
\begin{equation}
    \Delta A \cdot \Delta B \geq \frac{1}{2} \left| \langle \Psi | [\hat{A}, \hat{B}] | \Psi \rangle \right|.
\end{equation}
Questa relazione, nota come disuguaglianza di Robertson, evidenzia che il limite inferiore dell'indeterminazione dipende dal valore medio del commutatore calcolato sullo stato del sistema.
Se $[\hat{A}, \hat{B}] = 0$ (operatori compatibili), il limite inferiore è nullo e le grandezze sono simultaneamente determinabili.

\subsection{Relazione posizione-impulso}
Applicando la formulazione generale alla coppia di osservabili posizione $\hat{x}$ e impulso $\hat{p}_x$, e ricordando la relazione di commutazione canonica $[\hat{x}, \hat{p}_x] = i\hbar$, si ottiene la forma più celebre del principio.
Poiché il commutatore è una costante ($i\hbar$), il valore medio è indipendente dallo stato:
\begin{equation}
    \left| \langle [\hat{x}, \hat{p}_x] \rangle \right| = |i\hbar| = \hbar.
\end{equation}
Sostituendo nella disuguaglianza generale, si giunge alla relazione di Heisenberg:
\begin{equation}
    \Delta x \cdot \Delta p_x \geq \frac{\hbar}{2}.
\end{equation}

\subsection{Significato fisico e natura ondulatoria}
È cruciale interpretare correttamente tale disuguaglianza:
\begin{itemize}
    \item \textit{proprietà intrinseca:} l'indeterminazione non è frutto di errori sperimentali o disturbi incontrollabili introdotti dallo strumento (sebbene tali effetti esistano). Essa è una proprietà intrinseca dello stato quantistico. Non esiste uno stato fisico in natura che sia "più definito" di quanto concesso da questo limite;
    
    \item \textit{dualismo di fourier:} matematicamente, il principio è una diretta conseguenza della relazione di trasformata di Fourier che lega lo spazio delle posizioni allo spazio degli impulsi (numeri d'onda $k = p/\hbar$).
    Un pacchetto d'onde molto localizzato nello spazio ($\Delta x \to 0$) richiede necessariamente la sovrapposizione di un ampio spettro di frequenze spaziali ($\Delta k \to \infty$, quindi $\Delta p \to \infty$). Viceversa, un'onda con impulso perfettamente definito (onda piana monocromatica, $\Delta p = 0$) è totalmente delocalizzata nello spazio ($\Delta x \to \infty$).
\end{itemize}


\subsection{Lo stato di minima indeterminazione}
Esiste una classe di stati particolari che saturano la disuguaglianza, ovvero per i quali vale il segno di uguaglianza $\Delta x \Delta p = \hbar/2$.
Tali stati sono i \textit{pacchetti d'onda gaussiani}. La funzione gaussiana è l'unica funzione matematica che minimizza il prodotto delle larghezze nelle due rappresentazioni coniugate (spazio e frequenza).

\section{Il Principio di Complementarietà}
Il Principio di Indeterminazione di Heisenberg trova la sua controparte concettuale nel \textit{Principio di Complementarietà}, formulato da Bohr nel 1927.
Mentre Heisenberg stabilisce il limite matematico della precisione ($\Delta A \Delta B \ge \dots$), Bohr definisce la natura logica ed epistemologica della misurazione quantistica.

\subsection{Enunciato e significato fisico}
Il principio afferma che aspetti diversi di un sistema quantistico (come la natura ondulatoria e quella corpuscolare, o la posizione e l'impulso) rappresentano descrizioni \textit{complementari}: esse sono mutualmente esclusive in un singolo esperimento, ma entrambe necessarie per una descrizione completa del fenomeno.
In termini operativi:
\begin{quote}
    \textit{Non è possibile osservare simultaneamente aspetti complementari dello stesso sistema. La scelta dell'apparato sperimentale determina quale aspetto della realtà quantistica si manifesterà, precludendo inevitabilmente l'osservazione dell'altro.}
\end{quote}

\subsection{La scelta degli operatori e la perdita di informazione}
Tradotto nel linguaggio dell'algebra lineare finora sviluppato, il principio di complementarietà si manifesta nella scelta del set di osservabili che si decide di misurare.
L'atto di misura non è una semplice registrazione passiva, ma una \textit{scelta attiva del contesto}:
\begin{enumerate}
    \item \textit{selezione della base:}
    quando lo sperimentatore decide di misurare una grandezza $F$, sta arbitrariamente scegliendo di proiettare lo stato $\Psi$ sulla base degli autovettori di $\hat{F}$. Se decide di misurare un set di grandezze, deve assicurarsi che i relativi operatori commutino (formando un CSCO), definendo così una base comune;
    
    \item \textit{esclusione e cecità:}
    scegliendo di "interrogare" il sistema rispetto alla base di $\hat{F}$, si rinuncia \textit{a priori} alla possibilità di ottenere informazioni precise su qualsiasi osservabile $\hat{G}$ che non commuti con $\hat{F}$ (osservabile complementare);
    
    \item \textit{distruzione dell'informazione:}
    non si tratta solo di ignorare il valore di $\hat{G}$, ma di renderlo fisicamente \textit{indefinito}. Se il sistema possedeva un'informazione definita su $\hat{G}$ prima della misura, l'atto di misurare $\hat{F}$ (e il conseguente collasso della funzione d'onda) distrugge tale informazione, disperdendola nella sovrapposizione statistica della nuova base.
\end{enumerate}
\textbf{In sintesi:} il sistema quantistico possiede una ricchezza di informazioni potenziali che non possono essere attualizzate tutte contemporaneamente. Lo sperimentatore è costretto a scegliere quale "finestra" aprire sul sistema (quale set di operatori commutanti misurare); tale scelta oscura irrimediabilmente la vista attraverso le finestre complementari.





\chapter{Cenni di fisica dello stato solido}

La trattazione svolta nei capitoli precedenti si è focalizzata sulla descrizione quantistica di sistemi isolati o costituiti da un numero ridotto di particelle interagenti. Si estende ora l'indagine allo studio della \textit{materia condensata}, ovvero di sistemi macroscopici in cui il numero di costituenti è dell'ordine del numero di Avogadro ($N \approx 10^{23}$). \\

\noindent
In questo regime, la complessità del problema aumenta drasticamente. L'interazione elettromagnetica tra un numero così elevato di nuclei ed elettroni rende analiticamente intrattabile l'equazione di Schrödinger per la funzione d'onda globale del sistema (problema a molti corpi).
Tuttavia, lo stato solido, e in particolare lo stato \textit{cristallino}, presenta una proprietà fondamentale che ne semplifica radicalmente la descrizione: l'\textit{ordine a lungo raggio}. La disposizione periodica e regolare degli atomi all'interno del reticolo cristallino permette di semplificare il potenziale subito dagli elettroni. Questo passaggio concettuale, supportato dal Teorema di Bloch, porta a una delle conseguenze più rilevanti per l'ingegneria fisica: la trasformazione dei livelli energetici discreti (tipici dell'atomo isolato) in \textit{bande di energia} continue intervallate da regioni proibite (gap).
L'obiettivo di questo capitolo è comprendere come le proprietà microscopiche di simmetria del reticolo determinino le proprietà macroscopiche di conduzione elettrica, permettendo di classificare i materiali in conduttori, isolanti e semiconduttori. Questa classificazione costituisce il fondamento fisico su cui poggia l'intera elettronica moderna.

\subsubsection{Semplificazioni necessarie alla trattazione}
Stante la difficoltà nella risoluzione del'equazione di Schrödinger per un ensamble di elettroni nel reticolo cristallino, risulta fondamentale apportare delle \textit{semplificazioni} al modello:
\begin{enumerate}
\item il potenziale a cui è soggetto l’elettrone è costituito da una schiera di buche di potenziale;
\item il problema è monodimensionale;
\item il solido ha dimensione infinita;
\item il solido contiene un solo elettrone;
\item il solido si trova alla temperatura di $T = \SI{0}{\kelvin}$, cioè i nuclei sono fissi.
\end{enumerate} 
Si procederà, una volta completata l'analisi formale, a rimuovere una a una le approssimazioni per ottenere una trattazione completa.

\section{Il caso dell'elettrone in un reticolo cristallino}
Si consideri, per semplicità, un reticolo monodimensionale di passo $d$ (costante reticolare).
Un elettrone che si muove all'interno di tale struttura non risente del campo di un singolo atomo isolato, bensì della sovrapposizione dei potenziali elettrostatici generati da tutti gli ioni del reticolo.
Assumendo che i nuclei siano fissi nelle loro posizioni di equilibrio $x_n = n \cdot d$ (approssimazione di Born-Oppenheimer), il potenziale effettivo $V(x)$ è dato dalla somma dei contributi coulombiani:
\begin{equation}
    V(x) = -\sum_{n=-\infty}^{+\infty} \frac{Ze^2}{4\pi \varepsilon_0 |x - x_n|}.
\end{equation}
Si noti il segno negativo, indicativo della natura attrattiva dell'interazione nucleo-elettrone.
La caratteristica fondamentale di questa funzione è la sua \textit{periodicità}. Il potenziale, stante la distribuzione spaziale dei nuclei, riproduce se stesso traslando di un passo reticolare:
\begin{equation}
    V(x) = V(x+d).
\end{equation}

\begin{figure}[h!]
    \centering
    \begin{tikzpicture}[scale=1.5, >=stealth]
        % --- CLIP AREA ---
        % Fondamentale per gestire gli asintoti che vanno a -infinito
        \clip (-0.5, -3.5) rectangle (5.5, 1.5);

        % --- DEFINIZIONI ---
        \def\d{2.0} % Passo reticolare

        % --- ASSI ---
        \draw[->] (-0.2, 0) -- (5.2, 0) node[right] {$x$};
        \draw[->] (0, -3.2) -- (0, 1.2) node[above] {$V(x)$};
        \node[below left] at (0,0) {$0$};

        % --- NUCLEI (POSIZIONI) ---
        \foreach \i in {1, 2} {
            % Disegno il nucleo sull'asse
            \filldraw[black] (\i*\d - \d/2, 0) circle (3pt);
            \node[above=5pt] at (\i*\d - \d/2, 0) {\small Nucleo};
            % Asse verticale tratteggiato
            \draw[dotted, gray] (\i*\d - \d/2, 0) -- (\i*\d - \d/2, -3.5);
        }
        % Terzo nucleo parziale a destra
        \filldraw[black] (2.5*\d, 0) circle (3pt);

        % --- CURVA DEL POTENZIALE (IPERBOLICO) ---
        % Usiamo due funzioni asintotiche simulate con curve di Bezier
        \draw[thick, blue!80!black] 
            % Parte sinistra (prima del primo nucleo)
            (-0.2, -0.5) to[out=-10, in=95] (0.9, -3.5);
            
        \draw[thick, blue!80!black]
            % Tra nucleo 1 e nucleo 2
            (1.1, -3.5) to[out=85, in=180] (2.0, -1.2) % Risalita
                        to[out=0, in=95] (2.9, -3.5);  % Discesa

        \draw[thick, blue!80!black]
            % Dopo nucleo 2
            (3.1, -3.5) to[out=85, in=180] (4.0, -1.2) % Risalita
                        to[out=0, in=100] (5.2, -3.0); % Uscita

        % --- ETICHETTE ---
        \draw[<->] (1.0, 0.5) -- (3.0, 0.5) node[midway, fill=white] {$d$};
        \node[text=blue!80!black, align=center] at (4.0, -2.5) {Potenziale\\Coulombiano\\Effettivo};

    \end{tikzpicture}
    \caption{Andamento qualitativo del potenziale reale $V(r)$. Le buche di potenziale coulombiano (iperboliche) divergono negativamente in corrispondenza dei nuclei.}
\end{figure}


\subsection{Dal potenziale reale al modello a buche quadre}
L'andamento reale del potenziale (curva solida in figura) presenta singolarità in corrispondenza dei nuclei ($V \to -\infty$ per $x \to x_n$) e barriere di altezza finita tra un atomo e l'altro.
Tuttavia, la risoluzione analitica dell'equazione di Schrödinger con un potenziale di tipo $1/r$ è estremamente complessa.
Per procedere nello studio delle proprietà elettroniche, si introduce un'idealizzazione nota come \textit{Modello di Kronig-Penney}. Si approssima il potenziale reale con una successione periodica di barriere rettangolari (potenziale a "muri verticali"):
\begin{itemize}
    \item \textit{in prossimità dei nuclei:} si assume che il potenziale sia nullo, rappresentando la "buca" in cui l'elettrone è confinato.
    \item \textit{tra i nuclei:} si assume una barriera di potenziale costante $V_0$ e spessore $b$. La larghezza della buca di potenziale è detta $a$ e vale $a+b = d$.
\end{itemize}
Questa semplificazione geometrica, giustificabile nell'assunzione per cui, all'equilibrio termodinamico\footnote{Perché un elettrone salga nella buca, e perciò non la percepisca come una coppia di rette parallele e verticali, è necessario che venga eccitato. Bisogna quindi ionizzare il materiale, violando l'assunzione di equilibrio termodinamico.}, il generico elettrone si trova nel suo stato fondamentale, molto in profondità nella buca iperbolica di Coulomb, permette di risolvere l'equazione di Schrödinger a tratti, conservando però la caratteristica fisica essenziale del sistema: la periodicità spaziale.

\begin{figure}[h!]
    \centering
    \begin{tikzpicture}[scale=1.5, >=stealth]
        % --- PARAMETRI ---
        \def\a{1.2}     % Larghezza buca (zona classica permessa, V=0)
        \def\b{0.5}     % Larghezza barriera (V=V0)
        \def\Vzero{1.5} % Altezza barriera
        \def\N{2}       % Indice max (0, 1, 2 => 3 celle)

        % --- ASSI ---
        % Calcolo lunghezza totale per l'asse x
        \pgfmathsetmacro{\Ltot}{(\N+1)*(\a+\b) + 0.5}
        \draw[->] (-0.5, 0) -- (\Ltot, 0) node[right] {$x$};
        \draw[->] (0, -0.5) -- (0, \Vzero + 0.5) node[above] {$V(x)$};

        % --- DISEGNO LOOP ---
        % Ciclo da 0 a N (per 3 celle totali)
        \foreach \i in {0, ..., \N} {
            
            % Calcolo la x di inizio di questa cella: x = i * (a + b)
            \pgfmathsetmacro{\startX}{\i * (\a + \b)}
            
            % Disegno il profilo del potenziale per una cella
            \draw[thick, red!80!black] 
                (\startX, \Vzero)          % Parte dall'alto (fine barriera prec)
                -- (\startX, 0)            % Scende a 0 (inizio buca)
                -- (\startX + \a, 0)       % Fondo buca
                -- (\startX + \a, \Vzero)  % Risale a V0
                -- (\startX + \a + \b, \Vzero); % Cima barriera
            
            % --- NUCLEO ---
            % Posizionato al centro della buca (startX + a/2)
            \filldraw[black] (\startX + \a/2, 0) circle (2.5pt);
            \node[below=3pt, font=\scriptsize] at (\startX + \a/2, 0) {Nucleo};
            
            % --- LINEE GUIDA PER IL PERIODO ---
            % Disegno linee tratteggiate sotto i nuclei per indicare il passo d
            \draw[dotted, gray] (\startX + \a/2, 0) -- (\startX + \a/2, -0.8);
        }

        % --- QUOTE E ETICHETTE ---
        
        % Livello V0 e 0
        \draw[dashed] (0, \Vzero) -- (\Ltot, \Vzero) node[right] {$V_0$};
        \node[left] at (0,0) {$0$};

        % Quote a e b (sulla prima cella, i=0)
        \draw[<->] (0, 0.3) -- (\a, 0.3) node[midway, fill=white, inner sep=1pt] {$a$};
        \draw[<->] (\a, 0.3) -- (\a+\b, 0.3) node[midway, fill=white, inner sep=1pt] {$b$};
        
        % Quota Periodo d (tra il primo e il secondo nucleo)
        % x1 = a/2, x2 = a+b + a/2
        \pgfmathsetmacro{\xUno}{\a/2}
        \pgfmathsetmacro{\xDue}{\a + \b + \a/2}
        \draw[<->] (\xUno, -0.6) -- (\xDue, -0.6) node[midway, fill=white] {$d = a+b$};

    \end{tikzpicture}
    \caption{Modello idealizzato di Kronig-Penney. Il potenziale è schematizzato con barriere rettangolari positive di altezza $V_0$. I nuclei sono posizionati al centro delle buche di potenziale (regione a $V=0$).}
    \label{fig:kronig_penney}
\end{figure}

\subsection{Trattazione analitica}
L'equazione di Schrödinger stazionaria per l'elettrone si scrive:
\begin{equation}
\begin{aligned}
&-\frac{\hbar^2}{2m}\frac{d^2\psi(x)}{dx^2} = E\psi(x) \implies \frac{d^2\psi(x)}{dx^2} = -\overbrace{\frac{2mE}{\hbar^2}}^{\alpha^2}\psi(x), \quad V(x) = 0 \\
&-\frac{\hbar^2}{2m}\frac{d^2\psi(x)}{dx^2} + V_0 \psi(x) = E\psi(x) \implies \frac{d^2\psi(x)}{dx^2} = \underbrace{\frac{2m(V_0-E)}{\hbar^2}}_{\gamma^2}\psi(x), \quad V(x) = V_0,
\end{aligned}
\end{equation}
dove si sono nuovamente definiti i parametri $\alpha^2$ e $\gamma^2$ per rendere i calcoli più agevoli, ottenendo le due seguenti equazioni differenziali del secondo ordine:
\begin{equation}
\begin{dcases}
\frac{d^2\psi(x)}{dx^2} = -\alpha^2\psi(x), \quad V(x) = 0 \\
\frac{d^2\psi(x)}{dx^2} = \gamma^2\psi(x), \quad V(x) = V_0.
\end{dcases}
\end{equation}
Si noti come non si è fatto alcun accenno alla periodicità spaziale, che si introduce a questo punto della trattazione sfruttando i risultati del \textit{Teorema di Bloch}.

\subsubsection{Il teorema di Bloch}
Il teorema di Bloch costituisce il fondamento matematico per la descrizione degli stati elettronici in un potenziale periodico.
Esso stabilisce che, in un reticolo cristallino caratterizzato da un potenziale $V(x)$ con periodicità spaziale $d$ (ovvero $V(x+d) = V(x)$), le soluzioni dell'equazione di Schrödinger assumono la forma:
\begin{equation}
    \psi_k(x) = e^{ikx} u(x),
\end{equation}
dove $k$ è un parametro il cui significato fisico verrà chiarito più avanti e $u(x)$ una funzione, detta \textit{funzione di Bloch}, che possiede la \textit{stessa periodicità del reticolo}:
    \begin{equation}
        u_k(x+d) = u(x).
    \end{equation}
Fisicamente, questo risultato implica che gli elettroni in un cristallo non sono localizzati attorno a un singolo atomo, ma sono \textit{delocalizzati} sull'intero solido sotto forma di onde modulate periodicamente dalla struttura atomica.
Una formulazione equivalente del teorema afferma che una traslazione di un passo reticolare modifica la funzione d'onda solo per un fattore di fase costante:
\begin{equation}
    \psi(x+d) = e^{ikd}\psi(x).
\end{equation}

\subsubsection{Conseguenze del Teorema di Bloch}
Sostituendo l'espressione della funzione d'onda dettata dal Teorema di Bloch, le differenziali diventano:
\begin{equation}
\begin{dcases}
\frac{d^2}{dx^2}(e^{ikx} u(x)) = -\alpha^2e^{ikx} u(x), \quad V(x) = 0 \\
\frac{d^2}{dx^2}(e^{ikx} u(x)) = \gamma^2e^{ikx} u(x), \quad V(x) = V_0.
\end{dcases}
\end{equation}
L'incognita diventa allora la funzione di Bloch $u(x)$:
\begin{equation}
\begin{dcases}
\frac{d^2 u(x)}{dx^2} +2ik\frac{du(x)}{dx} -(k^2-\alpha^2)u(x)= 0, \quad V(x) = 0 \\
\frac{d^2 u(x)}{dx^2} +2ik\frac{du(x)}{dx} -(k^2+\gamma^2)u(x)= 0, \quad V(x) = V_0.
\end{dcases}
\end{equation}
Le equazioni caratteristiche associate sono:
\begin{equation}
\begin{aligned}
&\lambda^2 +2ik \lambda - (k^2-\alpha^2) = 0 \implies \lambda_{1,2} = -ik\pm i\alpha \\
&\lambda^2 +2ik \lambda - (k^2+\gamma^2) = 0 \implies \lambda_{1,2} = -ik\pm \gamma,
\end{aligned}
\end{equation}
da cui le soluzioni delle equazioni risultano:
\begin{equation}
u(x) =
\begin{dcases}
e^{-ikx}(Ae^{i\alpha x} + Be^{-i\alpha x}), \quad V(x) = 0 \\
e^{-ikx}(Ce^{\gamma x} + De^{\gamma x}), \quad V(x) = V_0.
\end{dcases}
\end{equation}
Per trovare le costanti di integrazione si dovrebbero imporre condizioni di continuità e derivabilità a $\psi(x)$. Dal Teorema di Bloch si evince che è sufficiente le sfruttare suddette condizioni al contorno sulla funzione $u(x)$\footnote{La funzione di Bloch differisce dalla funzione d'onda per un fattore moltiplicativo esponenziale, ininfluente in questi termini derivativi e di continuità (si eliderebbe a entrambi i membri dell'equazione).}. Sfruttando un sistema di riferimento come quello in Figura \eqref{fig:kronig_penney}, si ottiene il seguente sistema $4\times 4$:
\begin{equation}
\begin{dcases}
A+B-C-D = 0 \\
Ai(-k+\alpha) -Bi(k+\alpha) = C(-ik+\gamma)-D(ik + \gamma) = 0 \\
Ae^{i(-k+\alpha)a} + Be^{-i(k+\alpha)a} - Ce^{-(-ik+\gamma)b}-De^{(ik+\gamma)b} = 0 \\
Ai(-k+\alpha)e^{i(k+\alpha)a} -Bi(k+\alpha)e^{-i(k+\alpha)a} - C(-ik+\gamma)e^{-(-ik+\gamma)b}+D(ik + \gamma)e^{(ik+\gamma)b} = 0.
\end{dcases}
\end{equation}
In termini matriciali, si riscrive come:
\begin{equation}
\begin{bmatrix}
1 & 1 & -1 & -1 \\[10pt]
i(\alpha-k) & -i(\alpha+k) & -(\gamma-ik) & (\gamma+ik) \\[10pt]
e^{i(\alpha-k)a} & e^{-i(\alpha+k)a} & -e^{-(\gamma-ik)b} & -e^{(\gamma+ik)b} \\[10pt]
i(\alpha-k)e^{i(\alpha-k)a} & -i(\alpha+k)e^{-i(\alpha+k)a} & -(\gamma-ik)e^{-(\gamma-ik)b} & (\gamma+ik)e^{(\gamma+ik)b}
\end{bmatrix}
\begin{bmatrix}
A \\
B \\
C \\
D
\end{bmatrix}
= \bm{0}
\end{equation}
che ammette soluzioni non banali se e solo se il determinante della matrice dei coefficienti è nullo. Svolgendo i calcoli, si ottiene la seguente equazione (trascendente):
\begin{equation}
    \frac{\gamma^2-\alpha^2}{2\alpha \gamma}\sinh(\gamma b)\sin(\alpha a) + \cosh(\gamma b)\cos(\alpha a) = \cos[k(a+b)].
\end{equation}
Per semplificare l'espressione, si introduce l'approssimazione della \textit{barriera a delta di Dirac}. Si assume che l'energia dell'elettrone sia trascurabile rispetto all'altezza della barriera ($V_0 \gg E$, quindi $\gamma \gg \alpha$) e che la barriera diventi infinitamente sottile ($b \to 0$) pur mantenendo finita l'area del potenziale ($V_0 b = \text{cost}$).
In questo limite valgono le stime asintotiche $\sinh(\gamma b) \simeq \gamma b$ e $\cosh(\gamma b) \simeq 1$, mentre il passo reticolare si riduce alla sola buca ($a+b \simeq a$).
Sostituendo e semplificando il termine pre-fattore ($\frac{\gamma^2-\alpha^2}{2\alpha\gamma} \gamma b \simeq \frac{\gamma^2 b}{2\alpha}$), l'equazione diviene:
\begin{equation}
    \frac{m V_0 b}{\hbar^2 \alpha} \sin(\alpha a) + \cos(\alpha a) = \cos(ka).
\end{equation}
Moltiplicando e dividendo il primo termine per $a$, si definisce il parametro adimensionale $P$ (potenza della barriera\footnote{Il termine $P$ rappresenta l'effetto globale della barriera sulla particella.}):
\begin{equation}
    P = \frac{m V_0 b a}{\hbar^2}.
\end{equation}
Si giunge infine alla celebre \textit{equazione di Kronig-Penney}:
\begin{equation}
    P \frac{\sin(\alpha a)}{\alpha a} + \cos(\alpha a) = \cos(ka),
\end{equation}
espressione che stabilisce quali siano le coppie energia (contenuta nel fattore $\alpha$) - fattore $k$ per cui un elettrone può esistere nel solido.   

\subsection{L'elettrone libero, significato del termine $\bm{k}$}
Nel limite di elettrone libero, le barriere scompaiono e si ha $V_0 = 0$:
\begin{equation}
 \cos(\alpha a) = \cos(ka),
\end{equation}
vera se e solo se $k = \alpha$, da cui:
\begin{equation}
 k^2 = \alpha^2 \implies k^2 = \frac{2mE}{\hbar^2} \implies E = \frac{\overbrace{k^2\hbar^2}^{p}}{2m} = \frac{p^2}{2m}
\end{equation}
da cui è evidente che $k$ deve essere il vettore d'onda associato all'elettrone ($k = 2\pi/\lambda$ con $\lambda$ la lunghezza d'onda di De Broglie). Si ottiene così lo spettro parabolico di energia continuo, dove tutte le energie sono permesse.

\subsection{Il limite della barriera infinita}
Nel caso di barriere infinite e impenetrabili, con $V_0 \to \infty$ (quindi $P \to \infty$), l'unica possibilità per cui il membro sinistro dell'equazione di Kronig-Penney non diverga è che il suo fattore moltiplicativo si annulli: 
    \begin{equation}
        \frac{\sin(\alpha a)}{\alpha a} \to 0 \implies \sin(\alpha a) = 0 \implies \alpha a = n\pi.
    \end{equation}
    Sviluppando il parametro $\alpha$, si ottiene:
        \begin{equation}
         E_n = \frac{n^2h^2}{8ma^2},
    \end{equation}
    esattamente i livelli discreti di energia permessi\footnote{Gli autovalori dell'hamiltoniano.} a una particella in una buca infinita di potenziale.
    
    \subsection{Il caso generale di $\bm{P}$ finito}
L'equazione di Kronig-Penney lega in generale l'energia dell'elettrone (contenuta nel parametro $\alpha$) al vettore d'onda $k$:
\begin{equation}
    f(\alpha a) = P \frac{\sin(\alpha a)}{\alpha a} + \cos(\alpha a) = \cos(ka).
    \label{eq:KP_transcendental}
\end{equation}
Poiché il membro di destra è una funzione coseno, esso può assumere esclusivamente valori compresi nell'intervallo $[-1, 1]$. Di conseguenza, le uniche soluzioni fisicamente accettabili per l'energia $E$ sono quelle che soddisfano la condizione:
\begin{equation}
    -1 \leq f(\alpha a) \leq 1.
\end{equation}

\begin{figure}[H]
    \centering
    \begin{tikzpicture}[scale=1.2, >=stealth]
        % --- CONFIGURAZIONE ---
        \def\P{6}        % Forza della barriera
        \def\xmax{13.5}  % Estensione asse x
        \def\ymin{-2.2}  % Limite inferiore
        \def\ymax{4.5}   % Limite superiore (tagliamo il picco iniziale a 7 per chiarezza)
        
        % --- CALCOLO BANDE (INTERSEZIONI ESATTE per P=6) ---
        % Banda 1: incrocio f(x)=1 a x~2.38, fine a pi (f(x)=-1)
        \fill[green!15] (2.38, \ymin) rectangle (3.1415, \ymax);
        
        % Banda 2: incrocio f(x)=-1 a x~4.90, fine a 2pi (f(x)=1)
        \fill[green!15] (4.90, \ymin) rectangle (6.2831, \ymax);
        
        % Banda 3: incrocio f(x)=1 a x~7.55, fine a 3pi (f(x)=-1)
        \fill[green!15] (7.55, \ymin) rectangle (9.4247, \ymax);
        
        % Banda 4: incrocio f(x)=-1 a x~10.50, fine a 4pi (f(x)=1)
        \fill[green!15] (10.50, \ymin) rectangle (12.5663, \ymax);

        % --- GRIGLIA E LINEE LIMITE ---
        % Linee tratteggiate a +1 e -1 (zona permessa)
        \draw[dashed, thick, black!60] (0, 1) -- (\xmax, 1) node[right, font=\footnotesize] {+1};
        \draw[dashed, thick, black!60] (0, -1) -- (\xmax, -1) node[right, font=\footnotesize] {-1};
        
        % Asse X
        \draw[->] (0, 0) -- (\xmax, 0) node[right] {$\alpha a$};
        % Asse Y
        \draw[->] (0, \ymin) -- (0, \ymax) node[above] {$f(\alpha a)$};

        % Ticks asse X (multipli di pi)
        \foreach \n in {1, 2, 3, 4} {
            \draw (\n*3.1415, 0.1) -- (\n*3.1415, -0.1);
            \node[below, font=\footnotesize] at (\n*3.1415, -0.1) {$\n\pi$};
        }

        % --- PLOT FUNZIONE ---
        % Clipping per gestire i picchi alti senza uscire dal grafico
        \clip (0, \ymin) rectangle (\xmax, \ymax);
        
        % Funzione: P*sin(x)/x + cos(x)
        % domain parte da 0.1 per evitare divisione per zero
        \draw[thick, blue, samples=300, domain=0.1:\xmax, variable=\x] 
            plot (\x, {\P*sin(\x r)/\x + cos(\x r)});

        % --- ANNOTAZIONI ---
        % Etichette bande
        \node[green!40!black, font=\small, rotate=90] at (2.76, 2.5) {banda 1};
        \node[green!40!black, font=\small, rotate=90] at (5.6, 2.5) {banda 2};
        \node[green!40!black, font=\small, rotate=90] at (8.5, 2.5) {banda 3};
        
        % Etichetta gap (con freccia orizzontale)
        \draw[<->, red, thick] (3.1415, 0) -- (4.90, 0); 
        \node[red, font=\small, fill=white, inner sep=1pt] at (4.0, 0) {gap};

        % Formula nel grafico
        \node[blue, right, fill=white, inner sep=2pt] at (1, 3.5) {
            $f(\alpha a) = P \frac{\sin \alpha a}{\alpha a} + \cos \alpha a$
        };
        
    \end{tikzpicture}
    \caption{Analisi grafica del modello di Kronig-Penney ($P=6$); le zone colorate indicano le bande di energia permesse, corrispondenti ai valori di $\alpha a$ per cui la funzione cade nell'intervallo $[-1, 1]$. Si noti la perfetta corrispondenza tra i limiti delle bande e le intersezioni della curva con le rette orizzontali.}
    \label{fig:kp_plot_revised}
\end{figure}

\subsubsection{Interpretazione dello spettro energetico}
Si analizzi l'andamento della funzione $f(\alpha a)$ riportato in Figura (\ref{fig:kp_plot_revised}).
La funzione è somma di un coseno oscillante e di un termine $\text{sinc}(x)$ modulato dal parametro di forza $P$.
Si osservano due regimi distinti sull'asse delle ascisse (che rappresenta l'energia, essendo $E \propto (\alpha a)^2$):

\begin{itemize}
    \item \textit{bande permesse:} sono gli intervalli di $\alpha a$ in cui la curva $f(\alpha a)$ cade all'interno della fascia compresa tra $-1$ e $+1$ (zone evidenziate in verde); per questi valori di energia, esiste un $k$ reale che soddisfa l'equazione e l'elettrone può propagarsi nel cristallo come un'onda di Bloch;
    
    \item \textit{gap proibiti (band gap):} sono gli intervalli in cui $|f(\alpha a)| > 1$; in queste regioni non esiste alcun $k$ reale soluzione dell'equazione e ciò corrisponderebbe matematicamente a un $k$ immaginario, ovvero a una funzione d'onda evanescente che si smorza esponenzialmente.
\end{itemize}
Si nota inoltre una proprietà fondamentale: al crescere dell'energia (ovvero per $\alpha a$ grandi), il termine $P \frac{\sin(\alpha a)}{\alpha a}$ tende a zero. Di conseguenza, le oscillazioni della funzione si smorzano, rientrando sempre più spesso nella fascia $[-1, 1]$.
I gap di energia diventano progressivamente più stretti all'aumentare dell'energia e lo spettro tende al continuo (come per l'elettrone libero) per $E \to \infty$: si forma una \textit{banda di energia continua}, il cui livello più basso è il limite di energia libera dell'elettrone. Fisicamente, questo risultato indica che gli elettroni ad alta energia si comportano come \textit{particelle libere}: la loro lunghezza d'onda di De Broglie è molto piccola rispetto al passo reticolare e le barriere di potenziale rappresentano solo una perturbazione trascurabile sul loro moto. 

\subsection{Il modello a bande di energia}
Fissando un valore di $E$ che per ogni banda derivante dall'equazione di Kronig-Penney renda possibile l'esistenza di un $k$ e rendendo graficamente per ogni bande energetica si ottiene.
\begin{figure}[h!]
    \centering
    \begin{tikzpicture}[scale=1, >=stealth]
        % --- PARAMETRI ---
        \def\K{1.57} % pi/a (circa 1.57 per scala visiva)
        \def\gap{0.4}
        \def\ymax{6.5}
        \def\xmax{5.5}

        % --- ASSI ---
        \draw[->] (-\xmax, 0) -- (\xmax, 0) node[right] {$k$};
        \draw[->] (0, -0.5) -- (0, \ymax) node[above] {$E$};

        % --- GRIGLIA VERTICALE (Zone di Brillouin) ---
        % Disegno linee a n*pi/a
        \foreach \n in {-3, -2, -1, 1, 2, 3} {
            \draw[dotted, gray] (\n*\K, 0) -- (\n*\K, \ymax);
        }
        % Etichette asse x
        \node[below, font=\scriptsize] at (0, 0) {$0$};
        \node[below, font=\scriptsize] at (\K, 0) {$\frac{\pi}{a}$};
        \node[below, font=\scriptsize] at (-\K, 0) {$-\frac{\pi}{a}$};
        \node[below, font=\scriptsize] at (2*\K, 0) {$\frac{2\pi}{a}$};
        \node[below, font=\scriptsize] at (3*\K, 0) {$\frac{3\pi}{a}$};

        % --- FUNZIONI PERIODICHE (Bande) ---
        % Usiamo coseni traslati e scalati per simulare le bande piegate
        
        % BANDA 1 (Minimo a k=0, Max ai bordi)
        % Periodo 2*K. Forma: convessa verso l'alto
        % y = 1 - cos(x) scalato
        \foreach \shift in {-2, 0, 2} {
             \draw[thick, blue, domain=(\shift-1)*\K:(\shift+1)*\K, samples=50] 
                plot (\x, {1.0 - 0.8*cos(\x/\K*180)});
        }
        
        % BANDA 2 (Minimo ai bordi, Max a k=0)
        % Sopra la banda 1 con un gap. Concava verso il basso a k=0.
        \foreach \shift in {-2, 0, 2} {
             \draw[thick, red!80!black, domain=(\shift-1)*\K:(\shift+1)*\K, samples=50] 
                plot (\x, {2.8 + 0.8*cos(\x/\K*180)});
        }
        
        % BANDA 3 (Minimo a k=0, Max ai bordi)
        % Sopra la banda 2 con un gap.
        \foreach \shift in {-2, 0, 2} {
             \draw[thick, green!60!black, domain=(\shift-1)*\K:(\shift+1)*\K, samples=50] 
                plot (\x, {4.6 - 0.6*cos(\x/\K*180)});
        }
        
        % --- EVIDENZIAZIONE 1° ZONA BRILLOUIN ---
        % Rettangolo di sfondo per evidenziare la zona ridotta
        \fill[gray, opacity=0.1] (-\K, 0) rectangle (\K, \ymax);
        \node[gray, font=\footnotesize, anchor=north] at (\K/2, \ymax) {1\textsuperscript{a} Z.B.};

        % --- GAP ENERGETICO ---
        % Disegno il gap tra banda 1 e 2 a pi/a
        \draw[<->, black, thick] (\K, 1.8) -- (\K, 2.0); % (fine banda 1 a inizio banda 2)
        % Nota: con i parametri sopra, banda 1 finisce a 1.8, banda 2 inizia a 2.0
        \node[right, font=\footnotesize] at (\K, 1.9) {$E_g$};

        % --- ETICHETTE BANDE (nella zona centrale) ---
        \node[blue, font=\scriptsize, fill=white, inner sep=1pt] at (0.8, 0.5) {1};
        \node[red!80!black, font=\scriptsize, fill=white, inner sep=1pt] at (0.8, 3.0) {2};
        \node[green!60!black, font=\scriptsize, fill=white, inner sep=1pt] at (0.8, 4.3) {3};

    \end{tikzpicture}
    \caption{Schema a zone ripetute. Poiché l'energia è una funzione periodica nel reticolo reciproco ($E(k) = E(k + G)$), è possibile rappresentare l'intera struttura a bande ripetendo la prima zona di Brillouin.}
    \label{fig:band_structure_repeated}
\end{figure}

\noindent
La proiezione dei valori assunti sull'asse delle ordinate dipinge le varie \textit{bande di energia permesse}. Gli intervalli che non sono immagine di alcuna relazione $E(k)$ corrispondono alle \textit{bande di energia proibite}. \\
L'insieme delle relazioni $E(k_x)$ (in Figura) $E(k_y)$ e $E(k_z)$ (perfettamente simmetrici al caso monodimensionale discusso) costituiscono il cosiddetto \textit{schema a bande di energia} e lo spazio identificato dalla terna di vettori d'onda $(k_x,k_y,k_z)$ viene detto \textit{spazio reciproco}.

\subsubsection{Periodicità nel reticolo reciproco e Prima Zona di Brillouin}
La struttura periodica della relazione di dispersione $E(k)$ emerge analiticamente osservando il termine dominante nell'equazione trascendente, ovvero $\cos(ka)$. Sfruttando la proprietà di periodicità delle funzioni trigonometriche, si nota che:
\begin{equation}
    \cos(ka) = \cos(ka + 2\pi n) = \cos\left[\left(k + \frac{2\pi}{a}n\right)a\right], \quad \text{con } n \in \mathbb{Z}.
\end{equation}
Questa invarianza suggerisce l'introduzione di un vettore fondamentale del reticolo reciproco, definito come $G = 2\pi/a$.
Fisicamente, ciò implica che i vettori d'onda $k$ e $k' = k+G$ non descrivono stati fisici distinti, bensì corrispondono al medesimo stato elettronico all'interno del cristallo. Ne consegue che l'autovalore dell'energia deve soddisfare la condizione di periodicità:
\begin{equation}
    E(k) = E(k+G).
\end{equation}
Tale proprietà rende ridondante l'analisi dello spettro su tutto l'asse reale dei vettori d'onda\footnote{Con un abuso di notazione, il fatto che l'autovalore energia dell'equazione stazionaria di Schrödinger sia periodico, implica necessariamente che l'autofunzione associata sia anch'essa "periodica" $\psi(k) = \psi(k+G)$.}. Poiché l'informazione fisica si ripete periodicamente ogni intervallo di ampiezza $G$, è sufficiente restringere lo studio a una singola cella primitiva nello spazio $k$.
Considerando inoltre che la relazione di dispersione è una funzione pari, ovvero $E(k) = E(-k)$ (simmetria di inversione), si sceglie convenzionalmente l'intervallo simmetrico intorno all'origine:
\begin{equation}
    -\frac{\pi}{a} \leq k \leq \frac{\pi}{a}.
\end{equation}
Tale regione è definita come \textit{Prima zona di Brillouin} (= 1ZB). Ogni vettore d'onda $k$ (più propriamente detto \textit{pseudo-quantità di moto}) esterno a questo intervallo può essere ricondotto al suo equivalente interno tramite una traslazione di un vettore reticolare $nG$. Confinare il vettore d'onda nella 1ZB significa assumere come valori possibili della quantità di moto quelli compresi fra:
\begin{equation}
    -\frac{h}{2a} \leq p \leq \frac{h}{2a}.
\end{equation}

\section{Introduzione della finitezza del solido}
Mantenendo un approccio monodimensionale (eventualmente estensibile) e assumendo una schiera infinita di atomi equidistanti $a$, la si sezioni in $L$ parti uguali, ciascuna contenente un numero $\mathcal{N}_a = L/a$ di atomi. Essendo le componenti della divisione uguali, la funzione d'onda che descrive un elettrone che vive nella schiera deve essere $L$-periodica: $\psi(x) = \psi(x+L)$. Per il Teorema di Bloch, poi:
\begin{equation}
    \psi(x) = u(x)e^{ikx}
\end{equation}
con la funzione di Bloch $a$-periodica\footnote{Si ricordi che la funzione di Bloch ha la stessa periodicità del potenziale scalare in cui è immersa la particella. In questo caso, stante l'equispaziatura del reticolo, deve essere che $V(x) = V(x + a)$.}. Essendo $L$ multiplo di $a$, si scrive:
\begin{equation}
     \psi(x) = u(x)e^{ikx} = \psi(x+L) = u(x+na)e^{ik(x+na)} = u(x+L)e^{ik(x+L)}. 
\end{equation}
La funzione di Bloch a membro destro e sinistro dell'equazione deve essere uguale se valutata in traslazioni multiple del periodo fondamentale, da cui:
\begin{equation}
e^{ikx} = e^{ik(x+L)} = e^{ikx}e^{ikL} \implies 1 = e^{ikL}.
\end{equation}
Continuando con l'identità di Eulero, poi:
\begin{equation}
1 = \cos(kL)+i\sin(kL) \implies k = \frac{2\pi}{L} n, \quad n \in \mathbb{Z}.
\end{equation}
Il numero d'onda può assumere quindi esclusivamente valori discreti nello schema a bande, in particolare nella 1ZB. Il numero di valori di $k$ disponibili (proporzionale al numero di \textit{stati quantistici disponibili}) nella 1ZB, poi:
\begin{equation}
\mathcal{N}_k = \frac{2\pi/a}{2\pi/L} = \frac{L}{a} = \mathcal{N}_a
\end{equation}
che corrisponde al numero di atomi nel solido.


\section{Sistemi a più elettroni}
Fino a questo punto si è analizzato il comportamento di un singolo elettrone immerso nel potenziale periodico. In un cristallo macroscopico reale, tuttavia, è presente un numero di elettroni dell'ordine del numero di Avogadro ($\sim 10^{23}$). Poiché gli elettroni sono fermioni, essi obbediscono al \textit{Principio di esclusione di Pauli}: ogni stato quantico può essere occupato al massimo da un singolo elettrone.
Considerando anche il grado di libertà di spin (che può assumere due valori, \textit{up} e \textit{down}), la capacità totale di una singola banda di energia viene determinata moltiplicando per due il numero di stati orbitali disponibili nella prima zona di Brillouin.
Poiché, come dimostrato precedentemente, il numero di vettori d'onda $k$ ammissibili nella 1ZB coincide con il numero di celle primitive $\mathcal{N}_a$ (o atomi, nel caso di reticolo monoatomico), si deduce che:
\begin{equation}
    \text{Capacità di una banda} = 2 \times \mathcal{N}_a.
\end{equation}
Allo zero assoluto ($T=0$ K), il riempimento degli stati segue il principio di minima energia: gli elettroni occupano progressivamente tutti i livelli disponibili partendo dal fondo della buca di potenziale fino a esaurimento del numero totale di particelle.
L'ultimo livello energetico occupato definisce l'\textit{Energia di Fermi} ($E_F$). La configurazione elettronica fondamentale del solido dipende dunque crucialmente dalla posizione di $E_F$ rispetto alla struttura a bande.

\subsection{Riempimento delle bande}
In un cristallo elementare e perfetto, il numero di elettroni $\mathcal{N}_e$ è legato al numero totale di atomi $\mathcal{N}_a$ dal numero atomico $Z$ caratteristico dell'elemento:
\begin{equation}
    Z\mathcal{N}_a = \mathcal{N}_e, \quad Z \in \mathbb{N}
\end{equation}
da cui il numero di bande occupate:
\begin{equation}
    \mathcal{N}_B = \frac{\mathcal{N}_e}{2\mathcal{N}_a} = \frac{Z}{2}, \quad Z \in \mathbb{N}.
\end{equation}
 Possono dunque verificarsi due condizioni di riempimento complementari:
\begin{itemize}
\item \textit{ultima bande piena}: se $Z$ è pari, gli elettroni riempiono completamente l'ultima banda;
\item \textit{ultima bande semi-piena}: se $Z$ è dispari, gli elettroni riempiono per metà l'ultima banda.
\end{itemize}
Si segnala che l'ultima banda occupata da elettroni allo zero assoluto è detta \textit{Banda di Valenza} (BV), la prima banda vuota è detta \textit{Banda di Conduzione} (BC).

\subsection{Classificazione elettrica dei solidi}
La distinzione fenomenologica tra materiali conduttori e isolanti non risiede nella capacità degli elettroni di muoversi (la velocità di gruppo è non nulla in ogni banda), ma nella possibilità di modificare la distribuzione degli elettroni sotto l'azione di un campo elettrico esterno.
Sulla base del grado di riempimento delle bande a $T=0$ K, si distinguono tre categorie fondamentali:

\begin{itemize}
    \item \textit{conduttori (metalli):} si hanno quando l'ultima banda contenente elettroni è piena solo parzialmente oppure quando vi è sovrapposizione energetica tra una banda piena e una vuota; in questo scenario, il livello di Fermi cade all'interno di una banda permessa e vi sono stati vuoti accessibili a distanza energetica infinitesima sopra $E_F$, permettendo agli elettroni di guadagnare energia da un campo esterno e generare corrente;
    
    \item \textit{isolanti:} si verificano quando un numero intero di bande è completamente riempito e la successiva banda vuota è separata da un ampio gap di energia ($E_g \gg k_B T$, tipicamente $> 3-4$ eV); una banda completamente piena non può trasportare corrente elettrica netta poiché, per ogni elettrone con velocità $v$, ne esiste uno con velocità $-v$ e l'assenza di stati liberi vicini impedisce l'alterazione di questo equilibrio;
    
    \item \textit{semiconduttori:} presentano una configurazione qualitativamente identica a quella degli isolanti (banda di valenza piena e banda di conduzione vuota a $T=0$ K), ma sono caratterizzati da un gap di energia ridotto ($E_g \sim 1$ eV); questa piccola separazione energetica consente, a temperature finite ($T > 0$), l'eccitazione termica di una frazione di elettroni dalla banda di valenza a quella di conduzione, rendendo il materiale debolmente conduttivo.
\end{itemize}


\subsection{Approssimazione parabolica delle bande}
Dall'analisi del modello di Kronig-Penney, si osserva che nel limite in cui il parametro di forza della barriera $\alpha$ tende a infinito, lo spettro energetico diviene discreto. Tuttavia, per valori finiti di $\alpha$, si ottiene una relazione di dispersione continua che, in prossimità dei minimi, presenta un andamento \textit{parabolico}.
Questa osservazione suggerisce la possibilità di trattare l'elettrone in un cristallo periodicamente perturbato in analogia a una particella libera\footnote{La cui energia è esclusivamente quadratica in $k$.}, a patto di rinormalizzarne la massa inerziale introducendo il concetto di \textit{massa efficace}.

\subsubsection{Espansione energetica in prossimità degli estremi}
Si consideri il caso monodimensionale e sia $E(k)$ la funzione che descrive l'energia in funzione del vettore d'onda (o pseudo-quantità di moto) $k$. Analizzando la prima zona di Brillouin (1ZB), si assume che la banda di interesse (ad esempio la banda di conduzione) presenti un minimo per $k = 0$.
Sviluppando $E(k)$ in serie di Maclaurin (Taylor attorno all'origine) e troncando al secondo ordine, si ottiene:
\begin{equation}
    E(k) \approx E(0) + \eval{\frac{d E}{d k}}_{k = 0}k + \frac{1}{2}\eval{\frac{d^2 E}{d k^2}}_{k = 0}k^2.
\end{equation}
Ponendo lo zero dell'energia in corrispondenza del fondo della banda, si ha $E(0) = 0$. Inoltre, essendo $k=0$ un punto stazionario (minimo), la derivata prima deve annullarsi:
\begin{equation}
    \eval{\frac{d E}{d k}}_{k = 0} = 0.
\end{equation}
L'espressione dell'energia si riduce quindi al solo termine quadratico:
\begin{equation}
    E(k) \approx \frac{1}{2}\eval{\frac{d^2 E}{d k^2}}_{k = 0}k^2.
    \label{eq:parabolic_approx}
\end{equation}
Questa approssimazione risulta valida per valori di $k$ piccoli, ovvero in prossimità del centro della 1ZB.

\subsubsection{Definizione di massa efficace}
L'energia cinetica di una particella libera di massa $m_0$ è data da:
\begin{equation}
    E_{free} = \frac{p^2}{2m_0} = \frac{\hbar^2 k^2}{2m_0}.
    \label{eq:free_electron}
\end{equation}
Confrontando l'equazione (\ref{eq:free_electron}) con l'espansione parabolica (\ref{eq:parabolic_approx}), si impone l'uguaglianza tra i coefficienti del termine $k^2$. Affinché l'elettrone nel cristallo possa essere descritto formalmente come una particella libera, si introduce una massa fittizia $m^*$ tale che:
\begin{equation}
    \frac{\hbar^2}{2m^*} = \frac{1}{2} \eval{\frac{d^2 E}{d k^2}}_{k=0}.
\end{equation}
Da questa uguaglianza discende la definizione formale di massa efficace:
\begin{equation}
    m^* \equiv \hbar^2 \left( \eval{\frac{d^2 E}{d k^2}}_{k=0} \right)^{-1}.
    \label{eq:effective_mass_def}
\end{equation}
La massa efficace è dunque il parametro che "consta" delle interazioni reticolari dell'elettrone, risultando inversamente proporzionale alla curvatura della banda\footnote{Le bande di energia, con l'aumentare della curvatura vedono aumentare anche la loro ampiezza. Da qui la diminuzione della massa efficace all'aumentare della curvatura: in bande più ampie, un elettrone è più libero di muoversi. La sua inerzia (rappresentata dalla massa) deve necessariamente essere minore.}: maggiore è la curvatura (banda "stretta"), minore sarà la massa efficace.  

\subsubsection{Dinamica semiclassica e interpretazione fisica}
L'introduzione della massa efficace permette di preservare la validità della seconda legge di Newton per un pacchetto d'onde elettronico sottoposto a una forza esterna $F_{ext}$ (ad esempio un campo elettrico).
La velocità di gruppo del pacchetto d'onde è definita come:
\begin{equation}
    v_g = \frac{1}{\hbar} \frac{dE}{dk}.
\end{equation}
L'accelerazione subita dall'elettrone nel tempo è quindi:
\begin{equation}
    a = \frac{dv_g}{dt} = \frac{1}{\hbar} \frac{d}{dt} \left( \frac{dE}{dk} \right) = \frac{1}{\hbar} \frac{d^2E}{dk^2} \frac{dk}{dt}.
\end{equation}
Ricordando che la variazione del momento cristallino è legata alla forza esterna dalla relazione $\hbar \frac{dk}{dt} = F_{ext}$, si ottiene:
\begin{equation}
    a = \frac{1}{\hbar^2} \frac{d^2E}{dk^2} F_{ext}.
\end{equation}
Sostituendo la definizione (\ref{eq:effective_mass_def}), si recupera la legge del moto classica:
\begin{equation}
    F_{ext} = m^* a.
\end{equation}
Questo risultato dimostra che la massa efficace ingloba in sé l'interazione dell'elettrone con il potenziale periodico del reticolo. In particolare:
\begin{itemize}
    \item se la curvatura è positiva (minimo di banda), $m^* > 0$ (comportamento tipo elettrone);
    \item se la curvatura è negativa (massimo di banda, es. banda di valenza), $m^* < 0$. In questo contesto è conveniente trattare la lacuna come particella a massa positiva e carica opposta.
\end{itemize}

\subsection{Densità degli stati nelle bande}

La naturale conseguenza della descrizione energetica delle bande è la determinazione del numero di stati quantistici disponibili per unità di volume a una data energia $E$. Si introduce a tal fine la funzione \textit{densità degli stati} $g(E)$, definita in modo tale che il numero di stati disponibili per unità di volume nell'intervallo infinitesimo $[E, E+dE]$ sia:
\begin{equation}
    dn = g(E) dE
\end{equation}
dove $\mathcal{N}_{tot}$ rappresenta il numero cumulativo di stati quantistici ed $V_{cristallo} = L^3$ è il volume fisico del materiale.

\subsubsection{Conteggio degli stati nello spazio dei momenti}
Per determinare la distribuzione dei vettori d'onda permessi, si consideri un cristallo di forma cubica con lato $L$. Le condizioni al contorno impongono che le componenti del vettore d'onda $\bm{k} = (k_x, k_y, k_z)$ siano quantizzate:
\begin{equation}
    k_i = \frac{2\pi}{L} n_i, \quad \text{con } i \in \{x,y,z\} \text{ e } n_i \in \mathbb{Z}.
\end{equation}
Questa quantizzazione individua una griglia di punti nello spazio $k$, dove la distanza minima tra due stati adiacenti è $\Delta k = 2\pi/L$. È possibile associare a ogni punto della griglia un volumetto elementare cubico $\Omega_k$:
\begin{equation}
    \Omega_k = (\Delta k)^3 = \left( \frac{2\pi}{L} \right)^3 = \frac{8\pi^3}{L^3}.
\end{equation}
Geometricamente, ogni cubetto $\Omega_k$ possiede 8 vertici. Tuttavia, in una griglia estesa (considerando $\bar{k} \gg \Delta k$), ogni vertice è condiviso contemporaneamente da 8 cubi adiacenti. Di conseguenza, il contributo netto di ogni vertice al singolo cubetto è $1/8$. Poiché ci sono 8 vertici, si ha che $8 \times (1/8) = 1$: in media, ogni volume elementare $\Omega_k$ "contiene" un solo vettore d'onda $\bm{k}$ permesso.
Il numero di stati $\mathcal{N}_{k}$ con vettore d'onda inferiore a un certo raggio $\bar{k}$ si ottiene quindi dividendo il volume della sfera di raggio $\bar{k}$ per il volume elementare $\Omega_k$:
\begin{equation}
    \mathcal{N}_{k}(\bar{k}) = \frac{V_{sfera}}{\Omega_k} = \frac{\frac{4}{3}\pi \bar{k}^3}{\frac{8\pi^3}{L^3}} = \frac{L^3}{6\pi^2} \bar{k}^3.
\end{equation}

\subsubsection{Inclusione dello spin e densità energetica}
Per ottenere il numero totale di stati quantistici disponibili $\mathcal{N}_{tot}$, è necessario considerare la natura fermionica degli elettroni. Secondo il Principio di Esclusione di Pauli, ogni stato orbitale definito dal vettore $\bm{k}$ può essere occupato da al massimo due elettroni, aventi numero quantico di spin opposto ($s = \pm 1/2$).
È quindi necessario introdurre un fattore moltiplicativo pari a 2:
\begin{equation}
    \mathcal{N}_{tot}(\bar{k}) = 2 \cdot \mathcal{N}_{k}(\bar{k}) = \frac{L^3}{3\pi^2} \bar{k}^3.
\end{equation}
Per esprimere tale numero in funzione dell'energia, si inverte la relazione parabolica $E = (\hbar^2 k^2)/(2m^*)$, ottenendo $\bar{k} = \frac{\sqrt{2m^* E}}{\hbar}$. Sostituendo nella relazione precedente:
\begin{equation}
    \mathcal{N}_{tot}(E) = \frac{L^3}{3\pi^2} \left( \frac{2m^* E}{\hbar^2} \right)^{3/2}.
\end{equation}
Il numero di stati nell'intervallo di energia $[E, E+dE]$ lo si valuta come:
\begin{equation}
    \mathcal{N}_{tot}(E+dE) - \mathcal{N}_{tot}(E) \approx \pdv{\mathcal{N}_{tot}}{E}dE.
\end{equation}
Essendo il numero di stati all'energia $E$ dato da $dn = g(E)dE$, si evince che la densità degli stati disponibili $g(E)$ si ottiene infine derivando rispetto all'energia $\mathcal{N}_{tot}$ e normalizzando per il volume del cristallo $L^3$:
\begin{equation}
    g(E) = \frac{1}{L^3} \frac{d \mathcal{N}_{tot}}{dE} = \frac{1}{2\pi^2} \left( \frac{2m^*}{\hbar^2} \right)^{3/2} \sqrt{E}.
\end{equation}

\section{Sistemi a temperature superiori allo zero assoluto}
L'imposizione della temperatura $T = \SI{0}{\kelvin}$ permette di trascurare fenomeni di agitazione termica e di posizionare gli elettroni nelle bande seguendo il criterio di minima energia potenziale, giungendo così a una distinzione formale fra conduttori e isolanti. Che un sistema si trovi però allo zero assoluto è impossibile: è necessario rimuovere dunque quest'ultima ipotesi semplificativa. 

\subsection{Processo di riscaldamento}
In un cristallo sottoposto a un processo di riscaldamento, l'energia termica proveniente dall'esterno eccita i nuclei atomici. L'agitazione termica casuale derivante dall'aumento di temperatura rende la probabilità che un nucleo urti un elettrone diversa da zero: anche gli elettroni, inizialmente confinati nelle buche di potenziale, si agitano. Il movimento disordinato e casuale degli elettroni può essere tale, in determinate condizioni, da permettere a un corpuscolo in banda di valenza di raggiungere la banda di conduzione: viene così a crearsi un "buco" in BV (detto \textit{lacuna}) mentre la BC consta di un nuovo occupante. Questo processo caotico potrebbe, in linea di principio, mutare quindi le caratteristiche elettriche di un materiale? Può un isolante diventare un conduttore e viceversa, casualmente? 

\subsubsection{L'ipotesi di Fermi}
La risposta all'ultimo fantomatico quesito è \textit{no} e la motivazione viene suggerita da Fermi. Deve infatti esistere una distribuzione elettronica negli stati disponibili del cristallo tale per cui la probabilità che si realizzi è significativamente maggiore rispetto a tutte le \textit{configurazioni spurie} che implicherebbero una variazione dello stato del materiale. 

\subsection{La distribuzione di Fermi-Dirac}
\subsubsection{Assunzioni teoriche}
Dette $g(E)$ e $n(E)$ rispettivamente la \textit{densità degli stati disponibili} e la \textit{densità degli stati occupati} all'energia $E$, deve in generale valere:
\begin{itemize}
\item $g(E) \geq n(E)$: il principio di esclusione di Pauli garantisce che non possono esistere più elettroni che condividono lo stesso stato quantistico;
\item l'insieme degli elettroni viene interpretato come un \textit{gas di fermioni} che risponde alla meccanica statistica e alla termodinamica: l'energia totale del sistema deve quindi essere $E_{tot} = \sum_E n(E) E$;
\item il numero totale di elettroni è $\mathcal{N}_e = \sum_E n(E)$.
\end{itemize}


\subsubsection{Analisi combinatoria}
Detti $g_i = g(E_i)dE$ il numero di stati disponibili e $n_i$ numero di stati occupati all'energia $E_i$, il numero $W_i$ di combinazioni realizzabili elettrone - stato disponibile sono date dal coefficiente binomiale\footnote{Si sta essenzialmente risolvendo il problema del sistemare $n_i$ palline \textit{indistinguibili} in $g_i$ scatole.}:  
\begin{equation}
W_i = 
    \begin{pmatrix}
    g_i \\
    n_i
    \end{pmatrix}
    = \frac{g_i!}{n_i!(g_i-n_i)!}.
\end{equation}
La totalità delle realizzabili sistemazioni elettrone - stato disponibile nell'intero schema a bande è ottenuta dal prodotto delle singole $W_i$, essendo indipendenti le une dalle altre: 
\begin{equation}
W = \prod_i W_i.
\end{equation}

\subsubsection{Analisi termodinamica}
Dal primo principio della termodinamica, lo scambio di energia $Q$ fra nuclei e ambiente esterno comporta:
\begin{equation}
\delta Q = dU + \delta W.
\end{equation}
Assumendo una variazione di volume trascurabile nel caso del cristallo, si ha $\delta W = 0$, da cui:
\begin{equation}
\delta Q = dU.
\end{equation}
Per definizione, la variazione infinitesima di entropia corrisponde a:
\begin{equation}
dS = \eval{\frac{\delta Q}{T}}_{rev},
\end{equation}
dunque al calore scambiato rapportato alla temperatura su un percorso reversibile. Essendo l'entropia definita da un differenziale esatto, ai fini del calcolo della sua variazione è sufficiente conoscere stato finale e iniziale: assumendo che siano dati, si riscrive il primo principio:
\begin{equation}
dU = \delta Q = TdS, 
\end{equation}
che, in equilibrio a una temperatura $T$ costante, equivale a:
\begin{equation}
d(U-TS) = 0.
\label{eq:prim_prin}
\end{equation}

\subsubsection{L'entropia}
Non esiste un modo univoco per definire l'entropia, l'unica costante è che rappresenta il "grado di disordine" del sistema. In questo senso, è legata alle combinazioni degli stati realizzabili da un sistema, in numero $\xi$ generalmente maggiore o uguale a 1. Se $\xi = 1$, il sistema può trovarsi esclusivamente in un singolo stato: il "disordine è nullo". \\

\noindent
Boltzmann ipotizzò dunque che l'entropia $S$ di un sistema deve essere proporzionale, a meno di una costante, al logaritmo\footnote{Il logaritmo garantisce la nullità nel caso di argomento unitario.} degli stati assumibili\footnote{Microstati.} da un sistema:
\begin{equation}
S = k_B \ln( W),
\end{equation}
 con $k_B$ la costante di Boltzmann. Sostituendo nella \eqref{eq:prim_prin}, si ottiene:
 \begin{equation}
d[U-Tk_B\ln( W)] = 0 \implies d{\left[ \sum_i n_i E_i  - Tk_B \ln\left( \prod_i W_i \right)\right]} = 0.
\end{equation}
Sviluppando il logaritmo della produttoria, ancora:
 \begin{equation}
d{\left[ \sum_i n_i E_i - \sum_iTk_B \ln (W_i)\right]} = \sum_id\underbrace{[n_iE_i- Tk_B\ln(W_i)]}_{\mathcal{F}(n_i)} = \sum_id\mathcal{F}(n_i) =  0.
\end{equation}
 dove si è sfruttata la linearità dell'operatore differenziale rispetto alla somma e si è definito il termine differenziato come $\mathcal{F}$ \textit{energia libera}. Dalla definizione di differenziale, poi: 
\begin{equation}
	\sum_id\mathcal{F}(n_i) = \sum_i \pdv{\mathcal{F}}{n_i} dn_i
\end{equation}
con $dn_i$ che rappresenta la variazione della popolazione elettronica nello stato di energia $E_i$. Dal momento che il numero totale di elettroni si deve conservare in un sistema isolato, $\sum_i dn_i = 0$. Perché quest'ultima condizione valga per qualsiasi variazione $dn_i$ e affinché la somma dei differenziali allora sia nulla si impone che il termine moltiplicativo $\mu = \partial_{n_i} \mathcal{F}$, detto \textit{potenziale chimico}, sia costante. 
\begin{equation}
\begin{aligned}
    \mu &= \pdv{\mathcal{F}}{n_i} = \pdv{}{n_i} \left[  n_i E_i - k_B T \ln\left( \frac{g_i!}{n_i!(g_i-n_i)!} \right)\right] \\
    &= \pdv{}{n_i} \left\{ n_i E_i- k_B T \left[ \ln(g_i!) - \ln(n_i!) - \ln((g_i-n_i)!) \right] \right\}.
\end{aligned}
\end{equation}
Sfruttando l'approssimazione asintotica di Stirling\footnote{Valida in questi casi essendo $n_i$ e $g_i$ molto grandi.} $\ln(n!) \approx n\ln(n) - n$ si ottiene:
\begin{equation}
\begin{aligned}
    \mu &= \pdv{}{n_i} \left\{n_i E_i - k_B T \left[ g_i\ln(g_i)- n_i\ln(n_i) - (g_i-n_i)\ln(g_i-n_i) \right]  \right\}. \\
     &= E_i - k_B T[ -\ln(n_i) - 1 + \ln(g_i-n_i) + 1]  \\
     &=  k_B T \ln\left(\frac{n_i}{g_i-n_i}\right) + E_i.
\end{aligned}
\end{equation}
Detta $P = n_i/g_i$ la probabilità che uno stato all'energia $E_i$ sia occupato da un elettrone, si ottiene la \textit{distribuzione di Fermi-Dirac} (f(E) = P):
\begin{equation}
\begin{aligned}
    \mu - E_i &= -k_B T \ln\left(\frac{g_i -n_i}{n_i}\right) \implies \frac{E_i-\mu}{k_BT} = \ln\left( \frac{1}{P}-1 \right) \\
    f(E) &= \frac{1}{1+e^{\frac{E_i-\mu}{k_BT}}}.
\end{aligned}
\end{equation}

\subsection{Analisi della distribuzione: il Livello di Fermi}
Si riporta in Figura (\ref{fig:fermi_dirac}) il grafico della distribuzione di Fermi-Dirac. Nel tentativo di dare una definizione più concreta del potenziale elettrochimico $mu$, si analizzi il caso di temperatura assoluta prossima allo zero $T \to 0$:  l'argomento dell'esponenziale va a $\infty$ se $E_i > \mu$ mentre a $-\infty$ nel caso contrario. La distribuzione assume dunque il valore unitario per valori di energia inferiori al potenziale elettrochimico e zero per energie superiori. In questo contesto, $\mu$, si configura come il \textit{massimo} livello energetico accessibile a un elettrone nel reticolo allo zero assoluto: questo livello prende il nome di \textit{Livello di Fermi} $\mu = E_F$. Si noti poi come, per $T > 0$, il Livello di Fermi corrisponda al livello energetico per cui si ha probabilità $1/2$ di occupazione. Infatti:
\begin{equation}
	f(E_F) = \frac{1}{1+e^{\frac{E_F-E_F}{k_BT}}} = \frac{1}{2}.
\end{equation}
Si definisce allora \textit{Livello di Fermi} anche il livello energetico che ha probabilità $1/2$ di occupazione a prescindere della temperatura.


\noindent
Il Livello di Fermi è una caratteristica intrinseca del sistema e, condizione necessaria perché si verifichi in un solido l'equilibrio, è la sua costanza spaziale. 

\subsubsection{Calcolo del livello di Fermi}
In un conduttore, la probabilità di occupazione degli stati in banda di valenza con energia inferiore a quella di Fermi allo zero assoluto è unitaria. Scegliendo come zero dell'energia il fondo della BV, si ottiene che la densità di elettroni $n$ che si trovano in stati energetici di energia $E>0$ è dato dal prodotto della densità di stati disponibili moltiplicata per l'annessa probabilità di occupazione:
\begin{equation}
\begin{aligned}
n = \int_0^{+\infty} f(E)g(E)dE =& \int_0^{E_F} \underbrace{f(E)}_{=1}g(E)dE + \int_0^{+\infty} \underbrace{f(E)}_{=0}g(E)dE = \int_0^{E_F}g(E)dE \\
n =& \int_0^{E_F} \frac{\sqrt{E}}{2\pi^2}\left(\frac{2m^*}{\hbar^2}\right)^{\frac{3}{2}} dE = \left(\frac{2mE_F}{\hbar^2}\right)^{\frac{3}{2}}\frac{1}{3\pi^2} 
\end{aligned}
\end{equation}
da cui il Livello di Fermi $E_F$:
\begin{equation}
E_F = \frac{\hbar^2}{2m}(3\pi^2n)^{\frac{2}{3}}.
\end{equation}


\end{document}